% 当前TeX为2026年10月12日补充公共品效率与政治风险敏感性后的编辑源稿；旧Markdown尚未同步，勿以旧稿重建覆盖。
\documentclass[UTF8,fontset=windows,12pt,a4paper,oneside,openany]{ctexbook}
\usepackage[margin=25mm,headheight=16pt]{geometry}
\usepackage{amsmath,amssymb,booktabs,longtable,array,graphicx}
\usepackage{setspace,fancyhdr,xcolor,tikz,pgfplots,xurl}
\usetikzlibrary{arrows.meta,positioning}
\pgfplotsset{compat=1.18}
\usepackage[colorlinks=true,linkcolor=blue!45!black,urlcolor=blue!55!black,breaklinks=true]{hyperref}
\hypersetup{pdftitle={新技术对资本主义基本矛盾的影响},pdfsubject={学位论文体例研究稿；理论、公开证据与条件模型}}
\setstretch{1.32}
\setlength{\emergencystretch}{3em}
\setlength{\parskip}{0.25em}
\setcounter{tocdepth}{1}
\setcounter{secnumdepth}{2}
\ctexset{chapter={name={第,章},number=\chinese{chapter}},section={format=\Large\bfseries}}
\pagestyle{fancy}\fancyhf{}\fancyhead[L]{\small 新技术与资本主义基本矛盾}\fancyfoot[C]{\thepage}
\begin{document}
\begin{titlepage}\centering
\vspace*{25mm}{\Large 学位论文体例研究稿\par}
\vspace{18mm}{\Huge\bfseries 新技术对资本主义\\基本矛盾的影响\par}
\vspace{12mm}{\Large 从生产能力、市场实现与收益分配\\到社会再生产和决定权\par}
\vfill
{\large 作者与所属机构：待填写\par}
\vspace{8mm}{\large 2026年10月12日修订\par}
\vspace{8mm}\begin{minipage}{0.8\textwidth}\small
本文结合理论辨析、公开资料与未经验校准的条件模型。作者、学校、资助与利益冲突事实尚未提供；本稿不声明通过学位审核。中文排版使用CTeX通用类，不冒称特定院校模板。
\end{minipage}
\end{titlepage}
\frontmatter


\chapter*{摘要与研究透明说明}


\section*{中文摘要}

新技术扩大生产能力，却不自动改变生产条件的归属和使用决定。本文结合理论辨析、跨行业公开证据与条件模型，研究技术节约如何转为实际收益，以及共同生产者能否决定其用途。得到四项条件性结论：第一，技术改变劳动力使用，可变资本实际使用额取决于产量、岗位工时与报酬，预付额另受周转约束；减员与在岗劳动强化可以并存，单位省工亦不必减少工资总额，不能按性能增幅推定增收或缩时。第二，知识公开与协作扩展可以同关键生产条件集中控制并存，平台规则、模型权限和制造接口影响参与者的实际选择。第三，企业降本应区分完整投入减少与家庭、公共部门新增的未补偿承担。第四，补偿责任、公共品效率和政治成本会反馈到制度选择；条件算例中，效率提高与新增任务均可反转资本方政策排序，风险厌恶的作用则取决于两种方案的失败风险。竞争算例进一步表明，个别采用有利可以与普遍采用后利润下降并存；采用动机、私人利润和社会收益因而需要分别判断。资本主义基本矛盾在上述过程中表现为共同生产能力与生产支配的分离。竞争、劳动保障和公共服务能够改善具体收益；占有关系转化则要求改变生产资料归属及投资、劳动组织和剩余用途的实际决定权。现有材料支持局部事实与条件解释，尚不足以证明新技术普遍强化资本控制，或所有智力工作必将转为资源依赖型工作。

关键词：生产社会化；资本主义占有；可变资本；技术收益；劳动分配；生产决定权。

\section*{Abstract}

New technologies expand productive capacity without automatically changing ownership of production assets or authority over their use. Combining theoretical analysis, public evidence across industries, and conditional models, this study examines how technical savings become realized benefits and whether collective producers can determine their use. It develops four conditional conclusions. First, technology changes the use of labor-power. The amount of variable capital employed depends on output, paid hours, and the composition of pay, while the required advance also depends on turnover. Workforce reductions can coexist with intensified work among those retained, while lower labor input per unit can coexist with a larger wage bill. Performance gains alone do not establish higher incomes or shorter hours. Second, public knowledge and wider collaboration can coexist with concentrated control over essential production conditions. Platform rules, model permissions, and manufacturing interfaces shape participants' feasible choices. Third, lower firm costs must be separated into reduced total inputs and additional uncompensated burdens on households and public institutions. Fourth, compensation obligations, public-good efficiency, and political costs feed back into policy choices. Conditional examples show that efficiency improvements and new labor-intensive tasks can reverse capital-owners' policy rankings. The effect of risk aversion depends on the failure risks of both policies. A competition example further shows that individually profitable adoption can coexist with lower profits after all firms adopt. Adoption incentives, private profits, and social gains therefore require separate assessment. The contradiction between socialized production and capitalist appropriation takes concrete form in the separation between collective productive capacity and authority over production. Competition, labor protection, and public services can improve specific benefits; transforming appropriation relations additionally requires changes in asset ownership and actual authority over investment, work organization, and surplus use. The evidence supports local findings and conditional explanations, but does not establish that new technologies universally strengthen capital control or that all intellectual work will become resource-dependent.

Keywords: socialized production; capitalist appropriation; variable capital; technological gains; labor distribution; production decision rights.

\section*{数据、代码与材料可得性}

本目录保存完整章节、逐章来源核查、模型参数和Python代码、CSV/JSON原始计算输出、PNG/SVG图、概念节点边表及候选DAG文件。公开外部研究的资料获取范围分别见台账；本报告没有取得、处理或重新估计相关研究的个人微观数据。已保存的公开原文快照只按实际可访问范围用于核查，正文提供原出版方或作者入口。第三方材料的开放访问不被理解为无限转载许可。

篇幅工具排除标题、书目、代码、图片标签与URL目标，仅统计本目录实际章节正文中文汉字。旧报告、重复版本、来源台账和本说明不累计。统计与链接核对不能证明学术质量，最终完成仍须内容和证据审查。

\section*{AI使用、作者与责任}

报告由Codex在用户要求下辅助检索、核查、组织论证、撰写、翻译、构建模型和绘图，并使用协作代理承担部分章节和审读。AI输出与引用均可能有误；实际核查范围、访问失败、版本差异及未解决问题保留在台账与审校记录中，没有宣称逐句人工核读全部引用文献。

人类作者姓名、所属机构和最终作者署名尚未提供，本报告不自动为用户填写学术作者身份。工具生成的综合框架、条件命题和探索模型属于本项目生成内容，不冒充某项已发表研究。若用于正式学术提交，人类作者需完成适用的独立审查、来源核读和机构要求；本文不标为已通过学位审核的博士论文。

\section*{作者贡献说明}

用户提出研究主题、篇幅和证据要求。Codex辅助完成研究设计与文本整合、公开资料检索、证据登记、情景模型、概念图和初步审校。协作代理辅助理论、劳动、制造、能源、交通、半导体和社会再生产章节及审读。以上为实际工具分工说明，不自动对应人类作者的正式贡献认证或作者顺序。供后续人类作者核定的CRediT活动类别包括Conceptualization、Methodology、Software、Validation、Visualization及Writing；工具不作为学术作者署名。

\section*{伦理、利益冲突与资助}

本报告使用公开文本、公开统计说明和自拟情景，没有进行访谈、干预、招募或新采集个人敏感数据；没有以机器转写视频替代准确原文引证。不存在可供本报告声称的新人体研究审批编号。涉及医疗与法律的材料仅分析原研究和明确版本的制度安排，不提供个人诊疗或具体法律资格判定。

人类作者的利益冲突及资助情况尚未提供，因此标记为“未提供，待作者确认”，不擅自填写“无利益冲突”或“无资助”。第三方临床研究、企业披露和机构报告的资助、作者角色与利益关系在其读取范围可识别时分别记录，不能因存在机构来源就视为无偏。

\tableofcontents\mainmatter

\chapter{导论：共同生产能力、收益分配与生产决定权}\label{ch:1}

\section{生产率提高，劳动者未必同时受益}

技术进步可以减少单位产品所需的劳动，却不规定节省的工时怎样使用。企业可以缩短工作日，也可以扩产或减少用工；增加的产出可以通过降价惠及购买者，也可以用于积累。\textbf{技术提供了新的生产可能，市场条件与生产决定权规定这些可能怎样实现。}生产能力由劳动、知识和设施共同形成，而投资、劳动组织和成果用途未必由参与者共同决定。这种分离构成本文的研究问题。

生成式AI辅助客服的研究为考察任务生产率提供了具体证据。Brynjolfsson、Li与Raymond在2025年正式发表的论文中，利用一家企业5,172名客服人员的分批部署资料，估计AI辅助使每小时成功解决的问题数平均提高约15\%。该指标反映既定客服工时内可完成任务数量的增加。由此进一步分析劳动后果，还需考察企业对新增生产能力的使用方式，以及劳动者参与收益分配的具体条件。\href{https://doi.org/10.1093/qje/qjae044}{正式论文}

工时需求与任务总量的数量关系可以说明这一分析的必要性。假如某项完整工作流程的生产率确实提高15\%，而任务总量保持不变，那么所需工时变为原来的\(1/1.15\)，减少约13.04\%。若任务总量也增加15\%，工时便维持不变；若任务增长更快，总工时还会增加。这些数字是基于比例关系的情景演算，不能据此认定上述客服研究已估计出相同的就业效应。生产率反映单位任务所需劳动，需求则影响任务总量；工资、就业岗位和闲暇的变化，还取决于两者的共同作用及相应劳动安排。

类似的分离也见于电力、医疗和制造。设备降本要经过融资、电网和交易，才可能降低用户支出；疗法有效要经过筛查、支付和诊疗供给，才可能使患者受益；分散加工要具备设计、数据和认证条件，才可能形成独立交付能力。技术效果与社会结果之间的这些环节，是跨行业比较的对象。

本文依次考察企业采用、市场实现、劳动分配和持续生产条件，说明技术收益怎样形成、由谁取得、由谁决定用途。这一分析承认技术改善的实际价值，同时解释改善为何可能与既有占有关系的延续并存。

\section{技术节约可以转为工资、闲暇或经营余额}

设一家工厂生产同一种产品。改造前，每工时生产10件，每工时工资30元，材料等其他投入为每件2元，产品售价8元。生产并售出100件，需要10工时，支付工资300元、其他投入200元，余下300元。此处的余额尚未扣除固定设备费用、折旧、利息和税，因此称为经营余额。

改造使每工时产出由10件提高到15件，其他物理投入保持不变。表1-1暂不计入设备购置费用，分别考察生产率、价格和工资的变化。全部数值都是解释机制的假设算例。

\begingroup\footnotesize\setlength{\tabcolsep}{4pt}\renewcommand{\arraystretch}{1.25}
\begin{longtable}{>{\raggedright\arraybackslash}p{\dimexpr 0.23636\textwidth-8.0pt\relax}>{\raggedright\arraybackslash}p{\dimexpr 0.12727\textwidth-8.0pt\relax}>{\raggedright\arraybackslash}p{\dimexpr 0.12727\textwidth-8.0pt\relax}>{\raggedright\arraybackslash}p{\dimexpr 0.12727\textwidth-8.0pt\relax}>{\raggedright\arraybackslash}p{\dimexpr 0.12727\textwidth-8.0pt\relax}>{\raggedright\arraybackslash}p{\dimexpr 0.12727\textwidth-8.0pt\relax}>{\raggedright\arraybackslash}p{\dimexpr 0.12727\textwidth-8.0pt\relax}}
\toprule
\textbf{情形} & \textbf{销量：件} & \textbf{时薪：元/工时} & \textbf{售价：元/件} & \textbf{总工时} & \textbf{工资总额：元} & \textbf{经营余额：元} \\
\midrule\endfirsthead
\toprule
\textbf{情形} & \textbf{销量：件} & \textbf{时薪：元/工时} & \textbf{售价：元/件} & \textbf{总工时} & \textbf{工资总额：元} & \textbf{经营余额：元} \\
\midrule\endhead
\bottomrule\endfoot
改造前 & 100 & 30 & 8 & 10.00 & 300 & 300 \\
改造后，销量与工资价格不变 & 100 & 30 & 8 & 6.67 & 200 & 400 \\
改造后，降价但销量不变 & 100 & 30 & 7 & 6.67 & 200 & 300 \\
改造后，时薪与生产率同比提高 & 100 & 45 & 8 & 6.67 & 300 & 300 \\
改造后，降价并扩大销量 & 150 & 30 & 7 & 10.00 & 300 & 450 \\
\end{longtable}\endgroup


\textbf{表1-1　技术收益的不同去向。} 工资总额等于时薪乘以总工时，经营余额等于销售收入减去工资和每件2元的其他投入。最后一行假设降价后的市场能够吸收150件产品，尚未由需求方程证明。第十四章将放开这一假设，通过需求和竞争求出销量。

第二行以更少工时完成原有订单：企业经营余额增加100元，工资总额减少100元。工时减少若集中于部分人员，可能表现为岗位收缩；若分摊到全体人员，则可能表现为人人缩时、人人减收。表中总量不能辨认具体安排。

第三行在销量不变时降价，购买者少付100元，经营余额回到原水平。第四行以提高时薪维持工资总额，同时减少总工时。第五行在降价后扩大销量，用完原有10工时，工资总额不变，经营余额增加。三行分别展示价格传导、劳动分享和需求扩张的不同作用。

同一技术可以产生不同的工时、工资和经营余额，但表中各行只是给定条件下的分配方案。采用费用限制预算，市场需求限制扩产，劳动合同与协商限制报酬调整。因而，算例不仅说明结果可以分化，也引出后文要解释的条件：哪些方案可行，谁能从中作出选择。

用符号表示，设\(Q\)为产出，\(a\)为每工时产出，\(H\)为总工时，则

\[
H=\frac{Q}{a},\qquad
W=wH,\qquad
\Pi^{\mathrm{op}}=pQ-W-mQ.
\tag{1-1}
\]

其中，\(W\)为工资总额，\(w\)为时薪，\(p\)为售价，\(m\)为每单位其他投入，\(\Pi^{\mathrm{op}}\)为经营余额。第一式算完成产量需要多少工时，后两式算工资和经营余额。公式列清了账，价格、工资和产量为什么取这些值，还要到市场竞争和劳动安排中去找原因。生产关系的分析，就是要说明增加了哪些选择，又由谁作选择。

\section{共同参加生产，不等于共同决定生产}

资本主义基本矛盾是社会化生产与资本主义占有之间的矛盾。生产社会化体现为劳动过程的结合、社会分工的扩展和生产条件的相互依赖；资本主义占有则使生产资料作为资本运行，产品及其剩余由资本关系支配。问题不仅在于收益分得是否均匀，还在于直接生产者与生产条件相分离，共同劳动形成的能力受资本增殖目的组织。恩格斯在《反杜林论》第三编第二章据此分析生产方式与占有形式的冲突。\href{https://www.marxists.org/archive/marx/works/1877/anti-duhring/ch24.htm}{原著}

本文以“共同生产与生产支配的分离”考察这一矛盾的具体表现。多家机构参与设备研发，说明能力依靠协作；某家机构掌握关键部件的供给条件，则涉及特定交易权力。二者不能混同：协作不必然产生排他控制，排他控制也不能仅凭参与者数量推断。后文分别检验资本对雇佣劳动的组织权与关键供应者对其他生产者的入口控制。

马克思对协作与机器体系的讨论，为区分生产力的形成和它在资本关系中的表现提供了起点。共同劳动增加的能力，需要一定的组织和协调；协调由谁实施、服从什么目标，则构成进一步的问题。\href{https://www.marxists.org/archive/marx/works/1867-c1/ch13.htm}{《资本论》第一卷“协作”}、\href{https://www.marxists.org/archive/marx/works/1867-c1/ch15.htm}{“机器与大工业”}的中文与英文版本章号差异在第二章说明。

价格、工资和利润反映不同结果。降价可能扩大消费，工资上涨可能改善生活，利润增加可能来自节约或收费变化。确认这些变化之后，占有关系的判断还需回答生产条件由谁掌握、劳动与剩余用途由谁决定。分配指标能够显示得失，不能独自说明支配关系。

全文围绕资产归属、使用准入、生产控制、剩余处置和风险责任展开。公共出资不必然带来公共决策，员工分红不必然带来表决权，免费文件也不必然带来独立运行能力。五类事实的区分，将笼统的“控制”落实为可以核查的权利和行为。

第二章先把价值、财富、价格与利润分清。本章和后文的数量模型，算的是产量、成本和货币收入，用来说明具体的经济过程。模型变量不能因为名称相近，就替代马克思主义范畴的论证；理论概念也不能免去对每一笔投入和收入的核算。

\section{技术进步扩大选择，生产关系决定谁能作出选择}

本文的中心判断是：\textbf{新技术扩大共同生产能力，却不自动改变生产条件及其使用决策的归属。}这使效率提高、所得改善与占有关系延续能够同时发生。论文据此分析基本矛盾在三个层面的表现：劳动过程中，节省的时间受用工安排支配；生产组织中，参与者的选择受关键条件及准入规则约束；社会再生产中，收益取得与维持生产的费用承担可能分属不同主体。

在劳动过程层面，单位产品省时经过销量、报酬和排班形成实际结果。需求扩张可以吸收节省的工时，降价可以把收益传给购买者，工资与工时协商可以把节约转为收入或闲暇。因此，生产率是收益形成的条件，具体分配还取决于市场与劳动安排。

在生产组织层面，知识、设施和接口的互补性规定了生产需要什么；产权、合同与治理规定了谁能使用和改变这些条件。只有必要投入由特定主体控制、使用者缺少可行替代时，技术依赖才可能成为入口支配。雇主对劳动过程的支配则另有雇佣关系基础，不以关键投入垄断为前提。

在社会再生产层面，企业账户不涵盖全部照护、训练和维护投入。同等产出下，完整投入减少构成资源节约；投入转由家庭或公共部门承担，则改变费用分配。两者可以同时发生，需要按相同服务、时期和补偿条件比较。

\section{本文解释收益形成与生产支配的分离}

本文解释技术节约怎样形成参与者的实际收益，以及收益形成为什么不必然改变生产支配。经济分析追踪需求、竞争和报酬，生产关系分析追踪投资、劳动组织、准入和剩余处置；社会再生产分析进一步追踪维持上述过程的必要投入。

本文的工作在于建立这三个层次之间的可核查联系。行业资料提供不同环节的证据，条件模型检验若干推论，制度比较区分改变分配、扩大准入和改变决定权的措施。本文不主张发现了贯通全部环节的新因果效应，也不将跨行业材料拼接为同一条已识别的经验链。

\section{现有理论的分工与本文的连接}

任务框架解释技术怎样改变劳动与资本承担的任务。Acemoglu与Restrepo区分自动化的置换作用和新任务的恢复作用，使劳动需求分析超越单一“生产率提高”的描述。本文借助这一框架组织第五章的证据，再进一步追问新增任务的入口、学习条件和收益决定权。\href{https://doi.org/10.1257/jep.33.2.3}{原论文}

通用技术研究将下游互补创新纳入技术扩散分析。Bresnahan与Trajtenberg的研究问题提示：一种基础能力要产生广泛作用，需要其他部门进行调整和创新。本文由此研究算力与电力、制造与认证、疗法与医疗设施之间的互补关系；这些关系的存在并不要求把十项技术都归为同一种通用技术。\href{https://www.nber.org/papers/w4148}{原工作论文及发表信息}

产权理论讨论不完全合同下的剩余控制权，以及所有权配置对投资激励的影响。Grossman与Hart据此解释企业边界。本文使用这一理论辨认具体资产控制，却不把剩余控制权等同于马克思主义的剩余价值或资本主义占有；后两者涉及资本与劳动的生产关系，解释范围不同。\href{https://doi.org/10.1086/261404}{DOI}，\href{https://dash.harvard.edu/server/api/core/bitstreams/7312037c-527a-6bd4-e053-0100007fdf3b/content}{作者机构公开原文，期刊第691—694页}

共同资源治理研究则为制度比较提出另一项问题：多个使用者如何形成、执行和修改共同规则？Ostrom的研究议题有助于说明，公共性需要实际治理安排和问责条件。本章依据出版社的研究说明定位这一问题，具体制度效果在第十五章依靠相应经验材料分析。\href{https://doi.org/10.1017/CBO9780511807763}{出版社资料}

制度还可能成为经济变化的结果，而非始终给定的外部条件。Acemoglu、Gitmez与Shadmehr（2026）的\emph{Automation and Repression}为连接技术选择与政治回应提供理论参照。\href{https://doi.org/10.3386/w35336}{工作论文}。本文据此补问：收入和补偿责任变化后，资本方为什么接受协商、限制部署或试图绕过政治约束？第十四章先以自有成本与参与约束推导选择，再检验公共品效率和政治风险变化能否改变排序。所得为条件推论；工作论文定理、修改假设后的反例和现实制度变化分别判断。

任务框架、通用技术研究与产权理论分别解释任务配置、互补投入和资产控制，马克思主义政治经济学则考察这些过程所处的资本关系。本文按问题使用相应分析工具，并明确连接条件；理论并置本身不构成统一证明。

\section{案例选择与比较逻辑}

十项技术按机制差异选取。案例比较旨在考察解释框架在不同物质约束和组织条件下的适用范围，不据此估计全部新技术的平均影响。表1-2给出各案例在论证中的任务。

\begingroup\footnotesize\setlength{\tabcolsep}{4pt}\renewcommand{\arraystretch}{1.25}
\begin{longtable}{>{\raggedright\arraybackslash}p{\dimexpr 0.23077\textwidth-8.0pt\relax}>{\raggedright\arraybackslash}p{\dimexpr 0.30769\textwidth-8.0pt\relax}>{\raggedright\arraybackslash}p{\dimexpr 0.46154\textwidth-8.0pt\relax}}
\toprule
\textbf{技术及主要章节} & \textbf{首先改变的生产条件} & \textbf{需要进一步解释的转换} \\
\midrule\endfirsthead
\toprule
\textbf{技术及主要章节} & \textbf{首先改变的生产条件} & \textbf{需要进一步解释的转换} \\
\midrule\endhead
\bottomrule\endfoot
工业机器人，第五章 & 部分任务由设备承担，单位直接工时改变 & 成本节约、规模扩张和新任务怎样共同决定劳动需求 \\
生成式AI，第五、七章 & 知识调用与任务辅助成本改变 & 局部任务收益怎样经过组织调整、算力与接口实现 \\
数字劳动平台，第六章 & 信息匹配和远程协调范围扩大 & 匹配收益怎样受到收费、评价规则和退出条件影响 \\
增材制造，第八章 & 部分产品的成形方式和小批量成本改变 & 分散加工能否同时扩大设计、材料和认证的取得能力 \\
工业物联网与数字孪生，第八章 & 生产状态可见性和过程协调能力提高 & 数据互通带来的效率怎样与专用接口依赖并存 \\
光伏，第九章 & 部分发电的设备与运行成本改变 & 低设备成本如何经过融资、电网和交易成为可用低价电力 \\
锂电池与储能，第九章 & 能量在时段之间转移的能力增强 & 储能收入、系统价值和生命周期成本如何对应 \\
L4自动驾驶，第十章 & 特定运行范围内的驾驶任务改变 & 车辆效率怎样经过调度、道路和公共交通形成城市效果 \\
CRISPR基因编辑，第十一章 & 特定疾病的治疗可能性扩大 & 疗效怎样穿过制造、医疗容量和支付条件形成可及性 \\
先进EUV光刻，第十二章 & 先进芯片制造中的部分加工能力提高 & 高度互补的国际协作怎样产生关键节点和进入约束 \\
\end{longtable}\endgroup


\textbf{表1-2　案例分工。} 技术边界及数据口径在各章定义。工业物联网与数字孪生作为一个关联技术组，锂电池与储能作为一个技术组；生成式AI跨章讨论，仍只计一项。

案例比较保留各领域的有效测量：医疗使用患者分母、临床终点和随访期，能源区分投资、供电与收费，劳动区分任务、企业和地区结果。跨行业归纳比较的是节约实现、准入控制和费用承担的机制，不能合成一个没有共同分母的平均技术效应。

部分案例存在直接技术联系，如AI运行需要芯片与电力。其他案例主要承担机制对照。论文因而同时观察企业内部与相关生产网络，不以章节顺序假定十项技术组成一条统一因果链。

\section{用事实查明结果，用模型检验推论}

本文从原著、论文和公开原始资料中整理事实，比较案例，再用可以复算的模型核对比例、费用和推理条件。没有开展新的访谈、企业调查或因果估计。视频和机器文字稿曾提供问题线索，但正文中的判断须由原著和实际研究支持。

一种工具能做什么，企业是否采用，采用后各方得到什么，要分别查证。外部研究的人员、处理办法、时期和估计对象要说明，统计的单位、分母和版本要保留。任务做得更快，还没有证明工资提高；运行规模扩大，还没有证明成果归谁；采用成功的企业，也不能代表所有采用者。第三章据此划清每项证据能够说明的事情。

数理分析分成本与竞争、条件谈判、支出与现金核算三个模块，分别说明企业为什么采用、收益怎样分享、付款怎样落实。几个模块尚未联合校准，也没有联立成一般均衡。它们能查明设定条件下结果怎样产生、怎样改变；要解释现实行业和实际权限，还须取得相应资料。

\section{收益改善和占有关系改变须分别查证}

具体冲突的变化应由可比较结果判断：生产条件是否更难取得，净所得是否下降，必要负担是否增加，替代与纠错是否更困难。价格、服务或安全改善应按受益群体分别确认。生产规模扩大和单项平均数改善，都不足以独自确定冲突的方向。

判断占有关系有没有转化，要进一步查生产条件归谁、谁真正作决定。投资由谁提出和批准，劳动怎样组织，经营剩余怎样使用，错误安排能否纠正，都要有事实。收入提高、部分事项能够协商、生产条件共同所有，各有不同作用，不能彼此代替。第十五章比较这些制度，第十六章据已查明的结果形成结论。

\section{全文从劳动收益推进到共同条件与生产决定权}

全文按“节约怎样形成收益、参与者为何不能同等改变安排、生产怎样持续”推进。第1—4章界定理论与证据；第5—6章从工时和报酬进入规则权；第7—8章比较知识工具与制造条件；第9—12章并列检验共同设施、外部影响、服务进入和关键接口；第13章补进持续生产的必要投入；第14章检验经济推论；第15—16章形成制度回应和结论。表1-3列出各组交出的判断。

\begingroup\footnotesize\setlength{\tabcolsep}{4pt}\renewcommand{\arraystretch}{1.25}
\begin{longtable}{>{\raggedright\arraybackslash}p{\dimexpr 0.22414\textwidth-8.0pt\relax}>{\raggedright\arraybackslash}p{\dimexpr 0.50000\textwidth-8.0pt\relax}>{\raggedright\arraybackslash}p{\dimexpr 0.27586\textwidth-8.0pt\relax}}
\toprule
\textbf{章节} & \textbf{本组必须确立的判断} & \textbf{交给后文的未决问题} \\
\midrule\endfirsthead
\toprule
\textbf{章节} & \textbf{本组必须确立的判断} & \textbf{交给后文的未决问题} \\
\midrule\endhead
\bottomrule\endfoot
第1—4章：问题、概念、证据与机制 & 共同生产、收入分享和共同决定是不同关系；各项判断需要相应证据 & 这些区别在实际生产中怎样表现 \\
第5—6章：劳动与平台 & 省时不自动转为缩时或增收；订单和规则影响劳动者所得 & 规则权依托哪些生产条件 \\
第7—8章：AI与制造 & 开放材料、分散加工可以与上游条件控制并存 & 获得工具之后能否持续交付服务 \\
第9—12章：共同条件的四类比较 & 网络、公共空间、诊疗容量和跨国互补投入分别约束技术收益 & 这些共同条件及其长期费用由谁维持 \\
第13章：社会再生产 & 时间、现金与共同服务形成持续参与能力，责任安排影响完整收益 & 共同费用与市场收益怎样由不同主体承担或取得 \\
第14章：条件模型 & 采用有利、行业利润下降、产量增加和工时不变可以在给定条件下并存 & 改变哪些条件能够改变结果 \\
第15—16章：制度与结论 & 收益改善与生产决定权变化需要分别识别 & 对基本矛盾能够作出何种有范围的判断 \\
\end{longtable}\endgroup


\textbf{表1-3　章节承担的论证及其衔接。} 第5—8章由劳动结果逐步追踪控制条件；第9—12章是围绕共同命题的并列比较，不构成能源、交通、医疗和半导体依次引发的因果链。



第十四章以条件反例检验经济推论，第十五章将制度措施对应到所得、使用机会和实际决定权。两章的衔接依靠变量与机制的一致性，不把模型情景当作前述案例的现实预测。

\section*{本章参考文献}

Acemoglu, D., \& Restrepo, P. (2019). Automation and new tasks: How technology displaces and reinstates labor. \emph{Journal of Economic Perspectives, 33}(2), 3–30. \url{https://doi.org/10.1257/jep.33.2.3}

Bresnahan, T. F., \& Trajtenberg, M. (1992). \emph{General Purpose Technologies “Engines of Growth?”} NBER Working Paper 4148. \url{https://doi.org/10.3386/w4148} 。本章核查原工作稿摘要与发表信息，引用此版本，未把正式1995稿逐页核查与工作稿混同。

Brynjolfsson, E., Li, D., \& Raymond, L. (2025). Generative AI at work. \emph{Quarterly Journal of Economics, 140}(2), 889–942. \url{https://doi.org/10.1093/qje/qjae044}

Grossman, S. J., \& Hart, O. D. (1986). The costs and benefits of ownership: A theory of vertical and lateral integration. \emph{Journal of Political Economy, 94}(4), 691–719. \url{https://doi.org/10.1086/261404}

Ostrom, E. (1990). \emph{Governing the Commons: The Evolution of Institutions for Collective Action}. Cambridge University Press. \url{https://doi.org/10.1017/CBO9780511807763} 。本章只核查出版社书目与研究说明。

恩格斯（1878）。《反杜林论》，第三编第二章。\href{https://www.marxists.org/archive/marx/works/1877/anti-duhring/index.htm}{公开档案与版本说明}。档案路径标1877，正文按全书发表年1878引；版本辨析见第二章。

马克思（1867）。《资本论》第一卷，第十一、十三章；英文档案第十三、十五章。\href{https://www.marxists.org/archive/marx/works/1867-c1/index.htm}{版本说明}。中文译本页码未指定，采用章节定位。

\chapter{理论基础：共同生产与生产支配的分离}\label{ch:2}

本章解释技术能力为何不能直接决定所得与占有关系。分析从共同劳动形成的生产力出发，依次讨论劳动过程控制、资本预付、商品销售、收益分配与再生产。贯穿各环节的问题是：生产者怎样取得生产条件，生产由什么目标组织，成果和持续费用由谁处置。



第一章算例中的价格、工资和产量尚属给定条件，本章说明这些变量所处的关系。资本取得生产条件并购买劳动力，以增殖为目的组织劳动；生产者则通过工资取得生活资料。技术改变可行的生产方式，资本与劳动的关系规定这些方式怎样进入实际生产。

\section{共同劳动形成生产力，资本关系规定其使用}

恩格斯在《反杜林论》第三编第二章讨论现代生产力的社会性与资本主义占有之间的矛盾，联系了资本与劳动的对立，以及企业内部有组织生产和社会生产无政府状态的并存。这一论证所指的社会性，是产品和生产条件越来越需要共同劳动才能形成，而成果仍按照资本占有的关系处理。（恩格斯，1878，第三编第二章；\href{https://www.marxists.org/archive/marx/works/1877/anti-duhring/ch24.htm}{原著}）

生产社会化包含两种相互联系的形式：劳动者在一个过程内协作，以及不同生产过程通过分工和交换相互依赖。车间工序、数字系统和供应链体现这些联系，教育、科研与基础设施又提供其持续条件。“社会化”在此描述生产的组织性质，不意味着参加者已经平等，也不意味着资产已归公共所有。

马克思在《资本论》第一卷第十一章分析协作产生的集体生产力，以及这种生产力在资本组织下表现为资本自身能力的现象。集体生产力确实存在：有些任务只有共同完成才可能实现，某些设施可以由多人共用，从而减少每件产品的投入。将这种能力归于共同劳动，并不否定组织管理和工程改进的贡献，而是强调各项投入只有在相互配合中才能发挥作用。（马克思，1867，中文第十一章、英文第十三章；\href{https://www.marxists.org/chinese/marx-engels/23/021.htm}{中文原著}、\href{https://www.marxists.org/archive/marx/works/1867-c1/ch13.htm}{英文原著}）

协作增加以后，参与者能做什么，要看协作怎样组织。假定一家工厂把加工、检测和交货记录接入同一系统，返工减少了。采购方若能单方修改订单，供应者只能接收指令，供应者就可能要更频繁地被动调整。供应者若能保留质量记录、把记录用于多家客户，并与采购方约定变更补偿，也可能更容易独立经营。两种安排都连接数据，差别在于谁能利用数据作决定。这个情景用来说明机制；要判断现实中发生了哪一种，还需第八章讨论的合同、权限和迁移资料。

共同劳动所形成的能力，可以在资本组织下作为资本的生产力发挥作用。生产者参与生产，却未因此取得生产资料和产品的共同处置权。这一关系不以工资持续下降为成立条件；生活改善与资本对生产条件的支配可以并存。

\section{谁控制劳动过程，谁就能影响节省时间的用途}

在《资本论》第一卷第四章第三节，劳动力成为商品有特定条件：劳动者能够以法律上的自由主体出售劳动力，同时没有足够的生产条件使劳动独立实现。第五章继而讨论资本取得劳动过程的控制和产品所有权。工资支付与生产中的支配由此相连，即使买卖形式平等，生产者和生产条件之间仍有结构性分离。（马克思，1867，第四章第三节、第五章；\href{https://www.marxists.org/chinese/marx-engels/23/014.htm}{劳动力的买卖}、\href{https://www.marxists.org/chinese/marx-engels/23/015.htm}{劳动过程和价值增殖过程}）

完成同样的合格产出，技术省下的工时可以用于缩短工作日，也可以用于减少付酬工时或承接更多订单。新添的核验和维护还会吃掉一部分节省。市场和预算限制了可选办法；在这些限制以内由谁选择，仍要看资本与劳动之间的权利安排。

生产控制体现于可执行的权限：谁安排工时，谁规定绩效，谁批准投资，谁能修改或拒绝决定。知情、建议、协商和否决的作用不同。员工得到奖金，可能改善分配；取得工时否决权，则改变特定事项的决策程序。权限是否有效，还取决于拒绝后的就业与替代条件。

为查清这些决定，本文分别考察所有权、使用准入、生产控制、剩余处置与风险责任。每项都对应具体事项，不能把它们相加成为一个权力分数。

\begingroup\footnotesize\setlength{\tabcolsep}{4pt}\renewcommand{\arraystretch}{1.25}
\begin{longtable}{>{\raggedright\arraybackslash}p{\dimexpr 0.18919\textwidth-8.0pt\relax}>{\raggedright\arraybackslash}p{\dimexpr 0.43243\textwidth-8.0pt\relax}>{\raggedright\arraybackslash}p{\dimexpr 0.37838\textwidth-8.0pt\relax}}
\toprule
\textbf{关系} & \textbf{在技术采用中的具体问题} & \textbf{可能取得的资料} \\
\midrule\endfirsthead
\toprule
\textbf{关系} & \textbf{在技术采用中的具体问题} & \textbf{可能取得的资料} \\
\midrule\endhead
\bottomrule\endfoot
所有权 & 谁能转让设备、专利和其他资产 & 产权、许可、股权和章程 \\
使用准入 & 谁能接入服务，费用与资格怎样规定 & 接入合同、收费、资格要求 \\
生产控制 & 谁安排任务、工时、核验与停机 & 作业权限、排班、申诉处理 \\
剩余处置 & 谁决定增收用于工资、分红或再投资 & 预算权、协商安排、收益去向 \\
风险责任 & 失败和转换费用由谁承担 & 赔偿、维护、保险和保障条款 \\
\end{longtable}\endgroup


上述权利在同一生产网络中可能分属不同主体。设备所有者可能没有独立维修许可，操作者可能不能修改安全参数，公共出资者可能无权调整企业绩效。因此，所有权调查必须延伸到使用和决策，同时避免把持有少量股份直接理解为共同管理。

协调是共同劳动的必要职能，资本支配则是这一职能在特定生产关系中的组织形式。工序排序、质量记录和安全停机并不因为需要管理，就必须由资本所有者决定。马克思第三卷第二十三章区分共同劳动的协调与资本、劳动对立产生的监督。（马克思，1894，第二十三章；\href{https://www.marxists.org/archive/marx/works/1894-c3/ch23.htm}{原著}）相同系统可以减少等待，也可以增加纪律约束。评价这类系统，需要分别考察工序改进与指挥权限的变化，并分析协调收益在执行者与管理者之间的分配。

承担风险和支配收益，不能直接画等号。投资者可能损失本金，劳动者可能失去收入和多年积累的职业能力，社会可能承担故障和环境损失。各方分散风险、取得保险和改变决定的能力并不相同。确认投资有风险以后，还要说明谁决定风险大小、谁取得收益；确认劳动者受损以后，还要查损失在哪个采用或分配环节形成。某一方承担风险，并不足以证明全部剩余都应由这一方处置，仍须考察具体制度。

资本对雇佣劳动的支配与关键供应者的市场权力具有不同基础。竞争充分的企业仍可以依据资产和雇佣关系组织劳动、处置剩余；市场竞争会约束其销售与回报，但不自动把这些权限交给劳动者。关键供应者对独立生产者的入口控制，则取决于必要投入、替代条件和准入规则。前者不能用垄断程度衡量，后者不能仅由协作依赖推断。

\section{完整任务与资本预付限定采用回报}

技术能力只有与生产条件结合，才能成为稳定交付能力。《资本论》第一卷第十三章把机器放在相互连接的机器体系中分析，区分其在劳动过程中的全部效能和价值逐步转入产品的方式。（马克思，1867，中文第十三章第一、二节，对应英文第十五章；\href{https://www.marxists.org/chinese/marx-engels/23/023.htm}{中文}、\href{https://www.marxists.org/archive/marx/works/1867-c1/ch15.htm}{英文}）这一区分有助于考察现代技术的采用条件：软件或设备取得之后，人员配置、接口适配与质量保障共同影响潜在资源节约的实现。

设工具把文书草拟从8分钟缩至4分钟，核验却从2分钟增至3分钟，完整任务便从10分钟缩至7分钟。草拟用时减少50\%，完整任务用时只减少30\%。这是自拟算例，说明局部省时与完整任务节约的差别。核验负担还可能随版本和人员学习变化，采用评价应覆盖实际交付流程。

资本分类为区分各类投入的经济职能提供了依据。不变资本与可变资本按价值增殖职能区分，固定资本与流动资本按价值周转方式区分，不能把“固定资本与可变资本”合成同一组分类。前一组来自第一卷第六章，后一组见第二卷第八章。机器通常既是不变资本又是固定资本，原料属于不变资本并以流动方式周转，购买劳动力的资本为可变资本。这里分类的是资本在生产中的存在形式与职能，不是把劳动者本人当作一件归企业所有的资产。这两组分类回答不同问题。（马克思，1867，第六章；\href{https://www.marxists.org/chinese/marx-engels/23/016.htm}{中文}；马克思，1885，第八章；\href{https://www.marxists.org/archive/marx/works/1885-c2/ch08.htm}{英文}）

“不变”指生产资料在该理论中的价值转移角色，设备的市场价格当然会变；“固定”指跨周期使用和分批转移价值，亦不表示永远固定在一个地点。软件、研发和服务的分析还要指定边界。企业购买云服务与云服务商购买服务器不是同一项支出，前者价格可能包含后者的折旧、劳动和利润。会计上资本化或费用化，也不能直接决定政治经济学分类。须先辨明谁组织生产、生产什么、投入怎样周转，再研究其增殖职能。

采用门槛具有先支付、后回收的时间结构。沿第四章约定，将当期新增固定采用支出记为$F$，单位货币成本记为$c$。暂时给定价格和产销量，若$c$由80元/件降至60元/件，采用$F=300$元，10件产量只节省200元，20件节省400元。两者扣除采用支出之后分别为负100元、正100元。相同单位能力的改进，对不同订单量的企业具有不同财务意义。

上述比较暂时固定订单量，第十四章再引入价格与竞争。预期收益为正，也不意味着企业能够先支付采用费用；流动资金、抵押和信用会限制投资。共享设施若减少重复准备，可降低实际采用成本；若仅由财政代付，则主要改变付款承担。二者的作用不能混为一谈。

资本的技术构成、价值构成和有机构成由此需要分别理解。技术构成涉及劳动运用多少、什么生产资料；价值构成用不变与可变资本价值的比例表示；有机构成讨论由技术构成决定并反映其变化的价值构成。更强的设备若因制造改进而便宜，物质规模可以上升而价值投入未按同一比例增长。马克思第一卷第二十三章对积累和资本构成的讨论提供这一辨析基础。（马克思，1867，第二十三章；\href{https://www.marxists.org/chinese/marx-engels/23/034.htm}{中文}、\href{https://www.marxists.org/archive/marx/works/1867-c1/ch25.htm}{对应英文第二十五章}）

\subsection{可变资本购买劳动力，技术改变其使用条件}\label{sec:variable-capital}

\textbf{可变资本的“可变”，指劳动力在生产中创造的新价值超过资本为购买它所支付的价值，不是指工资随产量浮动。}第一卷中文第六章、对应英文第八章按价值增殖职能区分不变资本与可变资本。自动化改变两者的数量关系，并不使机器在这一理论中取得活劳动的价值创造职能。\href{https://www.marxists.org/chinese/marx-engels/23/016.htm}{中文原著}、\href{https://www.marxists.org/archive/marx/works/1867-c1/ch08.htm}{对应英文}

这一区分把分析推进到劳动力购买与使用：企业购买哪些人的劳动力、支付多少、安排多长时间、以何种强度组织工作。技术可能减少某种任务，也可能扩大生产、增加维护与核验，或改变招聘和培训要求。所谓生产资料替代劳动，是相同交付中的部分任务转由设备完成；原料本身不能像机器一样执行原有操作，其用量增加也可能来自劳动者在同样时间内加工更多材料。只有把投入变化接到完整任务与实际用工，才能说明可变资本怎样变化。设备投入增加或人员数量减少，本身还不能给出其变动额。

脑力劳动与体力劳动也不是资本分类标准。第一卷英文第十六章以资本经营的学校说明，生产性劳动不以制造物质产品为唯一条件；第三卷第十七章又讨论商人用于购买商业雇员劳动力的可变资本，同时区分纯粹流通劳动与直接创造剩余价值的劳动。\href{https://www.marxists.org/archive/marx/works/1867-c1/ch16.htm}{生产性劳动原著}、\href{https://www.marxists.org/archive/marx/works/1894-c3/ch17.htm}{商业利润原著}。因此，不能把软件、教育和诊疗雇员全部排除，也不能把所有领工资者都放进同一种增殖过程。本文先辨认生产、纯粹流通、公共服务和独立经营等关系，再按各自职能分析技术影响。

\subsection{工资支出、相对比重和资本预付分别变化}

在已经确认由资本购买生产性劳动力的范围内，按劳动归属期间核算的工资支出记录这一购买的货币额。记该范围为$P_t$，岗位组为$j$，付酬工时为$H_{jt}$，相应平均时薪为$w_{jt}$，则

\[
W_t^{P}=\sum_{j\in P_t}w_{jt}H_{jt}.
\tag{2-V1}
\]

本式以归属本期劳动的工资为观察指标，延期或提前支付另列现金时间表；福利、社会缴费与其他雇主劳动支出须另列并说明是否纳入，不能把包含这些项目的劳动成本与只含工资的账户直接比较。工资记录也不是劳动力价值的独立估计。全年工资支出是流量，维持这一支出所需垫付的资本是存量；二者虽有联系，不能直接相等。

技术对这笔支出有三条直接路径：减少完成既定交付的工时，改变实际交付规模，改变岗位构成和报酬。在单一产品、质量不变且付酬工时覆盖完整流程的简化下，$W=wQ/a$，因而$W_1/W_0=(w_1/w_0)(Q_1/Q_0)/(a_1/a_0)$。产出与平均时薪增长因子的乘积若超过生产率增长因子，工资支出仍会增加。多岗位生产则回到式（2-V1），分别计算减少的操作工资与新增的工程、维护、核验工资，不能沿用原时薪乘以新工时。

可变资本绝对额与其在资本中的比重还要分开。不变资本增加更快时，可变资本可以绝对增加而相对减少。若讨论$c_M/v_M$，两项必须采用相同资本边界、估值时点及预付口径，不能拿设备存量除以全年工资流量。设备价格下降、工资变化和周转加快都可能改变价值构成；技术构成趋向更密集的设备使用，尚不能独自证明有机构成按同一幅度上升。

\subsection{周转加快可以减少预付，而不减少全年用工}

第二卷第十六章区分预付可变资本与一定期间实际使用的可变资本。\href{https://www.marxists.org/archive/marx/works/1885-c2/ch16.htm}{原著}。在每次等额投入、回流后重复使用、全年完成$n_v$次周转的简化下，全年使用额$V_{M,\mathrm{年}}^{\mathrm{用}}$与预付额$v_M^{\mathrm{预}}$满足

\[
V_{M,\mathrm{年}}^{\mathrm{用}}=n_vv_M^{\mathrm{预}}.
\tag{2-V2}
\]

缩短生产、验收或销售回款周期，可以让同一笔资本更快再次用于劳动力购买。全年使用额不变时，所需预付额随周转次数增加而下降。这种下降不意味着劳动者少领了工资，也不意味着单位劳动创造了更多价值。生产率节约与周转节约须分别识别。

预付是资本循环中的概念，不表示劳动者一律在工作前拿到工资。第一卷英文第六章明确讨论先使用劳动力、以后支付工资的情形。\href{https://www.marxists.org/archive/marx/works/1867-c1/ch06.htm}{原著}。真实企业还存在工资支付、产品交付与收款日期的差异。数字结算若仅把资金更早交给劳动者，未必减少企业垫款；若缩短企业应收账款，则可能减少。具体需要用支付和回流时间表检验，不能只看应用显示的结算速度。

\subsection{生活资料生产率影响劳动力价值，工资调整另有条件}

另一条路径发生在劳动者生活资料及其生产投入中。第一卷英文第十二章说明，相对剩余价值的生产需要降低再生产劳动力所需的必要劳动时间；有关生产率提高须作用于必要生活资料或生产这些资料的投入部门。\href{https://www.marxists.org/archive/marx/works/1867-c1/ch12.htm}{原著}。某企业节省程序员工时，不能仅凭这一事实推定全社会劳动力价值下降；它还须通过相关商品与服务的社会生产条件发生作用。

这条关系以必要生活资料的范围和标准为条件。第一卷英文第六章指出，劳动力价值包含历史和道德因素。\href{https://www.marxists.org/archive/marx/works/1867-c1/ch06.htm}{原著}。本文据此将名义工资、实际购买力与劳动力价值分开：同额工资可以买到更多生活资料，可能改善生活；维持既定生活标准所需支出下降，也未必立即导致工资下调，合同、最低标准、劳动供求和协商仍在起作用。

工作日不变时，必要劳动时间缩短可以改变必要劳动与剩余劳动的划分；人数减少、工资压低、劳动强度提高则是另外的变化。经验上的利润增加或工资份额下降不能单独识别哪条路径，更不能直接测出剩余价值率。后文的行业分析依次检验劳动力购买、任务配置、周转及再生产条件，并保留现有资料的识别范围。

\subsection{机器替代与劳动强化可以同时推进}\label{sec:replacement-intensity}

\textbf{资本既可以减少某些劳动力的购买，也可以加强对留任人员的使用；减员与劳动强化能够在同一次技术改造中发生。}学院派Academia（2021）的《【硬核社会学】资本主义的本质是什么？》在约32—35分钟，将竞争推动机械化同资本加强劳动使用联系起来。\href{https://www.bilibili.com/video/BV1Jt4y1B72f/?t=1945}{视频相关段落}\footnote{本地保存的相关文字稿由faster-whisper机器转写，未经逐句听校，时间戳用于定位讨论；本文转述其问题线索，不将转写句子当作马克思原文。资本分类、价值增殖与利润率的依据分别回到原著核对。}。本文据此分别分析资本购买多少、什么劳动力，以及怎样安排其任务、时间和责任，避免把全部技术后果归结为工资总额增减。

机器替代与劳动强化的结合已有原著依据。第一卷“机器与大工业”第三节分别讨论工作日延长和劳动强化，第四节又讨论劳动者同机器体系的关系，因而不能将这种结合表述为马克思未曾考虑的后来现象。\href{https://www.marxists.org/archive/marx/works/1867-c1/ch15.htm}{英文第十五章，第三、四节}。本文据此区分技术效率与劳动耗费：工具使同等劳动完成更多产出，属于生产率改善；缩短休息、提高操作速度或并行承担更多责任，则可能增加同一钟表时间内的劳动耗费。实际变化必须查流程和工作规则，不能从每小时产出增加直接推定其中一项。

假设同样的合格产出原需10人，改造后6人即可完成，且其余条件不变。节省的任务时间可以用于降低工时、扩大产量，或减少岗位；如果企业又提高6人的任务定额、覆盖范围和响应要求，就同时调整了劳动购买与劳动使用。这是组织选择的条件例子，不是实际裁员比例。机器承担了哪些任务，与留任者是否获得休息、收入和决定权，需要分别查明。

剩余价值率说明资本怎样使用已经购买的劳动力。在劳动力按价值交换、劳动强度与复杂程度给定的简化下，以$v_M$表示同期实际使用的可变资本、$s_M$表示剩余价值，以$t_N$和$t_S$表示必要劳动时间和剩余劳动时间，则

\[
e_M=\frac{s_M}{v_M}=\frac{t_S}{t_N}.
\tag{2-V3}
\]

这属于价值理论中的关系，不是把工资账直接换算成时间。剩余劳动的无偿性指其创造的价值没有作为劳动力价值的等价物返还，不仅指未登记或未支付的加班；资本主义生产中正常付薪的工作日也包含必要劳动和剩余劳动。第一卷“剩余价值率”提供这一划分。\href{https://www.marxists.org/archive/marx/works/1867-c1/ch09.htm}{英文第九章}。因此，工资按时支付、名义工资提高或可变资本绝对增加，都不能单独证明剩余劳动比例下降；工资减少也不能单独测出它提高了多少。

高技能、高工资也不能独自决定这一关系。技术可能减少部分操作岗位，同时增加研发、维护和核验岗位；劳动者为了取得条件更好的岗位参与竞争，是其依靠出售劳动力生活的具体表现。较高工资能够改善生活，但剩余价值率还取决于新创造价值与劳动力价值的关系，生产支配则取决于谁安排劳动和处置剩余。学历、工资绝对额、剩余劳动比例与决定权分别回答不同问题，不能相互代替。

由此形成两条需要共同检验的路径：技术替代改变资本购买多少、什么劳动力；技术与工作规则结合，改变资本如何使用这些劳动力。机器人中的多机看管、平台中的派单考核、AI中的交付定额和先进工厂中的连续运行，分别给出可观察的组织环节。后文把它们与岗位、报酬和责任对应；涉及劳动力价值和剩余价值率时，则保留现有数据不能直接估计的范围。

\section{产量、价值、价格与利润具有不同含义}

资源节省首先改变有用产品的生产条件。《资本论》第一卷第一章将使用价值与价值、具体劳动与抽象劳动区分。在一定条件下，同一劳动时间可以生产更多有用产品，单位商品的社会必要劳动时间却下降。社会必要劳动时间是社会通常生产条件下所需时间，不能直接以某一家先锋企业的实测工时替代。（马克思，1867，第一章第一、二节；\href{https://www.marxists.org/archive/marx/works/1867-c1/ch01.htm}{原著}）

第一章的产量和工时例子，只能说明一定条件下可以生产多少，并没有测出商品价值。企业抢先降低成本，可能先取得优势；技术普及以后，同行的生产条件和售价又会变化。第一卷第十章关于相对剩余价值的论述，涉及个别劳动生产率与社会生产条件的这种联系。（马克思，1867，中文第十章、英文第十二章；\href{https://www.marxists.org/archive/marx/works/1867-c1/ch12.htm}{英文原著}）所以，要分别算抢先采用和全行业采用的结果。第四章也据此区分单家采用的激励与行业同步改进的后果。

产品完成生产后，还需通过销售实现货币收入。《资本论》第二卷第一章将货币资本循环写为$M-C\ldots P\ldots C'-M'$，其中购买生产条件、生产和销售属于不同环节。商品资本回到货币形式，才能支付后续费用并继续循环。（马克思，1885，第一章，尤其第一至第三阶段；\href{https://www.marxists.org/archive/marx/works/1885-c2/ch01.htm}{原著}）当代采用研究因而需要区分最大产能、合格产量、已售数量和回款。能够生产150件，并不保证有150件订单，也不保证买方按期付款。

价值与价格之间还存在数量偏离。第一卷第三章第一节讨论价格形式，第三卷第九章讨论平均利润和生产价格；单个部门利润与该部门创造的剩余价值不能直接画等号。（马克思，1867，第三章第一节；\href{https://www.marxists.org/archive/marx/works/1867-c1/ch03.htm}{价格形式}；马克思，1894，第九章；\href{https://www.marxists.org/archive/marx/works/1894-c3/ch09.htm}{生产价格}）本文用货币成本、收入和企业利润解释竞争时，分析的是货币表现及其配置。劳动价值论中的价值量另有理论含义，不能把两套量直接相减。

接入费体现货币支付与资源消耗的区别。下游每件支付10元，是下游成本，也是供应者收入；其中可能补偿能源、维护、研发和其他投入。生产网络的资源核算应继续识别这些实际投入，避免把服务投入与同一笔付款重复计入，同时保留费用和收益在主体之间的分配。

估值、营业利润与剩余价值亦须区分。估值包含对未来收益的预期，营业利润按当期会计口径形成，剩余价值则是资本主义生产关系中的理论范畴。三者可能相互影响，却不能因数值同为货币表现就直接相等；比较必须说明时期、边界和资本估价。

利润额与利润率的差别则可以用条件公式说明。此处以$c_M$、$v_M$、$s_M$专指马克思资本范畴中的不变资本、可变资本和剩余价值，和第四章单位货币成本$c$区别。在单周期、忽略周转与生产价格差异的简化下，$r_M=s_M/(c_M+v_M)$。令$e_M=s_M/v_M$、$k_M=c_M/v_M$，便有$r_M=e_M/(1+k_M)$。

若$c_M=80$、$v_M=20$、$s_M=30$，简化利润率为30\%。只把$c_M$提高至120，其余不变，利润率降至约21.4\%。若其他条件使$s_M$增加到45，利润率又升至约32.1\%。这些数字使用同一种抽象价值单位，展示分母增长与剩余价值变化的相互作用，没有对应现实企业。第三卷第十三、十四章分别讨论下降趋势和反作用因素，可与这类条件比较对照。（马克思，1894，第十三、十四章；\href{https://www.marxists.org/archive/marx/works/1894-c3/ch13.htm}{趋势}、\href{https://www.marxists.org/archive/marx/works/1894-c3/ch14.htm}{反作用因素}）

利润率下降还不等于利润总额减少。一个仅用于区分比例与总量的同口径算例是：预付资本100、利润20时利润率为20\%；预付资本扩大到300、利润45时，利润率降至15\%，利润总额却增加25。该例未设定资本构成和市场行为，不解释下降的原因。它只说明，讨论采用回报、积累和危机时，不能用利润率下降代替利润额、可支配现金或偿付能力的变化。

可变资本相对减少与利润率变化，还须同剩余价值率联立考察。在上述单周期、同一价值估价和忽略周转差异的条件下，若$e_0>0$，由$r=e/(1+k)$得到

\[
\frac{r_1}{r_0}=\frac{e_1}{e_0}\frac{1+k_0}{1+k_1}.
\tag{2-V4}
\]

若资本构成比值由4升至6，要保持原利润率，剩余价值率须达到原来的$7/5$。这只是条件阈值：它不说明企业已能或必将通过延长劳动、提高强度或压低工资取得这一增幅。第三卷第十三、十四章将下降趋势同反作用因素分别讨论，后者包括加强劳动使用、不变资本要素便宜等。\href{https://www.marxists.org/archive/marx/works/1894-c3/ch13.htm}{第十三章}、\href{https://www.marxists.org/archive/marx/works/1894-c3/ch14.htm}{第十四章}。视频所提示的竞争、替代和劳动强化，因而可组织成待检验的关系，不能缩成每一项自动化都造成利润率下降和危机的预测。企业竞争模型中的售价、收费与会计净利润变化，又不同于这里的价值构成关系。

研究现代企业资本收益率，还要查资本怎样估价、多久周转一次，以及税收和融资条件。多年设备的原值、当年折旧和当年工资，不能混作同一期的投入。统一预付资本和观察期之后，才可以比较利润额与利润率。设备性能提高多少，也不能按同一比例推出利润率提高多少。

\subsection{价值生产与市场实现具有不同条件}\label{sec:production-realization}

\textbf{生产能力扩大，不保证商品能够以支持资本继续运行的价格售出。}第三卷第十五章区分剩余价值的生产条件与实现条件：商品完成生产以后，还须通过销售转回货币形式。\href{https://www.marxists.org/archive/marx/works/1894-c3/ch15.htm}{原著相关讨论}。这使技术收益的判断包含两个环节：生产中是否减少必要投入、增加可交付产出；市场中是否形成相应订单、售价和回款。

商品过剩与剩余价值不是同一概念。前者涉及一定时期、产品和价格条件下的交付超过可实现需求，后者涉及资本主义生产中活劳动创造的新价值怎样划分。商品库存增加不能独自证明剩余价值总量增加；企业亏损也不能倒推劳动者没有提供剩余劳动。社会仍有需要而销售困难，是需要未能通过收入、支付资格或公共购买转化为相应需求的一种可能情形；产品不适用、质量不足和部门比例失调也可能阻碍销售，不能全部归因于居民收入。

技术竞争由此可能产生不同的压力。资本构成变化通过前述价值关系影响利润率；同行降价、订单不足和接入收费则通过具体交易影响经营利润。企业减少工资支出还可能改变消费需求，但投资、公共购买、外需、价格和信贷会调整这种作用。本文将这些路径分别核算，再考察它们是否在特定时期汇合；不把机器增加、工资下降和商品卖不出去排列成无需查证的必然链条。

\section{省下的成本怎样分，要看工资、工时和谈判}

劳动分配并非销售之后才开始。工资和用工安排在生产前及生产中形成，销量与回款再约束其维持和调整。成本节约能否成为奖金、缩时或培训投入，取决于预算和双方权限。第十四章的条件谈判模型将收入与工时暂时给定，用于分离外部机会和谈判权重的作用。

劳动者靠出售劳动力取得工资，工资涨不涨，还要看市场和制度。一项技术使任务更容易完成，劳动者可能更容易换工作，雇主也可能更容易换人。前一种变化增强劳动者的外部机会，后一种可能增强雇主的外部机会。因此，技能能否带到其他岗位、还有哪些职位、能否集体协商、生活有无保障，都会影响省下的成本如何分配。资本与劳动的关系，要经过这些具体条件才表现为工资结果。

有没有更多自由时间，要另查一笔时间账。产量保持100件，劳动时间减少而原收入不减，劳动者就可能多休息；付酬工时减少，工资总额也减少，生活反而可能更困难。名义工作时间即使缩短，如果待命、核验和家庭工作增加，真正可支配的时间也未必增加。第一卷第十三章第三节已讨论机器使用对工作日和劳动强度的影响。现代研究还要把待命、外包和家庭劳动计入，不能只记直接操作时间。

劳动者也要购买生活用品。产品降价可以提高工资购买力，住房、能源和照护涨价又可能抵消工资增加。因而，劳动条件、实际购买力和生产决定权须分别报告。同一项技术可能让消费者受益，也让某些劳动者收入减少，公共服务还可能补偿一部分损失。先查受益者与受损者是不是同一群人，再作评价；消费者平均受益，不能代替对特定劳动群体的调查。

分配还要看眼前与以后。设想一位新人借助工具更快完成任务，企业立即受益；如果原来供新人练习的简单工作全部取消，新人可能难以学会独立判断。工具也可能提供反馈、降低训练材料费用，让更多人得到练习。哪一种结果会发生，不能凭当期效率决定。第四章H4分别考察训练、独立能力和岗位入口，后续劳动研究再核查相应时期的证据。

劳动份额、实际收入与决定权可能不同向变化。份额下降时，收入仍可能随总产出增加；工资提高时，工作规则仍可能由管理层单方修改。因而，劳动分配应分别报告时薪、总收入、实际购买力与工时，生产关系则由可执行的投资、劳动组织和剩余处置权判断。

\section{再生产需要持续投入与支出衔接}

生产不能只做一期。资本要继续取得原料、维护设备、支付劳动、卖出产品；劳动者要有生活、照护、教育和健康条件，才能继续工作。所谓社会再生产，就要把这些持续维持的条件查清。当期成本下降以后，哪些费用确实省了，哪些换了付款人，哪些只是推到以后，都关系到生产能力能否维持。

例如，企业取消内部培训，当期费用可以下降。同行若仍在培训，它还可能招到熟练人员，暂时维持产能；多家企业都不培训，今后的技能供给却可能受影响。这种后果要经过人员流动和多年积累才能出现，不能拿一家企业一年的利润来证明，也不能只凭担忧认定它已经发生。要查培训投入、人员去了哪里、能否独立完成任务，以及以后还能否招到所需人员，岗位数量本身不够。

再生产还要求生产与支出衔接。工资、利润和公共收入分别进入消费、投资或储蓄，分配改变可能影响产品销售，但消费下降可以被投资、公共购买或外需补偿。第十四章的局部支出模型检查给定规则下的需求、预算与现金一致性；它不是对现实跨期循环的估计，不能据此把劳动收入下降直接判为必然销售危机。

马克思在《哥达纲领批判》第一部分先讨论生产资料补偿、扩大生产和储备，再讨论共同需要与个人消费。这一顺序表明，共同占有仍需决定积累和支出安排，社会总产品并非全部可以立即分给个人。（马克思，1875，第一部分；\href{https://www.marxists.org/archive/marx/works/1875/gotha/ch01.htm}{原著}）公共或合作方案同样必须支付设备、维护、培训和风险准备。其区别在于这些用途由什么程序决定，以及受影响者能否监督和纠正。

家庭照护等活动使人能够继续劳动，但不能把这种社会作用直接并入某家企业的剩余价值核算。可以从具体时间和费用查起：通勤少了，随时响应是否多了；照护工具究竟省了哪段时间；公共服务减少以后，谁多花了钱。这样既保留价值核算的范围，也把企业账目以外的维持条件纳入分析。

再生产比较应延长到维护、转换和损失责任。一项技术当期降本，可能同时改变未来停机、环境修复与转岗负担。保险和赔偿改变损失承担，预防投入改变风险发生条件，参与治理则可能改变风险决策。三种作用应分别识别，不能把企业账外发生的一切费用都判作成本外移。

\subsection{危机风险来自生产扩张与实现条件的失配}

个别资本有动力提高效率、扩大生产和争取市场，社会总体却需要各部门供给、支付与持续投入衔接。各企业分别作出的有利选择，可能汇合为重复投资、产能闲置或销售价格不足以覆盖持续费用。这构成从个别采用推进到总体风险的研究问题；能证明某项技术省时，尚不能证明上述失配已经发生。

市场扩张和信用可以支持新增需求，也可能使投资依赖尚未实现的未来收入。是否形成金融危机，须继续核对借款、偿债、资产价格和付款链；生态压力则须核对实际规模、单位资源耗费、排放与修复责任。销售困难、利润率压力、债务风险和环境负担可以相互影响，但并非同一指标。需要有对应过程和资料，才能说明它们怎样扩大或相互传递。

本文据此把危机讨论落实为可检验的条件：社会能够生产什么，资本在什么价格与回款条件下愿意继续生产，劳动者和其他主体取得的收入能支持哪些支出。生产能力与增殖条件的冲突具有制度含义，但不能仅由技术替代推定制度必然崩溃，也不能将一次市场调整视为全部矛盾已经消失。第十四章验证局部竞争和现金关系，尚未建立跨期债务与宏观危机模型。

\section{取得公开知识，还要有条件用起来}

在《政治经济学批判大纲》中，一般社会知识、固定资本与自由时间的讨论把知识积累和生产组织联系起来。“一般智力”可用于追问知识怎样成为现实生产能力，对当代大模型制度性质的判断仍需结合具体生产关系展开。（马克思，1857—1858，网站分段ch14，固定资本与一般社会知识段；\href{https://www.marxists.org/archive/marx/works/1857/grundrisse/ch14.htm}{原著}）

公开文件、开放许可与运行能力是三个不同条件。文件可读不必然允许修改或再分发；许可允许使用，也不保证材料、设备、技能和验证条件齐备。开放减少的是许可明确解除的限制，互补投入决定这一机会能否转为实际生产。

使用和换用技术，各有必要的费用。学会操作、完成任务需要劳动；换供应者以后，重新适配、验证和建立信用也需要劳动。两类费用可能重叠，不能把所有迁移费用都说成人为锁定。共同标准若让测试记录得到互认，可以少做重复验证；新服务若仍须原平台单方认可，接口即使公开，使用者也可能离不开原平台。第八、十二章将结合产品和流程查明这种差别。

获得免费使用许可可以减少许可支出，整合、托管和认证仍可能收费。采用者需要比较总费用，资源核算需要确认服务补偿的实际投入。开放知识与商业服务可以并存，生产自主性取决于可保留的知识、可替换的服务和实际决策权限。第七章的数学发布据此区分取得研究材料与持续调用产生材料的模型。（\hyperref[ch:7]{第七章}）

共同设施可以分摊固定投入，实际作用取决于容量、维护和分配规则。使用者能够接入，不代表能够决定研究方向、服务标准或成果用途。共同治理要落实为信息取得、议程提出、表决和监督程序；管理仍可分工，关键在于授权和问责由谁决定。

\section{分到收益、参与决定与共同占有，各有判断依据}

马克思在第三卷第二十七章讨论股份公司中的资本筹集、所有与经营分离，以及工人合作工厂在旧制度内部形成的新组织形式。股份社会化和生产者共同组织具有不同关系，合作工厂又受到既有环境约束。（马克思，1894，第二十七章；\href{https://www.marxists.org/archive/marx/works/1894-c3/ch27.htm}{原著}）这一论述提示，判断占有关系应同时考察生产条件的归属、组织内部的实际决定和外部依赖。

先看收益分享。企业增加工资或奖金，劳动者的生活可以改善；技术预算、岗位调整和再投资若仍由资本单方决定，劳动者与生产条件的基本关系就仍然延续。资本支配生产，并不要求工资一直下降。承认劳动者所得改善，同时判断资本关系仍在延续，两者并不冲突。

再看共同决定。协议若让劳动者能共同决定安全、工时或培训预算，就使原来由管理层决定的事项受到约束，超出了单纯分钱。这是规定范围内的生产关系调整。知情、建议、协商和否决不是同一种权利；发生争议时谁能改变决定，才是判断依据。“参与”或“民主”的名称不能代替执行结果。

占有关系的整体转化涉及生产者与生产资料的分离是否改变，投资、劳动组织和剩余用途是否形成共同支配。共同支配可以通过授权与分工运行，不要求所有参与者直接决定每项操作；判断依据是授权、监督和规则修正的实际权利。局部协商权、少量员工股份或公共补贴，只能按其确实改变的范围评价。

资产归谁，提供了权利基础；权利怎样用，还要看治理程序。合作组织拥有设备，成员有表决权，却仍可能离不开单一客户、封闭认证或昂贵融资。公共资产也要有预算、监督和使用程序，才能由公众实际控制。内部共同决定可以确实存在，外部交易条件也可以同时限制它。研究须据此划清转化范围：已经发生的内部变化应承认，还未改变的全社会关系也须说明。

集中度是产业结构指标，不是决定权的直接测量。第一卷第二十三章区分积聚与集中；现实中的市场份额、资产规模和规则权限也有不同含义。供应者较少，使用者仍可能具有稳定准入和可行替代；企业数量较多，也可能共同依赖单一入口。产权、合同与实际决策资料才说明这种结构怎样影响生产选择。

收益分享、共同决定和共同占有，相互联系，却不是改革必经的三个阶段。工资提高可以不改变治理，取得某项表决权可以暂时不增加收入；共同所有还要接受实际运行和外部依赖的检验。本文分别查所得、决策过程和生产条件归属。第三章说明各项判断要用什么证据，第四章再解释同样的能力为何在不同权利安排下产生不同结果。

\section*{本章参考文献}

1. 恩格斯（1878）。《反杜林论》，第三编第二章。访问版本：Emile Burns英译，Progress Publishers，1947；\href{https://www.marxists.org/archive/marx/works/1877/anti-duhring/index.htm}{MIA版本说明与目录}。
2. 马克思（1867）。《资本论》第一卷。访问版本：Samuel Moore与Edward Aveling英译，恩格斯编辑，1887；\href{https://www.marxists.org/archive/marx/works/1867-c1/index.htm}{英文版本说明}，\href{https://www.marxists.org/chinese/marx-engels/23/index.htm}{《马克思恩格斯全集》第二十三卷中文转录目录}。章号对应见正文及理论核查台账。
3. 马克思（1885）。《资本论》第二卷，恩格斯整理出版。访问版本：I. Lasker英译，Progress Publishers，1956，依据1893年修订第二版；\href{https://www.marxists.org/archive/marx/works/1885-c2/index.htm}{版本说明}。
4. 马克思（1894）。《资本论》第三卷，恩格斯整理出版。访问文本为MIA英文转录，所列来源包括Institute of Marxism-Leninism，1959及International Publishers；\href{https://www.marxists.org/archive/marx/works/1894-c3/index.htm}{版本说明}。
5. 马克思（1857—1858年草稿；德文首次出版1939—1941）。《政治经济学批判大纲》。访问版本：Martin Nicolaus英译，1973；\href{https://www.marxists.org/archive/marx/works/1857/grundrisse/index.htm}{版本说明与草稿目录}。一般智力段落按网站分段定位，不将网页编号冒充作者正式章号。
6. 马克思（1875年写作；1891年刊布）。《哥达纲领批判》，第一部分。访问版本来自Progress Publishers，1970，《马克思恩格斯选集》第三卷英文转录；\href{https://www.marxists.org/archive/marx/works/1875/gotha/index.htm}{版本说明}。

补充讨论材料：学院派Academia（2021）。《【硬核社会学】资本主义的本质是什么？》，2021年2月2日发布，P1“资本主义三幻神”，BV1Jt4y1B72f。\href{https://www.bilibili.com/video/BV1Jt4y1B72f/}{视频入口}。本文参考本地机器文字稿约32:25—34:43的相关论述，并结合20:56—25:20的劳动力购买与社会关系讨论；未逐句听校，不以其替代原著或因果证据。第一卷英文第九、十五章和第三卷第十三、十四章另于2026年10月11日核对相关段落；2026年10月12日补核第一卷英文第十五章第三节及第三卷第十五章的生产与实现相关段落。

各章页链接已嵌入相应论述。原有章节访问日期为2026年10月7日；此前新增核对第一卷英文第十二章、第二卷第一章，访问日期为2026年10月8日。本轮可变资本讨论另核对第一卷英文第六、八、十二、十六章、第二卷第十六章及第三卷第十七章，访问日期为2026年10月11日。未核查译本印页，不补写页码或DOI。


\chapter{证据比较：生产能力、收益分配与决定权分别验证}\label{ch:3}

\section{能力、收益和决定权，各要有相应证据}

本文区分三个结果：技术能力、实际所得与生产决定权。完整流程资料衡量能力，价格、工资和费用记录衡量所得，合同与决策记录衡量权限。\textbf{一项结果的证据不能替代另一项结果的证据。}本章据此规定资料比较与因果解释的范围。

第五章的客服研究能够说明既定工作流程内的任务产出提高，不能直接说明节余时间被用来缩短工作日；第九章的储能资料能够确认装机与使用同时扩大，不能直接确认消费者电费或项目利润的变化。两个例子表明，能力与利用证据有独立价值，收益判断还须找到实际付款、领取者及其成本。（\hyperref[ch:5]{第五章}；\hyperref[ch:9]{第九章}）

第六章的报酬规则与第七章的模型发布进一步涉及组织权限。报送收入上升确认报告账户中的收入改善，个人净额仍需扣除经营成本；平台能否继续决定订单与账户规则，又是不同的观察对象。数学成果公开扩大可供研究和核查的材料范围，研究者能否自主调用产生结果的模型，则影响其持续研究自选问题的能力。收入、知识取得与生产自主性由此进入不同的观察范围。（\hyperref[ch:6]{第六章}；\hyperref[ch:7]{第七章}）

结果变化与变化原因还需区分。采用者工资较高，可能来自技术，也可能来自原有管理、融资和订单优势。可信比较解决的是归因问题；缺少工资或权限资料，则是结果尚未被观测的问题。两种证据缺口需要不同补充。

本文没有贯通全部环节的新建微观面板。经验结果保留原研究的主体、处理、时期和识别条件，跨环节联系则明确所需假设。任务省时与工资变化若来自不同样本，可以支持机制讨论，不能据此认定前者造成后者。

权限证据同样具有层次。合同能够确认形式权利，执行和争议记录能够确认权利是否有效；只有可比的采用前后资料，才有条件进一步判断技术是否改变权利。本文分别报告这些层次，避免以某项控制存在推断其已因技术而增强。

\section{替代条件与实际权限界定生产依赖}

生产依赖应由必要投入及其替代条件描述。参与者数量反映协作规模，备份、接口兼容、迁移费用与转换时间反映替代能力。供应商众多的系统可能仍有唯一关键部件，供应商较少的系统也可能提供可替换服务。因此，规模与依赖不能互相代替。

谁支配生产，要从权利和行动查起。法律所有权、技术接入权、生产决策权、收益权和风险责任须分别记录。企业能否改变接口收费，劳动者能否拒绝特定绩效安排，公众能否审查投资用途，这些具体事项比一个笼统的“控制程度”分数更容易核查。表3-1据此列出所需资料。

\begingroup\footnotesize\setlength{\tabcolsep}{4pt}\renewcommand{\arraystretch}{1.25}
\begin{longtable}{>{\raggedright\arraybackslash}p{\dimexpr 0.12963\textwidth-8.0pt\relax}>{\raggedright\arraybackslash}p{\dimexpr 0.42593\textwidth-8.0pt\relax}>{\raggedright\arraybackslash}p{\dimexpr 0.44444\textwidth-8.0pt\relax}}
\toprule
\textbf{待解释环节} & \textbf{主要观察量或材料} & \textbf{比较时必须保持的口径} \\
\midrule\endfirsthead
\toprule
\textbf{待解释环节} & \textbf{主要观察量或材料} & \textbf{比较时必须保持的口径} \\
\midrule\endhead
\bottomrule\endfoot
生产能力 & 合格产出、完整流程工时、能源和材料投入、质量结果 & 同类任务或服务；质量与安全要求；直接和辅助投入 \\
采用门槛 & 部署时间、固定费用、融资条件、培训、互补设施 & 计划采用与实际使用；可获得与已投运；单期支出与存量 \\
市场实现 & 成交价格、销量、订单、客户流动、替代供应者 & 相同产品质量、市场边界、名义与实际价格 \\
劳动分配 & 总工资、实际时薪、就业人数、总工时、岗位入口 & 在岗者与全体受影响者；人均与总量；工时与人数 \\
控制关系 & 合同、表决权、接口规则、协商和申诉记录 & 形式权限与实际执行；规则改变前后；谁具有否决能力 \\
社会再生产 & 照护时间、培训机会、环境负担、公共支出 & 成本发生地点、承担群体、时间范围及服务质量 \\
\end{longtable}\endgroup


\textbf{表3-1　概念与证据的对应。} 表中是研究所需观察项，并非宣称已经取得每一个变量的完整数据。

平均数的解释取决于群体范围和权重。工资总额等于按工时加权的平均时薪乘以总工时；留任者平均工资上升时，退出者与失去工时的劳动者仍可能减收。判断全体受影响者的结果，应同时登记人员进入退出、工时与所得，不能只追踪留任者。

查明资产归谁以后，还要查权利怎样使用。股份分散，表决权未必分散；资产归公共所有，用途也可能仍由少数机构决定。要判断生产者能否共同支配条件，就追踪投资方案谁提出、谁批准、谁能否决、各方能取得哪些信息，以及决定能否执行。收入增加可以与这些权限原样保留同时发生。因此，本文分别报告收入、时间、服务和权限，不把它们合成一个没有权重依据的“矛盾指数”。

\subsection{工资账户识别购买支出，任务数据识别使用变化}

可变资本的经验分析需要连接三组资料。第一组记录谁购买劳动力：雇佣主体、生产职能、合同、外包和独立经营关系。第二组记录怎样使用：岗位工时、任务、强度、质量、合格产量及人员进入退出。第三组记录何时支付与回流：工资支付、在制品停留、验收和收款日期。第一组限定资本关系，第二组解释实际使用额，第三组解释预付和周转。缺少工资记录的任务实验，只能推进第二组的局部分析。

边界也要随外包调整。企业将原有岗位交给供应商后，自身工资支出可能下降，但供应商仍购买劳动力；外包价还包括设备费用、税费和利润。核算生产网络时，应分别登记各雇主的工资，不把外包合同全额当作工资，也不再把它与合同内工资重复相加。工时可以向另一个地区转移，不能仅由采购方减员认定社会劳动总量减少。

机器人安装、研发费用、产业增加值与电站装机各有用处，却不是可变资本的直接观测值。本文已有材料尚未形成跨行业一致的工资、生产职能和周转面板，因而不能估计所有部门统一的“可变资本替代率”。条件公式用于列明待检验关系，部门案例用于确认具体环节，缺失变量须保留为空，不能由设备价格或模型成绩补出。

\section{事实、因果估计与条件演算分别使用}

正文把数字分清来源和用途。装机、成本与利用记录说明规定范围内发生了什么；随机实验或准实验在其识别条件下，估计特定处理造成的变化；均衡、阈值与反事实演算则检查给定假设能推出什么。读者由此可以分清事实记录、因果估计和可能机制。

观测资料首先需要单位与时间一致。功率以GW计，储存能量以GWh计，两者相除才得到小时；全年发电量是流量，年末装机是某个时点的存量。若大量设备在年中投运，将全年发电量除以年末装机容量，会低估按实际投运时长计算的利用强度。其原因在于流量与存量指标的时间口径不一致，需要依据设备实际运行期间调整比较。

对百分数也应保留基数。设某指标由\(x_0\)变为\(x_1\)，增长率为

\[
g_x=\frac{x_1-x_0}{x_0}.
\tag{3-1}
\]

若指标本身是占比，变化还可以用百分点表示。占比由20\%提高到30\%，是提高10个百分点，也可以说相对提高50\%；二者分别描述比例指标的绝对变动与相对变动。本文在数据表中同时保存原值和计算式，避免只留下一个脱离分母的百分数。

使用外部研究的系数，要把对象和单位带上。每千名劳动者多一台机器人所对应的地区工资变化，不能写成全国每增加一台机器人的工资变化；任务生产率提高多少，也不能当成企业总成本下降多少。研究比较了谁、用了什么设计、按什么单位计量，都限制了结论可以推广到哪里。

演算中的观测输入与假设输入分别标明。第一章给定价格、工资和产量，第十四章给定需求与行为规则，均用于检验条件推论。第九章将报告中的安装成本与另行假定的寿命、折现率和运行条件分列，计算结果不冒充原报告估计。数据锚定部分输入，不能使其余设定自动成为观测事实。

数字怎样产生，就应怎样说明不确定性。抽样估计报告原研究给出的标准误、置信区间等信息；确定性演算检查参数改变后结果怎样变；行政统计说明覆盖范围、修订和定义。把三类数字都画成一种误差带，会混淆抽样误差、参数假设与统计覆盖。

\section{选择差异限制采用效应的归因}

本节先固定其他主体的采用状态；同行响应和干扰的记号扩展见第3.8.2节。企业局部处理效应不直接等于行业净效果。

因果问题需要比较同一单位在不同状态下会怎样。用\(Y_i(1)\)表示单位\(i\)采用技术时的结果，\(Y_i(0)\)表示它未采用时的结果，\(D_i\)表示是否实际采用。对同一时期的同一单位，只能观察其中一种状态。采用者的平均处理效应定义为

\[
ATT=E[Y_i(1)-Y_i(0)\mid D_i=1].
\tag{3-2}
\]

这里的\(E[\cdot]\)表示相应群体的平均值。式（3-2）所问的是：实际采用者如果没有采用，结果会相差多少。它没有把所有潜在使用者都作为同一群体。

现实中容易取得的是两组实际结果的差异。给这项差异加上再减去采用者的未采用潜在结果，就得到

\[
\begin{aligned}
E[Y_i\mid D_i=1]-E[Y_i\mid D_i=0]
={}&ATT\\
&+\{E[Y_i(0)\mid D_i=1]-E[Y_i(0)\mid D_i=0]\}.
\end{aligned}
\tag{3-3}
\]

第二行是选择差异：即使两组都不采用技术，结果也可能原本不同。这个分解说明，实际组间差异不全是采用的效果。采用者工资较高，只有在选择差异被消除或得到可信处理后，才可以归因于采用本身。

哪些企业先采用，往往正由本研究讨论的条件决定。大企业可能更容易融资，管理好的企业可能更容易完成改造，订单增长快的企业可能先投资设备。这些条件既影响采用，也影响后果。理论既然认为固定成本和共同设施限制进入，经验比较就不能又把采用当成随机事件。

随机分配针对的是被分配的处理，未必是实际采用。若随机开放工具，有人不使用或通过其他渠道取得工具，按原分组比较首先对应开放资格的影响；直接比较实际使用者与未使用者，可能重新引入选择。转述实验时因此保留处理定义、依从情况与样本范围，不把开放、培训和实际使用混为同一干预。

成功案例适合追踪流程、费用和权限，不能代表采用失败、融资受阻或退出者。本文用它们说明机制如何运行，群体平均结果则依靠覆盖相应对象的研究；两类材料不互相代替。

\section{双重差分依赖平行趋势与处理定义}

如果能同时观察采用组和比较组的前后变化，就可以先扣除两组共同经历的时间变化。表3-2是纯粹用于解释方法的假设数据，结果变量为同口径产出指数。

\begingroup\footnotesize\setlength{\tabcolsep}{4pt}\renewcommand{\arraystretch}{1.25}
\begin{longtable}{>{\raggedright\arraybackslash}p{\dimexpr 0.25000\textwidth-8.0pt\relax}>{\raggedright\arraybackslash}p{\dimexpr 0.25000\textwidth-8.0pt\relax}>{\raggedright\arraybackslash}p{\dimexpr 0.25000\textwidth-8.0pt\relax}>{\raggedright\arraybackslash}p{\dimexpr 0.25000\textwidth-8.0pt\relax}}
\toprule
\textbf{组别} & \textbf{采用前} & \textbf{采用后} & \textbf{前后变化} \\
\midrule\endfirsthead
\toprule
\textbf{组别} & \textbf{采用前} & \textbf{采用后} & \textbf{前后变化} \\
\midrule\endhead
\bottomrule\endfoot
采用组 & 100 & 124 & 24 \\
比较组 & 80 & 88 & 8 \\
两组差异 & 20 & 36 & 16 \\
\end{longtable}\endgroup


\textbf{表3-2　双重差分的算术例子。} 数值为假设，不对应任何真实企业。

采用后的两组差异为36，包含了采用前已经存在的20；采用组自己的增量为24，又包含可能共同经历的8。双重差分为\(24-8=16\)。若采用组在没有采用时也会增加8，它的反事实结果便是108，实际124比反事实高16。这个计算之所以具有因果含义，依靠的是“未采用时两组会有相同变化”的假设，而不是差分运算本身。

若采用组原本就因订单增长会增加12，则真实采用效应是\(124-112=12\)，双重差分会多算4。这一例子解释了为什么研究需要考察采用前趋势、需求冲击和预期采用。采用前趋势相近可以提供支持，但不能观察到采用后未发生的反事实，因此也不能单独完成证明。

企业在不同年份采用时，比较对象还会变化。Callaway与Sant’Anna以“某一采用批次在某一时期的平均效应”为基本单位处理这种问题，并讨论条件平行趋势、比较组选择和效应汇总。\href{https://doi.org/10.1016/j.jeconom.2020.12.001}{正式论文}，\href{https://pedrohcgs.github.io/files/Callaway_SantAnna_2020.pdf}{作者稿第5—7页}的定义和识别设定为本节提供方法依据。

借用这一记法，令\(G_i=g\)表示单位在\(g\)期首次采用，\(Y_{it}(g)\)表示从该期采用时的结果，\(Y_{it}(0)\)表示在观察期间始终未采用的结果，则

\[
ATT(g,t)=E[Y_{it}(g)-Y_{it}(0)\mid G_i=g].
\tag{3-4}
\]

式（3-4）区分刚采用与采用三年后的动态效应，同时保留早期和晚期采用群体之间的异质性。对于一旦采用就持续使用的研究，已经采用的单位未必适合继续作为未处理比较组；若技术会反复停用、升级或退出，则还需要重新定义处理过程。

本研究没有据此实施新的双重差分估计。公式的作用，是为阅读现有采用研究提供明确问题：作者比较了谁，比较期有多长，采用之前是否存在预期调整，以及估计结果适用于哪个批次和时期。这些识别条件和样本范围，比方法名称本身更能说明估计结果的适用性。

\section{混杂、中介与反馈需要不同识别条件}

因果比较需区分采用前混杂与采用后变化。原有管理能力可能同时影响采用和生产率，任务重组则可能是采用产生的结果。将后者直接加入回归，不仅可能截断部分传导，还可能引入处理后偏差；剩余系数不能自动解释为直接效应。直接效应需要明确干预对象和额外识别条件。Montgomery等以分析、模拟和经验重分析讨论了这种偏差。\href{https://doi.org/10.1111/ajps.12357}{正式论文}，\href{https://cpb-us-e1.wpmucdn.com/sites.dartmouth.edu/dist/5/2293/files/2021/03/post-treatment-bias.pdf}{作者公开稿，正文前五页}

本报告的候选时序图将采用前条件、采用、传导环节和后续结果分开。它的用途是把假设写出来：哪些因素被认为发生在采用前，哪些变化可能属于中介，哪些反馈只能在下一时期发生。Textor的DAGitty手册讨论了总效应、直接效应与调整路径之间的关系；本文使用这种区分组织问题，尚未以该软件为所有技术案例求出充分调整集。\href{https://www.dagitty.net/manual-3.x.pdf}{手册第8—10页}

例如，研究AI采用对劳动收入的总影响，需记录既有管理与需求条件，同时把任务重组、销量扩张和工资调整作为候选传导环节。若要解释给定销量下的作用，则需另行定义估计对象并说明识别条件。事后按销量分组与外部设定销量不同，不能仅靠变量进入回归就认为问题已经解决。

只看留下来的人，也会漏掉结果。研究采用后仍在企业的劳动者，可能遗漏离职和裁员；研究仍在运营的项目，可能遗漏融资失败和亏损退出。因此，正文要注明“留任者”或“存续项目”，涉及全体受影响者时，再说明还缺哪些结果。

相互影响也要有先后。工资变化影响消费，消费影响下一期订单，订单再影响投资和雇佣。机制图可以提出这个循环；若把各项都放在同一期，又声称每条箭头都已识别，就无法交代谁先改变。第十四章核对单期需求反馈中的收支，同时区分这类核算与跨期因果研究。模型中独立比较的反事实，不能当作连续年度预测。

\section{企业、行业和地区的结果，须按各自范围比较}

一家企业多卖了，不等于全行业多卖了。它可能得到新增需求，也可能抢走同行订单。若是后者，把这家企业的销量增量乘以全部企业数，就把转来的订单重复计算。第十四章分别计算单家采用和共同采用的均衡结果，以说明这种差别。

影响还可能越出企业。制造和维护设备，会增加其他企业的劳动需求；失业者可能转入别的行业；产品降价可能改变消费者的支出。任务实验看局部流程，企业资料看组织调整，地区研究纳入更多人员流动和替代，结论可能因范围不同而不同。先把单位、结果和时期对齐；范围接近而方向冲突，才更可能是需要解释的经验争议。

不能把客服速度、地区就业人口比、发电成本和临床终点的百分数合成一个平均技术效应，它们没有共同分母。跨案例要比较具体机制：固定投入在什么条件下限制进入，接口有没有替代者，降本是否传给用户，转换费用由谁承担。先保留各领域的有效测量，再讨论共同关系。

上游开发与下游使用还具有不同观察范围。OECD将法国和葡萄牙2021／2023年企业ICT调查连接2011—2022年企业资料，未发现非生成式AI使用与市场权力代理指标系统性上升相一致的普遍关联。这一证据削弱“所有使用AI的企业都会增强市场权力”的无条件说法，同时留下上游开发和设施控制的不同问题。\href{https://www.oecd.org/content/dam/oecd/en/publications/reports/2026/07/competition-in-the-age-of-ai_e9f49e20/6f88a1ea-en.pdf}{原报告第9、52页}，\href{https://doi.org/10.1787/6f88a1ea-en}{DOI}

比较结果应当改变解释。上游设施即使高度集中，下游若出现更多替代服务、价格下降、迁移更容易，就须分别说明哪些控制仍在、哪些已经减弱。相反证据限制了原有判断的适用范围，不能只用一句“存在例外”把它搁下。

\section{先统一口径，再说明规模和利用怎样变化}

大量数值的罗列本身不能解释机制。以第九章重新整理的新型储能数据为例，国家能源局等发布的同一份2026年报告同时提供装机规模和分组利用资料。本文依据图19与图24逐项核对原图，保留原单位，换算后计算增长率和比值。\href{https://www.nea.gov.cn/20260730/0a1e9823b7a14e8ba16e82c34b3bea2b/202607300a1e9823b7a14e8ba16e82c34b3bea2b_270e42e93bb1a647609c851527f42364ee.pdf}{官方报告}

2021年至2025年，报告所列年末新型储能功率由4GW增至135.93GW；2023年至2025年，全国平均利用小时由611增至1195。两组数值共同表明，观察期内扩张发生在装机规模和利用强度两个方面，因而“规模扩大但完全没有更多利用”的描述与这组资料不符。进一步判断资源配置是否有效，仍须取得服务价值、成本和收益。

分组资料又提出更具体的问题。2025年独立储能平均利用1246小时，新能源配建储能为917小时，两者相差329小时。差异值得解释，但它并未把项目位置、时长、投运月份、调度资格与其他条件固定。若直接将329小时全部归因于组织形式，就回到了式（3-3）的选择问题。这项描述分析确认了统计差异，并明确了进一步解释差异所需的项目特征与制度资料。

来源与计算分别保存在装机数据、利用数据和复算程序。这种安排使读者能够区分报告原值、单位换算和本文计算。图表不提供原资料没有的置信区间，也不从年末功率与全年利用均值反推未经口径核实的全国总放电量。

同样的规则适用于其他章节。对临床结果，应说明接受治疗、进入主要分析和满足终点的不同人数；对光伏，应把建设成本、折现成本和售价分别列出；对劳动，应同时解释生产率的质量口径和工时范围。明确指标定义与数据口径，有助于核查机制判断的经验依据。


\subsection{行业总量变化，还要查哪些项目进入和退出}

行业扩大，要查哪些项目在变化。原有项目多运行了，新项目进来了，低效项目退出了，都可能改变总量和平均数。年末装机与全年运行的时间不同，组织类型的均值也混合了位置、设备年龄和需求。现有汇总资料能说明规模及使用变化；解释原因还需项目标识、投运退出日期和运行条件。企业名称不能代替项目标识，一家公司可以经营多座设施，资产转移也可能只改变名称。

本文的整理保留“年份—统计范围—指标—原单位—转换规则—原文定位”，将原值与派生值分别保存；缺失年份不以插值冒充观测。储能装机与利用统一采用同一2026年报告中的历史系列，减少跨版混用。若以后取得项目资料，可在这套记录上补入组织方式、调度资格、时段电量、停机、价格与费用，逐步连接设施使用和收益实现。

同一项目怎样变化，地区总体服务怎样变化，要分别报告。持续运行的项目便于前后比较，却漏掉新进入和退出者；全部项目反映总体，却含有项目构成的变化。把两者并列，才能查明改善来自原有设施用得更好、规模扩大，还是样本换了，也才能确定制度解释须处理哪些选择差异。

\subsection{自己采用的效果，要与同行响应分开}

式（3-2）至式（3-4）将同行状态固定，或采用不受其他单位处理影响的设定。降价、订单转移和公共设施使用可能违背这一条件。企业自身采用效应与同行响应必须分别定义，才能解释局部效果怎样联系到市场结果。

把同行采用状态的一个明确定义的汇总记为\(z\)，则企业结果可写成\(Y_i(d,z)\)。例如\(z\)可以是在事先确定的相关市场中，竞争者已经采用的产能份额。这里仅示范如何定义研究对象，没有假设这一份额足以概括现实全部网络作用。在给定同行状态\(z\)时，企业自身采用的差别是

\[
Y_i(1,z)-Y_i(0,z).
\tag{3-5}
\]

若同时将同行状态从\(z_0\)改到\(z_1\)，整体变化可以写成

\[
Y_i(1,z_1)-Y_i(0,z_0)
=\left[Y_i(1,z_0)-Y_i(0,z_0)\right]
+\left[Y_i(1,z_1)-Y_i(1,z_0)\right].
\tag{3-6}
\]

第一项改变自身采用，第二项改变同行状态。该分解是潜在结果的代数恒等式，不提供识别方法；自身与同行状态存在交互时，换一种分解顺序，两项数值也会改变。因此，它区分研究对象，却不能给出唯一的“技术贡献”。

第十四章的两企业模型使这个问题可以直接复算。没有企业采用时，每家利润约177.778；一家先采用并支付固定费后，其利润约411.111；对方随后也采用时，两家各约100。同一企业从没有采用到抢先采用是利润增加，从原状态到双方都采用却是利润减少。若只收集成功先行者，便会漏掉其后竞争变化；若把原本未采用的对手当作完全不受影响的对照，对照组利润变化又会把订单转移混入估计。

有相互影响，仍可以作比较，但要先选对对象。研究任务辅助，可以限定任务环境；研究产品市场，要同时查价格、销量和同行；研究电网服务，要按物理连接与调度范围确定谁可能受影响。距离近未必有经济联系，行政边界也未必隔得开供应链。范围须靠行业和制度资料确定，不能只靠回归设定。

\subsection{价格、产量和所得，要放在同一范围内核对}

假定未来资料允许估计采用使单位成本下降，仍需问这种变化经过了什么渠道。若成交价下降而销量扩大，收入结果取决于两者的乘积；若平均时薪提高而总工时减少，劳动总收入也需重新核算。仅报告一个系数的正负，不能交代共同形成的收益如何在产品、工资、利润和接入收费之间流动。

最基本的核对是让相互关联的量使用相同期间和对象。企业收入等于同范围成交价格与数量的加总，异质产品不能直接把简单均价乘总数量。劳动收入的核对同样需要工时权重：一人工作1小时、时薪10元，另一人工作3小时、时薪30元，总工资100元、总工时4小时，按工时加权平均为25元。用两个人时薪的简单平均20元相乘，只得到80元。这个算例说明，经济含义取决于统计定义，单位看似相同也可能无法直接相乘。

成本、价格和工资同时变化，说明多个环节发生了变化，不能独自确认传导原因；需求冲击也可能同时影响它们。总效应、给定条件的效应和中介效应须事先定义。尤其不能把处理后的销量直接作为控制，就声称已识别成本渠道。识别条件不足时，账户核对只能检查数值一致性，不能完成因果归因。

平均结果还要分人看。新进入和原有企业、留任者和离岗者、不同收入消费者，所得变化可能不同。报告分组结果，须带上各组范围、人数或权重以及缺失情况；小组的大幅增长不能说成多数人的改善。没有所需微观资料，就把群体效果保留为待检验的问题。

本文的描述整理、账户核对和条件模型分别承担测量、核算与推理检验。项目、合同与人员资料未贯通的部分保留为机制假说，不能因可计算或可解释而升级为经验结论。

\section{相反结果限制主张，模型核验检查推理}

文献选择围绕机制而非结论方向展开。从主张的原始出处追踪作者版本、方法和更新，再比较相近问题中的不同样本与相反结果。本文属于定向的机制综述，没有穷尽检索的完整记录，不据此报告系统综述式覆盖率。来源的作用由其生成方式决定：企业材料可确认自身披露和运行安排，实验可估计限定处理，行政记录可描述规定总体，合同则说明权利内容。

不同结果应改变解释，而非仅增加一句例外。程序生成的省时与真实代码库中的减速，要求把核验和整合计入完整任务；制造技术扩散后贸易增加，限制了技术必然带来普遍就地生产的判断；AI使用未普遍伴随市场权力代理指标提高，则要求把下游采用与上游控制分开。（\hyperref[ch:5]{第五章}；\hyperref[ch:8]{第八章}；第3.7节）本文保留这些结果，因为它们决定中心主张能够覆盖的条件。

版本变了，判断也要重新核对。正式稿可能更新样本和估计，运行统计可能修订历史值，公开数学稿件也可能撤回。精确说法要对应实际读过的版本与位置：登记了链接，不等于查过全文；公开稿件有多少，不等于独立验证了多少。第七章据此使用2026年数学发布材料，使事件在论证中承担与证据相称的作用。

模型审查分别检查数学解、经济核算和现实适用。竞争模型将解析解与独立最优反应比较，谈判模型与数值优化比较，社会收益和现金关系从不同账户核对；融资例子将年金闭式与逐年折现比较。这些检查能够发现推导和遗漏，却不能证明真实市场的需求曲线、谈判权重或资产安排已符合假设。第十四章将数学核验与经验检验明确分开。

改变参数，要找出结论在哪里反转。固定费能改变采用激励，销量扩张能改变总工时，外部机会和谈判条件能改变劳动所得，支出与资金条件能改变交付。阈值告诉研究者下一步该测什么；参数网格中某种结果占的面积，却不是它在现实中发生的概率。资料表明条件不成立，就应修改解释，不能换一组参数来保住原方向。

\section{现有证据支持局部结果与条件解释}

当前证据的强度并不均匀。任务研究、运行统计和部分因果研究，对技术能力及局部收益提供了较直接的支持；合同、许可和组织资料，可以确认特定主体拥有的权限；关于技术采用是否使这些权限普遍增强，材料仍不足。全文结论按这三个层次表述，不把第三层的缺口用第一层的数据填满。

因此，本文能够说明技术改善与资本支配可以并存，并给出收益分化的具体过程。它尚未证明全部新技术加强控制，也没有证明公共所有或合作所有在全部部门取得更优结果。未检出工资提高、设施不足或供应者较少，都须解释原因，不能直接等同于剥夺、垄断或占有关系的变化。

第四章据此提出条件解释。第五至十三章分别检验它们在劳动、组织、共同设施和再生产中的适用范围；第十四章检查经济推理；第十五章讨论与已识别机制相对应的制度措施。缺少直接证据的环节保留为研究问题，不写成已完成的结论。

\section*{本章参考文献}

Montgomery, J. M., Nyhan, B., \& Torres, M. (2018). How conditioning on posttreatment variables can ruin your experiment and what to do about it. \emph{American Journal of Political Science, 62}(3), 760–775. \url{https://doi.org/10.1111/ajps.12357} 。本轮核查出版社摘要与作者公开稿正文前五页，用于处理后调整偏差的说明；未重做原文模拟或实证分析。

Callaway, B., \& Sant’Anna, P. H. C. (2021). Difference-in-Differences with multiple time periods. \emph{Journal of Econometrics, 225}(2), 200–230. \url{https://doi.org/10.1016/j.jeconom.2020.12.001} 。核查作者2020-12-01公开稿第1—7页及定义；未在本报告实施估计。

OECD. (2026). \emph{Competition in the age of AI: Initial evidence from microdata}. OECD Artificial Intelligence Papers, No. 64. \url{https://doi.org/10.1787/6f88a1ea-en} 。核查公开PDF第9、52页的范围与讨论。

Textor, J. (2023). \emph{Drawing and Analyzing Causal DAGs with DAGitty}, version 3.1 manual, July 18. \url{https://www.dagitty.net/manual-3.x.pdf} 。核查第8—9页；本报告未运行DAGitty效应识别算法。


国家能源局能源节约和科技装备司、国家发展和改革委员会创新和高技术发展司、电力规划设计总院（2026）。《中国新型储能发展报告2026》。\url{https://www.nea.gov.cn/20260730/0a1e9823b7a14e8ba16e82c34b3bea2b/202607300a1e9823b7a14e8ba16e82c34b3bea2b_270e42e93bb1a647609c851527f42364ee.pdf} 。图19、图24；本章计算为描述统计。

\chapter{作用机制：技术节约经过市场和生产安排形成收益}\label{ch:4}

\section{技术节约经过交付、市场与分配形成收益}

技术节约通过四个环节形成社会结果：完整流程实现节约，市场将节约转为价格与销量，劳动和交易安排分配收入，持续投入维持生产。投资、用工、准入和剩余用途的权限贯穿这些环节，既影响节约能否实现，也影响由谁取得。

以下命题说明同一技术为何可能产生不同结果。它们要求同时观察相关投入、行为与结果，而非仅按技术名称归因。例如扩张可能来自新增需求，也可能来自转移同行订单；平台增收可能来自匹配改善，也可能来自收费调整。解释需要在相同范围内区分这些变化。

下文将解释集中到完整生产条件、竞争与劳动分配、持续生产投入及制度反馈四组关系。H1、H2、H5讨论条件与准入，H3、H4讨论劳动与技能，H6、H7讨论费用与需求反馈；H8是区分收益分享和占有关系的概念判据，H9补充技术选择与制度回应之间的反馈。现有材料支持其中若干环节，各命题不能累计成整条链已经得到经验识别的次数。

\begin{center}\resizebox{\textwidth}{!}{\begin{tikzpicture}[>=Stealth, every node/.style={font=\small}]
\node[draw=gray!65,fill=blue!7,rounded corners,align=center,inner sep=6pt,text width=3.1cm] (TC) at (2.24,5.94){技术能力\\任务效率／信息／能量与材料\\版本、质量与部署条件};
\node[draw=gray!65,fill=blue!7,rounded corners,align=center,inner sep=6pt,text width=3.1cm] (SD) at (7.36,8.14){共同生产依赖\\协作、标准与公共设施};
\node[draw=gray!65,fill=blue!7,rounded corners,align=center,inner sep=6pt,text width=3.1cm] (ET) at (7.36,5.94){效率与任务结构\\核验、替代、增强与学习};
\node[draw=gray!65,fill=blue!7,rounded corners,align=center,inner sep=6pt,text width=3.1cm] (XC) at (7.36,3.74){可排他控制\\资产、接口与准入规则};
\node[draw=gray!65,fill=blue!7,rounded corners,align=center,inner sep=6pt,text width=3.1cm] (ED) at (12.96,8.47){经济分配\\工资／盈余／价格};
\node[draw=gray!65,fill=blue!7,rounded corners,align=center,inner sep=6pt,text width=3.1cm] (LO) at (12.96,6.489999999999999){劳动与入口\\工时／强度／技能／机会};
\node[draw=gray!65,fill=blue!7,rounded corners,align=center,inner sep=6pt,text width=3.1cm] (DE) at (12.96,4.51){决策与退出\\参与／迁移／申诉};
\node[draw=gray!65,fill=blue!7,rounded corners,align=center,inner sep=6pt,text width=3.1cm] (CR) at (12.96,2.5300000000000002){成本与风险\\投资／转换／公共责任};
\node[draw=gray!65,fill=blue!7,rounded corners,align=center,inner sep=6pt,text width=3.1cm] (DF) at (18.08,8.47){需求与产出实现\\消费／投资／公共购买};
\node[draw=gray!65,fill=gray!10,rounded corners,align=center,inner sep=6pt,text width=18.7cm] (INST) at (10.4,0.2){制度条件：产权与治理、竞争与替代、议价与分享、开放与公共设施\\H8：区分收益分享、局部共同决定与占有关系转化};
\draw[->,dashed,gray!70!black,line width=0.6pt] (TC.north east) -- node[midway,fill=white,inner sep=2pt,font=\scriptsize]{H1} (SD.west);
\draw[->,dashed,gray!70!black,line width=0.6pt] (TC.east) -- node[midway,fill=white,inner sep=2pt,font=\scriptsize]{H1,H4} (ET.west);
\draw[->,dashed,gray!70!black,line width=0.6pt] (TC.south east) -- node[midway,fill=white,inner sep=2pt,font=\scriptsize]{H2,H5} (XC.west);
\draw[->,dashed,gray!70!black,line width=0.6pt] (ET.east) -- node[midway,fill=white,inner sep=2pt,font=\scriptsize]{H3} (ED.west);
\draw[->,dashed,gray!70!black,line width=0.6pt] (ED.east) -- node[midway,fill=white,inner sep=2pt,font=\scriptsize]{H7} (DF.west);
\draw[->,dashed,gray!70!black,line width=0.6pt] (ET.east) -- node[midway,fill=white,inner sep=2pt,font=\scriptsize]{H3,H4} (LO.west);
\draw[->,dashed,gray!70!black,line width=0.6pt] (XC.east) -- node[midway,fill=white,inner sep=2pt,font=\scriptsize]{H2} (DE.west);
\draw[->,dashed,gray!70!black,line width=0.6pt] (SD.north) -- (7.36,9.8) -- node[midway,fill=white,inner sep=2pt,font=\scriptsize]{H6} (15.6,9.8) -- (15.6,2.53) -- (CR.east);
\end{tikzpicture}}\end{center}


图4-1表示H1至H7的候选通道，箭头不代表已估计的因果效应；H7单列经济分配通向需求与产出实现的通道，区别于效率影响分配；H8列在制度判据中，不作为因果箭头。本章保留机制与关键结果，第十四章集中给出成本、竞争、谈判、支出和制度选择的完整推导。H9的制度反馈另在后文展开，未包含在本图H1—H7的候选箭头中。各模型均为未校准的条件分析，不以构造参数代替行业资料。

\section{互补投入与开放许可共同限制交付}

局部技术能力须与互补工序结合，才能形成合格交付。生成、加工或发电加快时，核验、后处理或接入容量可能成为限制。评价完整流程节约，需要同时记录质量、辅助劳动和交付数量。

\textbf{H1：在工艺、质量与需求给定时，当前瓶颈限制完整交付；仅提高非瓶颈环节能力，产出不会同比增加。}发电接网、细胞制备与诊疗容量提供不同对象。若补足被判定的瓶颈仍未增加交付，应检查是否另有约束，或原瓶颈判断有误；工艺替代和需求变化也可能改变必要条件。

瓶颈首先是技术与组织约束，其来源可能是建设周期、物理资源或排他准入。只有明确经营者的规则权限、使用者的替代条件及实际后果，才能进一步分析入口支配。设施不足本身不能证明私人控制。

配套投入还决定谁用得起技术。企业先付调试、培训、认证和改造费，再靠销售收回。同样节省单位成本，订单稳定、容易融资的企业更可能做成投资；订单不足、借不到钱的企业，则可能仍用不上。所以，资源原本怎样分布，会影响技术怎样扩散；采用者表现更好，可能既有技术作用，也有原有优势。

\textbf{H5：开放许可降低其明确解除的使用限制；许可与互补投入共同决定独立部署和持续交付的范围。}文件公开不必然允许修改、再分发或商用。权限改善后，计算、技能、质量验证与销售条件仍可能限制使用；这些条件同时改善，则可能扩大独立供给。检验应以许可内容与实际使用为准，不把下载量当作自主生产能力。

共同设施通过分摊投入或减少重复准备降低门槛，容量分配和维护规则再决定机会如何落实。免费接入、收费服务与成员共同维护应按持续使用、完整费用和治理权限比较，不能仅按收费是否为零判断开放程度。

检验H1和H5，先在采用前查明需要哪些配套条件，再比较补足投入以后，完整任务产出和持续使用怎样变化。已确认的设施缺口若并未限制产出，开放以后小组织若已能自主部署，就应缩小瓶颈解释的范围。管理改善和需求扩张也可能解释结果，须另行比较。

\section{谁掌握难以替代的入口，谁就可能改变交易条件}

\textbf{H2：在关键投入难以替代、迁移费用高且使用者缺乏有效约束权时，提供者具有改变交易条件的空间；替代供给和可执行的迁移权会压缩这一空间。}这是入口权力的条件解释，不是所有资本支配的来源。是否形成更高收费或其他收益，还取决于服务成本、需求、规则约束和提供者的实际行为。

第十四章将企业单位成本分为劳动、其他真实投入和接入费三项。在其构造例子中，生产率提高使每件劳动费用减少20；若接入费也提高20，下游企业的总单位成本便保持原值。这个比较显示一种可能的收益转移：技术确实节省了劳动，下游却没有取得降价所需的成本空间。它没有证明供应者实际上能提高收费20，更没有估计这类转移在现实中的普遍程度。

入口权力的检验需区分成本补偿与规则优势。服务量、维护或质量要求提高，可能使收费增加；同等服务下用户难以替换、提供者能够单方改约，则更有根据分析权力作用。要判断超出竞争条件的收益，还需相应竞争基准，不能把收费或利润本身称为垄断租金。

控制也可以不涨价。平台能排订单、停账户，采购者能改验收条件，设施经营者能排接入顺序。这些决定会改变参与者取得多少收入、做多少工作、付多少转换费用。劳动者或供应者名义上独立，工作仍可能受到外部规则支配；供应者虽少，规则也未必不对称。实际权限要从这些事项查明。

OECD对法国和葡萄牙2011—2022年非生成式AI资料的研究，未发现使用AI与市场权力代理指标系统提高的平均关联；开发活动另行分析。这一关联证据要求把下游采用与上游控制分开，限制了“技术采用普遍增强市场权力”的概括。（OECD，2026，\href{https://doi.org/10.1787/6f88a1ea-en}{正式来源}）因此，本文关于资本控制的判断，落在可核查的费用、准入和规则权限上，不以采用规模或集中度替代。

\section{企业各自争先采用，行业利润仍可能下降}

资本控制生产条件，并不保证每家企业都能长期多赚。企业采用技术，一方面为取得额外收益，另一方面也为保住经营地位。先降成本的企业可以多接订单，同行跟进以后，先行优势会缩小，降价和重复投入又会改变行业净收益。

第十四章的两企业模型把这种关系明确化。给定线性需求、固定工资、同质产品和每家300的当期采用支出，两家均未采用时各获净利润177.778；一家先采用时，其净利润为411.111，对手降至44.444；两家均采用后，各获100。对手无论是否采用，自己采用都更有利，因而采用成为各自的严格占优选择；最终利润却低于双方最初都未采用的状态。以上均为同一期、同一货币单位下的构造情景。

该例区分个别采用激励与共同采用结果。每家按自己的竞争位置选择，双方反应却消除了部分成本优势，并支付采用费用。因而，实际采用不证明高回报。模型中降价又增加消费者剩余，企业利润下降同样不足以证明社会净损失。

结果由固定支出限定。在该算例中，支出处于约222.222至355.556之间时，才同时出现“采用具有严格占优性”与“双方采用后的净利润低于原状”。较低支出可以使双方都获益；足够高的支出又会抑制采用。第十四章提供阈值推导、边界与社会账户核算。这些条件限制了模型结论，也指明经验研究应调查采用费用和竞争者反应。

这一机制把资本间竞争接到基本矛盾：相互依赖的生产由各资本分别按自身回报决定，个别选择与整体结果因而可能背离。模型只证明这种背离的可能性；现实判断还需需求、进入退出、资产寿命和投资顺序。多期设备应比较多期收益，不能由当期费用推断永久损失。

\section{劳动者能否增收缩时，要看订单、用工和谈判}

效率首先改变单位产品所需工时，市场需求决定生产多少，工资和工作规则决定劳动者如何分享结果。总工时满足$H=Q/a$：产量增长慢于生产率，总工时减少；二者同比增长，总工时不变；产量增长更快，总工时增加。总工时又可分为岗位人数和每人时长。由此，技术省时可以表现为缩短工作日、减少岗位、扩大生产或多种变化的组合。

\textbf{H3：在效率变化给定时，销量决定所需总工时，报酬与排班决定劳动者分享收入还是闲暇；外部机会和协商权限影响这种分享。}工时关系提供数量约束，工资和岗位分配仍需行为与制度解释。销量、平均时薪、总工时及人员流动应共同观测。

第十四章另设谈判模型，将供双方分配的净收入记为$V$，劳动方协议破裂后的所得为$bH$，企业对应所得为$\bar\pi$，劳动谈判权重为$\eta$。其解为$W=bH+\eta(V-bH-\bar\pi)$。在工时和外部条件固定时，净收入增加，劳动所得按权重分享增量；若技术同时提高雇主的替代能力，或削弱劳动方的议价条件，劳动所得也可能不增反降。纳什原论文提供破裂基点与增益乘积解的理论来源；加权货币形式和本报告的情景数值是本文设定。（Nash，1950，原刊pp.157—159；\href{https://doi.org/10.2307/1907266}{DOI}，\href{https://www.haverford.edu/sites/default/files/Nash1950.pdf}{公开全文}）

劳动者能否拒绝不利条件，关系到能分到多少。能换到其他岗位，失业期间有生活保障，能共同协商工资和强度，都会改变拒绝的能力。技能和记录若只能留在原企业，离职又立即危及生活，法律上能够选择，现实中却可能难以选择。这些条件通过外部机会与协商影响分配，具体效果仍需合同、劳动记录和制度变化来识别。

工资提高，要承认分配改善；投资、裁员和剩余用途若仍由资本决定，就不能同时宣称劳动者已取得这些控制权。协议若让劳动者能共同决定缩时、安全或培训，则又改变了部分劳动控制关系。因此，后文分别记录收入与权限。竞争模型固定工资，谈判模型固定收入和工时，两者还未联立求解，不能把各自结果直接拼成完整均衡。

\subsection{劳动节约经扩产、岗位和周转形成资本后果}

H3还需要明确资本一侧的结果。给定产量和报酬，完整付酬工时减少才会降低工资支出；省时但保薪保岗，首先增加的是可用时间或可完成任务。扩产和新增岗位能够吸收这种能力，也能够扩大可变资本实际使用额。若工作与工资不变而验收、回款加快，变化则主要发生在所需预付资本。

因此，技术影响可变资本不能归结为“机器代人”。任务替代解释每单位用工，市场实现解释总量，岗位与协商解释购买价格和结构，周转解释资本需要垫付多久。生活资料生产条件再把这些企业变化接到劳动力价值与社会再生产。四个层次的连接，才构成后续各行业比较的共同依据。

劳动使用还构成这一机制的另一面。企业可以在少招操作人员的同时提高留任者的岗位覆盖，在采购AI服务的同时增加核验和交付定额，或在实现远程运行后扩大每名支持人员负责的范围。工资支出未下降时，这些变化仍可能改变劳动密度和责任；工资支出下降时，也不能独自说明劳动强度上升。将岗位购买与实际使用分开识别，才能讨论第二章所述技术替代同劳动强化并存的条件。

\section{工具提高当前表现，长期能力仍须练习和反馈}

\textbf{H4：工具提高当期表现，并不足以保证长期独立能力；练习、反馈和职责进阶影响技能积累。}两种能力需用不同结果衡量。若撤去工具后能力持续提高，且训练和岗位入口扩大，应修正工具必然削弱技能的判断。

新人先作判断，再用工具核对、学习错误原因；或者直接提交工具输出，原有练习任务也被取消。这两种使用方法都可能加快当期交付，对以后技能的作用却可能不同。因此，要查撤去工具后的表现、培训时间和职责进阶。即时任务实验没有观察这些后果，不能据此认定整个职业入口已经收缩。

谁培养人，谁用培养出来的人，也可能不是同一家企业。提供培训的企业未必留住全部熟练者，其他企业可以招聘他们。如果各家都少培养，未来技能供给就可能不足。公共教育、共同培训和内部晋升可以改变这些选择。判断这种跨期机制是否发生，还须连续记录培训、人员流动和能力。

技能既用于当前生产，也决定劳动者以后能否换雇主、能否自主工作。工具扩大高质量学习机会，可能增强这些能力；知识若只保存在不能带走的系统里，劳动者只按规定步骤执行，则可能更依赖组织。第五、六章按研究实际覆盖的时期，分清已经观察到的任务效果与尚待检验的职业后果。

\section{完整投入与承担变化区分节约和转移}

\textbf{H6：在产出、质量和观察期可比时，企业降本若伴随他人新增必要投入而未获补偿，就可能包含费用转移；完整投入减少则构成资源节约。}企业停止培训、劳动者增加待命或公共设施增加维护，只是候选机制。证明转移还需对应的投入、责任与补偿记录；公共支出增长本身不构成证据。

资源节约与费用转移可以并存。设备减少危险作业，家庭却增加培训投入，需要分别记录；扩产导致维护增加，也不能直接归为外移。成本比较应保持服务范围一致，或明确处理规模变化。环境、健康和时间缺乏可靠估值时，分别报告，不强行合成为一个净收益数。

付款和资源消耗，还要分账。下游付接入费，上游同时取得收入；全社会核算应继续查服务用了多少劳动、能源和设施，不能重复算同一付款或同一资源。分清以后，再查共同条件由谁出资，收费是否覆盖成本，剩余归谁，以及承担负担者有没有相应决定权。

\textbf{H7：收入分配通过消费、投资和公共支出影响技术产出的实现；在其他需求来源未作相应补偿时，收入转向消费倾向较低的主体会削弱消费需求。}第十四章用局部支出式说明这一条件关系，并检查商品、收入和现金账户。它不把劳动收入减少直接等同于总需求不足，因为新增利润可以投资，价格下降也可以扩大购买数量。

产能扩大，只说明企业能多生产；产品卖出去，还要有适用需要、相应支付和可接受的交易条件。企业追求减支与扩产时，工资分配又参与形成购买力，生产决定因而可能反馈到销售条件。消费、投资、公共购买和外部需求须共同核算；特定产品的销售不足不能直接外推为总需求不足。第二章据此区分价值生产与市场实现，第十四章检查局部商品和现金账户。宏观危机还需跨期投资、信用和部门联系的证据。

\section{改善所得与改变生产支配，须分别检验}

制度比较分别考察分配、准入与生产决定权。竞争和开放可能扩大使用机会，保障与公共服务可能改善所得和负担，所有权与治理可能改变实际决策。三类效果可以相互支持，但需各自的结果与执行证据。

\textbf{H8：分配改革改变劳动者获得多少收入与时间；判断占有关系变化，还须查谁控制生产条件、谁决定投资和劳动组织、谁处置剩余。}这是制度比较的概念判据。奖金增加、收费降低、公共服务扩大，都有实际作用；资产和决策权限是否改变，须继续核查。

劳动者取得安全、工时或培训的共同决定权，说明特定生产关系已调整。若进一步判断资本主义占有关系整体转化，就还要查生产者与生产条件的分离怎样改变，生产目标和社会剩余是否由共同需要支配。企业内部的一项参与制度，不能单独证明生产已从资本增殖转向这种目的。

公共所有或合作所有改变了权利基础，实际运行仍要逐项查。成员能否决定技术投资、取得信息、监督管理、修改剩余用途，才能说明名义权利有多少实际内容；融资、客户和设施依赖，又会限制内部选择。因此，要把资产归属、决定程序能否执行和生产能否持续，放在一起比较。

以上机制将技术收益分解到完整交付、市场实现、劳动分享和持续投入，并说明具体决定权怎样介入这些过程。第五至十三章分别提供环节证据，第十四章检验经济与制度选择的条件推论，第十五章据此比较制度措施。各章支持的范围仍由资料与假设限定。

\subsection{制度回应反过来改变自动化与任务创造}

\textbf{H9：当劳动参与约束可执行时，自动化引起的补偿责任会影响资本方部署选择；绕过约束是否更有利，取决于补偿效率、政治成本、两种方案的失败风险、风险偏好与新增任务。} 本命题把制度回应作为待解释结果，并允许它反馈到技术选择。原H1—H7说明经济和组织通道，H8区分收益与占有关系，H9则追问生产决定权为什么被分享、维持或限制。

第十四章的参与约束模型先比较各制度下的最优方案，再计算转换阈值，避免只看一种技术的工资结果就判断政治走向。若公共品效率提高、有效补偿、新任务或较高制度成本改变收益排序，就不能把自动化与镇压的互补性推广为必然关系；若实施机构不按资本方收益作决定，也不能继续套用本节目标函数。工资低、监控密集或存在冲突，均不足以独立证明整条政治反馈链。

\section{机制命题的证据与反驳条件}

后续案例按照上述关系分工，既考察技术的实际改善，也核查控制机制成立的条件。

\begingroup\footnotesize\setlength{\tabcolsep}{4pt}\renewcommand{\arraystretch}{1.25}
\begin{longtable}{>{\raggedright\arraybackslash}p{\dimexpr 0.15254\textwidth-8.0pt\relax}>{\raggedright\arraybackslash}p{\dimexpr 0.37288\textwidth-8.0pt\relax}>{\raggedright\arraybackslash}p{\dimexpr 0.47458\textwidth-8.0pt\relax}}
\toprule
\textbf{机制及命题} & \textbf{主要证据} & \textbf{削弱或限定解释的事实} \\
\midrule\endfirsthead
\toprule
\textbf{机制及命题} & \textbf{主要证据} & \textbf{削弱或限定解释的事实} \\
\midrule\endhead
\bottomrule\endfoot
共同条件与开放，H1、H5 & 完整任务产出、互补投入、许可内容和持续使用 & 已确认的设施缺口未限制交付；小组织已能独立部署和维护 \\
关键入口，H2 & 实际收费、服务成本、替代供应、迁移和规则变更 & 收费对应真实新增服务；替代便利；使用者能够共同改变规则 \\
劳动分配与技能，H3、H4 & 工资、总工时、人数、强度、协商记录和后续独立能力 & 需求与人员筛选解释收入变化；工具持续扩大训练与职业入口 \\
跨界费用，H6 & 企业、家庭与公共部门之间对应的投入和补偿 & 完整流程的资源消耗确实减少；其他主体的新增承担得到补偿 \\
需求反馈，H7 & 收入、价格、消费、投资和交付的对应变化 & 投资、降价、公共购买或外需补偿消费变化 \\
占有关系，H8 & 生产条件归属、实际决定程序和剩余用途 & 只发生收入或价格变化；形式参与未形成可执行权限 \\
制度反馈，H9 & 部署方案、补偿效率与执行、政治成本、两分支风险、资本方目标及新增任务 & 效率改善或新任务反转排序；成本和风险设定改变；实际目标不符合模型 \\
\end{longtable}\endgroup


\begin{center}\resizebox{\textwidth}{!}{\begin{tikzpicture}[>=Stealth, every node/.style={font=\small}]
\node[draw=gray!65,fill=blue!7,rounded corners,align=center,inner sep=6pt,text width=2.7cm] (B0) at (2.08,7.5){B\ensuremath{_0}  采用前条件\\资金／技能／需求};
\node[draw=gray!65,fill=blue!7,rounded corners,align=center,inner sep=6pt,text width=2.7cm] (G0) at (2.08,3.0){G\ensuremath{_0}  产权与治理\\采用前的规则结构};
\node[draw=gray!65,fill=blue!7,rounded corners,align=center,inner sep=6pt,text width=2.7cm] (A1) at (6.08,5.300000000000001){A\ensuremath{_1}  技术采用\\明确版本与部署方案};
\node[draw=gray!65,fill=blue!7,rounded corners,align=center,inner sep=6pt,text width=2.7cm] (P2) at (9.76,7.5){P\ensuremath{_2}  实际效率\\质量约束下的表现};
\node[draw=gray!65,fill=blue!7,rounded corners,align=center,inner sep=6pt,text width=2.7cm] (T2) at (9.76,5.300000000000001){T\ensuremath{_2}  任务结构\\替代／增强／学习};
\node[draw=gray!65,fill=blue!7,rounded corners,align=center,inner sep=6pt,text width=2.7cm] (C2) at (9.76,3.0){C\ensuremath{_2}  实际控制\\接口／劳动规则};
\node[draw=gray!65,fill=blue!7,rounded corners,align=center,inner sep=6pt,text width=2.7cm] (Y3) at (13.76,5.300000000000001){Y\ensuremath{_3}  后续结果\\工资／工时／入口／风险\\实证时应拆成具体变量};
\node[draw=gray!65,fill=blue!7,rounded corners,align=center,inner sep=6pt,text width=2.7cm] (D4) at (13.76,1.3){D\ensuremath{_4}  更晚的需求\\下一阶段反馈};
\draw[->,dashed,gray!80] (B0.east) -- (A1.west);
\draw[->,dashed,gray!80] (B0.north) -- (2.08,9.1) -- (13.76,9.1) -- (Y3.north);
\draw[->,dashed,gray!80] (G0.east) -- (A1.west);
\draw[->,dashed,gray!80] (G0.east) -- (C2.west);
\draw[->,dashed,gray!80] (G0.south) -- (2.08,0.15) -- (16.2,0.15) -- (16.2,5.3) -- (Y3.east);
\draw[->,dashed,gray!80] (A1.east) -- (P2.west);
\draw[->,dashed,gray!80] (A1.east) -- (T2.west);
\draw[->,dashed,gray!80] (A1.east) -- (C2.west);
\draw[->,dashed,gray!80] (P2.east) -- (Y3.west);
\draw[->,dashed,gray!80] (T2.east) -- (Y3.west);
\draw[->,dashed,gray!80] (C2.east) -- (Y3.west);
\draw[->,dashed,gray!80] (Y3.south) -- (D4.north);
\end{tikzpicture}}\end{center}


图中$B_0$为采用前资金、技能和需求，$G_0$为既有产权与治理，$A_1$为明确版本的采用，$P_2,T_2,C_2$分别为效率、任务和实际控制变化，$Y_3$为后续结果，$D_4$为更晚的需求反馈。图形程序仅检查无环，未识别现实因果；收入、工时和权限仍须作为不同结果分别研究。

后续案例先查整个流程的产出能否带来劳动收益，再查运行、交付和规则怎样使分配不同。比较共同设施，要分清实际短缺与排他准入；研究社会再生产，要补查企业以外谁承担费用。模型检查这些关系能否从明列假设推出，制度比较再针对已查明的费用、选择和决定权。各章补足不同环节，行业并列比较不能写成技术之间依次发生的因果链。

\section{本章参考文献与图源}

Nash, J. F., Jr. (1950). The bargaining problem. \emph{Econometrica}, 18(2), 155–162. \href{https://doi.org/10.2307/1907266}{DOI:10.2307/1907266}。\href{https://www.haverford.edu/sites/default/files/Nash1950.pdf}{Haverford公开原论文}，原刊pp.157—159核查破裂基点、公理与增益乘积解；权重、货币简式和劳动分配算例为本报告设定。访问日期：2026年10月8日。

OECD. (2026). \emph{Competition in the age of AI: Initial evidence from microdata}. OECD Artificial Intelligence Papers, No. 64（正式网页所列系列号；本轮读取的公开PDF封面、封底标为No. 62，以同一DOI识别报告）. \href{https://doi.org/10.1787/6f88a1ea-en}{DOI:10.1787/6f88a1ea-en}。本章只保留使用者／开发者区分及关联并非单向的背景，本轮重新核对原报告第9、25、33—34页及正式网页元数据，不把它当成模型估计。访问日期：2026年10月8日。

图4-1、图4-2：本报告原创研究设计；图形源代码与节点／边数据见“图/绘制概念图.py”“图/概念节点与边.json”“图/候选识别图.dag”。本轮修正图4-1的H7通道，PNG、SVG、节点边表及TeX内嵌图同步；图形提出假设，未报告经验因果结果。

\chapter{自动化与劳动：劳动收益取决于订单和用工安排}\label{ch:5}

自动化对劳动的影响经过三次转换：局部省时转为完整流程节约，单位劳动节约经过需求变化形成总工时，总工时再通过排班与报酬形成个人结果。机器人与AI研究分别观察其中一些环节。本章据此解释生产率提高为何可能伴随扩产、缩时或就业调整，并区分谁获得收益、谁承担转换。

下文先把工时账算清，再看竞争、工资和职业入口。机器人研究比较任务、企业、部门与地区的不同结果，AI研究比较局部任务效果与早期劳动市场变化。原研究的样本和版本见\hyperref[ch:5]{劳动研究核查}。假设算例说明条件怎样起作用，实际效果仍以经验研究为据。

\section{总工时随产量和生产率变化}

设一个时期内合格产出为Q，总劳动工时为H，每工时产出为a。相同边界下，a=Q/H，因此：

\[
H=\frac{Q}{a},\qquad H=N h.
\tag{5.1}
\]

N是劳动者人数，h是每人平均工时。产出可以是吨钢、合格零件或完成并通过验收的服务，a相应为吨／工时、件／工时或服务／工时。企业间产品质量不同时，先统一合格标准，或分别列出产量和质量。人数与平均工时须覆盖同一群体；同时包含全职与兼职时，总工时应与岗位数量分别报告。

对正值取对数，工时的变化可精确写为：

\[
\Delta\ln H=\Delta\ln Q-\Delta\ln a,
\qquad
\Delta\ln N=\Delta\ln Q-\Delta\ln a-\Delta\ln h.
\tag{5.2}
\]

在连续变化中，同一关系是dlnH=dlnQ\ensuremath{-}dlna。对数变化便于拆分；若使用通常的百分比，有限变化应按H\ensuremath{_1}/H\ensuremath{_0}=(Q\ensuremath{_1}/Q\ensuremath{_0})/(a\ensuremath{_1}/a\ensuremath{_0})计算，不能把较大的增长率简单相减。工时变化的方向由以下条件决定：产出增长超过生产率增长，工时增加；产出增长不足，工时减少。设备技术水平本身不足以确定产出增长率。

\textbf{表5.1是教学假设，所有情境独立比较同一基线。}基线每期产出1000件，每工时10件，需要100工时。采用后每工时12.5件，提高25\%。下表不设价格或工资，也没有假定实际企业一定能够取得这些订单。

\begingroup\footnotesize\setlength{\tabcolsep}{4pt}\renewcommand{\arraystretch}{1.25}
\begin{longtable}{>{\raggedright\arraybackslash}p{\dimexpr 0.20455\textwidth-8.0pt\relax}>{\raggedright\arraybackslash}p{\dimexpr 0.15909\textwidth-8.0pt\relax}>{\raggedright\arraybackslash}p{\dimexpr 0.15909\textwidth-8.0pt\relax}>{\raggedright\arraybackslash}p{\dimexpr 0.15909\textwidth-8.0pt\relax}>{\raggedright\arraybackslash}p{\dimexpr 0.15909\textwidth-8.0pt\relax}>{\raggedright\arraybackslash}p{\dimexpr 0.15909\textwidth-8.0pt\relax}}
\toprule
\textbf{情境} & \textbf{每期Q（件）} & \textbf{a（件／工时）} & \textbf{H=Q/a（工时）} & \textbf{相对基线工时} & \textbf{若维持10人，每人h} \\
\midrule\endfirsthead
\toprule
\textbf{情境} & \textbf{每期Q（件）} & \textbf{a（件／工时）} & \textbf{H=Q/a（工时）} & \textbf{相对基线工时} & \textbf{若维持10人，每人h} \\
\midrule\endhead
\bottomrule\endfoot
基线 & 1000 & 10 & 100 & 0 & 10小时 \\
产出不变 & 1000 & 12.5 & 80 & \ensuremath{-}20\% & 8小时 \\
产出增加20\% & 1200 & 12.5 & 96 & \ensuremath{-}4\% & 9.6小时 \\
产出增加25\% & 1250 & 12.5 & 100 & 0 & 10小时 \\
产出增加50\% & 1500 & 12.5 & 120 & +20\% & 12小时 \\
\end{longtable}\endgroup


例如1200÷12.5=96，说明生产率提高25\%、产出增加20\%时，工时下降4\%，并非下降5\%。最后一列展示维持人数的排班结果。如果每人10小时固定，80工时可以由8人完成，120工时需要12人；96工时则可能需要兼职、加班或整数班组调整。9.6人的计算值是固定每人工时下的用工当量，实际人员配置仍受劳动合同与排班单位约束。

同样的20\%工时节约可以表现为保薪缩时、减薪缩时或减少人数。原来10人各工作10小时、时薪50元，每人收入500元；每人改为8小时且时薪不变，收入降至400元；维持500元则需时薪62.5元。保薪缩时没有减少单位产品5元的工资费用，却减少了完成同产出的劳动时间。资源节约与企业现金节省由此分开。

Acemoglu与Restrepo（2019）的任务理论进一步区分置换、生产率及新劳动任务的恢复效应。公式5.1可以容纳它们的最终结果，但没有单独识别这些机制。机器接手某项任务可能降低H，也可能使新增检查成为必要；价格下降可能提高Q；新任务又会改变产出的组成。任务理论进一步解释了这一总比率背后的配置变化，即劳动退出既有活动与进入新增活动的过程。\href{https://doi.org/10.1257/jep.33.2.3}{原论文}

\section{完整交付工时决定实际省时幅度}

完整流程工时包括准备、制作、核验、返工、交接和必要故障处理。机器运行时间、个人操作时间与团队劳动时间分别计量；后台运行不与并行劳动重复相加。以交付合格报告为产出时，草拟省时只有扣除新增工作后，才成为流程节约。

以下假设例子区分任务速度与完整流程生产率。原来制作100份合格报告需起草60小时、查证30小时和交接10小时，共100小时，a为1份／小时。工具使起草耗时减半，但新增复核12小时，其他工作保持不变，合计30+30+10+12=82小时，a为100/82\ensuremath{\approx}1.220份／小时。起草速度翻倍，流程生产率只提高约22\%。若复核新增35小时，合计105小时，流程生产率反而下降约4.8\%。这些算例用于说明工时核算边界，不提供任何产品错误率的估计。

合格标准限制可以计作产出的数量，错误损失又可能超出平均评分。生成量增加而核验时间不足，可能提高强度或转移错误责任。流程评价因此需要返工、投诉、损失及劳动投入资料，工具被称为“辅助”不能确定这些后果。

Noy与Zhang（2023）的预注册实验提供453名大学教育专业人员的短时写作证据，GPT-3.5访问权限使所选任务平均耗时约减少40\%、评分约提高18\%。它较好识别这些任务中的访问效果，而非整个职业一年中的工作结果。耗时与评分分别测量，不能将两个百分比相加为“58\%生产率收益”。\href{https://doi.org/10.1126/science.adh2586}{正式论文}

Dell’Acqua等的正式2026年论文使用2023年GPT-4，对758名咨询人员进行随机实验。边界内任务的数量与速度改善，越界管理任务的正确率却降低约19个百分点。能力对不同任务的适配呈不规则边界：某个环节的有效表现，仍需企业识别适用任务、安排核验并明确错误责任，才能转化为流程收益。\href{https://doi.org/10.1287/orsc.2025.21838}{正式论文}

实际采用涉及技术能力与组织责任安排的整合。职业“暴露”指数通常判断工具可能帮助哪些活动；实际采用还需要可靠性、培训、数据接入、权限与预算。愿意测试一次、能够持续使用、已重组流程，是不同阶段。若原工时中适合改进的流程份额为s，其中真正采用的比例为z，采用流程省时比例为r，新增组织工时为E，在产出固定且各部分互不重复时有：

\[
\frac{H_1}{H_0}=1-szr+\frac{E}{H_0}.
\tag{5.3}
\]

假设s=20\%、z=50\%、r=40\%，节省基线总工时4\%；若新增核验与集成占基线工时2\%，净工时节省只有2\%，相应每工时产出提高1/0.98\ensuremath{-}1\ensuremath{\approx}2.04\%。s是已经界定完整流程的工时份额，不能直接填入某个职业暴露指数的百分比。这一算例解释了局部实验收益在完整流程中的缩减，但没有确定采用后的销量变化。

组织调整也可能减少重复记录和查证，使完整节约大于单纯替换起草的结果。比较需同时计入取消与新增的工作。工具是否形成可携带的个人能力，则是另一个跨期结果，需要独立能力和后续任务资料，不能仅由当期生产率判断。

\section{降本后的用工变化取决于降价和销量}

生产率提高首先降低单位劳动需要，工资成本是否同步下降还取决于时薪。沿第四章的记号，单位成本为：

\[
c=\frac{w}{a}+m+\rho.
\tag{5.4}
\]

w为每工时工资，m为其他真实单位投入，\ensuremath{\rho}为按产出收取的接入费用。a提高直接改变的是单位劳动投入这一成本项。假设w=50元／工时，原a=10件／工时，m=3元／件、\ensuremath{\rho}=2元／件，单位成本为10元。a升到12.5件／工时后，劳动成本从5元降到4元，整体成本降为9元，降幅10\%，小于生产率增幅25\%。

\textbf{表5.2 假设例子：工资与接入费用怎样改变单位成本}

\begingroup\footnotesize\setlength{\tabcolsep}{4pt}\renewcommand{\arraystretch}{1.25}
\begin{longtable}{>{\raggedright\arraybackslash}p{\dimexpr 0.28571\textwidth-8.0pt\relax}>{\raggedright\arraybackslash}p{\dimexpr 0.14286\textwidth-8.0pt\relax}>{\raggedright\arraybackslash}p{\dimexpr 0.14286\textwidth-8.0pt\relax}>{\raggedright\arraybackslash}p{\dimexpr 0.14286\textwidth-8.0pt\relax}>{\raggedright\arraybackslash}p{\dimexpr 0.14286\textwidth-8.0pt\relax}>{\raggedright\arraybackslash}p{\dimexpr 0.14286\textwidth-8.0pt\relax}}
\toprule
\textbf{同一假设技术条件下的分配与收费} & \textbf{w} & \textbf{a} & \textbf{m} & \textbf{\ensuremath{\rho}} & \textbf{单位成本c（元／件）} \\
\midrule\endfirsthead
\toprule
\textbf{同一假设技术条件下的分配与收费} & \textbf{w} & \textbf{a} & \textbf{m} & \textbf{\ensuremath{\rho}} & \textbf{单位成本c（元／件）} \\
\midrule\endhead
\bottomrule\endfoot
基线 & 50 & 10 & 3 & 2 & 10 \\
工资及收费不变 & 50 & 12.5 & 3 & 2 & 9 \\
时薪提高20\%，收费不变 & 60 & 12.5 & 3 & 2 & 9.8 \\
时薪不变，接入费用增加1元 & 50 & 12.5 & 3 & 3 & 10 \\
\end{longtable}\endgroup


表中第三行让劳动者分享一部分节省，第四行则让上游接入提供者取得这部分收益。采用企业的单位成本甚至可以保持不变，但工资、售价和后续需求的后果不同。所以，训练或设备便宜了多少，还不足以说明使用企业愿意投入多少，须把工资和收费一起查。

固定采用成本F进一步影响采用的经济可行性。在原产量10000件、售价与订单保持不变的假设下，单位节省1元产生10000元毛节省，若当期归属的采用支出为4000元，净节省为6000元。这一计算仅作既定价格和产量下的成本初步比较。市场中的价格和对手销量会变化；单家采用与全行业同时采用的利润比较，需要分别求解第十四章的竞争模块。把全行业降成本后的利润增加当作先采用者的收益，会遗漏竞争者反应。

成本下降对劳动需求的传导，可以通过局部弹性关系分析。\ensuremath{\varepsilon}表示需求量对价格下降的敏感程度，\ensuremath{\phi}表示成本下降向价格下降的传递程度，\ensuremath{\lambda}=(w/a)/c是单位成本中的劳动部分。在工资、其他成本、收费、产品质量和需求条件保持不变的一次小变化中，dlnc=\ensuremath{-}\ensuremath{\lambda}dlna；若dlnp=\ensuremath{\phi}dlnc、dlnQ=\ensuremath{-}\ensuremath{\varepsilon}dlnp，则：

\[
d\ln H=(\varepsilon\phi\lambda-1)d\ln a.
\tag{5.5}
\]

乘积大于1，价格和需求扩张通道才足以超过单位劳动节省；乘积小于1，工时在这些限定条件下下降。\ensuremath{\varepsilon}、\ensuremath{\phi}和\ensuremath{\lambda}不是普遍常数：企业是否竞争、需求是否接近饱和、劳动成本占比怎样，都影响它们。例如客户已经得到足够服务，价格降低也可能只换来较少新需求；如果客户原被价格挡在市场之外，销量响应可以更大。大幅技术变化、收费上调或收入变化时，应重新计算，不能把局部式子当长期预测。

企业扩张与行业扩张还需分开。采用者可能从竞争者夺得订单，Q\_i增长而行业Q不变。以两家企业为教学假设，原来各生产500件、a均为10，合计100工时；采用者a升到12.5，取得700件订单，需要56工时，较原先增加6；另一家失去订单，只生产300件，需要30工时，合计降至86。观察成功采用者会看到工时增加，观察行业却看到下降14\%。这解释了企业处理效应与地区结果为何可能异号，而不借由任何一项研究宣布哪一种永远占优。

\section{机器人引起再配置，群体承担不同调整}

\subsection{安装增长与地区集中同时发生}

部署规模提供劳动分析的背景，而不是就业效果。IFR于2026年9月24日发布的同版年度图显示，全球工业机器人新安装量由2015年的25.4万台升至2025年的60.3万台，按公布精度计算增长约137.4\%，十个年度间隔的年均复合增长约9.0\%。2023—2024年接近平台，2025年重新增长，说明长期扩散并不等于每年同速增长。\href{https://ifr.org/downloads/press_docs/Graphic_World_Robotics_2026_GLOBAL.png}{IFR同版全球年度图}

\begin{center}\begin{tikzpicture}\begin{axis}[width=.46\textwidth,height=6.2cm,title={A　新安装规模与构成},ylabel={千台},ybar stacked,bar width=7pt,ymin=0,ymax=650,xmin=2014.4,xmax=2025.6,xtick={2015,2017,2019,2021,2023,2025},tick label style={font=\scriptsize},label style={font=\small},title style={font=\small},axis lines=left,grid=major,grid style={gray!15},scaled ticks=false,legend style={font=\scriptsize,draw=none,at={(0.5,-0.24)},anchor=north},legend columns=2]
\addplot+[red!65!black,fill=red!65!black,draw=none] coordinates {(2015,69) (2016,97) (2017,156) (2018,155) (2019,145) (2020,176) (2021,275) (2022,290) (2023,276) (2024,295) (2025,354)};
\addlegendentry{中国}
\addplot+[blue!40!gray,fill=blue!40!gray,draw=none] coordinates {(2015,185) (2016,207) (2017,244) (2018,268) (2019,242) (2020,214) (2021,251) (2022,263) (2023,265) (2024,247) (2025,249)};
\addlegendentry{其他地区}
\end{axis}\end{tikzpicture}\hfill\begin{tikzpicture}\begin{axis}[width=.46\textwidth,height=6.2cm,title={B　中国占全球安装比重},ylabel={百分比},ymin=0,ymax=70,xmin=2014.5,xmax=2025.5,xtick={2015,2017,2019,2021,2023,2025},tick label style={font=\scriptsize},label style={font=\small},title style={font=\small},axis lines=left,grid=major,grid style={gray!15},scaled ticks=false,]
\addplot+[red!65!black,mark=*,mark size=1.5pt,thick] coordinates {(2015,27.1653543) (2016,31.9078947) (2017,39) (2018,36.643026) (2019,37.4677003) (2020,45.1282051) (2021,52.2813688) (2022,52.4412297) (2023,51.0166359) (2024,54.4280443) (2025,58.7064677)};
\end{axis}\end{tikzpicture}\end{center}

\textbf{图5-1　工业机器人安装规模与地区构成。}\label{fig:open-industry-robots} 原始值为IFR World Robotics 2026公开图的2015—2025年标签，单位千台；其他地区由全球减中国，比重由显示值计算。图中是年度新增流量，不是运行存量、机器人密度或替代岗位数。\href{https://ifr.org/downloads/press_docs/Graphic_World_Robotics_2026_CHINA.png}{同版中国年度图}；整理序列与制图代码保存在开放产业数据目录。

中国同期新安装量由6.9万台增至35.4万台，占全球比重由约27.2\%升至58.7\%。全球十年新增流量的端点差为34.9万台，其中中国端点差为28.5万台，占约81.7\%。这是安装地点的增长分解，不能读作中国供应商占全球销售的份额，更没有测出生产条件的所有权集中。产线部署集中与设备供应、系统集成、客户认证分别处于产业链的不同位置。

这一趋势扩大了研究劳动调整的必要性，却不能按安装量乘一个固定人数估计失业。同台设备的任务、班次、利用、合格率与订单不同，劳动可能经减少直接操作、增加维护或扩大产量重新配置。下面的地区、行业和人员研究因此承担不同任务：它们分别检验部署暴露如何进入就业、生产率与职业轨迹。

\subsection{地区、行业与人员记录呈现不同调整}

机器人接手体力任务以后，还会改变制造业生产、产品贸易和地区用工。下列三项研究各看一部分，不能合成一个统一“替代率”。按公式5.1，行业生产率接近a的变化，地区就业还受到产出、劳动供给和迁移影响，跟踪个人则能看到谁留下、谁离开。先分清观察范围，再比较结果。

\textbf{表5.3 机器人研究的估计对象与传导位置}

\begingroup\footnotesize\setlength{\tabcolsep}{4pt}\renewcommand{\arraystretch}{1.25}
\begin{longtable}{>{\raggedright\arraybackslash}p{\dimexpr 0.26549\textwidth-8.0pt\relax}>{\raggedright\arraybackslash}p{\dimexpr 0.24779\textwidth-8.0pt\relax}>{\raggedright\arraybackslash}p{\dimexpr 0.26549\textwidth-8.0pt\relax}>{\raggedright\arraybackslash}p{\dimexpr 0.22124\textwidth-8.0pt\relax}}
\toprule
\textbf{原研究与观测范围} & \textbf{方法和结果单位} & \textbf{核实的主要发现} & \textbf{在本章传导链中的位置} \\
\midrule\endfirsthead
\toprule
\textbf{原研究与观测范围} & \textbf{方法和结果单位} & \textbf{核实的主要发现} & \textbf{在本章传导链中的位置} \\
\midrule\endhead
\bottomrule\endfoot
Acemoglu—Restrepo（2020），美国722通勤区，1990—2007 & 行业机器人变化与地区基线结构构成暴露，欧洲采用作工具变量；就业人口比、工资 & 每千劳动者增1台，局部首选约\ensuremath{-}0.39个百分点、\ensuremath{-}0.77\%；纳入贸易溢出的全国模型约\ensuremath{-}0.2个百分点、\ensuremath{-}0.42\% & 现实采用暴露之后的地区需求及分配；局部回归与全国推算分列 \\
Graetz—Michaels（2018），17国行业，1993—2007 & 行业长差分与任务工具变量；劳动生产率、工时、技能份额 & 年度劳动生产率增长贡献约0.36个百分点；总工时下降不显著，低技能工时份额下降 & 行业能力增长及结构变化，未逐人证明转岗结果 \\
Dauth等（2021），德国行政记录，1994—2014 & 地区暴露工具变量及劳动者轨迹 & 制造业置换由服务业新增抵消；年轻进入者更多承担制造业入口调整，在职者更多内部转任务 & 部门再配置与招聘流量，解释总量稳定中的群体差异 \\
\end{longtable}\endgroup


相关原论文分别为\href{https://doi.org/10.1086/705716}{美国研究}、\href{https://doi.org/10.1162/rest_a_00754}{跨国研究}和\href{https://doi.org/10.1093/jeea/jvab012}{德国研究}，定位及排除限制见台账L01、L02、L09。表中保留与本章机制直接相关的结果，后续分析着重讨论估计对象之间的联系。

美国的0.39是就业人口比的百分点变化。若一个仅用于换算的地区有10万劳动年龄人口且人口边界固定，下降0.39个百分点对应390名就业者；它不是390人已被某一台机器人替代的实际记录。暴露单位是每千劳动者增一台，且地区效应包含企业和服务需求的联动。工资下降0.77\%也不是0.77个百分点。单位换算有助于理解量级，但不改变原研究的识别对象。

三项研究的共同信息是劳动投入的再配置，而非统一替代率。行业生产率提高可以与特定地区就业下降并存，服务业补足总就业也可以与制造业入口收窄并存。因此，净就业是否恢复与原受影响者是否恢复收入，是两个不同结果。

德国的人员轨迹补充了总量研究：在职者更多内部转任务，年轻进入者更多承担入口调整。制造业可以通过少补离职者、少招新人收缩，而不主要依靠集中裁员。服务业新增职位提供机会，却不能据此确认同一批年轻人取得等同收入和职业前景。

跨国结果还混合技术、贸易、需求和制度差异，不能将其全部归因于劳动保护。识别制度作用，需要相近技术冲击下的排班、培训与议价比较，并检验各研究的识别条件。现有材料显示结果受环境影响，尚未分离每项制度的独立贡献。

\section{AI任务效果分化，早期劳动收入变化有限}

生成式AI可以缩短取得文字和部分知识的时间，但交付责任通常仍由人承担。工具可能少让人起草，也可能帮助查证和沟通，同时增加监督输出的工作。因此，要先分清测量了什么。下表把研究结果与代数换算列开；换算假定是同一合格任务，其余条件不变。

\textbf{表5.4 AI研究观测值与固定产出的代数换算}

\begingroup\footnotesize\setlength{\tabcolsep}{4pt}\renewcommand{\arraystretch}{1.25}
\begin{longtable}{>{\raggedright\arraybackslash}p{\dimexpr 0.28302\textwidth-8.0pt\relax}>{\raggedright\arraybackslash}p{\dimexpr 0.16981\textwidth-8.0pt\relax}>{\raggedright\arraybackslash}p{\dimexpr 0.28302\textwidth-8.0pt\relax}>{\raggedright\arraybackslash}p{\dimexpr 0.26415\textwidth-8.0pt\relax}}
\toprule
\textbf{研究及原始测量} & \textbf{已核结果} & \textbf{固定产出下的数学换算} & \textbf{不是这项研究直接测得的结果} \\
\midrule\endfirsthead
\toprule
\textbf{研究及原始测量} & \textbf{已核结果} & \textbf{固定产出下的数学换算} & \textbf{不是这项研究直接测得的结果} \\
\midrule\endhead
\bottomrule\endfoot
Brynjolfsson、Li、Raymond（2025），5,172名客服、分批部署、微调GPT-3；每小时解决问题数 & 平均约+15\%，低经验者收益较大 & 若a\ensuremath{_1}/a\ensuremath{_0}=1.15，H\ensuremath{_1}/H\ensuremath{_0}=1/1.15\ensuremath{\approx}0.8696，即\ensuremath{-}13.04\% & 工资涨15\%、员工减13.04\%、全国生产率涨15\% \\
Noy、Zhang（2023），453名专业人员、GPT-3.5随机写作；完成时间与评定 & 耗时约\ensuremath{-}40\%，评分约+18\% & 若同任务时间比0.60，则每小时任务能力比1/0.60\ensuremath{\approx}1.6667，即+66.67\% & 整个职业产出每小时增加66.67\%，或质量可简单加成 \\
Dell’Acqua等（2026正式稿），758名咨询者、2023年GPT-4随机实验 & 边界内耗时约\ensuremath{-}25.1\%；越界正确率合并约\ensuremath{-}19个百分点 & 同一合格任务耗时比0.749的倒数约1.335；越界正确率须另列 & 全流程统一增加33.5\%或所有复杂任务都变差 \\
\end{longtable}\endgroup


这些换算揭示一个常见单位错误：每小时产出提高15\%与完成同产出省时15\%不是同一变化。时间减少40\%对应速度提高66.67\%，也不能直接当作质量调整后的职业生产率。不同研究的测量口径与研究设计并不一致，不能仅将参与者数量相加，便把这些结果合并为统一的“AI劳动样本”。

客服正式论文提供实际工作场所中的重要证据：工具帮助经验较少者更快完成服务，知识提示可以在企业内部扩散。研究采用分批部署的双重差分与事件研究，而非将全部人员随机分到通用ChatGPT；低经验收益还取决于该企业的任务与系统。正式版5,172人、约15\%与2023工作稿5,179人、14\%是版本差异，正文使用前者。作者没有个人工资记录，该效率结果不提供下一节收入分配机制所需的参数。\href{https://doi.org/10.1093/qje/qjae044}{正式论文}

写作与咨询实验识别限定任务中的工具访问效果；企业交付还包含任务选择、质量责任和组织费用。越界任务正确率降低，说明适用范围是采用条件的一部分。人工核验能够纠错，但其新增工时也须进入式（5.3）。

Peng等（2023）的Copilot实验为标准化编程任务提供了较大的省时结果，但其分母需要完整保留。实验实际发生于2022年5月15日至6月20日，从Upwork招募95名专业开发者，随机分为45名获准使用Copilot者与50名对照者；任务是在JavaScript中实现HTTP服务器，以预先给出的12项测试验收。两组各有35人完成任务及调查。成功完成者的平均耗时分别为71.17和160.89分钟，论文报告耗时减少55.8\%，95\%置信区间为减少21\%至89\%。这里比较的是成功完成任务者的耗时，不能把95名入组者都视为已经交付，也不能将省时比例写成同幅度的每小时产出增长。\href{https://arxiv.org/pdf/2302.06590v1}{固定预印本方法与结果，印刷第3—6页}

Copilot实验预先规定需求、验收和起点代码，少了跨团队沟通、理解长期代码约定的负担。因此，工具可能在目标清楚的代码实现中，减少搜索和输入时间。这说明某类任务适合工具，不能把它当成编程职业不变的技术效果。工具若也影响谁能完成，只比较完成者还可能改变两组构成；成功率差异的区间跨过零，不能证明完成者与未完成者可以互换。完整评价仍要算失败尝试、后续维护和安全核验。原论文未独立评价代码质量，测试通过也未覆盖全部长期质量要求。

按项目交付的人，有订单时省时可以用于多接，也可能用于少工作；按时计酬的人，速度与收入又是另一种关系。该实验没有估计订单需求、市场工资和平台收费，不能直接推出增收。放回公式5.1和5.3，可以先确认既定任务怎样变化；报酬和就业仍须结合合同与采用后的市场结果考察。

Becker等（2025）的METR研究把任务移到开发者长期参与的大型开源代码库。本文使用2025年7月25日的第二版预印本：16名开发者完成246个实际问题，任务预先确定后随机允许或禁止AI，主要采用Cursor Pro及Claude 3.5／3.7 Sonnet。控制随机分配前的预计任务耗时后，回归估计允许AI使完成时间增加19\%；这是经任务难度预测调整后的结果，不是两个组原始均值直接相除。原论文图1与附录D报告区间方法；作者后续更新将该早期研究的95\%区间列为耗时增加2\%至39\%。\href{https://arxiv.org/pdf/2507.09089v2}{固定原稿§2—3、附录D}，\href{https://metr.org/blog/2026-02-24-uplift-update/}{作者对原研究区间与后续识别问题的说明}

测量包括开发者提交修改以及回应审查所花的自身工作时间，接近实际交付流程，但不是整个项目团队的总工时；其中27个未收到审查的问题，其审查后工作时间由作者按条件组均值插补。经筛选的录屏子样本显示，编码和搜索时间份额下降，同时出现提示、等待和审查AI输出等投入。这支持检查工具是否将时间从制作转向监督，却未分别随机改变各项活动，不能证明某一种活动单独造成全部减速。与Copilot实验相比，既有代码的熟悉程度、隐含规范、任务类型、工具代际和验收方式均不同，两项研究的正负结果无法单独识别“经验越高收益越低”，也不适合按样本人数求一个平均AI效应。

两项编程研究要求用完整交付评价工具，而非仅用代码生成量。它们同时改变任务、人员经验、工具版本和代码库环境，不能据此单独识别哪项差异造成效果反转。失败尝试、审查和整合是下一步机制比较的对象，工资与权限则尚未在两项实验中得到检验。

新增案例的逐项数值、实读范围与版本见劳动与平台新增案例核查及数值长表。

Humlum与Vestergaard的2026年3月13日工作稿把丹麦11类高暴露职业的调查与行政记录相连，最新轮约25,000回应者、7,000工作场所，观察到2024年底。汇总双重差分对收入和登记工时接近零，区间排除超过约2\%的平均变化，而任务、企业倡议和职业调整已经发生。\href{https://www.rfberlin.com/wp-content/uploads/2026/03/26078.pdf}{固定工作稿}

平均结果同时遮蔽了职业流动。原论文§3.3和图5(c)显示，AI采用与更活跃的职业迁移存在关联，转岗采用者更多进入工资溢价和AI相关度较高的职位。作者保留两种解释：工具帮助劳动者进入新职业，或转岗者更需要工具适应新工作。因此，这一关联尚不能单独识别AI对转岗及其收入收益的净因果效应。\href{https://www.rfberlin.com/wp-content/uploads/2026/03/26078.pdf}{原论文§3.3，印刷第26—29页}

丹麦资料区分了任务变化与早期劳动市场结果：截至2024年底，所研究职业已发生组织和职业调整，平均收入与登记工时仍无相应的大幅变化。这与任务实验并不矛盾，二者观测环节不同。省时后的继续工作、采用费用和报酬调整滞后是候选解释，现有跨样本比较尚未识别各自贡献。

\section{时薪、总收入与劳动份额可以不同向变化}

工资总额W、时薪w与付酬工时H的关系是W=wH。这一核算要求H与工资账户覆盖同一批付酬劳动；时薪不同时，w使用按工时加权的平均数。流程效率若包括外包或未付酬劳动，应分别报告其工时及成本，不能把未付酬时间直接乘一个在岗时薪。以下算例把流程工时与付酬工时边界设为一致。结合公式5.1，有：

\[
\frac{W_1}{W_0}=\frac{w_1}{w_0}\frac{Q_1/Q_0}{a_1/a_0},
\qquad
\Delta\ln W=\Delta\ln w+\Delta\ln Q-\Delta\ln a.
\tag{5.6}
\]

假设生产率提高25\%、产出增加20\%、时薪提高10\%，工时比为0.96，工资总额比为1.10\ensuremath{\times}0.96=1.056，即增加5.6\%。时薪提高与工时减少可以同时存在；劳动者个人收入还要检查h和失去岗位的时间。若同口径增加值提高20\%，劳动份额的比值为1.056/1.20=0.88。原份额60\%会降至52.8\%。工资总额增长与劳动份额下降可以并存，因为同口径增加值的增长更快。

工具对分配的作用取决于双方替代能力。经验复制可以提高新人表现，也可以扩大雇主可选人员；复杂业务又可能提高判断和责任承担的需要。实际工资结果须结合外部岗位、人员供给和协商条件。第六章说明平台场景，第十四章给出固定工时与收入下的条件谈判推导。

时薪之外，缩时的收入规则同样关键。5.1中保持500元的8小时工作，是劳动者取得时间收益；降为400元的8小时工作，则把工资现金节省留在企业。岗位是否稳定并不能区分这两种分配。若同样的节省改成10人中2人失去工作，剩余者收入不变，损失便集中于退出者。平均在岗工资不变可能掩盖群体工资总额下降，持续跟踪原人群才能恢复完整结果。

总工时与劳动强度分别计量。同一小时处理更多任务，可能来自减少无效搜索，也可能来自压缩休息、并行服务或加强考核。需要同时观察间歇、异常处理和责任要求，才能区分组织节约与强度增加；技术版本和绩效规则也须分别记录。

工具使用共同经验，个人产出就更依靠资料、培训和维护协作。这些能力怎样使用，才把本章问题接到占有关系。企业决定扩产、减员或改标准，劳动者可以通过工资和工时协议分享收益；协议若又让劳动者共同修改标准、决定培训或安全事项，才改变相应的劳动控制。因此，收入、闲暇和权限须分别查明，不能一项改善就说全部改善。

\subsection{机器人减少部分用工，工资购买仍受规模和结构影响}

机器人首先压缩其适用任务的工时，企业随后决定减少招聘、裁员、内部调岗还是扩大生产。重复操作岗位减少，而编程、设备维修、质检和异常处置岗位增加时，可变资本的构成已发生变化；总额却要把各岗位的报酬和实际工时合起来看。少了人数而平均报酬提高，可以与工资支出不变甚至增加同时发生。留任者工资也不能代表退出者的结果。

表5.3的研究分别观察地区、行业和人员轨迹，恰好限定了这项推论。跨国研究未检出总工时显著下降，同时发现低技能工时份额下降；德国研究发现制造业置换与服务业新增并存。它们支持技术改变劳动配置，却未提供一套可直接相加的全社会可变资本账户。\href{https://doi.org/10.1162/rest_a_00754}{跨国研究}、\href{https://doi.org/10.1093/jeea/jvab012}{德国研究}。服务业新增还要辨认生产关系与报酬，不能按岗位数量等额抵消制造业工资损失。

上述假设算例已给出一个相反于“必然减少”的结果：生产率提高25\%，产出增加20\%，时薪提高10\%，工资总额增加5.6\%，劳动份额却下降。它说明绝对支出与相对分配可以反向变化，并未测量$c_M/v_M$。机器人安装图同样没有工资和资本预付数据，不能用安装增长直接计算有机构成。真正的部门检验需要把采用后的岗位工资、产量、设备估值及回流周期配对。

机器人还改变资本对留任人员的使用。一人看管多台设备、设备提速后跟随更密集的操作节奏，可能让操作人数减少与个人劳动负担增加同时发生；自动检测减少重复检查，则也可能减轻负担。决定差别的是设备配置、工作标准和休息安排。工资账户反映购买额，岗位覆盖和负担记录反映使用方式，二者共同检验技术节约是减轻劳动还是集中到较少人员身上。

\section{新增任务提供机会，转岗仍需要时间和保障}

Autor、Chin、Salomons与Seegmiller（2024）使用美国1940—2018年的职业细称、普查和专利资料，区分与职业产出互补的创新和自动化创新，显示新工作的来源及职业位置随时期改变。\href{https://doi.org/10.1093/qje/qjae008}{正式论文}。这提供长期新任务形成的证据，却不是每个被替代岗位都找到同等新工作的追踪记录。职称、岗位、工时与实际采用是不同变量。

新任务只有进入持续需求和报酬安排，才形成稳定岗位。监督任务可以由新增人员承担，也可以加给现有员工；异地或高资格职位也未必向原受冲击者开放。职业名称增加不能替代人员进入、培训时间和后续收入的追踪。

在职人数之外，还要查招了多少人、走了多少人。企业可以不裁员，却少补离职者、少招初级岗位，逐步改变劳动投入。德国机器人研究说明职业入口确实会调整；丹麦早期AI工作稿则要求谨慎，不能把暴露职业的入口变化全归于实际采用。人数暂时不变，进入退出仍可能改变；变化是谁造成的，还需识别。

新人做简单任务，既生产本期Q，也靠反馈学会以后处理例外。企业取消初级任务，却提供更好的导师和模拟练习，学习可能加快；只省掉新人岗位、同时少训练，未来判断能力的供给就可能减弱。今天的成本与以后的专业人才由此相连。要判断这条机制，须跟踪同一群体的能力、导师工时和晋升，单次工具评分不够。

培训的可行性还取决于停工收入、生活与照护安排。学费免除只减少一类付款，不能消除学习的时间和其他必要投入。家庭承担更多费用可能降低企业支出，但是否构成转移，仍需相同培训活动的承担与补偿记录。第十三章进一步分析这些条件。

\section{所得的支出方式影响技术扩产后的销售}

效率改善只有进入实际生产并找到需求，才成为可实现的收入。节省劳动所形成的工资减少与利润增加，会改变不同主体的支出。设工资所得W、资本所得\ensuremath{\Pi}，在一个仅展示当期支出通道的简化关系中，消费C=c\_wW+c\_k\ensuremath{\Pi}，c\_w、c\_k分别是各类所得的支出倾向。若总所得固定而W转给\ensuremath{\Pi}，则\ensuremath{\Delta}C=(c\_k\ensuremath{-}c\_w)\ensuremath{\Delta}\ensuremath{\Pi}。

假设c\_w=0.8、c\_k=0.4，工资减少1000元、资本所得增加1000元，第一轮消费少400元。这是明确的假设例子，尚未包括投资、公共支出、外部需求和后续价格变化。它说明效率收益若主要归于较少支出的主体，原本期望的销量扩张可能减弱；它不证明所有利润增加都造成需求收缩。利润用于新增投资、公共转移或出口需求扩大时，结果可以不同，第十四章将检查相应现金和资源约束。

局部很大的任务收益也受到经济权重限制。Acemoglu（2025，2024年先在线）依据任务权重、可采用份额和成本节省校准AI宏观效应，十年累计TFP约提高0.66\%，考虑难学习任务后约0.53\%。\href{https://doi.org/10.1093/epolic/eiae042}{正式期刊稿}。它是条件推算，0.66\%不是每年增长率，也不是本报告测得的GDP。其意义在于要求将任务收益加权并计入采用，而非接受某一固定长期上限。

第5.2的2.04\%算例也不是宏观TFP：它把特定工时份额重新计算为劳动生产率，尚无资本、材料和价格权重。宏观加总要考虑部门增加值、投入产出联系、市场再配置与质量，客服的15\%局部增幅不足以直接推算全国就业或生产率变化。较小总量收益与特定群体承受明显损失并不矛盾，保护与培训的需要可以来自分布，而非必须等待一个很大的总量冲击。



\section{单位劳动节约通过需求和劳动安排形成不同收益}

本章将劳动后果分为完整流程效率、总劳动需求与个人分享。工时关系说明需求扩张怎样改变单位节约的总量后果；机器人研究显示企业、部门和地区结果可以不同；AI研究显示任务省时可以先于平均收入与登记工时的明显变化。材料支持这种分层解释，未给出统一职业替代率。

省下同样的劳动，收益形式可以不同。保薪缩时让劳动者多休息，扩产让社会多得产品，少用人则可能让离岗者承担调整；降价让购买者受益，购买者与失去工作者又未必是同一群体。因此，技术收益的分配必须落实到具体群体和时期，不能以企业利润或平均时薪代替全部劳动后果。

这一结论将技术条件接到了合同条件。企业能够扩产、排班和招聘，并不等于劳动者能够改变这些决定。第六章进一步解释这种不对称怎样在分散劳动中形成：订单、评价和账户规则，使没有集中在同一车间的劳动者也受到同一组织安排。

\section*{本章参考文献}

IFR. (2026). World Robotics 2026: Industrial Robots. 2026年9月24日公开年度图和发布材料。\href{https://ifr.org/ifr-press-releases/news/five-million-robots-now-operate-in-factories-globally}{官方发布入口}。新增图5-1采用同版标签，未读取收费完整数据库。

1. Acemoglu, D., \& Restrepo, P. (2019). Automation and new tasks: How technology displaces and reinstates labor. \emph{Journal of Economic Perspectives, 33}(2), 3–30. \url{https://doi.org/10.1257/jep.33.2.3}
2. Acemoglu, D., \& Restrepo, P. (2020). Robots and jobs: Evidence from US labor markets. \emph{Journal of Political Economy, 128}(6), 2188–2244. \url{https://doi.org/10.1086/705716}
3. Graetz, G., \& Michaels, G. (2018). Robots at work. \emph{Review of Economics and Statistics, 100}(5), 753–768. \url{https://doi.org/10.1162/rest_a_00754}
4. Dauth, W., Findeisen, S., Suedekum, J., \& Woessner, N. (2021). The adjustment of labor markets to robots. \emph{Journal of the European Economic Association, 19}(6), 3104–3153. \url{https://doi.org/10.1093/jeea/jvab012}
5. Brynjolfsson, E., Li, D., \& Raymond, L. (2025). Generative AI at work. \emph{Quarterly Journal of Economics, 140}(2), 889–942. \url{https://doi.org/10.1093/qje/qjae044}
6. Noy, S., \& Zhang, W. (2023). Experimental evidence on the productivity effects of generative artificial intelligence. \emph{Science, 381}(6654), 187–192. \url{https://doi.org/10.1126/science.adh2586}
7. Dell’Acqua, F., McFowland, E., III, Mollick, E., Lifshitz, H., Kellogg, K. C., Rajendran, S., Krayer, L., Candelon, F., \& Lakhani, K. R. (2026). Navigating the jagged technological frontier: Field experimental evidence of the effects of artificial intelligence on knowledge worker productivity and quality. \emph{Organization Science}, advance online publication. \url{https://doi.org/10.1287/orsc.2025.21838}
8. Autor, D., Chin, C., Salomons, A., \& Seegmiller, B. (2024). New frontiers: The origins and content of new work, 1940–2018. \emph{Quarterly Journal of Economics, 139}(3), 1399–1465. \url{https://doi.org/10.1093/qje/qjae008}
9. Acemoglu, D. (2025). The simple macroeconomics of AI. \emph{Economic Policy, 40}(121), 13–58. \url{https://doi.org/10.1093/epolic/eiae042}
10. Humlum, A., \& Vestergaard, E. (2026). Still waters, rapid currents: Early labor market transformation under generative AI. \emph{RFBerlin Discussion Paper 078/26}, March 13 version, working paper. \url{https://www.rfberlin.com/wp-content/uploads/2026/03/26078.pdf}
11. Peng, S., Kalliamvakou, E., Cihon, P., \& Demirer, M. (2023). \emph{The impact of AI on developer productivity: Evidence from GitHub Copilot}. arXiv:2302.06590v1, February 13, preprint. \href{https://doi.org/10.48550/arXiv.2302.06590}{DOI：10.48550/arXiv.2302.06590}; \href{https://arxiv.org/pdf/2302.06590v1}{本章所读固定版本}。访问日期：2026-10-08。
12. Becker, J., Rush, N., Barnes, E., \& Rein, D. (2025). \emph{Measuring the impact of early-2025 AI on experienced open-source developer productivity}. arXiv:2507.09089v2, July 25, preprint. \href{https://doi.org/10.48550/arXiv.2507.09089}{DOI：10.48550/arXiv.2507.09089}; \href{https://arxiv.org/pdf/2507.09089v2}{本章所读固定版本}。论文题页将Barnes署名列为Beth，arXiv书目信息列为Elizabeth。访问日期：2026-10-08。
13. Becker, J., Rush, N., Cunningham, T., Rein, D., \& Mahamud, K. (2026). \emph{We are changing our developer productivity experiment design}. METR, February 24, research update. \href{https://metr.org/blog/2026-02-24-uplift-update/}{作者原始说明}。仅用于补充原研究区间与时效、选择限制，不替代2025年原论文结果。访问日期：2026-10-08。

\chapter{平台劳动：共同协作形成收益，规则制定影响分配}\label{ch:6}

平台通过匹配、评价与结算组织分散劳动。减少等待可以创造真实协调收益，收费和必要成本决定劳动者保留多少，规则及外部机会决定其能否拒绝不利条件。本章将效率、净所得和实际权限分别分析，解释收入机会扩大与规则权集中为何可以并存。

下文先算全部时间和净收入，再查替代工作与谈判条件，最后比较规则怎样执行、劳动者怎样回应。假设算例用于分清机制，纽约报送资料用于查收入与时间，调查和历史案例用于查工作经历与制度形式。各项材料观察的群体不同，不能合成同一批人的因果估计。

\section{少等待可以提高每小时的实际产出}

平台能够提高劳动生产率的一条途径，是减少寻找订单、空驶和等待，而不改变完成任务本身的速度。把一段观察期内完成的服务量记为 $q$，付费任务时间记为 $h_{paid}$，取得这些任务所需但没有单独付费的时间记为 $h_{required}$，则劳动者实际投入为：

\[
H=h_{paid}+h_{required},\qquad
u=\frac{h_{paid}}{H},\qquad
\kappa=\frac{q}{h_{paid}},\qquad
a_{eff}=\frac{q}{H}=\kappa u.
\]

$u$ 是时间利用率，$\kappa$ 是付费任务时段的服务速度，两者共同决定完整工时生产率 $a_{eff}$。它与前章的 $a=Q/H$ 具有一致的解释：数值越大，每工时完成的产出越多。各项时间须按同一观察期计量，重叠活动只计一次。例如，等待期间同时做其他已计入工作，不能再次把整段等待加进总工时；但为了取得收入而不能自由支配的待命，也应按实际投入纳入工时分母，不能仅因没有单独付费而遗漏。

假设每付费小时完成两项服务，即 $\kappa=2$。原来30小时中15小时有任务，$u=1/2$，完整工时生产率为每小时一项服务。匹配改善后，同样30小时中20小时有任务，$u=2/3$，完整工时生产率变成每小时 $4/3$ 项，增加33.33\%。这一增幅来自闲置时间减少，服务执行速度在假设中保持不变。若平台只公布每个任务用时或订单总数，就会漏掉这条效率渠道，也会漏掉等待增加的相反情形。

Baiocco等将算法管理理解为通过计算机编程程序协调组织中的劳动投入，并讨论其在平台和常规工作场所中的运用。这一概念允许区分市场撮合、劳动安排和人工智能：自动排班未必采用机器学习，采用人工智能的写作工具也未必控制派单。（Baiocco等，2022，\href{https://www.ilo.org/publications/algorithmic-management-work-and-its-implications-different-contexts}{官方出版页}，概念框架）本章考察具体规则对劳动的影响，重点分析其如何改变 $u$、$\kappa$ 和劳动条件。

匹配能省多少时间，要看原来有多少等待。订单密度、路程和客户时段会限制匹配，远程服务还要搜索项目、报价、跨时区联系。原来已很少等待，改进空间就小；更多劳动者进入而订单没有增加，人均利用率还可能下降。平台能连接更多交易，不等于每个人都能得到同样多的订单，仍须查需求和劳动供给。

\section{匹配改善与收入增加，受不同条件制约}

任务数量与劳动者可支配的工作收入应分别核算。令 $R$ 为观察期内订单服务的收入、按扣除平台费之前的口径计量，$\theta$ 为该口径下的收费比例，$C$ 为劳动者承担的其他必要成本，则净工作收入与有效时薪为：

\[
Y_{net}=(1-\theta)R-C,\qquad
w_{eff}=\frac{Y_{net}}{H}.
\]

$C$ 不再包含已单列的平台费。资料若只有扣费后的到账收入，就使用到账收入减其他成本，避免重复扣费。车辆、设备、通信及维护费用要归入相应观察期，长期设备支出须按使用期合理归属，既不宜全部计入单日成本，也不宜遗漏其购置费用。这里的 $Y_{net}$ 是扣除上述经营成本后的收入，并未自动包含税费、社会保障或风险补偿。不同资料采用不同净额口径时，比较应先校正口径。

\textbf{表6.1 假设例子：匹配、收费与完整工时收入的分开变化}

假设每项服务收入50元，每项必要成本10元，付费时段每小时完成两项服务。以下各行为独立比较情景，币值和收费率均非实际平台数据。

\begingroup\footnotesize\setlength{\tabcolsep}{4pt}\renewcommand{\arraystretch}{1.25}
\begin{longtable}{>{\raggedright\arraybackslash}p{\dimexpr 0.14667\textwidth-8.0pt\relax}>{\raggedright\arraybackslash}p{\dimexpr 0.09333\textwidth-8.0pt\relax}>{\raggedright\arraybackslash}p{\dimexpr 0.09333\textwidth-8.0pt\relax}>{\raggedright\arraybackslash}p{\dimexpr 0.09333\textwidth-8.0pt\relax}>{\raggedright\arraybackslash}p{\dimexpr 0.09333\textwidth-8.0pt\relax}>{\raggedright\arraybackslash}p{\dimexpr 0.09333\textwidth-8.0pt\relax}>{\raggedright\arraybackslash}p{\dimexpr 0.09333\textwidth-8.0pt\relax}>{\raggedright\arraybackslash}p{\dimexpr 0.09333\textwidth-8.0pt\relax}>{\raggedright\arraybackslash}p{\dimexpr 0.09333\textwidth-8.0pt\relax}>{\raggedright\arraybackslash}p{\dimexpr 0.10667\textwidth-8.0pt\relax}}
\toprule
\textbf{情景} & \textbf{服务量 $q$} & \textbf{付费时间} & \textbf{必要未单付费时间} & \textbf{总工时 $H$} & \textbf{收入 $R$} & \textbf{收费率} & \textbf{其他成本 $C$} & \textbf{净收入 $Y_{net}$} & \textbf{有效时薪} \\
\midrule\endfirsthead
\toprule
\textbf{情景} & \textbf{服务量 $q$} & \textbf{付费时间} & \textbf{必要未单付费时间} & \textbf{总工时 $H$} & \textbf{收入 $R$} & \textbf{收费率} & \textbf{其他成本 $C$} & \textbf{净收入 $Y_{net}$} & \textbf{有效时薪} \\
\midrule\endhead
\bottomrule\endfoot
原有匹配 & 30项 & 15小时 & 15小时 & 30小时 & 1500元 & 10\% & 300元 & 1050元 & 35元 \\
保持服务量、减少等待 & 30项 & 15小时 & 10小时 & 25小时 & 1500元 & 10\% & 300元 & 1050元 & 42元 \\
保持工时、增加服务 & 40项 & 20小时 & 10小时 & 30小时 & 2000元 & 10\% & 400元 & 1400元 & 46.67元 \\
增加服务、提高收费 & 40项 & 20小时 & 10小时 & 30小时 & 2000元 & 25\% & 400元 & 1100元 & 36.67元 \\
增加服务、进一步收费 & 40项 & 20小时 & 10小时 & 30小时 & 2000元 & 30\% & 400元 & 1000元 & 33.33元 \\
\end{longtable}\endgroup


第二行保留原有收入并释放五小时，改善主要表现为时间；第三行保持工时，净收入增加33.33\%。第四行的完整工时生产率同样增加33.33\%，但有效时薪只比原有匹配提高4.76\%；第五行的生产率仍然较高，有效时薪却比原来下降4.76\%。增加服务的几项情景具有相同技术效率，其收入差异来自收费规则对效率收益的重新分配。真实劳动者也未必能自由选择第二行：如果最低接单要求、奖励门槛或家庭收入需要推动其继续工作，节省时间可能转成更多订单。

表中固定服务价格、必要成本和执行速度，分离匹配与收费的作用。现实中的补贴、距离、需求和人员进入可能同时变化，需用订单价格、平台费、非平台成本及完整时间核对净结果。算例给出可能机制，不预测个人服务量必然同比增加。

ILO2021报告的资料来自约1.2万名劳动者及85家企业代表的调查与访谈，涉及远程和在地平台、不同国家及调查期，包括2019—2020年的在地调查。这一范围有助于识别核算中容易遗漏的活动，不能视为统一抽取的全球随机样本。（ILO，2021，\href{https://www.ilo.org/publications/flagship-reports/role-digital-labour-platforms-transforming-world-work}{出版入口}、\href{https://www.ilo.org/sites/default/files/2025-03/WESO_2021_Appendix.pdf}{方法附录}，附录4）本章计算表未依据该报告的平均数校准，其作用是明确需要分别测量的变量及核算口径。

纽约的最低报酬实施初期资料可以将时间与收入分解落实到真实账户。DCWP于2024年公布的第一季度报告，使用五家应用按月报送的资料，表6.2在同一报告版本内比较2023与2024年第一季度的周均值。其每小时总收入包含应用报酬与小费，未扣除劳动者经营成本；工时包含配送和在线待命，并未调整跨应用重复计数。因此这些是监管报送账户的变化，不能直接当作去重后的个人有效时薪或净就业。\href{https://www.nyc.gov/assets/dca/downloads/pdf/workers/Restaurant-Delivery-App-Data-Q1-2024.pdf}{原报告表1及Technical Note，第3、5页}

\textbf{表6.2 同版纽约外卖报送资料的收入与时间变化}

\begingroup\footnotesize\setlength{\tabcolsep}{4pt}\renewcommand{\arraystretch}{1.25}
\begin{longtable}{>{\raggedright\arraybackslash}p{\dimexpr 0.50000\textwidth-8.0pt\relax}>{\raggedright\arraybackslash}p{\dimexpr 0.25000\textwidth-8.0pt\relax}>{\raggedright\arraybackslash}p{\dimexpr 0.25000\textwidth-8.0pt\relax}}
\toprule
\textbf{指标及口径} & \textbf{2023年第一季度} & \textbf{2024年第一季度} \\
\midrule\endfirsthead
\toprule
\textbf{指标及口径} & \textbf{2023年第一季度} & \textbf{2024年第一季度} \\
\midrule\endhead
\bottomrule\endfoot
应用报酬：美元／报告工时 & 5.89 & 16.95 \\
小费：美元／报告工时 & 5.83 & 2.31 \\
两项合计收入：美元／报告工时 & 11.72 & 19.26 \\
全部报告工时：千小时／周 & 1699 & 1468 \\
配送工时：千小时／周 & 787 & 908 \\
待命工时：千小时／周 & 912 & 560 \\
完成配送：百万次／周，按送达计数 & 2.56 & 2.77 \\
\end{longtable}\endgroup


在报告账户内，配送时间占全部时间的比例由约46.3\%升至61.9\%；每配送小时的送达数却由约3.25降至3.05。依据本章的时间分解，全部工时上的产出改善主要伴随待命份额收缩，不能解释成骑手执行任务速度同比提高。路程、合单与项目构成也可能改变配送小时上的比率，且重复在线影响分母。报酬上升和小费下降同时出现，则表明提高一项支付标准之后，总所得仍要核对其他收入渠道的调整；这里的小费抵消并未消除报告所见的总收入增加。

总量核算进一步确认，小时收入改善并非单由报送工时分母缩小造成。用表中已舍入的小时总收入乘以周工时，报送周收入约由1991万美元增至2827万美元，增长约42.0\%；与此同时，报告总工时下降约13.6\%。这是同版账户的描述计算，包含小费、未扣经营成本，也没有去重跟踪个人。它支持承认所记录群体的总收入改善，同时留下上线机会、退出者和个体净额的问题。报告没有反事实比较地区，不能将全部变化归为最低报酬政策的因果效应。

同版表2还显示，每次送达的平均消费者应用费由4.96升至7.26美元，小费由3.87降至1.23美元，总支付由38.35升至39.11美元；应用费包含订阅收费。总价小幅变化与支付渠道重组同时出现，说明劳动报酬、消费者负担和平台收入需要共同核对。平台净利润仍须扣除完整成本，单项政策的贡献则须另作反事实识别。\href{https://www.nyc.gov/assets/dca/downloads/pdf/workers/Restaurant-Delivery-App-Data-Q1-2024.pdf}{原报告表2，第4页}

数值换算与未去重口径见新增案例核查，原始值和计算值见数值长表。

\subsection{平台重组劳动付款，接单收入须辨明资本关系}

平台技术可能同时改变等待、派单、报酬和付款时间，其结果不能只用完成配送的速度解释。直接雇佣骑手时，平台或承包企业购买劳动力，排班和报酬变化直接进入相应工资账户；劳动者独立经营时，收到的服务款还须承担车辆、能源和业务费用，不宜全部记作平台可变资本。合同名称也不足以决定关系，应结合实际组织、风险承担和控制方式核查。

纽约报告在总收入之外列出了付款来源：2023至2024年第一季度，周均应用报酬由约1000万美元增至2490万美元，小费由约990万美元降至340万美元，报告工时则由169.9万降至146.8万小时。\href{https://www.nyc.gov/assets/dca/downloads/pdf/workers/Restaurant-Delivery-App-Data-Q1-2024.pdf}{DCWP原报告表1，第3页}。这说明总工时下降可以与应用支付大幅增加并存，也说明消费者小费与应用支付不能混成同一笔资本支出。资料未扣经营成本、未去重，且不识别全部合同关系，仍不能直接将2490万美元认定为平台生产性可变资本。

结算技术还涉及谁垫款。平台先收到消费者付款再向服务者结算，与雇主先支付工资再等待客户回款，资金关系不同。更快向劳动者付款有利于其生活现金，却可能增加付款方的短期资金需要；压缩客户回款周期则可能减少。这两项都可称结算改善，其对预付资本的方向却不同。应将订单、收款与实际付款日期逐笔连接，而不是以到账速度代替资本周转。

算法组织还能在不改变合同单价时改变劳动投入。更紧的送达时限、同时接单要求和排名惩罚，可能迫使劳动者提高劳动密度或延长为获得订单而投入的时间；合理路线和及时异常处理也可能减少无效劳动。受雇关系中的工资支出须同这些规则联读，独立经营关系则分别分析收费、风险与实际控制。派单技术提高可计算性，为资本调整劳动使用提供条件，具体是否增加负担仍要由规则版本和劳动记录证明。

\section{网络节省协调成本，替代选择限制收费权}

平台提供匹配，要投入软件、支付、安全和客户获取。订单和劳动者集中，匹配机会可能增多，固定投入也可能由更多交易分担，大网络因而可以确实节省成本。收费若补偿这些工作，就有相应服务。要进一步判断是否取得排他收益，还须比较服务成本、替代办法与风险，不能把全部收费都当作无功能的租金。

省下的协调成本归谁，与双方能否换渠道有关。劳动者能带走客户和信誉，客户也容易换用，平台就更需要用质量或低收费留住双方，参与者更可能分享匹配收益。交易历史若只在一个渠道可用，离开就失去主要订单，平台成本虽低，收费空间却可能扩大。有效网络可以让人愿意使用，难以转用又可能使规则制定者更主动。要解释后果，须把网络规模与转换条件一起查，不能从规模直接推出垄断。

应用多，也未必有足够替代。劳动者可以注册多个账户，大部分收入却仍来自一个；市场需求下降时，各平台还可能一起少派单。要查实际收入来自哪里、换用后拿到多少订单、客户能否带走、新人多久才接到第一批任务。专业服务中客户多元化也可能扩大选择。账户数量增加，不能单独证明竞争加强。

算法管理还超出平台交易。OECD对六国至少20名雇员机构的2024年调查获得6047份雇主方答复，所指采用包括部分或全部自动化管理软件，响应率为2\%—6\%。（Milanez等，2025，\href{https://doi.org/10.1787/287c13c4-en}{DOI}，第2节、附录A）它说明常规组织也在用软件配置劳动；雇主感知的安排改善，仍需与劳动者的完整时间、净收入及申诉权限分别观察。劳动控制的分析单位应是决定如何形成和执行，而非企业是否被称为平台。

匹配效率与退出限制依靠不同证据：前者由必要等待和合格服务量衡量，后者由可实现替代、迁移负担和规则权限衡量。二者可以并存。下一节将转换负担接到外部收入，再比较分配条件。

\section{客户和信誉难以带走，劳动者就难以退出}

劳动者能否拒绝收费或任务条件，取决于可实现的外部收入。为说明转换成本的作用，设另一种工作在没有转换成本时可提供每小时净收入 $b_0$，转换需要一次性成本 $k$，预计在那里工作 $T$ 小时。把该成本在这段期限内摊开，外部选择的简化净额为：

\[
b=b_0-\frac{k}{T}.
\]

这一简化关系用于分析转换成本对替代收入的影响，未包括完整的动态求职过程。假设外部工作每小时40元，转换要支出1600元，预计工作160小时，则这一期限的净外部选择相当于每小时30元。若转换后的工作持续更久，摊销负担会下降；若机会本来就不稳定，短期负担更大。对未来收入的估计还需要概率和等待时间，招聘广告时薪仅是潜在报酬信息，不等于已经可实现的外部收入。

$k$ 可以包括重新建立客户联系、获得信誉、学习新规则和收入中断等代价。这些代价不是都由平台故意设置：有些来自技能差异、地域和真实服务需求，有些则来自记录不互认或客户关系被封闭。判断控制关系时，需要区分这两种来源。可以跨平台使用的设备降低转换损失，难以复用的声誉或合同限制可能提高损失。已经花掉且无法收回的历史支出，不宜重复作为未来每次退出的现金成本；分析应以转换会新增或失去什么为准。

Rosenblat与Stark依据2014年12月至2015年9月的Uber司机经验，分析五个论坛约1350项档案并访谈七名司机，追踪信息与权力不对称形成间接控制的过程。（Rosenblat与Stark，2016，\href{https://ijoc.org/index.php/ijoc/article/download/4892/1739}{期刊全文}，3759—3760页）它为“规则信息影响行动”提供案例证据，样本不足以确定全部司机的收入结果或2026年的市场规则。劳动者知道自己可以拒绝一个任务，却未必知道拒绝后未来机会怎样变化；这种信息缺口会使其难以比较当前收入与外部选择。

有效选择同时需要可理解的规则和可实现的替代工作。价格、任务范围、拒单后果和复核程序有助于比较留下的回报，可靠外部收入则支持实际拒绝。账户可以注销，只确认技术上的退出，不确认生活和职业上的替代。

\section{共同收益增加，劳动所得仍可能减少}

退出机会把交易规则与收入分配连接起来。沿第十四章的谈判模型，给定劳动工时$H$和扣除外部必要成本后的共同收益$V$，劳动者可取得外部收入$bH$，平台可取得外部所得$\bar\pi$。当合作增益$S=V-bH-\bar\pi>0$时，加权纳什解为：

\[
W=bH+\eta(V-bH-\bar\pi)=bH+\eta S.
\]

$W$为劳动者所得，$V-W$为平台所得，$0<\eta<1$为劳动方的谈判权重。完整推导见第十四章；本节比较这一机制在平台场景中的含义。外部收入越可实现，劳动者越有能力拒绝不利条件；共同收益增加时，劳动者能够保留的份额还受协商条件影响。

模型中的$V$已扣除定义范围内的必要成本，不能直接填入订单营业额；$\eta$也不是平台收费比例。它们作为条件变量使用，未由本章材料估计。下表选择同一工时，分别改变共同收益、劳动权重和外部选择，检验效率增长是否足以保证劳动者增收。

\textbf{表6.3 假设例子：在固定160小时下比较分配条件}

各情景均固定 $H=160$ 小时、平台外部选择 $\bar\pi=2000$ 元。平台所得包含其外部选择和增益份额，不是会计资料中的已估计利润。

\begingroup\footnotesize\setlength{\tabcolsep}{4pt}\renewcommand{\arraystretch}{1.25}
\begin{longtable}{>{\raggedright\arraybackslash}p{\dimexpr 0.15942\textwidth-8.0pt\relax}>{\raggedright\arraybackslash}p{\dimexpr 0.11594\textwidth-8.0pt\relax}>{\raggedright\arraybackslash}p{\dimexpr 0.11594\textwidth-8.0pt\relax}>{\raggedright\arraybackslash}p{\dimexpr 0.10145\textwidth-8.0pt\relax}>{\raggedright\arraybackslash}p{\dimexpr 0.11594\textwidth-8.0pt\relax}>{\raggedright\arraybackslash}p{\dimexpr 0.11594\textwidth-8.0pt\relax}>{\raggedright\arraybackslash}p{\dimexpr 0.11594\textwidth-8.0pt\relax}>{\raggedright\arraybackslash}p{\dimexpr 0.15942\textwidth-8.0pt\relax}}
\toprule
\textbf{情景} & \textbf{共同收益 $V$} & \textbf{外部选择 $b$} & \textbf{权重 $\eta$} & \textbf{增益 $S$} & \textbf{劳动者所得 $W$} & \textbf{平台所得 $V-W$} & \textbf{$W/H$} \\
\midrule\endfirsthead
\toprule
\textbf{情景} & \textbf{共同收益 $V$} & \textbf{外部选择 $b$} & \textbf{权重 $\eta$} & \textbf{增益 $S$} & \textbf{劳动者所得 $W$} & \textbf{平台所得 $V-W$} & \textbf{$W/H$} \\
\midrule\endhead
\bottomrule\endfoot
基准 & 16000元 & 40元/小时 & 0.50 & 7600元 & 10200元 & 5800元 & 63.75元/小时 \\
效率提高、分配条件保持 & 20000元 & 40元/小时 & 0.50 & 11600元 & 12200元 & 7800元 & 76.25元/小时 \\
效率提高、权重降低 & 20000元 & 40元/小时 & 0.25 & 11600元 & 9300元 & 10700元 & 58.13元/小时 \\
效率提高、权重与外部选择降低 & 20000元 & 30元/小时 & 0.25 & 13200元 & 8100元 & 11900元 & 50.63元/小时 \\
\end{longtable}\endgroup


共同收益增加且分配条件不变时，双方增收；第三行将劳动权重降低后，劳动者所得反比基准少8.82\%；第四行再降低外部选择，所得继续下降。算例证明效率提高不足以保证增收，不证明技术必然降低$b$或$\eta$。经验解释必须找到改变这些条件的规则与行为。

比较静态把这个关系写得更清楚：$\partial W/\partial V=\eta$，$\partial W/\partial b=(1-\eta)H$，$\partial W/\partial\eta=S$，$\partial W/\partial\bar\pi=-\eta$。在最后两行固定的 $V$、$H$ 和 $\eta$ 下，$b$ 每小时减少10元，劳动者所得减少 $0.75\times160\times10=1200$ 元，恰好等于9300元与8100元的差额。记录携带、降低转换负担和组织化协商因此具有不同作用：前两者可能改善外部机会，后者可能改变增益如何分配。

本节将外部机会与谈判权重作为不同条件，尚未求解订单、收费和人员进入的共同变化。它与第十四章固定工资的竞争模块独立，不能直接拼接成同一市场均衡。

\section{数据用途与执行权限形成算法控制}

经济依赖不会自动说明每一项算法都有控制效果。需要看到数据、评分与收入或准入决定之间的连接。Kellogg、Valentine和Christin的综述把算法控制分为限制、推荐、记录、评分、替换和奖励，分别涉及行动引导、评价和纪律。（Kellogg等，2020，\href{https://doi.org/10.5465/annals.2018.0174}{DOI}、\href{https://angelechristin.com/wp-content/uploads/2020/01/Algorithms-at-Work_Annals.pdf}{作者稿}，表2—4、方法附录）这是从多种研究中形成的机制分类，而非同一群体的六个效应系数。

控制怎样发生，可以沿“记录—评价—任务机会”查下去。日志用于纠正付款，能帮助劳动者核账；同一记录用于改变排序或停派单，则会影响后续收入。劳动者为了保住任务而改变当期行为，平台规则就对劳动形成持续约束。退出成本高，这种作用可能更强。因此，要同时查数据用途与能执行什么决定，仅有数据记录还不足以说明全部控制过程。

安全底线可以让所有人都达标，相对排名却总有人在末位，即使大家都更快。订单若按排名分配，一个人加大劳动强度，也不保证获得更多机会。由此可以提出待检验的情形：平均服务更快，个人收入却可能更不稳定。现实是否发生，要查规则版本与派单记录，排名界面本身不能算出收入损失。

评价还要查责任算给了谁。客户等待可能因为道路、商家或客户自己，系统却可能把时间记在最后执行者名下。评分影响报酬，又不分原因，劳动者就可能承担自己控制不了的风险。准确记录和有效复核则可以纠正归属，让交易更可信。须查评了哪些因素、能否纠错、收益与风险由谁承担，不能只问有没有评价。

Rani等的JRC—ILO报告研究法国、意大利、印度和南非物流及医疗工作场所，采用案例、访谈与观察，发现流程简化和效率收益也伴随监控与工作质量问题。各案例访问日期及安排不同，部分接触依赖管理层。（Rani等，2024，\href{https://doi.org/10.2760/712475}{DOI}，第2节、附录1—2）这类资料支持“不同目标可以同时变化”的过程判断。管理者感知的效率改善与劳动者经历的控制增强可以同时发生，应结合双方资料分析其形成过程。

申诉程序的作用取决于复核者能否修改记录、撤销评分、恢复机会及补偿损失。人工确认原结论与有权纠正结论不是同一种安排。实际事件的处理记录，比是否设有人工客服更能说明纠错权限。

\section{时间自主与收入风险可以并存}

自主选择地点和上线时间，确实是便利，但不等于收入也由自己决定。客户晚上有需求，劳动者晚上工作，可能是主动安排，也可能因没有别的收入；家庭照护还可能限制其空闲时段。因此，要把订单需求、对这份收入的依赖和家庭安排一起查，才说得清时间自主有多少实际空间。

Wood等以六国远程接单劳动者的107次访谈为主，辅以东南亚和撒哈拉以南非洲679人的非概率调查；访谈田野在2014年9月至2015年12月之间分七个月展开。作者同时记录地点、时间和任务自主性，以及不规则时间、低收入和过度劳动压力。（Wood等，2019，\href{https://doi.org/10.1177/0950017018785616}{DOI}、\href{https://pmc.ncbi.nlm.nih.gov/articles/PMC6380453/}{开放全文}，Methods、Findings）这项研究提供的是多个工作质量维度共存的证据，目的抽样和非概率样本不提供全体平台劳动者的统一比例。

表6.1的有效时薪已计入必要时间，但仍未表达收入波动。两名劳动者平均每小时都得到40元，一人每周有稳定订单，另一人时常连续无单，家庭承担的风险不同。支付房租、照护和债务需要现金在特定时间到达；平均数相同不等于维持生活的条件相同。平台提高即时匹配可以减少空闲，却也可能通过浮动奖励把需求波动更多交给劳动者。应同时观察平均收入、波动、支付延迟和对平台收入的依赖。

车辆和设备归劳动者，可以让劳动者服务多家客户，也让其承担折旧、维修和事故费用。平台少持有资产，整套生产系统却仍须有人付这些物质成本。设备能换到其他渠道使用，可能增强外部机会；客户和计价只在一个渠道，拥有设备仍未必能定价。因此，资产使用、客户联系和规则权限都要查，才能判断有多大经营独立性。

劳动关系身份关系到保障与责任，合同名称却不能代替事实。谁定价、谁掌握客户、谁能拒绝任务、谁承担失败和风险，都须调查。使用软件本身不能证明是雇佣关系，“合作伙伴”的名称也不能免去关系调查。收入与时间可以先核算；特定辖区现行资格和权利，则须另查法规与裁判，不能由模型推出。

收入能否维持，还要看劳动者以后能否继续工作。靠长时间劳动、延迟维修或挤掉照护取得的增收，可能难以持续；通勤减少、时间可选、付款稳定，则可能让同样收入满足更多生活需要。劳动者要维持生活、形成技能、维护设备，才能持续提供服务。这些再生产条件会反过来改变未来的 $u$、$\kappa$ 和外部选择，也是效率能否维持的一部分。

\section{取得新增收入，仍可能受平台规则支配}

判断平台新增了多少机会，要与劳动者原来能做的工作比较。没有本地客户、难以通勤或须照护家庭的人，可能靠远程接单增加收入；原有稳定客户的人，可能只多接一些补充订单。因此，平台比替代工作收入高，可以与平台掌握规则同时存在。前者说明机会改善，后者还要查合作收益和工作要求由谁决定。

ILO2021执行摘要报告，在其分析的国家中，部分发展中经济体应用载客和配送劳动者的收入倾向于高于传统部门。这是一个有用的反例，表明“平台收入必定低于传统工作”不适合作为起始假设。（ILO，2021，\href{https://www.ilo.org/sites/default/files/wcmsp5/groups/public/@dgreports/@dcomm/documents/publication/wcms_771672.pdf}{执行摘要}，在地平台收入相关段落）国家范围、收入口径和调查选择仍限制其含义：它不能单独决定完整工时净额、风险或长期福利，但足以要求分析认真对待劳动者的实际替代机会。

等待多了，还可能因为进来的人比新增客户更多。平台降低准入成本，客户和劳动者都可能增加；客户增长更快，人均利用率可能提高，劳动者增长更快，则可能增加等待。看到等待增加，要先查进入人数和需求，不能先认定是算法造成。收入下降也要分账查服务价格、收费、其他成本和工时，再判断各项作用。

只看活跃者，也可能把结果看错。满意度提高，可能是条件改善，也可能是不满意者退出了；高收入者更容易一直被记录，低收入的短期进入者则可能很快消失。因此，要纳入退出、停用和迁移，观察不到也须注明。企业数据库虽有很多交易，却可能漏掉平台外时间；劳动者提供记录能补经历，提供者又可能是经过选择的一群人。把资料接起来，要说明各自缺什么，数量多不能代替因果识别。

退出受限的解释也有反驳条件。任务信息充分，拒单不损害以后机会，客户与信誉能带走，替代收入可靠，收费又受竞争约束，依赖增强的判断就应减弱。合作平台和公共组织也须接受同样检验：共同所有可能改善决定权，客户不足、参与无效却仍可能使收入不稳。须比较净收入、退出机会和实际改规则的能力，不能只凭组织名称。

\section{信息披露帮助判断，迁移和共同决定扩大选择}

上述分解明确了不同治理措施所影响的变量。完整收入与时间披露改善对 $R$、$C$、$H$ 的了解；可携带记录和客户联系可能减少 $k$，提高可实现的 $b$；劳动者代表参与收费、等待支付及绩效规则，可能改变分配条件。三类措施不能互相替代。知道收费过高不等于有能力转换，能够转换也不等于有权改变共同使用的规则。

个体申诉修正付款、评分和停用，共同参与则可能修改造成重复争议的规则。两种程序相互补充：申诉揭示错误，规则修正处理共同原因。有效性应由证据取得、处理权限、时限与恢复结果衡量。

Hadwiger的ILO工作论文记录国家案例、专家访谈与既有资料，包括2021年西班牙Just Eat与工会协议中的算法委员会。（Hadwiger，2022，\href{https://doi.org/10.54394/VARD7939}{DOI}、\href{https://webapps.ilo.org/static/english/intserv/working-papers/wp080/index.html}{公开全文}，第2章相关部分）它证明劳动者代表参与算法治理存在具体制度形式，未测出一个通用于全部平台的改善幅度，也不是该协议2026年全部条款的现行证明。本章借此提出执行评价：代表提出了什么，哪些被采纳，如何执行，收入、安全或争议结果怎样改变。

共同治理还需处理成员差异。主业劳动者、补充接单者和专业服务者可能分别重视保障、时段或客户携带。代表产生、信息取得、议程提出和分歧处理的程序，决定这些需要能否进入规则；会议数量不能代替可执行的修改权。

改规则以后，还要重算费用与需求。支付必要等待可以增加劳动者收入，也可能改变平台成本和订单价格；收费降低可以吸引进入，必要服务仍须维持；开放记录还要处理客户与其他劳动者的隐私。表6.3只改变单项条件，说明一部分可能作用。制度评价须继续查其他反馈是否抵消预期改善，以及需要哪些结果资料来验证。

\section{平台靠共同协作运行，规则制定权影响劳动所得}

本章得到的组织判断是：匹配创造收益，计价与收费分配收益，账户和评价规则约束行动。自备设备、自选上线时间并不使这三类决定都归劳动者掌握。分散劳动可以通过共同规则被集中组织，控制应按具体事项和执行记录确认。

材料支持其中不同部分。纽约报送资料同时记录待命减少与收入增加，这种改善应当承认；Uber历史案例和远程接单研究，则说明信息与评级可以影响行动，时间自主也可以与收入依赖并存。这些材料不能合成一个全球平台效果，但要求把效率、所得和控制分开评价。



收入改善与组织权限因而应分别评价。第七章将这一分析用于智力生产：研究材料公开可以扩大机会，生成材料的工具和运行条件却未必同时开放。两章比较的是参与与支配的关系，并不假定平台普遍依赖生成式AI。

\section*{本章参考文献}

1. International Labour Office. (2021). \emph{World Employment and Social Outlook 2021: The role of digital labour platforms in transforming the world of work}. Geneva: ILO. PDF ISBN 9789220319413. \href{https://www.ilo.org/publications/flagship-reports/role-digital-labour-platforms-transforming-world-work}{官方出版入口}。本章实际核查执行摘要与方法附录，总报告全文访问限制见台账。
2. Milanez, A., Lemmens, A., \& Ruggiu, C. (2025). \emph{Algorithmic management in the workplace: New evidence from an OECD employer survey}. OECD Artificial Intelligence Papers, No. 31. \href{https://doi.org/10.1787/287c13c4-en}{DOI：10.1787/287c13c4-en}。
3. Kellogg, K. C., Valentine, M. A., \& Christin, A. (2020). Algorithms at work: The new contested terrain of control. \emph{Academy of Management Annals, 14}(1), 366—410. \href{https://doi.org/10.5465/annals.2018.0174}{DOI：10.5465/annals.2018.0174}。
4. Rosenblat, A., \& Stark, L. (2016). Algorithmic labor and information asymmetries: A case study of Uber’s drivers. \emph{International Journal of Communication, 10}, 3758—3784. \href{https://ijoc.org/index.php/ijoc/article/download/4892/1739}{期刊全文}。\href{https://doi.org/10.2139/ssrn.2686227}{DOI：10.2139/ssrn.2686227}对应SSRN版本，非期刊DOI。
5. Wood, A. J., Graham, M., Lehdonvirta, V., \& Hjorth, I. (2019). Good gig, bad gig: Autonomy and algorithmic control in the global gig economy. \emph{Work, Employment and Society, 33}(1), 56—75. \href{https://doi.org/10.1177/0950017018785616}{DOI：10.1177/0950017018785616}。
6. Baiocco, S., Fernández-Macías, E., Rani, U., \& Pesole, A. (2022). \emph{The algorithmic management of work and its implications in different contexts}. Background Paper No. 9, Joint EU–ILO Project “Building Partnerships on the Future of Work”. \href{https://www.ilo.org/publications/algorithmic-management-work-and-its-implications-different-contexts}{官方出版页}。
7. Rani, U., Pesole, A., \& González Vázquez, I. (2024). \emph{Algorithmic Management practices in regular workplaces: Case studies in logistics and healthcare}. Luxembourg: Publications Office of the European Union. EUR 31846 EN; JRC136063. \href{https://doi.org/10.2760/712475}{DOI：10.2760/712475}。
8. Hadwiger, F. (2022). \emph{Realizing the opportunities of the platform economy through freedom of association and collective bargaining}. ILO Working Paper 80. \href{https://doi.org/10.54394/VARD7939}{DOI：10.54394/VARD7939}。

9. New York City Department of Consumer and Worker Protection. (2024). \emph{Restaurant delivery app data: January–March 2024}. Q1 2024 report, July 2024 corrected version; Tables 1–3 and Technical Note. \href{https://www.nyc.gov/assets/dca/downloads/pdf/workers/Restaurant-Delivery-App-Data-Q1-2024.pdf}{官方报告}。观察期与出版、纠错时间分别保留。访问日期：2026-10-08。

原有资料的公开稿、方法定位与版本限制见本章核查台账，原核查访问日期为2026年10月7日；本轮新增纽约案例访问日期为2026年10月8日。

\chapter{AI与智力劳动：任务计算化与研究工具使用权}\label{ch:7}

模型复用使部分任务的新增供给可以依靠追加计算，而不必同比增加专家处理工时。本章区分两项问题：完整工作中的人员瓶颈是否减轻，运行条件的权利配置是否限制生产者选择。前者比较同等质量的劳动与计算投入，后者比较许可、容量和替代条件。资源需求增加并不直接证明私人控制。

技术报告提供任务能力及投入资料，许可提供允许行为，成本算例比较部署与迁移，竞争资料描述不同市场环节。下文按这些材料的职能推进，不把下载量、训练费用或容量集中度当作控制变化的直接测量。

\section{部分任务计算化尚不足以证明完整职业替代}

本文将“转向资源依赖”限定为一种投入结构变化：在质量和任务边界可比时，新增合格产出主要靠追加计算及配套投入取得，专家工时不再同比增加。资源在此指设备、能源、数据和运行设施；若把人才和知识也全部纳入，“所有工作依赖资源”就失去比较意义。检验须覆盖查证、返工、沟通和责任承担，基准成绩仅提供局部线索。

推理模型提供了减少逐步人工示范依赖的证据。Guo等的DeepSeek-R1-Zero在预训练模型基础上，以答案反馈进行强化学习，AIME 2024平均单次作答成功率由15.6\%升至77.9\%。这一步未先做监督微调，但预训练知识、题目及判分规则仍是必要条件。（Guo等，2025，\href{https://www.nature.com/articles/s41586-025-09422-z}{正式论文}，DeepSeek-R1-Zero节及图1）可验证任务中的部分解题策略，可以通过训练和搜索形成，不必全部由专家逐步传授。

推理阶段也可以用计算投入改善结果。Snell等在数学任务中比较搜索与答案修订；在较小模型已有一定成功率的题目上，增加推理计算可使其超过参数量约14倍的大模型，比较保持总计算量可比。最难题组则未获得有意义的进展。（Snell等，2024，\href{https://arxiv.org/html/2408.03314v1}{固定v1}，§5.3、§7）由此可见，计算是可以调配的解题投入，但效果受基础能力、验证器与任务难度约束，不能把增加计算等同于必然获得正确答案。

这种变化已延伸到部分研发活动。AlphaEvolve把候选程序生成、执行和评价连接起来，开发团队报告，其发现的调度规则投入运行后平均释放约0.7\%的Google全体数据中心计算资源。（Novikov等，2025，\href{https://arxiv.org/html/2506.13131v1}{固定v1白皮书}，§3.3.1）这项开发者报告支持算法搜索能够形成可用改进；任务定义、模拟环境、验证和部署仍由研发体系提供，不能据此声称整个科研职业已被自动化。

\textbf{OpenAI公开七百余篇数学稿件，使部分研究论证的计算化生成成为近期的重要案例。} 2026年8月1日，OpenAI公布十项数学与理论计算机科学成果，称数学论证由内部模型生成，人员参与稿件准备，并公开Lean形式化证书。这是开发者对研究过程的披露，本文没有逐项重跑证书或独立判定其原创性。（OpenAI，2026a，\href{https://openai.com/index/ten-advances-in-mathematics/}{发布说明}）它提供的新增证据是：模型应用已从求解既有试题延伸至提出研究论证，研究过程中的生成、编排与验证开始采用不同的人机分工。

2026年10月6日，OpenAI进一步公开内部模型生成的七百余篇数学稿件及证明材料。按10月8日读取的仓库目录和撤回记录核算，首发规模为722篇；10月7日撤回3篇后，现存719篇，分为372组相关成果。每组可以包含主结果、配套论证、推论或不同证明，因此稿件数不能换算成独立解决的开放问题数。仓库披露模型接受约4,000个问题的评估，并从输出中组织、筛选这批材料；这一流程也不允许以稿件数除以问题数计算解题成功率。（OpenAI，2026b，\href{https://github.com/openai/math}{仓库说明}；2026c，\href{https://github.com/openai/math/blob/main/history.md}{修订记录}）这里可以确认的是研究材料的批量生成与公开，全部成果的正确性和学术价值仍须分别核验。

计算投入与核验状态使这次事件具有更具体的经济含义。官方称，平均每项结果使用约相当于三小时ChatGPT Pro思考的计算量，同时公开十份推理摘要；这一等效口径不是研究者工时、真实历时或完整研发成本。（OpenAI，2026d，\href{https://openai.com/index/sharing-ai-progress-in-mathematics/}{10月6日发布说明}）10月7日的仓库记录将已有形式化的顶层结果口径报告为300/719，约42\%，并说明一个符号错误使三篇相互依赖的稿件被撤回，另有14篇接受证明修补等修订。（OpenAI，2026c，\href{https://github.com/openai/math/blob/main/history.md}{修订记录}）这揭示了供给变化的两个环节：模型计算能够集中产生候选研究结果，合格知识的形成仍需要证明检查、依赖核查与错误修复。形式化可以把部分核验也转为计算，但未覆盖的论证及问题表述仍需相应审查。三篇稿件因共同依赖而一并撤回，还说明批量生成的数量不能替代对证明链的检验；错误是否得到识别、修复和公开，同样属于知识生产能力。

研究治理也在同期发生变化。OpenAI于9月21日宣布与数学家建立联系；独立的数学与人工智能咨询组随后发布建议，并在10月6日声明，咨询不代表认可这批成果或其取得过程，研究者应能提出自己的问题，取得强大工具和足够计算资源。咨询组还明确表示，它在AI企业内部没有决策权。（OpenAI，2026e，\href{https://openai.com/index/advisory-group-on-mathematics-and-ai/}{成立说明}；AGMAI，2026，\href{https://agmai.org/}{独立性与发布声明}）本次公开可以确认的权利变化有明确范围：\textbf{研究者取得了阅读和检查已发布材料的机会，尚未因此取得调用同一内部模型、研究自己所选问题的条件，也没有通过咨询取得企业决策权。} 成果分享与研究工具的使用权由此分开；这比稿件数量更直接说明知识生产中的权限配置。

上述材料支持部分任务的计算化，普遍职业转化还需要三类证据：能力能够覆盖完整工作，新增合格产出主要依靠追加计算，完整费用能够持续承担。跨职业外推还须覆盖不同任务结构，不能用数学解题代表全部智力劳动。数学发布未提供贯通上述环节的劳动与成本资料，故不能据此证明普遍转化。

模型将部分既有知识转为可复用的工具，复制工具与持续交付服务却不同。后者仍需要设备、电力、存储和维护，算法改进又可以减少每项投入。任务计算化说明生产方式发生变化，不预先决定单位能耗上升或总能耗下降。

运行条件的控制影响供给的方式，取决于准入、改价、限量及版本权限。公共设施、开放模型和可迁移服务可能扩大选择，缺少合格替代而提供者又能单方改约，才构成具体入口依赖。下文分别比较许可、部署与迁移，以识别依赖产生在哪一环节。



\section{开放材料增加可做的工作，复现仍需投入}

“开放”应说明开放什么、允许做什么、由谁维持。OSI《开源AI定义》1.0列出使用、研究、修改和分享的自由，并要求参数、代码及足够的数据说明。说明可能涉及不能直接分享的数据，定义并未要求所有原始材料在任何情况下公开。（OSI，2024，\href{https://opensource.org/ai/open-source-ai-definition}{定义原文}）这是组织的规范标准，能帮助区分发布安排，不能自行证明市场竞争或社会福利。

上述数学发布还提供了区分成果开放与生产能力开放的具体案例。10月6日公开了稿件、部分证明代码和推理摘要，但产生结果的内部模型在该次发布时尚未开放；官方表示正在推进模型发布。（OpenAI，2026d，\href{https://openai.com/index/sharing-ai-progress-in-mathematics/}{发布说明}）社会可以取得和核查这些成果，却尚不能仅凭这些材料调用同一模型研究自选问题。公开知识扩大了后续研究的起点，模型的接入、运行与调配权则决定谁能继续生产同类成果。两种开放的经济作用应分别评价。

OLMo固定首版提供训练与评估代码、权重、中间检查点及Dolma数据，代码和权重采用Apache 2.0许可，部分日志仍列为后续计划。Llama 3.1社区许可则规定署名、使用政策，并要求在该版本发布时，使用者或关联主体产品、服务的前一个自然月活跃用户超过7亿时另行申请。（Groeneveld等，2024，\href{https://arxiv.org/html/2402.00838v1}{固定v1}，§5—6；Meta，2024，\href{https://raw.githubusercontent.com/meta-llama/llama-models/main/models/llama3_1/LICENSE}{3.1许可}，§1.b、§2）两个版本都提供具体的可用材料，但其范围和使用条件不能由“可下载”一个标签概括。

\textbf{表7.1 发布对象如何改变行动，以及仍需什么投入}

\begingroup\footnotesize\setlength{\tabcolsep}{4pt}\renewcommand{\arraystretch}{1.25}
\begin{longtable}{>{\raggedright\arraybackslash}p{\dimexpr 0.16364\textwidth-8.0pt\relax}>{\raggedright\arraybackslash}p{\dimexpr 0.23636\textwidth-8.0pt\relax}>{\raggedright\arraybackslash}p{\dimexpr 0.27273\textwidth-8.0pt\relax}>{\raggedright\arraybackslash}p{\dimexpr 0.32727\textwidth-8.0pt\relax}}
\toprule
\textbf{取得的对象} & \textbf{增加的具体行动} & \textbf{仍需完成的环节} & \textbf{主要经济作用} \\
\midrule\endfirsthead
\toprule
\textbf{取得的对象} & \textbf{增加的具体行动} & \textbf{仍需完成的环节} & \textbf{主要经济作用} \\
\midrule\endhead
\bottomrule\endfoot
论文与算法说明 & 理解设计、建立替代实现 & 代码、数据、试验和验证 & 减少独立摸索，不直接提供运行容量 \\
可用权重与许可 & 运行、测试、在允许范围内适配 & 硬件、维护、任务评测和接口 & 提供替代接口的基础，降低部分退出障碍 \\
代码、数据说明和检查点 & 追踪过程、重做部分试验 & 可用材料、预算、人员与可比设置 & 降低复现成本，改善原因核查 \\
稳定托管接口 & 低固定投入取得日常服务 & 持续付款、接入及质量责任 & 扩大下游使用，同时保留提供者依赖 \\
迁移材料与维护权 & 带走记录和工作流、保留版本 & 新环境适配、验证与持续维护 & 降低关系转换成本，提高实际选择 \\
\end{longtable}\endgroup


表7.1按行动区分开放范围。运行权重可以支持行为检查，却不等于重训或核实全部来源；独立部署也不等于独立开发前沿模型。开放效果应由增加的可行行动及其费用评价，不宜用一个缺乏依据的百分比概括。

Kapoor等把开放基础模型理解为广泛可得的权重，讨论本地运行、修改与难以撤回等特性，并要求以既有工具或封闭服务为比较基线评价边际风险。（Kapoor等，2024，\href{https://arxiv.org/html/2403.07918v1}{作者v1}，§2—5）收益评价同样需要明确基线，区分开放对既有收费接口的替代、对新进入者的支持，以及对已有大型使用者的便利。这些结果可能分别发生。保留比较对象，才能既承认开放的真实价值，又避免把许可安排提前解释成所有权改革。

\section{数据能否取得和带走，影响进入与退出}

原始材料可读、可下载、可用于训练和可再发布，是不同权利。公共研究、个人作品、商业数据库和业务记录也有不同的取得条件。本章讨论这些安排对进入和控制的经济作用，不对各辖区的当前版权争议作统一合法性判断。具体行为的法律结论需要按材料、合同、辖区和日期核对；材料的网络公开状态不足以确定相关行为的合法性。

取得文本和记录，只是开始；还要去重、整理、标注，让它们能用于任务。整理有真实劳动成本，独有业务记录可能提供别处没有的信息，持续客户反馈也可能帮助改进服务。资料多，不保证效果好，相关性、噪声和评测都会改变作用。独有材料若关系到任务能否合格，就可能形成优势；公开资料和可复用工具若已够用，它就未必是进入的关口。须看实际任务，不能只比数据规模。

这些材料怎样授权，直接影响采用者先花多少钱、以后能否换用。合法许可允许重复使用时，多家开发者可以共用整理成果；排他合同则可能迫使后来者另找或另做资料。处理记录能帮助查错，索引和标注不能带走，又会增加转换工作。收费中有多少补偿整理劳动、有多少来自无法替代的入口，须比较实际费用、替代来源和合同条件，不能见到收费便作同一种判断。

共同知识做成模型，仍有人承担工程、试错和维护。因而要同时查几方关系：材料贡献者怎样授权、怎样取得报酬，整理者的劳动怎样补偿，开发者能不能取得适用资料，使用者能不能保存查错与退出所需的记录。开放许可、限定用途访问和共享处理工具，各能改变其中一部分。知识来自社会协作，说明须讨论贡献与权利，却不能直接算出每种服务应该怎样收费。

判断数据控制是否限制进入，要看两件具体事：取得相近资料需要多少钱、多久，换用以后还能不能继续处理。独有资料明显影响质量而替代不可得，控制资料就能限制进入；公开资料与工具已经够用，这种限制便较弱。这里说的是条件怎样起作用。现有材料没有估计数据集中度增加以后，市场权力平均增加多少，不能把这一缺口用概括代替。

\section{使用已有模型，可以省去重复研发}

开发模型要付研究、训练和试错的钱；适配模型要把工具接到具体任务；使用模型则可以直接购买服务。三者所做的工作不同。开放权重让采用者省去重复训练的一部分投入，低价接口让没有维护团队的组织也能使用。它们都可以增加下游机会，但使用者取得工具，并没有因此取得开发工具的全部能力。

DeepSeek-V3固定首版报告给出正式训练278.8万H800 GPU小时，按每小时2美元假定折算557.6万美元，并明确排除前期研究和消融试验。其结论还讨论小团队实现高效部署的资源负担。（DeepSeek-AI，2024，\href{https://arxiv.org/pdf/2412.19437v1}{固定v1}，表1、§6）$2{,}788{,}000\times2=5{,}576{,}000$美元可直接复算，但这只验证报告的折算，不构成企业总研发成本审计。已发布权重的使用者，也不能把这个训练金额再次视为每次部署都要支付的费用。

这项结果说明，同样的设备可以因工程改进做更多工作。少做无效计算、改进通信、调整结构，都可能降低研发所需的投入，前沿门槛并非固定不变。新进入者能够创新，已经反驳了“只有原有巨头能够创新”的绝对判断。但门槛是否持续下降，还要查方法是否传播、硬件能否适配、人员能否取得、后续收入能否维持服务。正式训练费用较低，尚不能说明这些条件都已具备。

下游选择比较的是同质量服务的完整费用。免费权重仍需要评测、设备和维护，按次接口则可能减少初始投入。任务少或波动大时，托管有固定成本优势；任务稳定且具备维护能力时，自建更容易摊销。保密、停机与持续运行要求还会改变选择，下一节用阈值分离费用条件。

\section{任务量与固定投入决定自行部署的费用优势}

设某组织在一年内需要完成 $q$ 次同质量、同服务承诺的合格任务。接口方案总成本为 $TC_C=F_C+\rho q$，本地部署为 $TC_L=F_L+v_Lq$。$F_C$ 是接口接入与评测等当期固定成本，$F_L$ 是本地方案当期归属的部署固定成本，$\rho$ 是按合格任务计的接口成本，$v_L$ 是本地可变成本。共同且不随方案变化的成本可以同时省略；方案特有的人员、监测与故障处理则要分别计入。

本地与接口的差额是：

\[
TC_L-TC_C=(F_L-F_C)-(\rho-v_L)q.
\]

如果本地固定成本较高而可变成本较低，即 $F_L>F_C$ 且 $\rho>v_L$，本地方案的成本优势要求：

\[
q\ge q^{\dagger}=\frac{F_L-F_C}{\rho-v_L}.
\]

分子是本地多花的固定投入，分母是每项服务省下的可变成本，所得单位为每年任务数。每项只省一点，就要在相应期间做很多任务才能摊平；接口价格降到不高于本地可变成本，而本地固定投入仍更多，扩大任务量便不能使本地反超。若本地两项成本都更低，就无须等到大规模；固定少、可变多，则可能先有利、任务增加以后又失去优势。这些差别都来自费用结构。

\textbf{表7.2 假设例子：相同质量服务下的部署阈值}

设 $F_C=20000$ 元/年，$F_L=120000$ 元/年，$\rho=.08$ 元/合格任务，$v_L=.03$ 元/合格任务。它们是教学参数，不是厂商报价。阈值为 $(120000-20000)/(.08-.03)=200$ 万次/年。

\begingroup\footnotesize\setlength{\tabcolsep}{4pt}\renewcommand{\arraystretch}{1.25}
\begin{longtable}{>{\raggedright\arraybackslash}p{\dimexpr 0.20588\textwidth-8.0pt\relax}>{\raggedright\arraybackslash}p{\dimexpr 0.26471\textwidth-8.0pt\relax}>{\raggedright\arraybackslash}p{\dimexpr 0.26471\textwidth-8.0pt\relax}>{\raggedright\arraybackslash}p{\dimexpr 0.26471\textwidth-8.0pt\relax}}
\toprule
\textbf{年合格任务量} & \textbf{接口总成本} & \textbf{本地总成本} & \textbf{本地相对接口} \\
\midrule\endfirsthead
\toprule
\textbf{年合格任务量} & \textbf{接口总成本} & \textbf{本地总成本} & \textbf{本地相对接口} \\
\midrule\endhead
\bottomrule\endfoot
50万次 & 60000元 & 135000元 & 多75000元 \\
200万次 & 180000元 & 180000元 & 相等 \\
500万次 & 420000元 & 270000元 & 少150000元 \\
\end{longtable}\endgroup


这一计算说明，开放权重可以扩大选择范围，但独立运行是否经济可行仍取决于任务规模。接口若降至每合格任务.04元，其余条件保持，阈值升至1000万次。组织继续调用可能是较低全成本的结果，持续使用接口不足以否定开放的作用；开放保留的替代路径还可能约束接口价格。实际选择还取决于替代方案能否运行，以及其成本是否适合组织的任务规模。

先把质量单位统一，才能比较阈值。如果按尝试收费，一次尝试合格的比例为 $\gamma$，在独立重复且无额外复杂费用的简化条件下，每项合格任务的调用费用为 $\rho_{raw}/\gamma$。失败带来的查证、重做和责任费用还须另算。答案短些、响应快些或评分高些，不等于能交付相同服务。自行部署还须满足峰时需求，全年平均有空闲，不能证明忙时也够用。因此，API标价只是比较的一部分，还须计入达到同样质量和承诺所付的费用。

有成本优势，也未必立即做得成。预期500万次任务，不等于已有500万次订单；实际只有50万次，本地固定投入就难按原计划摊开。分期建设、共享设施和试用可以少作一次性承诺，却可能增加以后收费。没有技术替代、暂时筹不到资金、有替代但接口更合算，是三种不同情形。分别查清，才能判断需要新工具、融资支持，还是保留当前选择。

\section{上游投入集中，下游仍可能降价和进入}

模型开发者提供工具，采用企业再向自己的客户交付产品。上游少数企业供给，不意味着下游也只能剩少数企业；模型多，也不保证用户容易换掉办公软件内嵌的助手。须分别查模型的质量与价格、云的容量与适配、应用取得客户与迁移的条件，再看下游产品价格、进入和利润。不同环节有不同客户和替代，不能用一个集中度概括。

OECD2026年《Artificial Intelligence markets》同时记录模型创新和价格下降，以及硬件、云、数据与分发的结构风险。图2衡量云端模型质量调整价格，表1则按价值链讨论风险。（OECD，2026a，\href{https://doi.org/10.1787/d531d73f-en}{报告}）价格竞争的证据对“所有层都已经失去竞争”构成约束，但单次调用降价并未测得整个业务迁移成本。资本集中、价格变化和有效选择，分别反映不同经济条件，需要结合解释。

另一个OECD报告使用法国、葡萄牙ICT调查2021/2023波次连接2011—2022年至少10名雇员企业的资料，主要研究非生成式AI。表4.1控制规格样本分别14895和8144，未发现AI采用与市场权力代理指标提高之间存在系统性的平均关联，并明确缺少因果识别。（OECD，2026b，\href{https://doi.org/10.1787/6f88a1ea-en}{报告}，数据说明、表4.1—4.4）该报告虽于2026年出版，但不构成2026年生成式AI的观察结果，其技术范围及数据时期仍受原样本限定。它直接约束“任何企业使用AI都提高市场权力”，而非直接检验前沿开发和算力入口是否集中。

两份报告支持分层分析：模型价格竞争、下游采用效果和上游入口条件具有不同对象。廉价服务可以帮助下游进入，同时保留对上游的依赖。沿第十四章成本式$c=w/a+m+\rho$，接入费影响采用者保留的节约；是否形成控制，还需合同和替代资料。



\section{调用便宜，也可能换用困难}

CMA2024年基础模型技术更新提出控制关键投入、利用既有分发地位、通过伙伴关系延伸市场力量三类风险。它来自市场材料、利益相关方意见和工作坊，是监管评估，并非对每项合作的违法或因果认定。（CMA，2024，\href{https://assets.publishing.service.gov.uk/media/661e5a4c7469198185bd3d62/AI_Foundation_Models_technical_update_report.pdf}{技术更新}，§1.8—1.13、第五章）这些风险分别影响采用过程中的条件：优先容量影响谁能形成服务，适配工具影响迁移，默认入口影响谁取得客户。

合作也会解决实际困难。长期容量合同可以减少供给不稳，云工具可以少做维护，软件内嵌可以少花搜索时间。须继续查这些便利是否一定要附带排他条件，别的提供者能不能交付相近功能。实际换用时，文件、权限、检索索引和评测记录能否导出，旧信息能否恢复，人员要重学多久，服务会中断多久，都影响选择。带走文件副本，并不保证带走原来的业务能力。

为单独分析这一机制，假设下一年需完成固定的 $q$ 项任务，两个提供者功能与非转换成本相同，切换新增成本为 $K$。现提供者每项比替代者贵 $\delta$，继续使用与转换的成本比较取决于 $\delta q$ 和 $K$。如果 $\delta q<K$，一次较低替代价格仍不足以支持当期切换。设 $K=60000$ 元，$q=200$ 万次，则每项价差不超过.03元时，单看该年的价格收益难以覆盖转换成本。这个假设例子不计未来年份、风险和价格策略，只显示固定转换负担怎样形成收费空间。

$K$ 只算换用新增的投入和损失。过去已经花掉、换不换都不能收回的钱，不能再算作今后的退出费用。记录和接口能带走，便可少做转换工作、降低 $K$；自行保存版本与测试能力，也有助于核实替代。但多选几家供应商，不一定免除共同故障：它们若共用同一不可替代的芯片或电力节点，仍可能一起受限。因此，要分别查能否换商、能否适配，以及故障时能否继续供给。

低价服务与长期选择分别评价。若组织保有必要材料、能够验证替代且付得起转换费用，托管可以成为可调整的选择；若这些条件不足，当期低价也可能伴随依赖。控制的判断应观察提供者改约后使用者能否有效回应。

\section{融资与实体容量分别限制服务扩张}

下表分别列资金与电力，因为有钱投入和有能力持续运行，还隔着建设与交付。汇总金额不能换算成控制程度，未来情景也不能写成已经发生的增长。先保留原单位、统计范围和包含关系，再讨论它们可能限制哪一环。

\textbf{表7.3 既有统计怎样进入成本和能力分析}

\begingroup\footnotesize\setlength{\tabcolsep}{4pt}\renewcommand{\arraystretch}{1.25}
\begin{longtable}{>{\raggedright\arraybackslash}p{\dimexpr 0.26667\textwidth-8.0pt\relax}>{\raggedright\arraybackslash}p{\dimexpr 0.25714\textwidth-8.0pt\relax}>{\raggedright\arraybackslash}p{\dimexpr 0.25714\textwidth-8.0pt\relax}>{\raggedright\arraybackslash}p{\dimexpr 0.21905\textwidth-8.0pt\relax}}
\toprule
\textbf{原来源及范围} & \textbf{原始数字} & \textbf{本章复算或解释} & \textbf{对应的经济环节} \\
\midrule\endfirsthead
\toprule
\textbf{原来源及范围} & \textbf{原始数字} & \textbf{本章复算或解释} & \textbf{对应的经济环节} \\
\midrule\endhead
\bottomrule\endfoot
DeepSeek-V3固定2024首版，正式训练 & 278.8万H800 GPU小时，假设2美元/小时 & 557.6万美元，仅该训练折算，排除前期研究与消融 & 上游研发中的一部分投入；不是每个下游用户的部署费 \\
AI Index2026，2025全球企业AI投资交易 & 总额5816.9亿美元，私人3446.6亿美元 & 私人约占59.25\%；其余交易类别差额2370.3亿美元，私人不能再次加到总额 & 融资及交易活动，不是新增设备或研发资源总量 \\
IEA2026，全数据中心2025估计与2030情景 & 485TWh到约950TWh，2030约全球用电3\% & 情景水平约增加95.88\%；不是纯AI的观测增长 & 持续用电及实体扩容需求，限制实际服务规模 \\
\end{longtable}\endgroup


AI Index第四章图4.2.1—4.2.3的企业投资包括私人投资、并购、少数股权和上市交易，生成式AI私人投资又是私人投资子集；主要使用Quid交易资料，私人投资口径覆盖2013年以来融资超过150万美元的AI/ML公司，不能完整代表政府支持资金。（Stanford HAI，2026，\href{https://hai.stanford.edu/assets/files/ai_index_report_2026_chapter_4_economy.pdf}{经济章节}）表中差额只是总额减私人部分，未识别成任何一种资产。融资可以支付人员、购设备、留现金或购买企业；并购金额尤其不等于新增生产能力。资金流、实际资本形成、研发和收入要分别观察。

IEA2026年4月16日报告更新全部数据中心用电情景，并称AI重点数据中心用电在2025—2030年约增长三倍。485与950TWh属于全数据中心口径，不能通过差额465TWh直接取得AI新增用电。（IEA，2026，\href{https://www.iea.org/reports/key-questions-on-energy-and-ai/executive-summary}{执行摘要}）三倍是另一个子口径的预测，2030均非观察事实。数据中心容量的MW还包含支持设施，也不直接等于GPU计算能力。

资金转成服务，需要时间。模型可以较快更新，电源、输电、并网和冷却却须建设；钱筹到了，设备和接入也未必同时到位。富余供电、业务可调度的地区，新增负荷可能提高利用；本已紧张的节点，则可能另需扩网和峰时供给。谁付扩容费用，谁先取得容量，公共资金支持新增设施还是补贴原有容量，须逐项查项目成本和合同。全球用电比例没有记录这些权利和负担。

\section{人员工时限制减轻，算力与电力可能成为瓶颈}

设备、电力和人员须配合起来，任何一项不足，都可能限制服务。设一段时期可供推理使用的GPU工时为 $K_G$，每次合格任务所需GPU工时为 $g$；可供该业务使用的电量为 $B_E$，每次合格任务完整用电为 $e$；人员工时为 $H$，每个劳动工时所能支持的合格任务为 $a$；可实现的需求为 $D$。可交付数量的一个简化上界是：

\[
q\le\min\left\{\frac{K_G}{g},\frac{B_E}{e},aH,D\right\}.
\]

这里的 $a$ 仍是每工时产出。作事前容量比较时，要用其他投入充足条件下能维持的流程产出/人员工时；不能把事后实际 $q/H$ 填回去，再说预测已被证实。$B_E$ 是期间电量，不是瞬时功率。各分母须指同样质量、同样范围的任务。人和机器可以同时工作，时间不能机械相加。峰时并发、网络和安全还可能另加限制，使实际交付更少。因此，这不是已经估计的生产函数，只用来说明：多添一项投入，未必多交付一项产品。

假设一个服务期间的计算容量支持200万项，电力支持120万项，人员支持180万项，需求为150万项，则上界为120万项。将计算容量翻倍，在其他约束不变时，上界仍为120万。这里的各容量都是教学假设，没有使用厂商能耗参数。若某地电力足够而维护能力只有80万项，瓶颈就转到人员；免费提供更多权重也不能直接消除它。投资应针对限制交付的环节，而不是按设备数量评价进步。

当人员支持容量$aH$原为最小项，模型提高$a$可能减轻人员限制；计算或电力随后成为最小项，新增交付才更依赖相应投入。最终核验、沟通或现场判断仍是瓶颈时，追加设备不能代替这些劳动。经验检验应同时记录合格任务、专家工时、计算、核验与失败，识别瓶颈变化，而非只看模型答题数量。

单位能耗改善与总用电还需分开。对同口径任务，简化的运行电量关系为 $E=e q$，有限变化 $E_1/E_0=(e_1/e_0)(q_1/q_0)$。假设每项能耗下降20\%、合格任务增加50\%，运行电量增加20\%，而非下降20\%。若任务只增10\%，运行电量下降12\%。这些情景固定支持设施边界；新增冷却、建设和设备制造另列。所谓反弹应由任务需求与能耗的共同变化检验，反弹这一名称本身不确定总用电的变化方向。

容量怎样分配，又会影响采用与退出。长期电力和容量合同可以帮助保证服务；小组织若要排队、承受价格波动，就可能继续依赖托管商。共享设施稳定供给，可以降低独立部署的门槛；容量若已被长期包租，设施再大也未必让后来者取得服务。因此，既要看建了什么，也要看谁能申请、以什么条件使用，才能解释选择为何不同。

\subsection{AI重组脑力劳动购买，算力费用不能替代工资分析}

对知识工作企业而言，AI降低工资支出的条件是：同等质量的完整交付少用付酬工时，且释放的人员或时间确实减少了支出。固定工资下少写一份初稿、随后多接项目或多做核验，不能记为工资现金节省。招聘需求、工作时长和报酬的后续调整，才把任务性能变化转为可变资本变化。

这种调整还可能改变购买哪些劳动力。经验知识由工具提供，部分入门任务的训练和招聘要求可能下降；复杂任务的核验、客户责任和系统维护又可能增加对高技能人员的需求。前者不能直接证明工资降低，后者也不能保证原初级人员顺利转岗。第五章的客服研究没有个人工资记录，编程实验又显示完整交付效果依任务而变，现有证据不足以给出脑力劳动工资支出的统一降幅。

使用企业购买云服务，不等于直接购买云服务商工程师的劳动力。自身工资下降、云采购上升，可能包含工作向上游转移；采购价格还包含设备、电力及利润。研究生产网络时，要把模型研发、数据处理、设施维护、部署和最终核验的雇佣工资分别登记。开源贡献和用户无酬纠错则另列时间，不能凭一个假定时薪把它们全部算作购买劳动力的资本。

“脑力工作转为资源依赖”据此只能作为有待检验的趋势命题：部分任务对个别从业者经验的依赖减轻，同时对既有模型和运行设施的依赖增加。它不证明活劳动从生产网络消失，也不证明模型在马克思价值理论中创造了新的价值职能。若想进一步判断可变资本相对减少，须取得匹配的生产网络工资与资本投入资料；模型基准、数学解答数量和单次调用价格都不能代替这些账户。

AI还可能将知识节约转为更高交付定额。企业若在工资和人数不变时提高每日任务要求，工具省时被用于新增产出；若核验、纠错和客户责任同时增加，实际劳动甚至可能加重。相反，减少定额或保薪缩时可以把节约转为闲暇。这一分歧来自对劳动过程的安排权限，并不由模型性能决定。高技能人员收入上升与资本更紧密地组织其工作可以并存，所需观察因此包括工资、交付标准、核验工时和修改标准的权限。

\section{补贴扩大使用，设施和治理影响自主生产}

中国五部门“东数西算”实施意见签署于2023年12月25日、公布于12月29日，提出统一度量与交易结算、跨区域调度及收益协同，并提出中小企业算力券、运力券。（国家发展改革委等，2023，\href{https://www.ndrc.gov.cn/xxgk/zcfb/tz/202312/t20231229_1363000.html}{官方文件}，条款八、九、十八）该文件支持政策工具及目标的制度事实描述，不提供项目验收或使用者效果的证据。

公共支持改变什么，要按具体办法查。补贴使用费，让采用者当期少付钱；共享维护和设施，让组织少作重复投入；标准、迁移支持和使用者治理，则可能让组织更容易换用、更能参与改规则。它们分别改变费用、付款人和决定权。文件列有工具和目标，不能据此认定结果已经实现，还须看实际服务、负担和权限怎样变化。

第7.5的部署阈值能把区别算清。共享维护若降低本地当期固定投入，其他条件不变，$q^{\dagger}$就降低，小规模组织也较容易独立运行；只补贴接口价格，则可能多用服务，却未必更容易自行部署。两种办法都可能有用，各自解决的困难不同。公共设施还须有人维护、有办法排队、有明确故障责任。购置预算说明花了钱，能否交付仍要看运行和使用记录。

设备归公共机构，首先改变资产归属；谁分容量、谁定收费、谁担故障责任，还须看实际治理。同样大的设施，可以由少数机构长期占用，也可以让较多小组织稳定使用。因此，评价公共项目，要查合格服务、等待、完整费用、迁移和使用者权限。投资金额没有回答这些问题，公共所有的名称也不能代替具体决定权。

本章区分成果取得、工具使用和持续交付。数学发布提供批量生成与公开的案例，许可规定可采取的行动，部署和容量分析说明完整服务仍受规模、投入与运行条件限制。已支持的结论是部分任务计算化；普遍职业替代与私人控制增强仍缺直接证据。下一章比较制造中的加工、合格交付和独立经营，检验相近的条件分离。

\section*{本章参考文献}

1. OECD. 2026a. \emph{Artificial Intelligence markets: Recent developments and competition issues}. Published 10 July 2026. \href{https://doi.org/10.1787/d531d73f-en}{DOI:10.1787/d531d73f-en}.
2. OECD. 2026b. \emph{Competition in the age of AI: Initial evidence from microdata}. OECD Artificial Intelligence Papers, No.64. Published 30 July 2026. \href{https://doi.org/10.1787/6f88a1ea-en}{DOI:10.1787/6f88a1ea-en}.
3. Stanford HAI. 2026. \emph{The 2026 AI Index Report}, Chapter 4: Economy. \href{https://hai.stanford.edu/assets/files/ai_index_report_2026_chapter_4_economy.pdf}{官方章节PDF}.
4. International Energy Agency. 2026. \emph{Key Questions on Energy and AI}. Paris: IEA. Published 16 April 2026. \href{https://www.iea.org/reports/key-questions-on-energy-and-ai}{报告主页}.
5. Competition and Markets Authority. 2024. \emph{AI Foundation Models: Technical Update Report}. 16 April 2024. \href{https://assets.publishing.service.gov.uk/media/661e5a4c7469198185bd3d62/AI_Foundation_Models_technical_update_report.pdf}{官方PDF}.
6. Kapoor, S., Bommasani, R., Klyman, K., et al. 2024. \emph{On the Societal Impact of Open Foundation Models}. 本章使用arXiv固定v1作者稿。 \href{https://doi.org/10.48550/arXiv.2403.07918}{DOI:10.48550/arXiv.2403.07918}；\href{https://arxiv.org/html/2403.07918v1}{固定版本}.
7. Groeneveld, D., Beltagy, I., Walsh, P., et al. 2024. \emph{OLMo: Accelerating the Science of Language Models}. 本章使用arXiv固定v1作者稿，提交于1 February 2024。 \href{https://doi.org/10.48550/arXiv.2402.00838}{DOI:10.48550/arXiv.2402.00838}；\href{https://arxiv.org/html/2402.00838v1}{固定版本}.
8. DeepSeek-AI. 2024. \emph{DeepSeek-V3 Technical Report}. 本章使用arXiv固定v1技术报告。 \href{https://doi.org/10.48550/arXiv.2412.19437}{DOI:10.48550/arXiv.2412.19437}；\href{https://arxiv.org/pdf/2412.19437v1}{固定版本PDF}.
9. Open Source Initiative. 2024. \emph{The Open Source AI Definition—1.0}. \href{https://opensource.org/ai/open-source-ai-definition}{定义原文}.
10. Meta. 2024. \emph{Llama 3.1 Community License Agreement}. Version release date: 23 July 2024. \href{https://raw.githubusercontent.com/meta-llama/llama-models/main/models/llama3_1/LICENSE}{发布方许可文本}.
11. 国家发展改革委、国家数据局、中央网信办、工业和信息化部、国家能源局. 2023. 关于深入实施“东数西算”工程加快构建全国一体化算力网的实施意见. 发改数据〔2023〕1779号；签署2023年12月25日，公布2023年12月29日。\href{https://www.ndrc.gov.cn/xxgk/zcfb/tz/202312/t20231229_1363000.html}{官方文件}.

12. Guo, D., Yang, D., Zhang, H., et al. (2025). DeepSeek-R1 incentivizes reasoning in LLMs through reinforcement learning. \emph{Nature, 645}, 633–638. \href{https://doi.org/10.1038/s41586-025-09422-z}{DOI:10.1038/s41586-025-09422-z}；\href{https://www.nature.com/articles/s41586-025-09422-z}{正式论文}。本章引用正式稿中R1-Zero结果，不与2025年1月预印本的不同数字混用。
13. Snell, C., Lee, J., Xu, K., \& Kumar, A. (2024). \emph{Scaling LLM Test-Time Compute Optimally can be More Effective than Scaling Model Parameters}. arXiv:2408.03314v1，6 August 2024. \href{https://doi.org/10.48550/arXiv.2408.03314}{DOI}；\href{https://arxiv.org/html/2408.03314v1}{本章使用的固定v1}。未将其与改题后的ICLR 2025版本混作逐页核查。
14. Novikov, A., Vũ, N., Eisenberger, M., et al. (2025). \emph{AlphaEvolve: A coding agent for scientific and algorithmic discovery}. arXiv:2506.13131v1，16 June 2025. \href{https://doi.org/10.48550/arXiv.2506.13131}{DOI}；\href{https://arxiv.org/html/2506.13131v1}{固定v1白皮书}。系统开发团队的结果报告，不视为独立随机实验。

15. OpenAI. 2026a. \emph{Ten advances in mathematics and theoretical computer science}. 1 August 2026. \href{https://openai.com/index/ten-advances-in-mathematics/}{官方发布说明}。本章核对公开的过程披露，未独立验证十项数学证明。
16. OpenAI. 2026b. \emph{math: mathematical manuscripts and supporting proof artifacts}. \href{https://github.com/openai/math}{GitHub仓库及README}。访问日期2026年10月8日；目录为719篇稿件、372组相关成果。
17. OpenAI. 2026c. \emph{math/history.md}. 7 October 2026更新，\href{https://github.com/openai/math/blob/main/history.md}{撤回、修订与形式化记录}。本章首发722篇由当次目录719篇加已撤回3篇核算，保留版本差异。
18. OpenAI. 2026d. \emph{Sharing AI progress in mathematics}. 6 October 2026. \href{https://openai.com/index/sharing-ai-progress-in-mathematics/}{官方发布说明}。计算口径为ChatGPT Pro等效思考量，不是完整研发成本。
19. OpenAI. 2026e. \emph{Advisory Group on Mathematics and Artificial Intelligence}. 21 September 2026. \href{https://openai.com/index/advisory-group-on-mathematics-and-ai/}{成立说明}。
20. Advisory Group on Mathematics and Artificial Intelligence (AGMAI). 2026. \emph{Independence, Transparency, and Accountability; On OpenAI's Release of Mathematical Results}. 6 October 2026发布声明，\href{https://agmai.org/}{咨询组官网}。访问日期2026年10月8日。

\chapter{新制造与组织边界：工艺可以降本，独立交付另需条件}\label{ch:8}

增材制造可以减少部分产品的准备投入，数字化协调可以减少重复记录和过程损失。但加工地点、费用优势与经营自主性是不同结果。本章依次比较完整交付成本、订单和市场变化，再分析设计、质量证明与客户关系怎样限制生产者选择。

下文依次算工艺阈值，比较订单与贸易，再看现场信息怎样用于调整，以及记录能否随生产者换用。技术资料说明流程和费用，贸易与企业研究说明市场结果，权限则须按设计、验收和客户关系分别查。假设算例公开条件，不混作原研究测得的平均效果。

\section{同样的合格交付，才能比较工艺费用}

增材制造按数字模型逐层成形，传统加工也有多种办法。要比费用，先须做出同样功能、质量、寿命和交期的产品。内部流道的轻量构件可能少用装配件，普通标准件却未必需要这种设计。把不同功能的产品直接比价，就会把设计或质量差别当成工艺降本；只算打印用料，又会漏掉表面加工、热处理、检测和返工。先统一交付要求，再算完整费用，比较才有意义。

Thomas与Gilbert的2014年NIST报告整理材料、设备、构建空间、时间、能源、劳动及质量等成本项目，并指出单个零件研究对供应链费用覆盖不足。（Thomas与Gilbert，2014，\href{https://nvlpubs.nist.gov/nistpubs/specialpublications/nist.sp.1176.pdf}{全文}，§1.3、§3—4）这份历史资料用于确定成本核算项目，不提供2026年价格。新流程的经济性应由实际材料和功能确定，“打印”这一名称不足以消除聚合物与金属、普通用途零件与关键零件的成本及责任差异。

本章产量 $q$ 只算观察期内通过验收、能够交付的合格品。废品说明花了多少投入，不能与成品同算；下载设计、尝试打印和交付，也须分别记录。生产率仍为每劳动工时产出 $a=q/H$。机器在后台运行多久，不等于人工作多久；设计、准备、清理、后处理、检测和维护，却都可能需要人。完整算入这些劳动，才分得清少做了现场操作，还是把工作交给了别的岗位。机器旁边没人，不能据此认定整体劳动消失。

比较还须用同样期限，并查起点。新产品可能要买模具、作验证，既有产品却已有可用工具；共享设备的一件订单，也不应既承担整台机器购置费，又支付同一投入的全部折旧。下文固定成本指比较期内由该方案承担的投入，变动成本指多交付合格品所增加的费用。算现金流，就按时点列购置、残值和更换；算单位成本，就用一致办法分摊资本服务。两套核算不能把同一支出计两遍。

\section{固定投入与每件费用，决定批量优势}

设增材路线为A、传统路线为B。对同一种合格交付，两者简化总成本为：

\[
TC_A(q)=F_A+v_Aq,\qquad TC_B(q)=F_B+v_Bq.
\]

$F_A$、$F_B$ 为同一比较期的固定成本，$v_A$、$v_B$ 为每件合格品的变动成本。总成本差额是 $\Delta TC=F_A-F_B+(v_A-v_B)q$。令其等于零，在 $v_A\ne v_B$ 时得到：

\[
q^*=\frac{F_A-F_B}{v_B-v_A}.
\]

分子为元，分母为元/件，阈值便是件数。但只有 $q^*>0$，又没有超过设备容量、需求和线性近似的适用范围，才存在实际可比较的正产量交叉点。等式只找出哪处费用相等，哪边更便宜还要看固定与变动成本各自的方向。

增材若少作专用准备、每件却多花投入，即 $F_A<F_B$、$v_A>v_B$，小批量时便宜，超过阈值以后传统路线便宜。若增材固定投入多、每件投入少，方向便相反。两项都高，在简式中就都不利；两项都低，就始终较便宜。变动成本相同，增产只摊薄各自固定投入，不改变排序。因此，哪种工艺划算，要由实际费用决定，名称本身没有答案。

\textbf{表8.1 假设例子：低固定、高变动的增材路线}

取 $F_A=20000$ 元、$v_A=200$ 元/件，$F_B=100000$ 元、$v_B=40$ 元/件。参数仅用于推导，不是机器、模具或厂商报价。阈值为 $(20000-100000)/(40-200)=500$ 件。

\begingroup\footnotesize\setlength{\tabcolsep}{4pt}\renewcommand{\arraystretch}{1.25}
\begin{longtable}{>{\raggedright\arraybackslash}p{\dimexpr 0.14583\textwidth-8.0pt\relax}>{\raggedright\arraybackslash}p{\dimexpr 0.18750\textwidth-8.0pt\relax}>{\raggedright\arraybackslash}p{\dimexpr 0.18750\textwidth-8.0pt\relax}>{\raggedright\arraybackslash}p{\dimexpr 0.16667\textwidth-8.0pt\relax}>{\raggedright\arraybackslash}p{\dimexpr 0.16667\textwidth-8.0pt\relax}>{\raggedright\arraybackslash}p{\dimexpr 0.14583\textwidth-8.0pt\relax}}
\toprule
\textbf{合格交付量} & \textbf{增材总成本} & \textbf{传统总成本} & \textbf{增材单位成本} & \textbf{传统单位成本} & \textbf{较低成本路线} \\
\midrule\endfirsthead
\toprule
\textbf{合格交付量} & \textbf{增材总成本} & \textbf{传统总成本} & \textbf{增材单位成本} & \textbf{传统单位成本} & \textbf{较低成本路线} \\
\midrule\endhead
\bottomrule\endfoot
100件 & 40000元 & 104000元 & 400元/件 & 1040元/件 & 增材 \\
500件 & 120000元 & 120000元 & 240元/件 & 240元/件 & 相同 \\
2000件 & 420000元 & 180000元 & 210元/件 & 90元/件 & 传统 \\
\end{longtable}\endgroup


每条路线内部，增加批量都降低固定项的单位负担；两条路线之间，较低变动成本会在大批量中起更大作用。小批量下的单位成本优势不足以证明大规模生产中仍保有同样优势。需求是否稳定也重要：仅有100件订单时，扩大至2000件生产虽可摊薄单位固定成本，却会形成缺少需求支持的库存。

阈值对质量和材料费用敏感。其余参数保持，增材新增30000元固定验证支出，阈值降到312.5件，低批量优势范围缩小；若增材变动成本从200降到120元而功能不变，阈值升到1000件。对已有可用传统工具的订单，未来需要承担的 $F_B$ 可能较低，不能为证明增材优势把已经沉没的模具支出重新计入。新进入者与既有厂商面对不同起点，采用条件因此可能不同。

\section{废品与后处理，也要计入交付成本}

上节的 $v_j$ 已按合格品定义。真实资料若记录每次制造尝试的费用，需要先考虑合格率。令路线 $j$ 每次尝试费用为 $u_j$，合格比例为 $\gamma_j$，每件合格品后续费用为 $d_j$。在尝试可重复、比例稳定的期望成本近似下：

\[
v_j=\frac{u_j}{\gamma_j}+d_j,\qquad
TC_j(q)=F_j+\left(\frac{u_j}{\gamma_j}+d_j\right)q.
\]

这笔费用仍须按同样验收要求计算：为取得合格品花掉的失败投入，也要计入。质量要求改变了合格率、验证或加工时间，阈值就随之改变。失败若相互关联，返工若能修复，或每批至少要做一定数量，就须另按实际过程算；不能一律把不合格品当作互不相关、可以原样重来的尝试。

沿用表8.1的固定成本另设独立情景：增材每次尝试200元，合格率80\%，每件合格品的后处理费用50元，则 $v_A=200/.8+50=300$ 元。传统每件合格交付仍为40元，固定成本保持，阈值变成 $80000/(300-40)\approx307.69$ 件。只看成形尝试的200元会把增材低批量范围估得更大。对于整数订单，应直接比较307件与308件的总成本，实际订单由整数件构成，代数阈值只表示成本相等的理论位置。

Hernandez Korner等的FDM成本框架用ABS和PEEK两种材料进行两组模拟，显示机器与材料条件会改变费用结构。（Hernandez Korner等，2024，\href{https://pmc.ncbi.nlm.nih.gov/articles/PMC10880658/}{原论文}，Cost model、Test cases、Conclusion）它支持联合考察材料、机器及批量，并没有给全部增材产品一个统一阈值。本节200元与80\%的假设未取自该研究。

NIST零件质量合格评定项目指出内部缺陷、复杂表面与各向异性带来的测量挑战，研究后处理测量与无损检测。（NIST，2018创建、2025更新，\href{https://www.nist.gov/programs-projects/additive-manufacturing-part-qualification}{项目说明}）验证费用具有物理和质量依据，费用较高本身不足以证明存在排他安排；但同等质量下，封闭资料、重复测试与不互认也可能增加关系转换成本。前者涉及制造质量保障，后者涉及质量能力可服务的客户范围。

设备、设计、工艺与客户认可共同构成交付能力。共同检测若使一份证明得到多家客户接受，可以减少重复验证，也扩大经营选择；设备分散而验证和客户条件未变，则未必形成独立供给。空间分散不能替代交付权限的调查。

\section{能就地打印，不保证取代跨境供货}

车间加工便宜，还须把货交到客户手里。小批量、个体适配、复杂形状和停产备件，可能适合增材；运输、最低订货量、库存、缺货与延误，又会改变整笔费用。原料不到、设备停机、本地无从检测，也可能使就地生产失去优势。设计文件可以很快传过去，粉末、耗材和质量条件却未必随时到位。

可以在上节总成本之外列出两种方案各自的物流、库存及预期风险费用 $L_j(q)$，比较 $F_j+v_jq+L_j(q)$。若差额不再线性，就不必强求一个统一阈值；某些批量会因运输整批、备件保有和产能台阶产生多个区间。产品轻量化减少服役能耗时，还需要统一功能寿命、使用强度和回收边界，制造费用较高仍可能具有使用阶段收益。反过来，材料少但成形用能多，也可能削弱资源优势。

助听器贸易为生产布局提供了一个直接检验。Freund、Mulabdic与Ruta（2019年工作稿）使用1995—2016年贸易资料，以2007年为技术扩散节点；表1第4列双重差分系数为0.456（标准误0.157），对应出口相对比较基准约增加58\%。这一结果依赖平行趋势等识别条件，是贸易额效应，不能改写为产量、利润或劳动者所得。本文采用该固定工作稿，未将其数字冒充后续正式稿结果。\href{https://documents1.worldbank.org/curated/en/152701569432061451/pdf/Is-3D-Printing-a-Threat-to-Global-Trade-The-Trade-Effects-You-Didnt-Hear-About.pdf}{原稿第8—9页、表1}

助听器结果反驳了“能就地打印就会普遍取代跨境供货”的推论。按本章成本关系，单件定制也可以与批量分摊设备、检测和人员投入结合，在哪生产仍取决于完整交付费用与需求，不能只看文件传多远。这里解释可能机制，尚未分别估计各项费用的贡献。出口增加也没有证明控制集中；设计谁能改、谁定验收、客户能否更换，仍须另查。

制造业同样需要区分新增需求与订单转移。采用者交货更快，可能扩大总销量，也可能取得同行订单。最终售价、合格交付和客户需求才说明市场变化，设计下载或设备数量不等于销售收入。

工艺省下费用以后，车间能留住多少，也受采购关系影响。稳定订单让固定投入容易摊开，改设计却又要准备和验证。采购方单方变更而不补偿，便可能让供应商多担风险；双方共享预测、约定变更责任，则可以少做无效备货。这里的控制发生在订单和验收条件上，与机器归谁不同。采购价低了多少，也不能单独说明谁有决定权。

共享加工可以少付进入所需的固定投入，租到设备却不等于取得认证和客户。平台把文件派给多地车间、统一设计和验收，可以让交易可靠，也可能让车间只负责加工；车间保有可用工艺和证明、能为多家客户供货，则有更多换用和拒绝的余地。两者都可以分布在多个城市。因此，要分别记录加工在哪进行、事项由谁决定，不能把空间分散当成权力分散。

\section{信息经现场执行转为流程收益}

工业物联网采集、传输和连接设备记录，数字孪生则为特定机器或流程建立可分析的数字表示。它们有没有用，要看帮助解决了什么现场问题。记录哪台设备、多久更新、怎样核验、结果交给谁行动，都须说明。静态模型可以帮助设计，实时连接可以帮助监控；用途不同，所需连接和检查也不同。技术名称不能替代用途。

Shao围绕ISO 23247框架设计机器健康、调度与路线、虚拟调试三个情景，其中虚拟调试对象尚未与已安装实体机器实时相连。（Shao，2021，\href{https://nvlpubs.nist.gov/nistpubs/ams/NIST.AMS.400-2.pdf}{报告}，§3—4、§4.3）这些情景说明实施结构，不是三家工厂的投资回报。它们还表明，投产前的虚拟验证可能减少现场试验，不能因缺少当时的实时实体而否定该用途；也不能以仿真结果证明投产后长期运行已经可靠。

信息产生收益需要可执行的行动链：单位与时间一致，机器与材料批次对应，误差满足用途，现场具备调整资源和权限。批次错配会妨碍缺陷追踪，准确预测也不能替代备件和停机安排。模型评分与损失减少因此分别测量。

Kim、Shao与Jo的线材电弧增材制造孪生架构使用OPC UA和PMML连接数据与模型，分析异常检测与反馈。其表1比较功能，结论把进一步实际案例列为后续工作。（Kim等，2022，\href{https://tsapps.nist.gov/publication/get_pdf.cfm?pub_id=934111}{作者公开稿}，§2—4）它支持说明接口与控制路径，未提供长期缺陷率或利润改善。架构能够连通，是产生收益的必要条件之一，尚不足以确定下节的成本节省参数。

谁能使用记录，也影响生产。现场人员能看数据、标异常、改错误设置，互联就可能帮助其判断；只能接远端指令，系统又可能增加远端干预。数据在边缘处理，并不自动把决定权交给工人；用集中服务器，也不等于每件事都由远端决定。须查实际权限、作业办法和异常记录，才知道信息让谁多作了什么决定。

\section{流程省了投入，新增产能仍要有订单}

Thomas的数字孪生经济报告强调，系统复杂性、非最优决策的损失及建模投入共同决定采用是否值得，更准确和灵活的模型也可能更昂贵。（Thomas，2024b，\href{https://nvlpubs.nist.gov/nistpubs/ams/NIST.AMS.100-61.pdf}{报告}，§3—4）这一分析支持根据成本和用途，比较完整模型、简单模型与有限数据跟踪方案。它的行业潜力推算不是本章参数，也不是采用者的平均效果。

为分析需求对工程改善收益的限制，设基线交付$q$，系统使每件既有合格品实现货币成本节省$s$，额外合格交付能力为$\Delta q_{max}$，同期可实现需求为$D$。价格和质量固定时，新增交付$\Delta q=\min\{\Delta q_{max},(D-q)_+\}$，其中$(D-q)_+=\max(D-q,0)$。令新增交付单位净贡献为$\mu$，维护费用为$M$，基线增量服务费为$\rho q$，当期固定支出为$F_S$，一期净差额为：

\[
B(q,D)=sq+\mu\min\{\Delta q_{max},(D-q)_+\}-M-\rho q-F_S.
\]

$sq$只计入比较期内可实现的货币成本节省。人员保薪而将省时用于其他任务，不构成工资现金支出的减少，时间收益应另列。$\mu$已扣新增交付的全部增量费用，不再重复扣收费；停机减少形成的同批新增产品也只能计一次。$\Delta q_{max}$衡量合格交付能力，不能用机器小时直接代替。备用和可靠性可以有价值，但未货币化的收益不直接进入本式。

\textbf{表8.2 假设例子：相同技术改善在不同需求下的净差额}

取基线10000件、每件节省8元、新增能力2000件、新增净贡献30元/件、维护25000元、基线服务费1元/件、当期固定支出50000元。各情景独立比较，忽略一期内付款时序，不是实际数字孪生合同。

\begingroup\footnotesize\setlength{\tabcolsep}{4pt}\renewcommand{\arraystretch}{1.25}
\begin{longtable}{>{\raggedright\arraybackslash}p{\dimexpr 0.22000\textwidth-8.0pt\relax}>{\raggedright\arraybackslash}p{\dimexpr 0.14000\textwidth-8.0pt\relax}>{\raggedright\arraybackslash}p{\dimexpr 0.16000\textwidth-8.0pt\relax}>{\raggedright\arraybackslash}p{\dimexpr 0.16000\textwidth-8.0pt\relax}>{\raggedright\arraybackslash}p{\dimexpr 0.16000\textwidth-8.0pt\relax}>{\raggedright\arraybackslash}p{\dimexpr 0.16000\textwidth-8.0pt\relax}}
\toprule
\textbf{可实现总需求 $D$} & \textbf{实际新增交付} & \textbf{原有交付节省} & \textbf{新增交付贡献} & \textbf{维护、收费和固定支出} & \textbf{一期净差额} \\
\midrule\endfirsthead
\toprule
\textbf{可实现总需求 $D$} & \textbf{实际新增交付} & \textbf{原有交付节省} & \textbf{新增交付贡献} & \textbf{维护、收费和固定支出} & \textbf{一期净差额} \\
\midrule\endhead
\bottomrule\endfoot
10000件 & 0 & 80000元 & 0 & 85000元 & \ensuremath{-}5000元 \\
11000件 & 1000件 & 80000元 & 30000元 & 85000元 & 25000元 \\
12000件及以上 & 2000件 & 80000元 & 60000元 & 85000元 & 55000元 \\
\end{longtable}\endgroup


工程能力相同，收益方向因需求不同而变化。即使没有新订单，材料或工时节省仍存在；即使能够多生产，也未必全部卖出。若全部已交付产品的系统费率每件增加1元，基线收费和新增交付的净贡献均须调整，三行净差额分别减少10000、11000、12000元，说明流程效率可能被持续接入费用部分吸收。固定价格、需求与投入条件也是假设：行业普遍降本可能改变价格和净贡献，需要按第四章竞争机制重新比较。

长期项目应使用现金时序。若初始投入为 $F_0$，每期扣除维护与服务费后的增量现金为 $G_t$，期末残值为 $R_T$，折现率为 $r$，则 $NPV=-F_0+\sum_{t=1}^{T}G_t/(1+r)^t+R_T/(1+r)^T$。同一设备不能既全额列在初始支出，又把其全部折旧当作额外现金扣一次。模型校准、网络安全与规则更新也是运行成本；收益逐年稳定并非已证事实，方案比较应包括需求不足和费用上升的情景。

\subsection{增材制造和数字孪生改变岗位，也可能减少资本垫付}

增材制造减少了某些模具、装配和重复加工任务，却仍需要设计、材料处理、后处理、检测与认证。在相同合格交付边界内，只有减少的工资超过新增专业劳动的工资，才表现为工资支出下降。小批量中省去制模劳动有优势，大批量中新的后处理或检测负担可能抵消；因此不能把打印速度直接写成可变资本节省比例。

数字孪生还可能改变劳动安排，而非立即减员。预测维护若减少无效巡检、降低故障返工，人员可能转向预防性维护；数据清洗、模型校准和安全管理又产生新任务。NIST报告分别核算劳动、能源、设备和材料费用，提示完整成本不能只取人工一项。\href{https://nvlpubs.nist.gov/nistpubs/ams/NIST.AMS.100-61.pdf}{Thomas，2024b，成本结构及附录A}。本章净差额式中的$sq$也只能包含已实现的工资现金节省；保薪人员释放出的时间应列为能力收益。

生产和返工周期缩短，则可能减少用于工资的资本回流前需要的资金。这里要追踪从劳动投入到交付、验收、收款的全过程：设备停机少了，但客户账期延长，所需预付未必下降；更早交付并回款、工资安排不变时，周转改善才可能释放资本。这一机制解释了为什么企业可以在不减少全年工资的情况下，从流程技术中取得资金收益。

本章产业增加值与机器人图反映产出和采用的变化，没有匹配工资账户。中国和印度制造业实际产出增加，可以与产业份额或某些岗位份额下降并存；既不能以份额下降认定可变资本绝对下降，也不能以产出扩张认定所有劳动者分享了收益。下一步应追踪相同企业或生产网络的岗位支出与周期，而不是把这些宏观曲线当作工资证据。

数字孪生记录也可能被用于改变节拍和考核。预测故障能减少工人应急劳动；把预测节约全部变成更紧的排程，却可能压缩操作缓冲，使同一班组承担更多工位。增材制造减少重复加工，也可能将任务集中到少数设计、后处理和验收人员。因而，应把技术版本同排程、岗位覆盖及休息规则分别记录：工资总额稳定时，资本对劳动的使用仍可能加深，不能将保岗直接等同于劳动负担不变。

\section{采用规模、试验表现与经济收益分别衡量}

采用多少设备、机器能做多少、试验多快、支出少了多少，分别计量不同的事情。下表按成本和需求关系分别列出，先认清原数字记录了什么，再判断可以解释哪一步，不能把它们机械相加。

\textbf{表8.3 既有资料的计量对象及经济含义}

\begingroup\footnotesize\setlength{\tabcolsep}{4pt}\renewcommand{\arraystretch}{1.25}
\begin{longtable}{>{\raggedright\arraybackslash}p{\dimexpr 0.21978\textwidth-8.0pt\relax}>{\raggedright\arraybackslash}p{\dimexpr 0.29670\textwidth-8.0pt\relax}>{\raggedright\arraybackslash}p{\dimexpr 0.29670\textwidth-8.0pt\relax}>{\raggedright\arraybackslash}p{\dimexpr 0.18681\textwidth-8.0pt\relax}}
\toprule
\textbf{原资料与实际观察范围} & \textbf{保留的数字} & \textbf{可以进入的解释} & \textbf{仍需补的关键量} \\
\midrule\endfirsthead
\toprule
\textbf{原资料与实际观察范围} & \textbf{保留的数字} & \textbf{可以进入的解释} & \textbf{仍需补的关键量} \\
\midrule\endhead
\bottomrule\endfoot
NIST2016技术基础设施需求研究 & 弥补增材技术缺口的潜在年节约41亿美元、18.3\% & 共同计量与参考设施可能减少重复成本；数字是潜力估计 & 项目实际节省、实现费用及受益者 \\
Unilever2025 Tinsukia厂自述 & 包装试验时间减少84\%；每年试验2019年2次、2023年30次 & 次数增至15倍，显示特定试验活动扩大；时间与次数不是同一结果 & 试验边界、完整支出、上市需求及对照 \\
工信部2026通知的累计与名录口径 & 推动建设2.3万余工业5G专网、1260名录项目、100典型实践 & 描述基础设施与遴选范围，不能据三数建立统一增产率 & 每项目用途、工时成本、交付及退出 \\
NIST2024报告沿用的2021美国增材制品估计 & 出货27.612亿美元、增加值13.94亿美元 & 增加值约为出货50.49\%；其差额不是利润或重复的技术贡献 & 一致反事实及实际增加值变化 \\
\end{longtable}\endgroup


第一行来自Anderson2016年简报表1，访谈与需求估计支持技术缺口分析，而非已实现平均降本。（Anderson，2016，\href{https://www.nist.gov/system/files/documents/2024/12/05/NIST.EAB_.1_R_1.pdf}{原文}）NIST粉末计量实验室研究流动、铺展、粒度和层密度，正是材料到零件可重复性的一种基础设施。（NIST，2024，\href{https://www.nist.gov/laboratories/tools-instruments/additive-manufacturing-powder-metrology-laboratory}{设施说明}）设施对象具体，并未提供企业利润或议价的效果评价。

第二行是公司对包装试验场景的自述，时间减少84\%在同一可比试验的假设下表示耗时比.16，但与次数15倍之间不存在可直接相乘以推算整厂生产率的关系。（Unilever，2025，\href{https://www.unilever.com/news/news-search/2025/five-ways-unilevers-new-lighthouse-site-applies-ai-for-impact/}{案例}，第四项措施）更多试验可能用于更多方案而非更多已销售产品，需进一步考察试验结果对产品选择、实际销售与需求的影响。该来源未提供全部成本与独立对照。

工信部通知签署于2026年2月10日、网页发布于3月30日，附件为2025年实践；网络、工厂名录与典型案例不是同一分母。（工信部，2026，\href{https://www.miit.gov.cn/ztzl/rdzt/gyhlw/wjfb/art/2026/art_6b35987d001445d49487650795eac294.html}{通知}）NIST2024报告附录B表B.6.2没有更新2021估计，行业分布借用全球结构，美国份额依系统装机等假设推算。（Thomas，2024a，\href{https://tsapps.nist.gov/publication/get_pdf.cfm?pub_id=958865}{原表}，印刷页61）表中的50.49\%只是两项原估计的比值，不是净利润率，也不表示增材替代传统工艺的贡献。

\section{制造业产出、产业份额与企业扩张分别判断}

\subsection{实际产出增长可以伴随制造业份额下降}

制造业变化至少有三种量：当价金额受数量、价格和汇率共同影响，常价增加值描述产出变化，GDP份额描述制造业在经济总量中的相对位置。份额下降，不能独自证明生产能力萎缩；增加值增长，也不能独自证明劳动报酬或决定权限改善。

本文查询UNIDO国民账户2026版，选择中国、美国、德国、日本及印度，分别代表主要机器人部署市场与快速扩张的市场。这是有目的的产业比较，不是世界国家的代表性抽样。常价制造业增加值按各国2010年设为100；GDP份额由同版当价制造业增加值除以当价GDP计算。2025年的常价数值带有UNIDO估计标记，当价分子和分母尚未返回，因此两图保留不同终止年份。\href{https://stat.unido.org/unido-statistics-portal-api}{UNIDO官方API方法}；\href{https://stat.unido.org/data/table?dataset=national-accounts}{国民账户数据库}

\begin{center}\begin{tikzpicture}\begin{axis}[width=.94\textwidth,height=6.8cm,title={五国制造业常价增加值指数},ylabel={各国2010年=100},ymin=80,ymax=275,xtick={2010,2013,2016,2019,2022,2025},tick label style={font=\scriptsize},label style={font=\small},title style={font=\small},axis lines=left,grid=major,grid style={gray!15},scaled ticks=false,legend style={font=\scriptsize,draw=none,at={(0.5,-0.24)},anchor=north},legend columns=5]
\addplot+[red!65!black,mark=none,thick] coordinates {(2010,100) (2011,110.861736) (2012,121.526205) (2013,131.753338) (2014,143.814654) (2015,154.228056) (2016,163.939487) (2017,173.400927) (2018,183.958745) (2019,192.075635) (2020,196.557822) (2021,215.961118) (2022,220.754992) (2023,229.316979) (2024,242.413057)};
\addlegendentry{中国}
\addplot+[red!65!black,mark=none,thick,dashed,forget plot] coordinates {(2024,242.413057) (2025,253.351561)};
\addplot+[red!65!black,only marks,mark=o,mark size=2pt,forget plot] coordinates {(2025,253.351561)};
\addplot+[blue!65!black,mark=none,thick] coordinates {(2010,100) (2011,100.52164) (2012,99.9225287) (2013,103.090311) (2014,104.861071) (2015,106.20291) (2016,105.995374) (2017,104.350087) (2018,109.460117) (2019,110.043813) (2020,105.869153) (2021,113.752347) (2022,112.806196) (2023,113.368426) (2024,116.796364)};
\addlegendentry{美国}
\addplot+[blue!65!black,mark=none,thick,dashed,forget plot] coordinates {(2024,116.796364) (2025,118.180825)};
\addplot+[blue!65!black,only marks,mark=o,mark size=2pt,forget plot] coordinates {(2025,118.180825)};
\addplot+[green!45!black,mark=none,thick] coordinates {(2010,100) (2011,109.060137) (2012,107.1759) (2013,107.085636) (2014,112.997856) (2015,114.103576) (2016,118.673135) (2017,123.050885) (2018,124.348415) (2019,122.59957) (2020,112.828613) (2021,122.644702) (2022,123.400653) (2023,124.235602) (2024,120.613739)};
\addlegendentry{德国}
\addplot+[green!45!black,mark=none,thick,dashed,forget plot] coordinates {(2024,120.613739) (2025,118.365027)};
\addplot+[green!45!black,only marks,mark=o,mark size=2pt,forget plot] coordinates {(2025,118.365027)};
\addplot+[violet!65!black,mark=none,thick] coordinates {(2010,100) (2011,97.0209782) (2012,99.0541444) (2013,98.4812664) (2014,100.740026) (2015,103.615616) (2016,103.72413) (2017,108.013421) (2018,111.702634) (2019,109.678786) (2020,103.377505) (2021,112.719547) (2022,113.508889) (2023,114.632434) (2024,115.227551)};
\addlegendentry{日本}
\addplot+[violet!65!black,mark=none,thick,dashed,forget plot] coordinates {(2024,115.227551) (2025,117.178495)};
\addplot+[violet!65!black,only marks,mark=o,mark size=2pt,forget plot] coordinates {(2025,117.178495)};
\addplot+[orange!75!black,mark=none,thick] coordinates {(2010,100) (2011,103.125055) (2012,108.748551) (2013,114.148838) (2014,123.161653) (2015,139.245818) (2016,150.283585) (2017,161.595526) (2018,170.34036) (2019,165.272798) (2020,170.352641) (2021,187.31162) (2022,184.044537) (2023,206.686325) (2024,216.027017)};
\addlegendentry{印度}
\addplot+[orange!75!black,mark=none,thick,dashed,forget plot] coordinates {(2024,216.027017) (2025,231.026697)};
\addplot+[orange!75!black,only marks,mark=o,mark size=2pt,forget plot] coordinates {(2025,231.026697)};
\end{axis}\end{tikzpicture}\end{center}

\textbf{图8-1　五国制造业常价增加值趋势。}\label{fig:open-industry-real} 2010—2024年为该版历史序列，2025年虚线和空心点为UNIDO估计。各国分别以2010年为100，图形比较增长路径，不比较五国的绝对产出规模，也不是机器人采用的因果效应。

\begin{center}\begin{tikzpicture}\begin{axis}[width=.94\textwidth,height=6.8cm,title={五国制造业GDP份额},ylabel={当价增加值／GDP（\%）},ymin=0,ymax=36,xtick={2010,2012,2014,2016,2018,2020,2022,2024},tick label style={font=\scriptsize},label style={font=\small},title style={font=\small},axis lines=left,grid=major,grid style={gray!15},scaled ticks=false,legend style={font=\scriptsize,draw=none,at={(0.5,-0.24)},anchor=north},legend columns=5]
\addplot+[red!65!black,mark=none,thick] coordinates {(2010,31.6005411) (2011,31.5617177) (2012,31.0136801) (2013,30.1269059) (2014,29.7459797) (2015,28.3890015) (2016,27.5237555) (2017,27.5998607) (2018,27.3434229) (2019,26.1881102) (2020,25.5809825) (2021,26.6219353) (2022,25.9806695) (2023,25.0144232) (2024,24.9966588)};
\addlegendentry{中国}
\addplot+[blue!65!black,mark=none,thick] coordinates {(2010,11.9596217) (2011,12.0104627) (2012,11.9029382) (2013,11.8582127) (2014,11.7003714) (2015,11.7049055) (2016,11.2483906) (2017,10.7572253) (2018,10.9496635) (2019,10.5329155) (2020,10.0899356) (2021,10.1845222) (2022,10.2409095) (2023,10.127023) (2024,9.83249369)};
\addlegendentry{美国}
\addplot+[green!45!black,mark=none,thick] coordinates {(2010,19.4327729) (2011,20.134216) (2012,20.0571351) (2013,19.7677939) (2014,20.1171123) (2015,20.150892) (2016,20.4345908) (2017,20.2013743) (2018,19.8657585) (2019,19.3688329) (2020,18.6373282) (2021,18.6967626) (2022,18.2764141) (2023,18.3595465) (2024,17.8005974)};
\addlegendentry{德国}
\addplot+[violet!65!black,mark=none,thick] coordinates {(2010,20.7662009) (2011,19.5355543) (2012,19.6667084) (2013,19.3290317) (2014,19.5936092) (2015,20.4624704) (2016,20.2880386) (2017,20.4359461) (2018,20.621684) (2019,20.2241828) (2020,20.0515708) (2021,20.7796035) (2022,19.8229457) (2023,20.5810487) (2024,20.3945326)};
\addlegendentry{日本}
\addplot+[orange!75!black,mark=none,thick] coordinates {(2010,17.1023243) (2011,16.1393374) (2012,15.816923) (2013,15.2530227) (2014,15.0655701) (2015,15.5838546) (2016,15.1622376) (2017,15.0182388) (2018,14.8815307) (2019,13.4558076) (2020,14.1204854) (2021,14.3770289) (2022,13.1454236) (2023,13.018629) (2024,12.6085653)};
\addlegendentry{印度}
\end{axis}\end{tikzpicture}\end{center}

\textbf{图8-2　五国制造业GDP份额趋势。}\label{fig:open-industry-share} 2010—2024年，同版当价制造业增加值／当价GDP。份额包含相对价格和产业结构的变化，与图8-1的常价指数不同；2025年缺少完整当价资料，没有补齐或外推。

\textbf{表8.4　制造业实际增长与相对份额，2010—2024年。}

\begingroup\footnotesize\setlength{\tabcolsep}{4pt}\renewcommand{\arraystretch}{1.25}
\begin{longtable}{>{\raggedright\arraybackslash}p{\dimexpr .16\textwidth-8pt\relax}>{\raggedright\arraybackslash}p{\dimexpr .23\textwidth-8pt\relax}>{\raggedright\arraybackslash}p{\dimexpr .20\textwidth-8pt\relax}>{\raggedright\arraybackslash}p{\dimexpr .20\textwidth-8pt\relax}>{\raggedright\arraybackslash}p{\dimexpr .21\textwidth-8pt\relax}}
\toprule
国家 & 常价增加值累计增长（\%） & 2010年GDP份额（\%） & 2024年GDP份额（\%） & 份额变化（百分点）\\
\midrule\endfirsthead
\toprule
国家 & 常价增加值累计增长（\%） & 2010年GDP份额（\%） & 2024年GDP份额（\%） & 份额变化（百分点）\\
\midrule\endhead
\bottomrule\endfoot
中国 & 142.41 & 31.60 & 25.00 & -6.60\\
美国 & 16.80 & 11.96 & 9.83 & -2.13\\
德国 & 20.61 & 19.43 & 17.80 & -1.63\\
日本 & 15.23 & 20.77 & 20.39 & -0.37\\
印度 & 116.03 & 17.10 & 12.61 & -4.49\\
\end{longtable}\endgroup

五国在同一比较期均出现常价增加值增长与GDP份额下降。中国实际增长约142.4\%，份额由31.60\%降至25.00\%；印度实际增长约116.0\%，份额由17.10\%降至12.61\%。美国、德国和日本实际增长约16.8\%、20.6\%与15.2\%，对应份额分别下降约2.13、1.63与0.37个百分点。由此能够排除“制造业份额下降就意味着制造业产出下降”的直接推论。份额变化还受其他部门扩张和相对价格影响，不能把全部差额归因于制造效率或自动化。

产业判断由此回到具体环节。需要研究能力扩张，就查实际增加值、合格交付和产能利用；需要研究产业地位，就查同口径份额与投入产出联系；需要研究劳动受益，就接上人数、完整工时和报酬。机器人安装与制造业增长分别描述部署和产出，需求、出口、资本形成及其他技术同样可能改变后者，两条上升曲线不足以识别自动化贡献。

\subsection{安装量扩大与行业份额下降可以并存}

机器人行业结构提供另一种数量与份额的分离。按IFR2026发布演示第11页的台数，2015—2025年汽车组由97,533增至138,429台，电子电气组由91,660增至142,235台，两组均扩大。汽车份额却由约38.4\%降至22.9\%，电子电气由约36.1\%降至23.6\%，原因是其余统计组增加更快，而不是这两组的新安装量减少。\href{https://ifr.org/downloads/press_docs/Market_Presentation_WR_Press_Conference_2026.pdf}{IFR原始演示，第11页}

\begin{center}\begin{tikzpicture}\begin{axis}[width=.46\textwidth,height=6.2cm,title={A　三组安装数量均增加},xlabel={千台},xbar,bar width=7pt,xmin=0,xmax=375,ytick={0,1,2},yticklabels={汽车,电子电气,其他及未指定},y dir=reverse,tick label style={font=\scriptsize},label style={font=\small},title style={font=\small},axis lines=left,grid=major,grid style={gray!15},scaled ticks=false,legend style={font=\scriptsize,draw=none,at={(0.5,-0.24)},anchor=north},legend columns=2]
\addplot+[gray!55,fill=gray!55,draw=none] coordinates {(97.533,0) (91.66,1) (64.555,2)};
\addlegendentry{2015}
\addplot+[blue!65!black,fill=blue!65!black,draw=none] coordinates {(138.429,0) (142.235,1) (322.643,2)};
\addlegendentry{2025}
\end{axis}\end{tikzpicture}\hfill\begin{tikzpicture}\begin{axis}[width=.46\textwidth,height=6.2cm,title={B　其他统计组份额扩大},ylabel={百分比},ybar stacked,bar width=20pt,ymin=0,ymax=100,xmin=-.6,xmax=1.6,xtick={0,1},xticklabels={2015,2025},tick label style={font=\scriptsize},label style={font=\small},title style={font=\small},axis lines=left,grid=major,grid style={gray!15},scaled ticks=false,legend style={font=\scriptsize,draw=none,at={(0.5,-0.24)},anchor=north},legend columns=1]
\addplot+[blue!65!black,fill=blue!65!black,draw=none] coordinates {(0,38.4369532) (1,22.9450346)};
\addlegendentry{汽车}
\addplot+[red!65!black,fill=red!65!black,draw=none] coordinates {(0,36.1224522) (1,23.5758909)};
\addlegendentry{电子电气}
\addplot+[green!45!gray,fill=green!45!gray,draw=none] coordinates {(0,25.4405946) (1,53.4790745)};
\addlegendentry{其他及未指定}
\end{axis}\end{tikzpicture}\end{center}

\textbf{图8-3　机器人行业安装数量与份额。}\label{fig:open-industry-sectors} 台数来自同版演示，份额由各组台数重算。第三组为汽车和电子电气之外的余项，含未指定行业；不能把全部余项当作已确认的一般工业用途。

其余组由64,555增至322,643台，约为原来的五倍。这支持安装结构不再仅由汽车和电子两组主导的描述；第10页仍列出约9万台未指定行业，具体扩散到哪些工序和产业需要更细资料。更多行业使用设备，可以扩大专用集成、测试与维护需求，也可能改变劳动任务，这些结果不能仅由台数确定。

\subsection{采用工厂扩张仍须同竞争者变化比较}

WEF2025灯塔报告讨论主要来自2024年新入网的36个标杆，网络总数当时为189；同日发布说明称新增17成员，分类列13个第四次工业革命灯塔、5个可持续灯塔，其中一个已是工厂灯塔。报告案例与发布批次不是统一分母。发布把劳动生产率平均提高53\%、转换成本平均降低26\%归于最新批次，没有逐指标实际场址分母、统一基线期或加权算法，本章无法独立复算。（WEF，2025，\href{https://reports.weforum.org/docs/WEF_Global_Lighthouse_Network_2025.pdf}{报告}，第6、12—13页；\href{https://www.weforum.org/press/2025/01/global-lighthouse-network-2025-world-economic-forum-recognizes-companies-transforming-manufacturing-through-innovation/}{发布说明}）

灯塔材料说明入选工厂怎样配合技术、人员和流程。均值首先属于这些入选者，而且多项技术与组织调整同时发生。要推到所有采用者，就须补看没入选、延期和退出的项目；要推到行业劳动，还须看竞争工厂怎样回应。法国机器人研究提供了后一种比较，让“成功者之外发生了什么”不只停留在猜测。

平均数还存在变换问题。为说明而假设两家可比流程的每工时产出分别提高20\%和86\%，简单平均是53\%。固定各自基线产出和相同基线工时，工时比为1/1.20与1/1.86，平均约.6855，即下降31.45\%；直接用1/1.53得到的.6536、下降34.64\%，不是这两家的平均。这一教学例子未复原WEF样本，仅说明缺少原始分母与权重时，平均生产率增幅不足以直接推算总体就业变化或平均省时。

灯塔案例提供可实现路径及配套安排，未入选、延期和终止项目则提供失败与调整条件。前者可说明特定工厂的改进，缺少后者就不能估计总体成功率或平均回报。案例作用由其覆盖范围决定。



法国制造业机器人研究提供了相邻技术的经验比较。Acemoglu、Lelarge与Restrepo（2020）分析2010—2015年持续经营的55,390家企业，其中598家购置机器人。表1未加权模型中，采用企业总工时的对数变化系数为0.109（标准误0.020）；表2按基期工时加权后，竞争者采用暴露的系数为\ensuremath{-}0.250（标准误0.107），即暴露提高10个百分点对应总工时约减少2.5\%。这里衡量的是工时，不能把它改写为岗位人数；两表样本数也分别为55,390和55,388。（Acemoglu等，2020，第384—386页，表1第6列、表2第4列，\href{https://doi.org/10.1257/pandp.20201003}{DOI：10.1257/pandp.20201003}，\href{https://economics.mit.edu/sites/default/files/publications/Competing\%20with\%20Robots\%20-\%20Firm-Level\%20Evidence\%20from\%20F.pdf}{作者公开正式稿}）

作者使用表2同一加权模型合并自身与竞争者估计，将样本中按基期工时加权的20个百分点采用增幅联系到行业总工时约下降3.2\%；这是样本内汇总推算，未包含跨行业一般均衡调整。企业采用也不是随机安排，成长前景仍可能影响采用与随后扩张。因此，研究揭示的层次差异比一个通用就业系数更有解释价值：同一家工厂工时增加，与竞争企业劳动需求减弱可以同时发生。（同上，第386—387页及注7；采用选择讨论见第383页）

法国研究将采用者扩张与同行收缩放入同一市场，说明企业收益可能包含订单再分配。本章的增材和数字孪生模型仍需各自的需求与成本资料，不能直接借用机器人系数。工时资料也未测量产权或规则变化，不能据此确认资本集中。



\section{能改设计、保留质量证明和选客户，才有独立交付条件}

设计版本、批次、机器参数和检测结果对应起来，返修就较容易找范围，排产也较容易结合交期与设备状态。但这些能力未必能随换用保留下来。文件能导出，不等于含义也保全；缺对象编号、时间、校准和版本关系，就可能重做历史解释。工艺经验与质量信誉越难带走，换软件或换客户越可能损失原有能力。能否迁移，因而关系到生产者能否继续使用自己积累的成果。

ISO/ASTM52920:2023官方摘要把质量安排的对象放在工业增材流程和生产场所，而不只放在机器上。（ISO/ASTM，2023，\href{https://www.iso.org/standard/76911.html?browse=tc}{官方摘要}）本章只读取摘要和元数据，未据此推演完整认证步骤或费用。标准提出共同对象，可以减少沟通，也需要学习和测试；效果取决于结果能否被不同客户接受，成员是否能够负担以及参与修订。

信息能交换、质量证明能被多家客户接受、成员能参与改测试和责任规则，是不同能力。接口互通只解决前一部分。共享实验室做完检测，客户仍要求重测，便省不了全部费用；格式统一，工艺模型不能带走，也可能仍离不开原服务。判断实际选择扩大多少，应查换用以后还能保留哪些记录、证明和订单。

共同设施建成以后，还须组织日常工作。设备共用，要约定预约、维护和事故责任；工艺库共用，要有人更新、明确知识用途；检测共用，要让客户认可，并能追查责任。收费可以支付这些实际任务。须查费用和贡献能否核对、有没有替代、谁定规则，才知道协作条件。公共机构运行或采用开放许可，都是需要评价的安排，实际能用到什么还要由运行结果说明。

数字化让主导企业多看到供应商信息，也可能使供应商负担不同。共享预测并补偿变更，可以少担风险；只单向监控、不断改交期，则可能把调整交给供应商。共同使用数据以后，有没有共同作决定，须查谁提变更、谁同意、怎样执行、损失怎样分担。协作扩大，说明生产更依靠共同投入；生产者能否独立使用成果，还要看相应的占有和决定关系。

\section{工艺影响成本，交付权限影响经营选择}

本章的成本分析说明，工艺优势受合格率、后处理、批量与订单限制，数字信息须经过现场执行才能形成节约。制造地点增加可以与设计、认证和客户条件集中并存；加工能力扩大是否带来经营自主性，需要可保留记录、可互认证明和可替换客户的证据。

制造技术也可以扩大原有生产网络。助听器贸易研究反驳了普遍就地生产的预期，法国企业研究说明采用者扩张时，竞争者可能收缩。两者分别衡量贸易和企业工时，不能合成统一效果，却都要求把采用者放回市场中看。低成本可以帮助新进入，也可以让已有企业更有竞争力，哪种变化占主要地位还须由对应资料判断。

生产者能否自主，须看能否独立交付。加工设备之外，还要能改设计、验证工艺、提供质量证明、选择客户。接口互通让信息较易交换，质量互认让证明可用于更多订单，可执行的协商又影响谁改订单、谁担损失。这些权利不对称，便可能让分散车间受主导企业控制。现有工程与贸易资料说明部分条件，还没有估计这种控制变化的平均效果。

制造中的费用结果与权限结果仍须分别验证。第九至十二章将比较范围扩到电网、道路、诊疗和芯片体系，考察共同条件怎样影响交付、外部负担、服务进入与关键接口。它们分别提供不同机制证据。

\section*{本章参考文献}

UNIDO. (2026). National Accounts Database. 2026年10月11日查询2026版公开API，保留动态数据集元数据、请求体和完整JSON；常价及当价制造业增加值、当价GDP。\href{https://stat.unido.org/unido-statistics-portal-api}{官方API说明}。\par IFR. (2026). Market Presentation, World Robotics Press Conference 2026. September 24, 第10—11页。\href{https://ifr.org/downloads/press_docs/Market_Presentation_WR_Press_Conference_2026.pdf}{官方演示PDF}。

1. Thomas, D. S., \& Gilbert, S. W. (2014). \emph{Costs and Cost Effectiveness of Additive Manufacturing: A Literature Review and Discussion}. NIST SP 1176. \href{https://doi.org/10.6028/NIST.SP.1176}{DOI：10.6028/NIST.SP.1176}。
2. Hernandez Korner, M. E., Lamban, M. P., Albajez, J. A., Santolaria, J., Ng Corrales, L. C., \& Royo, J. (2024). Cost Model Framework for Pieces Additively Manufactured in Fused Deposition Modeling for Low to Medium Batches. \emph{3D Printing and Additive Manufacturing, 11}(1), 287—298. \href{https://doi.org/10.1089/3dp.2022.0044}{DOI：10.1089/3dp.2022.0044}。
3. NIST. (2018，2025年更新). \emph{Additive Manufacturing Part Qualification}. \href{https://www.nist.gov/programs-projects/additive-manufacturing-part-qualification}{官方项目说明}。
4. ISO/ASTM. (2023). \emph{ISO/ASTM 52920:2023: Additive manufacturing—Qualification principles—Requirements for industrial additive manufacturing processes and production sites}. \href{https://www.iso.org/standard/76911.html?browse=tc}{官方摘要}。
5. Anderson, G. (2016). \emph{The Economic Impact of Technology Infrastructure for Advanced Manufacturing: An Overview}. NIST Economic Analysis Briefs 1. \href{https://www.nist.gov/system/files/documents/2024/12/05/NIST.EAB_.1_R_1.pdf}{全文}。
6. NIST. (2024). \emph{Additive Manufacturing Powder Metrology Laboratory}. \href{https://www.nist.gov/laboratories/tools-instruments/additive-manufacturing-powder-metrology-laboratory}{官方设施说明}。
7. Shao, G. (2021). \emph{Use Case Scenarios for Digital Twin Implementation Based on ISO 23247}. NIST AMS 400-2. \href{https://doi.org/10.6028/NIST.AMS.400-2}{DOI：10.6028/NIST.AMS.400-2}。
8. Kim, D. B., Shao, G., \& Jo, G. (2022). A Digital Twin Implementation Architecture for Wire + Arc Additive Manufacturing Based on ISO 23247. \emph{Manufacturing Letters, 34}, 1—5. \href{https://doi.org/10.1016/j.mfglet.2022.08.008}{DOI：10.1016/j.mfglet.2022.08.008}。
9. Thomas, D. (2024b). \emph{Economics of Digital Twins: Costs, Benefits, and Economic Decision Making}. NIST AMS 100-61. \href{https://doi.org/10.6028/NIST.AMS.100-61}{DOI：10.6028/NIST.AMS.100-61}。
10. World Economic Forum. (2025). \emph{Global Lighthouse Network: The Mindset Shifts Driving Impact and Scale in Digital Transformation}. \href{https://www.weforum.org/publications/global-lighthouse-network-the-mindset-shifts-driving-impact-and-scale-in-digital-transformation/}{出版入口}。同日\href{https://www.weforum.org/press/2025/01/global-lighthouse-network-2025-world-economic-forum-recognizes-companies-transforming-manufacturing-through-innovation/}{新批次发布说明}为同一项目资料，不计独立重复验证。
11. Unilever. (2025). \emph{Five ways Unilever’s new Lighthouse site applies AI for impact}. \href{https://www.unilever.com/news/news-search/2025/five-ways-unilevers-new-lighthouse-site-applies-ai-for-impact/}{公司案例}。
12. 工业和信息化部办公厅（2026）。《关于推广5G工厂典型应用实践的通知》，工信厅信管函〔2026〕58号。\href{https://www.miit.gov.cn/ztzl/rdzt/gyhlw/wjfb/art/2026/art_6b35987d001445d49487650795eac294.html}{官方原文}。
13. Thomas, D. (2024a). \emph{Annual Report on the U.S. Manufacturing Economy: 2024}. NIST AMS 600-16. \href{https://doi.org/10.6028/NIST.AMS.600-16}{DOI：10.6028/NIST.AMS.600-16}。
14. Acemoglu, D., Lelarge, C., \& Restrepo, P. (2020). Competing with Robots: Firm-Level Evidence from France. \emph{AEA Papers and Proceedings, 110}, 383—388. \href{https://doi.org/10.1257/pandp.20201003}{DOI：10.1257/pandp.20201003}，\href{https://economics.mit.edu/sites/default/files/publications/Competing\%20with\%20Robots\%20-\%20Firm-Level\%20Evidence\%20from\%20F.pdf}{作者公开正式稿}。

访问日期统一为2026年10月7日，全文可得性、观察期及定位见本章核查台账。

新增第14项来源于2026年10月8日读取；上句访问日期适用于原有十三项来源。新增案例的原表定位及限制见《新增案例-制造与普惠》。

15. Freund, C., Mulabdic, A., \& Ruta, M. (2019). \emph{Is 3D Printing a Threat to Global Trade? The Trade Effects You Didn’t Hear About}. World Bank Policy Research Working Paper 9024, September. \href{https://documents1.worldbank.org/curated/en/152701569432061451/pdf/Is-3D-Printing-a-Threat-to-Global-Trade-The-Trade-Effects-You-Didnt-Hear-About.pdf}{本文采用的固定工作稿}。本轮核查表1及方法、结论相关段落，访问日期2026年10月8日。

\chapter{光伏与储能：设备降本以后，还要供得出、用得上}\label{ch:9}

\section{设备、资金和电网配合，才有可用电力}

光伏与储能的技术收益须通过融资、网络和调度形成有效供电。项目按售电与服务收入安排投资，系统则按地点、时段和可靠性需要组织供给，两种评价可能一致，也可能分离。本章从采用门槛推进到交付与系统效果，解释设备降本为何不直接决定用户和公众所得。

电网与调度是供电的共同生产条件。线路容量和运行时序规定能够交付多少，接入资格、合同与收费规定谁能使用及怎样分担费用。物理约束与权利安排分别识别，才能说明共同设施怎样进入生产关系。

下文先用融资算例说明，同样设备为何有的项目能投、有的不能投；再用发电时序和储能演算，说明项目回报为何不同于系统价值。装机与利用资料检查设备是否真正运行，接入实验检查用户能否取得服务。要进一步判断资本占有，还须查投资、接入和收益由谁决定。供电短缺本身没有说明原因，不能一律归于私人产权。

\section{安装成本与度电成本衡量不同条件}

IRENA于2026年发布的报告给出，2025年新投运独立大型光伏项目的全球加权平均平准化度电成本约为44美元/MWh，总安装成本约为667美元/kW。前者把建设与运行费用分摊到寿命期发电量，后者是在建设阶段投入的单位装机费用。两者对应不同问题：一座电站要筹集多少初始资金，以及按一定折现与运行条件回收这些资金需要怎样的单位收入。\href{https://www.irena.org/-/media/Files/IRENA/Agency/Publication/2026/Jul/IRENA_TEC_RPGC_in_2025_2026.pdf}{原报告第14、70页}

全球加权均值来自实际投运的项目组合，可以说明一定范围的技术经济水平，却不是每个地方都能拿到的报价。大市场权重较大，没筹到钱或没投运的项目，也不会同样进入样本。因此，要解释地区采用差异，先须查均值由哪些项目形成，再查当地设备、土地和融资条件。

成本下降，可以让同样发电少用一部分资源、少付一部分费用。但前期花得多，收入何时取得、能否持续就更重要。燃料型项目要随发电不断购料；前期投入重的项目，建成后即使少发电，融资和固定维护责任仍在。因此，技术少用燃料，也会让采用更依赖能否筹钱、收入能否回收初始投入。费用结构把设备与金融条件接了起来。

跨版历史值存在展示差异：2025版公布2024年43美元/MWh，2026版摘要将两年均写为44，图3.6仍将2024年标为43。本章据此使用2025年的水平值，不推算精确同比；版本核查保存在证据台账。\href{https://www.irena.org/Digital-Report/Renewable-Power-Generation-Costs-in-2024}{2024年成本报告}



\section{全寿命费用由交付电量分摊}

设一单位装机在建设时投入\(K\)，运行第\(t\)年的运维费用为\(O_t\)，交付电量为\(E_t\)，寿命为\(T\)年，实际折现率为\(r\)。若每MWh取得不变的实际售价\(P\)，净现值为

\[
NPV=-K+\sum_{t=1}^{T}\frac{PE_t-O_t}{(1+r)^t}.
\tag{9-1}
\]

净现值把不同年份的收入和支出折算到建设时点。这里的折现率是给定的资金回报要求；公式暂不单列税费、补贴、债务偿还和残值。令净现值为零，解出刚好回收所列成本的售价，即

\[
LCOE=\frac{K+\sum_{t=1}^{T}O_t(1+r)^{-t}}
{\sum_{t=1}^{T}E_t(1+r)^{-t}}.
\tag{9-2}
\]

式（9-2）的分子是折现费用，分母是折现电量，结果以货币/MWh计。安装成本以美元/kW计，发电成本以美元/MWh计；两者须通过设备寿命、运行费用和交付电量联系，不能直接比较。

在展示融资影响的简化情景中，设每年运维费用\(O\)和交付电量\(E\)不变。有限等比级数给出

\[
A(r,T)=\sum_{t=1}^{T}(1+r)^{-t}
=\frac{1-(1+r)^{-T}}{r},
\quad
CRF(r,T)=\frac{1}{A(r,T)}.
\tag{9-3}
\]

\(A\)是年金现值因子，\(CRF\)是资本回收因子。当\(r=0\)时分别取极限\(T\)与\(1/T\)。把式（9-3）代入式（9-2），得到

\[
LCOE=\frac{K\,CRF(r,T)+O}{E}.
\tag{9-4}
\]

其中\(K\,CRF\)是每年须回收的等额资本费用。折现率提高，以后收到同样一笔钱，折算到今天就更少；回收相同初始投入，便需要更高年度收入。变化发生在现金流现值和回收要求上，设备没有因此多发或少发电。

若以1kW装机为单位，假设容量因子为\(f\)，每年8760小时，则年交付电量为\(E=8760f/1000\)MWh。容量因子在这里已经概括了假定的年度发电条件；真实项目还须处理退化、限电、维护、天气和设备可用率。增加这些条件时应回到式（9-2），逐年列出电量和费用，而不是机械沿用常数版本。

\subsection{光伏与储能分阶段使用劳动，低电价另有再生产效应}

光伏和储能的可变资本分析须区分设备制造、工程建设与运行维护。运行环节自动监控减少某类巡检时，电站工资支出可能下降；新增电站、储能和电网建设却可能增加工程和制造劳动。即使建设期投入集中、运行期持续投入较少，也应区分建设企业购买劳动力与电站以后计提折旧的不同账户，不能只用电站在岗人数评价整个供应链。

式（9-2）的$O_t$包含运维费用，式（9-4）的$K\,CRF$是年度资本回收要求。两项都不能整体视为可变资本：运维还包含备件与采购服务，初始安装款也可能包含承包方工资和利润。应在各雇主账户中分离工资，并对应建设批次与运行年度。装机增长可能扩大总工资，远程监控又可能降低单位电量工资，两者并不矛盾；交付电量下降还会使单位工资费用上升，即使工资总额不变。

能源技术还有一条间接路径。电力既用于生活，也用于生活资料生产；若供给生产率提高和成本下降通过相关商品的社会生产条件发生作用，在生活标准给定时可以减少劳动力再生产所需劳动。这与某座电站减员不同。设备度电成本下降尚未证明终端价格和生活资料价值同幅下降，电网、交易、税费及其他投入都可能影响传导。因此，本章融资与利用数据解释供给和交付条件，不能单独给出劳动力价值或相对剩余价值的变化量。

储能提高设备利用，还可能增加调度、维护和安全任务。新增劳动是企业雇员、设备服务商人员还是公共机构承担，需要分别登记。以利用小时增加推定“少用人”，既漏了任务结构，也漏了运行责任。周期方面，工程验收和电费回款若缩短，可以减少垫付需要；增加融资授信则是资金供给改善，不等同于生产周转加快。

建设与运行的劳动组织还受设备利用要求影响。为了使电站、储能和电网持续交付，企业可以增加班组，也可以提高现有人员覆盖范围；前者增加工资购买，后者可能增加待命和责任。远程监控提高可见性，却不能独立决定人员配置。比较单位电量工资下降时，应查其来自自动化节约、交付量增加，还是更多负担由原班组承担，并将外包和公共机构接替的任务另列。

\section{融资费用改变同样设备的投资门槛}

取安装成本667美元/kW作为公开报告的数量锚点，再公开设定其他假设：寿命25年，每年运维13美元/kW，容量因子20\%，无退化和残值。全年电量为1.752MWh/kW。分别给定3\%、6\%与10\%的实际折现率。表9-1展示融资条件如何改变结果。

\begingroup\footnotesize\setlength{\tabcolsep}{4pt}\renewcommand{\arraystretch}{1.25}
\begin{longtable}{>{\raggedright\arraybackslash}p{\dimexpr 0.14894\textwidth-8.0pt\relax}>{\raggedright\arraybackslash}p{\dimexpr 0.19149\textwidth-8.0pt\relax}>{\raggedright\arraybackslash}p{\dimexpr 0.14894\textwidth-8.0pt\relax}>{\raggedright\arraybackslash}p{\dimexpr 0.14894\textwidth-8.0pt\relax}>{\raggedright\arraybackslash}p{\dimexpr 0.19149\textwidth-8.0pt\relax}>{\raggedright\arraybackslash}p{\dimexpr 0.17021\textwidth-8.0pt\relax}}
\toprule
\textbf{假设实际折现率} & \textbf{资本回收因子} & \textbf{等额年度资本费用：美元/kW} & \textbf{演算LCOE：美元/MWh} & \textbf{售价40美元/MWh时NPV：美元/kW} & \textbf{该售价下可承受最高建设费：美元/kW} \\
\midrule\endfirsthead
\toprule
\textbf{假设实际折现率} & \textbf{资本回收因子} & \textbf{等额年度资本费用：美元/kW} & \textbf{演算LCOE：美元/MWh} & \textbf{售价40美元/MWh时NPV：美元/kW} & \textbf{该售价下可承受最高建设费：美元/kW} \\
\midrule\endhead
\bottomrule\endfoot
3\% & 0.05743 & 38.30 & 29.28 & 326.94 & 993.94 \\
6\% & 0.07823 & 52.18 & 37.20 & 62.67 & 729.67 \\
10\% & 0.11017 & 73.48 & 49.36 & \ensuremath{-}148.88 & 518.12 \\
\end{longtable}\endgroup


\textbf{表9-1　公开成本锚点下的融资机制演算。} 除667美元/kW外，输入均为假设；三种利率不代表具体国家或项目。表中结果不是IRENA的44美元/MWh估计的复现。售价40也是假设的固定实际价格，未包含储能、电网增量费用和税费。完整参数与结果见演算参数、计算表。

以6\%一行为例，667乘以0.07823，得到每年约52.18美元的资本费用；加上13美元运维，再除以1.752MWh，得到37.20美元/MWh。3\%与10\%之间，设备和发电量完全不变，演算度电成本却由29.28升至49.36，约增加68.57\%。它量化了一条明确机制：资金的取得条件可以显著改变相同物理能力的经济可用性。

采用门槛也能直接推导。固定实际售价为\(P\)时，每年净经营现金为\(PE-O\)，所能支持的最高初始投入为

\[
K_{\max}=(PE-O)A(r,T).
\tag{9-5}
\]

当\(K\leq K_{\max}\)时，所列假设下净现值非负。表9-1中，相同667美元/kW项目在3\%和6\%下满足这一条件，在10\%下不满足。相同购置价格下，项目仍可能因资金成本不同而具有不同投资可行性；融资差异由此在生产和售电之前影响采用机会。

融资条件还会由制度与经验改变。IRENA的2023年报告以2019—2021年为主要观察期，结合专家调查、访谈和基准工具研究资金成本。Egli等对德国133个光伏和陆上风电项目的约18年资料进行分析，也将融资条件作为随时间变化的对象。\href{https://www.irena.org/-/media/Files/IRENA/Agency/Publication/2023/May/IRENA_The_cost_of_financing_renewable_power_2023.pdf}{IRENA融资报告}，\href{https://www.nature.com/articles/s41560-018-0277-y?error=cookies_not_supported}{Egli等原始论文}支持研究这一通道，但其资料没有用于给本表的3\%、6\%、10\%赋予国别含义。

融资改善可能来自交易摩擦减少，也可能来自风险转由公共部门承担。公共担保降低私人支出时，需同时记录担保责任、收费、收益分享和问责。资金成本下降确认项目条件改善，是否减少社会风险则是另一项判断。

\section{设备便宜了，发电时的售价也可能下降}

前面的净现值推导假定实际售价固定。现实中发电时间与市场价格相关，光伏项目得到的是按自身发电量加权的价格。若把等长时段记为\(t\)，每段电量为\(E_t\)，价格为\(p_t\)，捕获价格为

\[
p_{\mathrm{cap}}=\frac{\sum_t p_tE_t}{\sum_t E_t}
=\bar p+\frac{\operatorname{Cov}(p_t,E_t)}{\bar E}.
\tag{9-6}
\]

横线表示所列等长时段的算术均值，协方差按同一组n个时段以1/n为分母计算，且总发电量为正。第二个等号来自\(\overline{pE}=\bar p\,\bar E+\operatorname{Cov}(p,E)\)。不同电源即使面对相同的全年平均市场价格，也可能因发电时段分布不同而取得不同的单位电量售电收入。

以两个等长时段为假设例子，价格分别为20和80美元/MWh，简单平均为50。一座电站分别发出90和10MWh，其捕获价格为\((20\times90+80\times10)/100=26\)美元/MWh。如果另一种交付安排分别提供30和70MWh，同样100MWh电量的捕获价格就变为62。差别来自电量在时段之间的分布，尚未扣除改变这种分布所需的设备与运行费用。

大量同类光伏项目在相近时段增供，可能压低这些时段的价格，使协方差更为负向。设备成本下降与捕获价格下降由此可以同时出现。IEA在《Electricity 2026》的灵活性章节将捕获价格与系统灵活性联系分析，为这一实现问题提供了行业背景。\href{https://www.iea.org/reports/electricity-2026/flexibility}{原始章节}

捕获价格下降可以减少电站收入，同时使部分用户受益；实际传导取决于零售合同。负荷调整、储能与网络扩建又可能提高受限电量的利用。系统比较需将新增投入与供电收益对应，不能以维持所有既有项目利润为评价标准。

用户电费不只付组件，还付网络、供电服务和其他费用。组件便宜，原有线路的维护和融资支出并不会立即消失。自发自用少购电以后，谁继续付网络费用，要看资费安排；有屋顶、能融资的住户与租户，得到的收益也可能不同。须查用户合同和实际费用，才知道降本怎样分配，不能只拿设备价格作结论。

\section{接入与调度决定潜在出力能交付多少}

发电与用电须在运行时满足网络输送限制和电力平衡要求。以没有本地消费和储能的单一发电项目为例，某时段可向外送出的功率，既不能超过当时可发功率，也不能超过获准的网络接入功率。假如可发功率100MW而接入仅60MW，剩余能力当时就无法全部通过该接口实现。提高组件效率并未改变这一限制。

加州运行资料提供了一个实际对照。CAISO的2024年12月31日报告将当年累计弃风弃光列为3,423,377MWh，并按调度原因拆分如下。这是运营商按预测与实际出力差额形成的初步运行记录，报告过滤小于1MW的量，数值可修订。美国能源信息署随后概述为约340万MWh、比2023年增加29\%，其中太阳能占93\%；这些经过舍入的年度概述与原始分类表应分别保留。（CAISO，2024，第1、3—4页；EIA，2025，\href{https://www.caiso.com/documents/wind-solar-real-time-dispatch-curtailment-report-dec-31-2024.pdf}{年终运行报告}、\href{https://www.eia.gov/todayinenergy/detail.php?id=65364}{官方年度分析}）

\begingroup\footnotesize\setlength{\tabcolsep}{4pt}\renewcommand{\arraystretch}{1.25}
\begin{longtable}{>{\raggedright\arraybackslash}p{\dimexpr 0.30000\textwidth-8.0pt\relax}>{\raggedright\arraybackslash}p{\dimexpr 0.18000\textwidth-8.0pt\relax}>{\raggedright\arraybackslash}p{\dimexpr 0.18000\textwidth-8.0pt\relax}>{\raggedright\arraybackslash}p{\dimexpr 0.34000\textwidth-8.0pt\relax}}
\toprule
\textbf{2024年累计调度分类} & \textbf{弃电量：MWh} & \textbf{占所列总量：\%} & \textbf{分类所指} \\
\midrule\endfirsthead
\toprule
\textbf{2024年累计调度分类} & \textbf{弃电量：MWh} & \textbf{占所列总量：\%} & \textbf{分类所指} \\
\midrule\endhead
\bottomrule\endfoot
LocalEconomic & 3,149,868 & 92.0105 & 为缓解本地拥塞而进行的经济调度 \\
LocalSelfSchCut & 2,387 & 0.0697 & 为缓解本地拥塞而削减自计划出力 \\
SystemEconomic & 271,122 & 7.9197 & 为缓解全系统供给过剩而进行的经济调度 \\
合计 & 3,423,377 & 100.0000 & 限于该报告所列风、光资源及口径 \\
\end{longtable}\endgroup


\textbf{案例表9-A　CAISO弃风弃光的运行分类。} 电量为原报告值，占比为本报告用各类电量除以合计复算，四舍五入可能使显示占比之和略偏离100\%。本地与系统分类描述调度原因，不等同于某项政策的因果贡献。原值与复算见新增能源数据。

这一分类使约束的位置更清楚：按报告口径，约92\%的弃电记在本地拥塞经济调度项下，不能全部解释为全系统总发电量超过需求。某节点仍有廉价可发电量，并不保证线路能够把它送到另一个节点；另一时段有需求，也不保证现有储能足以把电量转移过去。设备、电网和负荷需要协调，扩建哪一部分才有效，取决于具体约束的位置与时段。年度总量本身尚不能识别增加输电、储能或需求响应各能减少多少弃电。

弃电的机会成本应按当地时段、可行替代与新增接收费用判断。为避免最后一部分弃电而过度扩建，可能增加完整资源成本。案例确认运行约束的位置，未比较所有制效率；产权解释仍需合同、接入次序和可用替代证据。

用户侧的接入费用构成另一种约束。Lee、Miguel与Wolfram（2020）的肯尼亚农村实验，在150个变压器社区中对2,289户原未接电家庭实施社区层面的随机分组。下表将一次连接报价与接受情况分开列出；接受并不表示当天已经通电。研究提供的干预还包括免费基础室内接线板机会和办理安排，组间差异不能全部归为脱离这些条件的纯价格效应。\href{https://doi.org/10.1086/705417}{正式论文}、\href{https://emiguel.econ.berkeley.edu/assets/miguel_research/74/REPP_2019-04-25-FINAL-All.pdf}{作者最终稿设计部分及表B3}

\begingroup\footnotesize\setlength{\tabcolsep}{4pt}\renewcommand{\arraystretch}{1.25}
\begin{longtable}{>{\raggedright\arraybackslash}p{\dimexpr 0.23333\textwidth-8.0pt\relax}>{\raggedright\arraybackslash}p{\dimexpr 0.23333\textwidth-8.0pt\relax}>{\raggedright\arraybackslash}p{\dimexpr 0.30000\textwidth-8.0pt\relax}>{\raggedright\arraybackslash}p{\dimexpr 0.23333\textwidth-8.0pt\relax}}
\toprule
\textbf{组别} & \textbf{一次连接价格：研究口径美元} & \textbf{选择户数／分配户数} & \textbf{本报告计算的比例：\%} \\
\midrule\endfirsthead
\toprule
\textbf{组别} & \textbf{一次连接价格：研究口径美元} & \textbf{选择户数／分配户数} & \textbf{本报告计算的比例：\%} \\
\midrule\endhead
\bottomrule\endfoot
高补贴 & 0 & 363／380 & 95.53 \\
中补贴 & 171 & 90／379 & 23.75 \\
低补贴 & 284 & 27／380 & 7.11 \\
对照 & 398 & 15／1150 & 1.30 \\
\end{longtable}\endgroup


\textbf{案例表9-B　连接价格与接电选择。} 处理组记录报价接受，对照组记录常规价格下自行选择接电；两者均不能解读为即时完成通电。干预在2014年实施，处理组报价期限为8周；金额不是当前电价或全部建网成本。比例为原计数相除，未经协变量调整。随机化及标准误按社区处理，不能将每户当作独立分配单位。原值、页表定位与计算保存在新增能源数据。

研究在约接电后16个月和32个月随访。原文表3以随机补贴指示变量作为接电的工具变量，估计成年家庭成员就业或自营比例的影响为2.2个百分点，社区聚类标准误为2.5个百分点。按已舍入系数计算的正态近似95\%区间为\ensuremath{-}2.70至7.10个百分点，不能据此断言真实影响为零。这里的工具变量解释还依赖排除限制等条件，亦不能直接视为全部居民接电的平均效应。与此同时，照明替代和生活满意度有正向结果，不能把未检出明确收入改善概括为接电毫无收益。（同文表3及表注；近似区间为本报告复算，未使用微观样本重估，也不是多重检验校正区间。）

肯尼亚实验区分接电选择、生活改善与市场收入。减免连接付款解除一项进入约束，后续电器、可靠供电和经营条件未被分别随机化，不能确定哪一项解释收入结果。补贴还主要改变付款承担，未必减少施工投入。住宅连接与储能运行分别研究，不合并为统一能源效应。



两个案例提供互补证据：CAISO定位已建能力的交付约束，肯尼亚实验检验付款与服务进入。前者说明设施运行需要协调，后者说明进入改善不保证收入同幅变化。共同设施的评价由此包括交付能力、使用机会及实际所得。

在现实系统中，线路、变电设备、调度与负荷分布共同决定可接入能力。IEA《Electricity 2026》的电网章节讨论了连接队列、建设时序和释放接入能力的办法。这些资料使一个抽象命题具体化：电网是发电服务的共同生产条件，单个项目的建设决策无法替代全部系统协调。\href{https://www.iea.org/reports/electricity-2026/grids}{电网原始章节}

制造、发电投资和电网建设具有不同决策主体与周期。各方局部安排合理，也可能形成设备先到、接入滞后的组合。系统协调应按各阶段实际进度与可执行责任判断，增加组件数量不能直接解决时序错配。

接入稀缺时，能等得久、拿到长期合同、取得好位置的项目，就有设备效率之外的优势。资格若能转让，这种优势还可能算进项目资产价格。收益有多少补偿投资风险和协调工作，有多少来自别人无从进入，须查合同、排队、收费及替代。接入贵，并没有单独回答这个问题。

公共协调既要安排建设，也要安排使用。规划可以让设备、网络和需求相衔接，公开队列、进度要求和可执行责任可以减少无效占位。但线路经过的居民还要承受土地、生态和施工影响，也须参与相应决定。用电者、所在地居民和领取资产收益的人可能不同。谁参与、怎样补偿，因而也是共同生产怎样组织的一部分。

\section{中国储能装机和利用同增，收入仍须另算}

国家能源局公布的2025年运行资料显示，中国当年新增光伏约317GW，年末累计约1200GW，全年发电约1.17万亿kWh，利用率95\%。这些数字表明光伏已经交付大规模电力服务；功率、全年电量和运行利用统计则应分别解释，95\%不是全年满功率运行的时间比例。\href{https://www.nea.gov.cn/20260212/742b8c6a078347b0b39de676c05c5d58/c.html}{官方光伏运行统计}

为比较装机规模与实际运行，本章从《中国新型储能发展报告2026》的图19提取同版年度数据。原图以万千瓦和万千瓦时计量，表9-2换算为GW和GWh；名义时长由能量容量除以功率得到。

\begingroup\footnotesize\setlength{\tabcolsep}{4pt}\renewcommand{\arraystretch}{1.25}
\begin{longtable}{>{\raggedright\arraybackslash}p{\dimexpr 0.14894\textwidth-8.0pt\relax}>{\raggedright\arraybackslash}p{\dimexpr 0.14894\textwidth-8.0pt\relax}>{\raggedright\arraybackslash}p{\dimexpr 0.17021\textwidth-8.0pt\relax}>{\raggedright\arraybackslash}p{\dimexpr 0.14894\textwidth-8.0pt\relax}>{\raggedright\arraybackslash}p{\dimexpr 0.19149\textwidth-8.0pt\relax}>{\raggedright\arraybackslash}p{\dimexpr 0.19149\textwidth-8.0pt\relax}}
\toprule
\textbf{年末} & \textbf{功率：GW} & \textbf{能量容量：GWh} & \textbf{容量/功率：小时} & \textbf{功率同比增长} & \textbf{能量容量同比增长} \\
\midrule\endfirsthead
\toprule
\textbf{年末} & \textbf{功率：GW} & \textbf{能量容量：GWh} & \textbf{容量/功率：小时} & \textbf{功率同比增长} & \textbf{能量容量同比增长} \\
\midrule\endhead
\bottomrule\endfoot
2021 & 4.00 & 8.00 & 2.000 & — & — \\
2022 & 8.70 & 18.05 & 2.075 & 117.50\% & 125.63\% \\
2023 & 31.39 & 66.87 & 2.130 & 260.80\% & 270.47\% \\
2024 & 73.76 & 168.20 & 2.280 & 134.98\% & 151.53\% \\
2025 & 135.93 & 351.12 & 2.583 & 84.29\% & 108.75\% \\
\end{longtable}\endgroup


\textbf{表9-2　中国新型储能年末装机。} 来源：国家能源局等（2026），印刷第32页、PDF第42页图19；逐项核对图中标签。增速和比率为本文计算，使用报告发布精度，非单个电站面板。2021—2025年间只有四个年度增长间隔。\href{https://www.nea.gov.cn/20260730/0a1e9823b7a14e8ba16e82c34b3bea2b/202607300a1e9823b7a14e8ba16e82c34b3bea2b_270e42e93bb1a647609c851527f42364ee.pdf}{报告原件}

2025年功率约为2021年的33.98倍，能量容量约为43.89倍。能量增幅更大，使两者比值由2小时升至约2.58小时，表明报告所列装机组合的持续放电能力发生变化。该比值是按容量关系形成的系统名义时长，不能当成每座电站都具有2.58小时，也没有说明实际充放电安排。

同一报告图24提供了利用数据。表9-3把全国、独立储能和新能源配建储能分开，避免只用规模增长代表使用情况。

\begingroup\footnotesize\setlength{\tabcolsep}{4pt}\renewcommand{\arraystretch}{1.25}
\begin{longtable}{>{\raggedright\arraybackslash}p{\dimexpr 0.20000\textwidth-8.0pt\relax}>{\raggedright\arraybackslash}p{\dimexpr 0.20000\textwidth-8.0pt\relax}>{\raggedright\arraybackslash}p{\dimexpr 0.20000\textwidth-8.0pt\relax}>{\raggedright\arraybackslash}p{\dimexpr 0.20000\textwidth-8.0pt\relax}>{\raggedright\arraybackslash}p{\dimexpr 0.20000\textwidth-8.0pt\relax}}
\toprule
\textbf{年份} & \textbf{全国平均利用小时} & \textbf{新能源配建储能利用小时} & \textbf{独立储能利用小时} & \textbf{独立减配建：小时} \\
\midrule\endfirsthead
\toprule
\textbf{年份} & \textbf{全国平均利用小时} & \textbf{新能源配建储能利用小时} & \textbf{独立储能利用小时} & \textbf{独立减配建：小时} \\
\midrule\endhead
\bottomrule\endfoot
2023 & 611 & 383 & 680 & 297 \\
2024 & 911 & 766 & 995 & 229 \\
2025 & 1195 & 917 & 1246 & 329 \\
\end{longtable}\endgroup


\textbf{表9-3　新型储能利用强度。} 来源：同报告印刷第40页、PDF第50页图24；差额由本文计算。各列沿用报告的利用小时定义，不能解释为物理设备全年处于满负荷放电的小时。全国均值的完整加权构成不由后两列反推。

全国平均利用小时在2023—2025年增加约95.58\%，2024—2025年增加约31.17\%。装机统计与利用统计同时上升，因而将这一时期的储能增长解释为完全未伴随实际使用，缺乏资料支持。是否取得足以覆盖成本的收入，仍需要运行时段价格和服务支付；利用增长为收益实现提供了条件，却没有代替收益核算。

2025年独立储能比配建储能高329小时，值得进一步研究调度资格、服务组合和收益激励。但两类项目在地点、时长、投运月份和电网条件上可能不同；这些均值并未构成随机分配的组织实验。按照第三章的选择分解，组间均值差异同时可能包含组织安排与项目组成的影响；识别组织形式的作用，仍须取得可比项目的运行资料。

报告另列2025年全国平均等效充放电256次，2024年为221次，按发布值计算增长约15.84\%。其增速不同于利用小时，进一步说明两种统计不能互换。本章不拿年末功率乘全年均值推算未经口径核实的总放电量。原始数据与复算程序保存了表9-2、表9-3及计算输出。

\begin{center}\begin{tikzpicture}\begin{axis}[width=.44\textwidth,height=6cm,
title={A 装机功率},xlabel={年份},ylabel={GW},
tick label style={font=\scriptsize},label style={font=\small},title style={font=\small},
legend style={font=\scriptsize,draw=none,at={(0.5,-0.3)},anchor=north},legend columns=1,
axis lines=left,grid=major,grid style={gray!15},scaled ticks=false,xtick={2021,2023,2025},xticklabels={2021,2023,2025},]
\addplot+[blue!70!black,mark=none,thick,mark=*] coordinates {(2021,4) (2022,8.7) (2023,31.39) (2024,73.76) (2025,135.93)};
\end{axis}\end{tikzpicture}\hfill\begin{tikzpicture}\begin{axis}[width=.44\textwidth,height=6cm,
title={B 利用小时},xlabel={年份},ylabel={小时/年},
tick label style={font=\scriptsize},label style={font=\small},title style={font=\small},
legend style={font=\scriptsize,draw=none,at={(0.5,-0.3)},anchor=north},legend columns=1,
axis lines=left,grid=major,grid style={gray!15},scaled ticks=false,xtick={2023,2024,2025},xticklabels={2023,2024,2025},]
\addplot+[blue!70!black,mark=none,thick,mark=*] coordinates {(2023,611) (2024,911) (2025,1195)};
\addlegendentry{全国}
\addplot+[red!70!black,mark=none,thick,mark=*] coordinates {(2023,383) (2024,766) (2025,917)};
\addlegendentry{新能源配建}
\addplot+[green!50!black,mark=none,thick,mark=*] coordinates {(2023,680) (2024,995) (2025,1246)};
\addlegendentry{独立储能}
\end{axis}\end{tikzpicture}\par\vspace{1.2cm}\begin{tikzpicture}\begin{axis}[width=.6\textwidth,height=6cm,
title={C 融资情景},xlabel={假设实际折现率（\%）},ylabel={美元/MWh},
tick label style={font=\scriptsize},label style={font=\small},title style={font=\small},
legend style={font=\scriptsize,draw=none,at={(0.5,-0.3)},anchor=north},legend columns=1,
axis lines=left,grid=major,grid style={gray!15},scaled ticks=false,xtick={3,6,10},]
\addplot+[blue!70!black,mark=none,thick,mark=*] coordinates {(3,29.2833276) (6,37.2016102) (10,49.3619316)};
\addlegendentry{演算LCOE}
\addplot+[red!70!black,mark=none,thick,dashed] coordinates {(3,40) (10,40)};
\addlegendentry{假设售价}
\end{axis}\end{tikzpicture}\end{center}


\textbf{图9-1　技术部署、利用强度与融资门槛。} A、B为上述官方资料的重新绘图；C为表9-1的假设融资演算。图中把观测与演算分列，三部分共同说明需要考察哪些实现环节，并未估计相互之间的因果系数。


\subsection{2024—2025年新增容量主要来自功率扩张项}

储存电量的容量增加，可以多装充放电功率，也可以让每单位功率配更多电量。两种变化所需设备、场地和运行时段不同，须拆开算，才能知道增长主要来自哪项投入。一个总装机增速没有说明这种结构。

记年末装机功率为\(P_t\)，能量容量为\(E_t\)，总量之比为\(d_t=E_t/P_t\)。这里\(P\)不是前文售电价格，局部下标明确其为装机统计；三者单位分别为GW、GWh、小时。定义立即给出\(E_t=P_td_t\)。比较两个年末，有

\[
E_1-E_0=d_0(P_1-P_0)+P_0(d_1-d_0)+(P_1-P_0)(d_1-d_0).
\tag{9-6a}
\]

最后一项是功率和时长一起改变形成的。先加功率、后加时长，这一项全部计给时长；反过来，全部计给功率。两种算法都成立，份额却不同。因此，本节对两种顺序取平均，将这一项各分一半：

\[
\Delta E=\underbrace{\frac{d_0+d_1}{2}\Delta P}_{\text{功率扩张项}}
+\underbrace{\frac{P_0+P_1}{2}\Delta d}_{\text{总量时长变化项}}.
\tag{9-6b}
\]

该分解是在明确分摊规则下成立的精确恒等式，分解项并非回归系数，亦不要求两种投资决策相互独立。把式（9-6b）展开，交叉项正好相加为式（9-6a）的最后一项，故不留下无法解释的残差。复算所需的输入为同一统计范围的功率与能量容量；但若两项统计的设备范围不一致，比值和分解本身便失去含义。

\begingroup\footnotesize\setlength{\tabcolsep}{4pt}\renewcommand{\arraystretch}{1.25}
\begin{longtable}{>{\raggedright\arraybackslash}p{\dimexpr 0.25581\textwidth-8.0pt\relax}>{\raggedright\arraybackslash}p{\dimexpr 0.20930\textwidth-8.0pt\relax}>{\raggedright\arraybackslash}p{\dimexpr 0.18605\textwidth-8.0pt\relax}>{\raggedright\arraybackslash}p{\dimexpr 0.16279\textwidth-8.0pt\relax}>{\raggedright\arraybackslash}p{\dimexpr 0.18605\textwidth-8.0pt\relax}}
\toprule
\textbf{相邻年末} & \textbf{能量容量增加：GWh} & \textbf{功率扩张项：GWh} & \textbf{总量时长变化项：GWh} & \textbf{功率扩张项占增加量} \\
\midrule\endfirsthead
\toprule
\textbf{相邻年末} & \textbf{能量容量增加：GWh} & \textbf{功率扩张项：GWh} & \textbf{总量时长变化项：GWh} & \textbf{功率扩张项占增加量} \\
\midrule\endhead
\bottomrule\endfoot
2021—2022 & 10.050 & 9.576 & 0.474 & 95.28\% \\
2022—2023 & 48.820 & 47.706 & 1.114 & 97.72\% \\
2023—2024 & 101.330 & 93.440 & 7.890 & 92.21\% \\
2024—2025 & 182.920 & 151.181 & 31.739 & 82.65\% \\
\end{longtable}\endgroup


表中原始年末数据仍来自国家能源局报告图19；增加量、分解和比例是本文计算，精确输出见分解数据。显示值经过四舍五入，核对恒等式使用未舍入数值。每一行的分母是该年能量容量增加量，不能把最后一列解释为全部存量中某类设备的比重。

2024—2025年的容量变化具有明确的结构差异。能量容量增加182.92GWh，其中151.181GWh归于功率扩张项，31.739GWh归于总量时长变化项。按照上述分解规则，功率扩张仍占主要部分，同时总量时长变化项所占份额比前一年提高。其中31.739GWh表示统计分解项，并非某项政策创造容量的效应估计，也未识别每个存量项目是否增加电池。

总量时长上升可能来自原项目增容、新项目配置或退出与统计组成变化。式（9-6b）分解统计结构，未识别逐项目投资行为。要解释变化来源，需投运、改造、退役与技术路线资料。

能量容量增加提高可持续放电能力，也增加相应投入；其服务价值取决于受限时段与实际运行。短时紧张与长时紧张对额外容量的需要不同，容量增幅因而不能直接换算为生产率或收入增幅。

\subsection{两组利用不同，还要分清规则与项目组成}

2023—2025年独立储能组的报告利用小时高于新能源配建组，两组本身又都提高。组间差异可能来自接入与调度规则，也可能来自项目位置、投运时间、技术配置或进入各组的项目选择。同一组均值同样会受这些因素的共同影响，因而两组水平及其变化都需要与项目组成联系分析。

若有一致的项目分母，组平均可以写成\(\bar u_t=\sum_j\omega_{jt}u_{jt}\)，其中\(u_{jt}\)是项目指标，\(\omega_{jt}\)是相应权重。权重究竟按功率、容量、项目数还是运行期确定，必须遵循报告原定义。不能只看到“小时”这个单位，就假设它一定是对电站数量的简单平均。本报告当前保存的汇总图不足以恢复每个项目权重，因而不进一步估计全国均值由哪些站点贡献。

在仅改变项目权重的假设情景中，两个项目的利用小时一直分别为400和1000，第一年权重各半，均值700；第二年高利用项目权重升至80\%，均值便为880。平均数增加180小时，任何一个项目的运行都没有改善。若权重保持各半而两个项目各增加180小时，均值也会升到880。这两种过程对运行改革和建设选择的含义不同，汇总均值却完全相同。

相应地，图24中的分组差异不能单独成为“改成独立储能即可增加329小时”的效应估计。即使另有材料表明独立项目拥有更灵活的市场参与机会，仍需辨别项目条件与选择。合理的下一步是将接入权限、调度安排和收入规则与同一项目运行记录对应，比较规则变化前后，并寻找没有受到同一变化但原趋势可比较的项目。这些项目记录与比较对象是识别调度和准入规则作用所需的资料。



\subsection{运行强度、项目收支与系统节约须分别核算}

投资者先看能收回多少、要付多少；系统则须查储能在什么时候、什么地方多供了电、多作了调节，省掉哪种替代投入。装机功率、能量容量、利用小时和循环次数，记录同一过程中的不同环节。四项都增加，不能当成四笔独立社会收益，否则可能重复计同一运行活动，也漏掉增加活动所耗的能量。

两座同容量电站，全年放电相同，一座在供给宽松时循环，一座在紧张时供电，系统价值仍可能不同。备用设施很少放电，也可能可靠地减少停供风险和昂贵备用。但这要查可用率、无此设施时的风险与替代费用。低利用小时不能单独证明浪费，“备用”名称也不能免去效果核查。

因此，要把设备运行、项目收支和系统作用对应起来。运行记录列配置、充放电、损耗、停机与时段，财务记录列建设、维护、融资和收入，系统记录列受限时段的新增供电、调节及可用替代。使用同一项目和时段标识，才能查一笔收入对应了什么服务、真实省了什么资源、谁改为付钱。项目获补偿而盈利以后，还要看系统收益，不能把支付本身当作公共效果。

上述区分把规模统计接到下一节的项目核算：充放电时间、损耗、价差和服务支付共同决定收支，网络和调度再限制可交付服务。运行强度、项目盈利与系统节约因此分别评价。

\section{储能收入要扣损耗和费用，进入还会改变价差}

储能的物理功能是转移能量。设时段长度为\(\Delta t\)，充电功率为\(c_t\)，放电功率为\(d_t\)，充放电效率分别为\(\eta_c,\eta_d\)，储存电量为\(e_t\)，则忽略自放电时

\[
e_{t+1}=e_t+\eta_c c_t\Delta t-\frac{d_t\Delta t}{\eta_d},
\quad 0\leq e_t\leq E_{\max}.
\tag{9-7}
\]

充放电功率还分别受设备功率约束；简化单设备调度中，可要求同一时段不同时充放电。式（9-7）说明储能服务受能量和功率两种限制。高功率设备能迅速放出电量，却未必能持续很久；更大的能量容量也不自动提高瞬时输出功率。

若往返效率为\(\eta_{\mathrm{rt}}=\eta_c\eta_d\)，交付1MWh需购入\(1/\eta_{\mathrm{rt}}\)MWh。低价时购电价格为\(p_L\)，高价时售电价格为\(p_H\)，则交付每MWh的毛价差收入为

\[
M=p_H-\frac{p_L}{\eta_{\mathrm{rt}}}.
\tag{9-8}
\]

假设\(p_L=20\)、\(p_H=80\)美元/MWh，往返效率90\%，毛收入为57.78美元/MWh，而不是60。还须从中支付循环退化、运行和建设费用。如果每MWh退化与变动成本为\(v_s\)，全年交付电量为\(E_d\)，年化固定费用为\(F_s\)，其他已核实服务净收入为\(R_s\)，年度余额可写为

\[
\Pi_s=E_d(M-v_s)+R_s-F_s.
\tag{9-9}
\]

式（9-9）暂将全年电量按同一毛价差交付，实际核算应逐时进行。造价、运行量、调度与其他服务分别改变$F_s$、$E_d$、$M$与$R_s$；只提高一项不保证余额转正。

IEA关于电池与能源转型的报告区分了电量转移、短时灵活性和其他系统用途，本章据此分别评价电量转移、短时调节及其他用途。\href{https://www.iea.org/reports/batteries-and-secure-energy-transitions/executive-summary}{原始报告执行摘要}中的未来部署情景不作为已实现数量。本章关注的实际运行数据已在表9-2与表9-3单列。

新储能还会改变自己赖以盈利的价差。大家在低价时买电、高价时卖电，可能抬高低价、压低高价，使\(M\)缩小。单站套利收入可能减少，系统紧张却可能缓解。这也要考虑第四章所说的采用反馈：按原价差有利可图，不能保证全部进入以后还取得同样收益。须把各站行动带来的价格变化也算进去。

\section{套利回报与运行减排可以分离}

储能是否让社会受益，要比交付服务的价值与新增资源、风险费用。价格若反映真实稀缺，可以引导投资；若没有充分计入排放、拥塞、安全或可靠性，就可能让盈利与社会净收益分开。因此，项目收支算清以后，还须查它增加了什么、减少了什么，由谁承担。

以排放为例，设充电时边际电源每MWh排放为\(g_L\)，放电时所替代电源排放为\(g_H\)。交付1MWh带来的运行排放变化是

\[
\Delta g=\frac{g_L}{\eta_{\mathrm{rt}}}-g_H.
\tag{9-10}
\]

只有在\(g_H>g_L/\eta_{\mathrm{rt}}\)时，该次转移按所列边际电源比较减少运行排放。式（9-10）尚未计入制造和退役阶段，因此不是生命周期排放结论。低价充电可以有利于套利，却不由价格本身保证充电边际排放足够低。

Hittinger与Azevedo对当时美国20个区域的套利型储能建模，展示了盈利和减排可能不一致的情形。\href{https://hero.epa.gov/reference/2826816/}{原始研究摘要}在本章用于确定需要检查的条件，不将其效应量移作2025年中国的估计。现实检验需要小时与地点匹配的边际供给及实际调度，年度平均能源结构不足以单独回答这一问题。

公共支付要对应实际多交付的服务。没有支付也会按价差建成并提供的服务，再付同一笔钱，可能主要增加资产收益；系统确需备用或可靠性、却不能由电量交易补偿时，购买则可能补上缺口。须比较有无支付的交付，查性能、可用容量和违约责任，才能分清哪一种情况。

同一设备还可能同时签下多项义务。某时段已承诺全功率放电，又将相同可用功率全卖作随时可调用服务，就须证明两项能同时做到。核收入还要回到式（9-7）的物理约束。账面可以列多笔钱，设备却只有一套，所承诺的服务能否实际交付，仍须分别证明。

\section{使用地降本，还须查供应链与退役负担}

光伏和电池把部分持续燃料支出换成前期设备和材料需求。这改变了采用者的费用，也让开采、精炼、制造与回收在哪发生成为问题。使用地取得便宜电力，可能仍靠别处的劳动和环境投入。要比较整体收益，就须沿供应链查这些投入由谁承担。

IEA《Global EV Outlook 2026》的电池章节报告，2025年底全球锂离子电池名义制造能力超过4TWh/年；新登记EV所装电池的容量中，LFP相关口径份额超过55\%。这些数据描述制造能力和车辆电池组合，不能代替中国固定储能的运行电量。\href{https://www.iea.org/reports/global-ev-outlook-2026/electric-vehicle-batteries}{原始章节及图注}

名义制造能力要先变成合格产品，产品要符合用途，再经运输、安装和接入，才能运行。某处设备富余，不排除别处原料、施工或接入短缺。因此，讨论“过剩”须统一用途、质量和期限，比较产量、订单、库存与实际可交付能力，不能由某项名义容量直接判断整条链。

材料环节同样具有专门化。矿产资源、开采、精炼、正负极材料、电芯和系统集成属于不同阶段；某一阶段集中，不代表其他阶段具有相同结构。IEA关于关键矿产的市场资料将需求和供应环节分别讨论，为识别依赖位置提供依据。\href{https://www.iea.org/reports/global-critical-minerals-outlook-2026/market-overview}{2026年原始市场章节}中的长期预测则需保留情景条件。

集中可能降低费用，也可能让一次中断影响更多下游。规模和专业知识可以省投入，少数节点故障或限制又会扩大影响。因此，要比备用供货、库存和冗余花多少钱，停供损失由谁担。只选最低采购价，可能漏掉脆弱性；只增加供应者数，也未必以较少费用保证可靠性。

报价没付的费用，以后仍可能有人要付。项目未计拆解和处理，盈利结束后就可能留下公共支出；回收还要收集、运输、加工，并取得合格再生材料。设备未退役，也不能先当作已有回收材料使用。因此，须把责任写入资金和合同，让采用时就看见未来支出，而不等资产收益领走后再交给别的承担者。

\section{公共支持须同时安排收益、责任与决定权}

前文的推导使制度比较能够落到具体变量。改善融资可能降低\(r\)，共同电网建设和公平接入可能提高可交付的\(E_t\)，灵活负荷和储能可能改变捕获价格，公共服务合同可能增加与责任相匹配的\(R_s\)。这些改变分别作用于采用、实现和分配，评价时应追踪它们是否实际发生。

\begingroup\footnotesize\setlength{\tabcolsep}{4pt}\renewcommand{\arraystretch}{1.25}
\begin{longtable}{>{\raggedright\arraybackslash}p{\dimexpr 0.25000\textwidth-8.0pt\relax}>{\raggedright\arraybackslash}p{\dimexpr 0.29545\textwidth-8.0pt\relax}>{\raggedright\arraybackslash}p{\dimexpr 0.45455\textwidth-8.0pt\relax}}
\toprule
\textbf{制度安排} & \textbf{直接改变的条件} & \textbf{需要一并核查的问题} \\
\midrule\endfirsthead
\toprule
\textbf{制度安排} & \textbf{直接改变的条件} & \textbf{需要一并核查的问题} \\
\midrule\endhead
\bottomrule\endfoot
标准化合同、信用改善 & 资金回报要求与交易费用 & 是否减少摩擦，风险是否仍被正确计量 \\
公共担保或长期采购 & 收入与违约风险的承担 & 谁在失败时付款，收益如何分享 \\
网络扩建与接入改革 & 可交付功率、等待和准入 & 容量分配、土地影响与维护费用由谁决定和承担 \\
容量及辅助服务支付 & 可验证服务的收入 & 是否增加真实服务，是否重复承诺和收费 \\
社区或公共参与 & 投资、收益使用和规则修改权 & 谁有资格参与，低收入者和租户是否被覆盖 \\
\end{longtable}\endgroup


\textbf{表9-4　制度作用的识别位置。} 表中为机制对应关系，实际效果仍需项目及规则变化的证据。

所有权属于这一比较的重要部分，但需要与权限和受益范围一起解释。Walker与Devine-Wright区分社区能源的参与过程与收益结果；Jenkins等从分配、承认和程序讨论能源正义。\href{https://research.manchester.ac.uk/en/publications/community-renewable-energy-what-should-it-mean/}{社区能源原始资料}，\href{https://www.storre.stir.ac.uk/bitstream/1893/23236/1/Energy\%20Justice\%20A\%20Conceptual\%20Review\%20-\%20final\%20pre-pub\%20version.pdf}{能源正义作者稿}为分别分析受益分布与决策参与提供了概念依据。

相同设备与供电可以配合不同的收益和治理安排。仅向出资住户分红与让租户、缺资金者参与规则具有不同范围。参与程序若缺信息、预算或修改权限，也未必影响投资。制度评价需落实到谁受益、谁有权作出和纠正决定。

降价、服务改善，可以减轻具体冲突；谁能定投资、接入和剩余用途，又关系到生产支配。两种变化可以同时发生。合同公开，可以少花融资费用，也可能让使用者取得监督权；共同设施，可以扩大使用，也可能仍由少数管理者作主。因此，要沿具体事项查权限怎样变，才知道改善发生到哪一步。

\section{协调决定交付，收支与权限说明受益}

本章区分三层结果。融资与网络影响采用和交付，捕获价格与服务支付形成项目收入，可靠性、成本和排放决定系统效果。三层可同向改善，也可分化；运行统计与条件演算分别支持其中不同环节。

资料记录了供给与生活的实际改善：所列时期中国储能装机和利用同增，肯尼亚实验观察到照明与满意度提高。承认这些结果，并不免除项目收支、服务覆盖和公共费用的比较；未显著的收入结果也不能抹去已有生活改善。

共同电网要整体运行，投资、接入和收益却可以分别决定，本章就在这一联系中分析基本矛盾。公共设施可以帮助私人投资，私人投资也可以提供真实服务；若公共承担风险、少数主体仍支配收益和使用规则，收益与责任就可能不相称。协调建设、验证服务和公平接入，可以改善具体关系；要进一步转化占有，还须改变资产归属，以及投资、调度和剩余由谁决定。运行统计没有比较各所有制的整体效果，这一步仍须查合同和实际权限。

能源案例确立了共同设施的第一种作用：设备能力须经网络与时序协调才能成为有效交付，项目收入又不等于系统价值。这一结论解释了单项降本为何不能独自决定受益范围。第十章检验另一种情形：服务已经交付，新增道路活动的影响却落在交易之外的人身上。

\section*{本章参考文献}

Egli, F., Steffen, B., \& Schmidt, T. S. (2018). A dynamic analysis of financing conditions for renewable energy technologies. \emph{Nature Energy, 3}, 1084–1092. \url{https://doi.org/10.1038/s41560-018-0277-y}

Hittinger, E. S., \& Azevedo, I. M. L. (2015). Bulk energy storage increases United States electricity system emissions. \emph{Environmental Science \& Technology, 49}(5), 3203–3210. \url{https://doi.org/10.1021/es505027p}

International Energy Agency. (2024). \emph{Batteries and secure energy transitions}. \url{https://www.iea.org/reports/batteries-and-secure-energy-transitions}

International Energy Agency. (2026a). \emph{Global EV outlook 2026: Electric vehicle batteries}. \url{https://www.iea.org/reports/global-ev-outlook-2026/electric-vehicle-batteries}

International Energy Agency. (2026b). \emph{Electricity 2026}. \url{https://www.iea.org/reports/electricity-2026}

International Energy Agency. (2026c). \emph{Global critical minerals outlook 2026: Market overview}. \url{https://www.iea.org/reports/global-critical-minerals-outlook-2026/market-overview}

International Renewable Energy Agency. (2023). \emph{The cost of financing for renewable power}. \url{https://www.irena.org/Publications/2023/May/The-cost-of-financing-for-renewable-power}

International Renewable Energy Agency. (2025). \emph{Renewable power generation costs in 2024}. \url{https://www.irena.org/Digital-Report/Renewable-Power-Generation-Costs-in-2024}

International Renewable Energy Agency. (2026). \emph{Renewable power generation costs in 2025}. \url{https://www.irena.org/Publications/2026/Jul/Renewable-Power-Generation-Costs-in-2025}

Jenkins, K., McCauley, D., Heffron, R., Stephan, H., \& Rehner, R. (2016). Energy justice: A conceptual review. \emph{Energy Research \& Social Science, 11}, 174–182. \url{https://doi.org/10.1016/j.erss.2015.10.004}

Walker, G., \& Devine-Wright, P. (2008). Community renewable energy: What should it mean? \emph{Energy Policy, 36}(2), 497–500. \url{https://doi.org/10.1016/j.enpol.2007.10.019}

国家能源局. (2026a). \emph{2025年可再生能源并网运行情况}. \url{https://www.nea.gov.cn/20260212/742b8c6a078347b0b39de676c05c5d58/c.html}

国家能源局能源节约和科技装备司、国家发展和改革委员会创新和高技术发展司、电力规划设计总院. (2026). \emph{中国新型储能发展报告2026}. \url{https://www.nea.gov.cn/20260730/0a1e9823b7a14e8ba16e82c34b3bea2b/c.html}

California Independent System Operator. (2024). \emph{Wind and Solar Curtailment: December 31, 2024}. Posted January 1, 2025. \url{https://www.caiso.com/documents/wind-solar-real-time-dispatch-curtailment-report-dec-31-2024.pdf}

U.S. Energy Information Administration. (2025, May 28). \emph{Solar and wind power curtailments are increasing in California}. \url{https://www.eia.gov/todayinenergy/detail.php?id=65364}

Lee, K., Miguel, E., \& Wolfram, C. (2020). Experimental evidence on the economics of rural electrification. \emph{Journal of Political Economy, 128}(4), 1523–1565. \url{https://doi.org/10.1086/705417} \href{https://emiguel.econ.berkeley.edu/assets/miguel_research/74/REPP_2019-04-25-FINAL-All.pdf}{作者2019年4月最终稿及附录}。

\chapter{自动驾驶：车队少用驾驶劳动，道路仍要共同组织}\label{ch:10}

\section{驾驶自动化须经运营形成运输结果}

自动驾驶首先改变驾驶任务，完整运输仍需要维护、调度、应急和道路。车队成本、乘客服务与公共影响经过运营安排形成。本章分别评价限定范围的安全证据、完整服务费用与道路活动，解释局部自动化为何可能伴随不同社会结果。

本章只讨论限定道路、区域、速度和天气下的L4应用。等级说明谁承担驾驶任务，运行条件说明能够在哪里承担。安全材料查已运行范围的事件，成本核算查经营投入，车辆公里与替代方式查道路影响。三者分别回答问题，须有事实连接；等级提高或测试里程增加，没有直接证明服务更省、城市负担更小。

与电网限制交付不同，道路还承受服务扩大带来的交易外影响。乘客和经营者取得收益时，其他道路使用者可能承担延误和风险，公共机构可能增加投入。这些结果未必进入车费，公共评价须覆盖交易之外的人群。

\section{同类事件与相应里程，才能比较安全}

安全首先是采用和服务扩张的约束。设同一种事件共观察到\(k\)次，对应行驶里程为\(L\)，则每百万英里的事件率为

\[
\hat\lambda=\frac{k}{L}\times10^6.
\tag{10-1}
\]

算事件率，分子和分母须对应。同一种车辆涉及事件，要配相应车辆里程；保险索赔要配保险暴露；测试接管又是另一类活动。都叫“安全数据”，也不能把它们合成一个事故率。先分清记的是什么，才能比较。

Kusano等的研究比较Waymo截至2023年10月底714万英里的无人驾驶员载客运行与人类驾驶基准。其接受稿报告，合并地区的警方报告涉及车辆率为每百万英里2.1，对照基准为4.68；有伤报告涉及率约低80\%。这些结果对该车队和所研究条件提供了有利证据。\href{https://doi.org/10.1080/15389588.2024.2380786}{正式DOI}，\href{https://arxiv.org/html/2312.12675v3}{接受稿方法与表7}

按公布并舍入后的2.1和4.68直接计算，相对差为\(1-2.1/4.68\approx55.13\%\)，与作者所称约55\%一致。该复算核对公布比率，未重新估计置信区间；两组是否可比仍取决于样本构成，而非计算本身。原研究按部分地区、道路和车辆条件构建比较，但天气、时段和报告过程仍可能留下差异；作者均属Waymo，独立审查和后续资料仍有价值。

Scanlon等的基准研究进一步说明，车型、道路、地理范围和报告门槛会改变比较结果。\href{https://arxiv.org/html/2312.13228v2}{原研究接受稿}，\href{https://doi.org/10.1080/15389588.2024.2380522}{正式DOI}。若自动驾驶车队更容易记录轻微接触，而人类驾驶资料只包括警方获知的事件，直接比较全部事件会混入记录差异。把同一等级事件和暴露条件对齐，是让证据有解释力的必要步骤。

\begingroup\footnotesize\setlength{\tabcolsep}{4pt}\renewcommand{\arraystretch}{1.25}
\begin{longtable}{>{\raggedright\arraybackslash}p{\dimexpr 0.22581\textwidth-8.0pt\relax}>{\raggedright\arraybackslash}p{\dimexpr 0.40323\textwidth-8.0pt\relax}>{\raggedright\arraybackslash}p{\dimexpr 0.37097\textwidth-8.0pt\relax}}
\toprule
\textbf{材料} & \textbf{观察对象与时期} & \textbf{能支持的主要判断} \\
\midrule\endfirsthead
\toprule
\textbf{材料} & \textbf{观察对象与时期} & \textbf{能支持的主要判断} \\
\midrule\endhead
\bottomrule\endfoot
Kusano等，2024 & 特定城市714万英里、截至2023年10月底；事故涉及率 & 所列运行范围内若干事件率低于构建的人类基准 \\
Di Lillo等，2024作者稿 & 2018年1月至2024年7月Waymo2530万英里；第三方责任索赔 & 财产损害与人伤索赔频率的限定比较 \\
NHTSA事故报告 & 不同企业按报告要求上报的事件 & 具体事件与监管调查线索；缺少统一里程时不能直接排序事故率 \\
California DMV年度测试报告 & 2024年12月至2025年11月公共道路测试 & 许可测试活动量及脱离报告，不能替代收费运营结果 \\
\end{longtable}\endgroup


\textbf{表10-1　安全资料的比较边界。} 来源共享企业数据或作者时，不计作完全独立的重复验证。各资料详细读取位置见本章证据台账。

保险研究报告总体财产损害与人伤索赔频率相对基准低88\%和92\%，较新车型子样本为86\%和90\%。其人类基准来自2017—2022年保险经验，空间调整及事件形成过程与事故论文不同。原稿对较新车型年份存在2018—2021与2018—2022两处表述，本章不替其确定唯一边界。\href{https://waymo.com/research/do-autonomous-vehicles-outperform-latest-generation-human-driven-vehicles-25-million-miles/}{作者发布的研究与PDF}中的结果因此按各自定义保留，不能写成所有碰撞或损失金额减少同样比例。

\section{未见严重事件，不等于风险为零}

严重事件较少，有限里程下观察到零次，并不意味着真实风险为零。本节真实事件率以每英里计，L以英里计，展示上限时再乘以百万；它与上一节每百万英里的报告单位须换算。用一个明确的统计假设说明：若事件以固定率\(\lambda\)独立发生，行驶里程为\(L\)，事件数服从均值\(\lambda L\)的泊松分布，则零次事件的概率为

\[
\Pr(k=0\mid\lambda,L)=e^{-\lambda L}.
\tag{10-2}
\]

令该概率等于0.05，可以得到零事件情况下的单侧95\%上限：

\[
\lambda_U=\frac{-\ln(0.05)}{L}\approx\frac{2.996}{L}.
\tag{10-3}
\]

在这一模型下，100万英里零事件，对应上限约每百万英里3次；1000万英里零事件，上限约0.30次。观察里程扩大十倍，上限才降低到十分之一。上述假设计算反映观察里程与统计上限的关系，不构成现实车队是否达标的判断。

Kalra与Paddock的2016年研究采用历史基准讨论稀有事件所需里程，表明单靠道路暴露证明可靠性可能需要很大样本。\href{https://doi.org/10.7249/RR1478}{原报告}中的历史基准不能作为当前所有城市的统一许可门槛。实际风险还会随软件版本、道路、天气和服务范围改变，跨条件累计里程也未必符合式（10-2）的固定率假设。

道路事件、场景测试、事故调查和持续监测分别提供运行、稀有条件、原因和版本更新的证据。既有范围的低事件率支持相应条件下的判断，扩展服务范围仍需对应证据，不宜用累计里程代替条件可比性。

信息取得本身具有治理含义。NHTSA说明其公开报告未按车辆数、里程或运行设计域标准化，记录能力、重复报告和后续更正也影响数据。2026年10月8日查阅时，页面图表说明的数据截至8月17日。\href{https://www.nhtsa.gov/laws-regulations/standing-general-order-crash-reporting}{官方资料及限制}。公众若只能取得事件总数，便难以判断风险；要求相匹配的暴露和版本资料，是使监督具有实际能力的条件。

\section{完整服务成本包含支持劳动与设备}

自动驾驶可以减少部分车内驾驶劳动，但完整服务仍需车辆、能源、维修、清洁、调度和异常处置。将这些活动纳入同一期间，设车队固定费用为\(F_v\)，车辆总行驶公里为\(K_v\)，单位公里其他运行费用为\(c_v\)，远程与现场支持工时为\(H_s\)，相应每工时费用为\(w_s\)，其他已明确列出的服务费用为\(C_o\)，完成出行次数为\(R>0\)。平均服务成本为

\[
AC_{\mathrm{trip}}
=\frac{F_v+c_vK_v+w_sH_s+C_o}{R}.
\tag{10-4}
\]

这里的\(F_v\)须与观察期一致，不能把全部购车费用算在某一天；\(C_o\)不再重复包含前面已经列出的费用。若保险和事故责任支出另行计入，也要说明其预期与实现口径。式（10-4）是成本清单的核算，没有假定哪种车队必然更便宜。

完整服务降本可来自单位公里节约、设施利用改善或支持工时减少。提高订单密度也可能同时增加完成出行与空驶公里，需要在同一范围内核算分子和分母。减少车内驾驶劳动只是其中一项变化。

Waymo的2026年远程协助说明将其角色描述为向系统提供建议，2024年车队响应说明还包括必要时由现场人员取回车辆。\href{https://waymo.com/blog/shorts/advice-not-control-the-role-of-remote-assistance/}{远程协助说明}，\href{https://waymo.com/blog/2024/05/fleet-response/}{车队响应说明}。这两项公司自述说明需要研究哪些服务劳动，尚不足以算出全流程劳动替代率。不同企业的远程能力与人员职责也可能不同。

Markey办公室的2026年监督报告依据七家企业答复讨论远程支持，指出协助频率等披露缺口。\href{https://www.markey.senate.gov/imo/media/doc/remote_assistance_investigation_report.pdf}{原报告}。若没有请求次数、持续时间、待命和现场处理资料，只用工作人员人数除以车辆数会误读支持成本。一个时点的人员比例不是每次出行投入工时。

验证和组织先花一笔费用，订单多就较容易摊开，大网络因而可能少花单位成本。但乘客是否分享节约，还要看别的车队能否进入、乘客能否换交通方式。运行资料与接口若不能带走，换用又可能损失原有能力。规模降本和退出受限，是两种不同作用，现有支持人员资料没有测出各自强度。

\section{车队可以缩小，道路行驶也可能增加}

车辆利用提高，可能用更少车辆满足同样需求；道路承受的却是车开了多少、停靠多少，并非停车场里有多少车。因此，车数与道路活动须分开算。下式把乘客服务、载客率和空驶接在一起，说明两者为何可能反向变化。

设某日乘客出行总量为\(P\)人公里，载客路段按车辆公里加权的平均乘员数为\(o\)，空驶公里占全部车辆公里的比例为\(e\)，其中\(0\leq e<1\)。载客车辆公里为\(P/o\)，占总公里的\(1-e\)，所以

\[
K_v=\frac{P}{o(1-e)}.
\tag{10-5}
\]

\(P\)增加表示服务需求扩大，\(o\)提高表示拼车或载客率提高，\(e\)增加表示空驶份额提高。三者共同决定道路上的总公里。这里的空驶包括无乘客的调度等活动，统计时应明确充电、回场和其他行程是否计入。

两个独立情境的比率为

\[
\frac{K_{v1}}{K_{v0}}
=\frac{P_1}{P_0}\frac{o_0}{o_1}\frac{1-e_0}{1-e_1}.
\tag{10-6}
\]

因此，要让总公里下降，需要满足

\[
\frac{P_1}{P_0}
<\frac{o_1(1-e_1)}{o_0(1-e_0)}.
\tag{10-7}
\]

式（10-7）给出了明确的反转条件：拼车与空驶控制带来的有效装载改善，须超过服务需求扩张。需求增长本身可能带来可及性收益，车辆公里增加也不等于福利必然下降；但它会改变道路、能源和外部成本，不能从车数减少中抹去。

考虑一个假设算例。基线每天10万人公里，载客平均1.2人，空驶占10\%；新方案需求增加20\%，载客平均提高到1.5人，空驶却增至25\%。由式（10-6），总车辆公里变成原来的\(1.2\times1.2/1.5\times0.9/0.75=1.152\)，增加15.2\%。

再假设每辆车每日实际行驶时间由8小时增至16小时，行驶平均速度均为25公里/小时，则满足这些公里的平均运力车辆数为\(N_v=K_v/(25h_v)\)。表10-2展示两个指标的不同方向。

\begingroup\footnotesize\setlength{\tabcolsep}{4pt}\renewcommand{\arraystretch}{1.25}
\begin{longtable}{>{\raggedright\arraybackslash}p{\dimexpr 0.32258\textwidth-8.0pt\relax}>{\raggedright\arraybackslash}p{\dimexpr 0.22581\textwidth-8.0pt\relax}>{\raggedright\arraybackslash}p{\dimexpr 0.22581\textwidth-8.0pt\relax}>{\raggedright\arraybackslash}p{\dimexpr 0.22581\textwidth-8.0pt\relax}}
\toprule
\textbf{假设量} & \textbf{基线} & \textbf{新方案} & \textbf{变化} \\
\midrule\endfirsthead
\toprule
\textbf{假设量} & \textbf{基线} & \textbf{新方案} & \textbf{变化} \\
\midrule\endhead
\bottomrule\endfoot
乘客服务量：人公里/日 & 100000 & 120000 & +20\% \\
平均载客人数 & 1.2 & 1.5 & +25\% \\
空驶公里占比 & 10\% & 25\% & +15个百分点 \\
全部车辆公里/日 & 92592.59 & 106666.67 & +15.2\% \\
每车实际行驶小时/日 & 8 & 16 & +100\% \\
平均运力折算车数 & 462.96 & 266.67 & \ensuremath{-}42.4\% \\
\end{longtable}\endgroup


\textbf{表10-2　车辆存量与道路活动的分离。} 完全为假设演算，折算车数未包括峰时备用、维护缓冲和车辆整数约束，不是城市所需车队规模预测。

ITF的2015年里斯本模拟提供了相同逻辑的具体例子：在全替代、拼车并保留高容量公共交通的一个情景中，车辆数为基线10.4\%，汽车公里为106.4\%。\href{https://doi.org/10.1787/5jlwvzdk29g5-en}{原报告表4—5及设定}。这是一组模型结果，说明大幅减少车辆可以与行驶增加并存；它没有证明现实城市已实现相同比例，更没有直接估计拥堵或排放。

\section{接驳公交与替代公交，形成不同道路后果}

订单给车队带来收入，道路则由乘客和其他交通参与者共同使用。便宜、少等、不用自己开，可以实现原先无法完成的出行，也可以让人从公交、步行或骑行改乘车。新增公里的意义因而不同：可能增加必要服务，也可能用低载客车辆替代高容量交通。须查服务了谁、替代了什么，才能由经营收入转向社会收益。

接驳公交，可以让更多人进入公共交通；主要抢走密集走廊公交乘客，则可能用更多车辆活动运同样的人。驾驶技术相同，组织不同，道路后果便不同。服务合同、接驳、停靠和信息要求，可以影响订单与空驶，将经营接到城市交通目标。公共支持须说明服务怎样安排，技术先进本身没有确定这些结果。

车辆公里的外部影响取决于地点、时段和停靠。相同公里在空闲道路与高峰瓶颈产生的延误不同，停靠还影响公交、行人和应急。公共核算需记录新增设施、监管和应急投入，不将历史道路费用全部归给新车队，也不遗漏其新增使用成本。

固定费用在收入高、订单密的地区可能较容易摊开，出行迫切的地方却未必有同样经营条件。覆盖义务和公开补偿可以帮助后者，仍须核实实际服务与责任。效率、覆盖和道路负担会相互影响，须查真实运行，才知道公共支付增加了哪些服务，又由谁承担费用。

\section{任务替代与原劳动者补偿分别评价}

驾驶任务减少之后，一部分车辆检查、乘客支持、装卸、异常处理和调度活动仍需组织。美国交通部2021年关于长途货运与公交的报告将非驾驶职责及采用条件列入劳动影响分析。\href{https://www.transportation.gov/av/workforce/report}{官方报告}。它是前瞻研究，不是自动驾驶普及后的就业因果估计，但说明职业比单一驾驶动作包含更多内容。

劳动变化可以延续第五章的分解。设完成每次出行所需直接驾驶工时为\(h_d\)，支持工时为\(h_s\)，全年服务次数为\(R\)，则对应的服务劳动需求为

\[
H=R(h_d+h_s).
\tag{10-8}
\]

这只是给定服务范围下的核算。自动化降低\(h_d\)，可能增加或减少\(h_s\)，价格和需求又改变\(R\)。净变化须联合计算；另有研发、制造和其他行业的劳动，还应在更宽边界中研究。

总劳动需求稳定不保证原司机被补偿。支持岗位可能位于其他地区、要求不同技能并支付不同报酬。评价转岗需追踪同一受影响群体的进入、训练、收入中断与后续所得，不能仅列未来岗位名称。

培训与保障要对准这些困难。课程免费，学习期间仍要生活；换了岗位名称，技能也未必能带到别处。劳动者能否参与部署时间、转岗和风险评估，又影响调整怎样落到个人。学习与过渡收入还要同家庭照护、公共服务一起安排，第十三章继续查这些条件。

部署速度、工作标准、过渡费用与新增剩余的权限，规定调整怎样落到劳动者身上。劳动参与改变哪些事项，应按可执行决定确认；就业或服务改善不自动意味着上述权限一并转移。

\subsection{自动驾驶减少驾驶工资，支持劳动和服务扩张影响总额}

自动驾驶最直接的变化，是以设备和软件承担部分原由司机完成的任务。运输企业的工资结果仍要计入实际保留的驾驶劳动、远程支持与现场维护。设服务次数为$R$，每次出行各类付酬工时为$h_k$，对应时薪为$w_k$，固定付酬支持工时为$H_0$、平均时薪为$w_0$，则在相同运输服务边界内

\[
W^{\mathrm{运}}=R\sum_k w_kh_k+w_0H_0.
\tag{10-V1}
\]

驾驶工时减少，不保证加权工资支出同比减少。高报酬的技术支持、新增现场任务及更大的服务量都可能抵消节省；一人服务多辆车，也须保留轮班、待命和异常处理时间。Waymo于2026年2月公布的说明将远程协助界定为对系统请求提供建议，并说明其仍需运行人员。\href{https://waymo.com/blog/shorts/advice-not-control-the-role-of-remote-assistance/}{企业原始说明}。该说明支持辨认支持任务，未提供本式所需的完整工时与工资，不能用车数和同时值班人数相除，估计全年可变资本节省。

自营司机变为雇员或独立经营者退出，还会改变资本关系的范围。自营者收入含车辆成本和经营余额，不能与受雇司机工资直接对等。自动化运输公司的工资下降也不能证明原司机已被补偿；应另查原人群的收入中断、转岗和资产损失。最后，乘客更早付款或平台更晚结算可以改变垫款，却不改变驾驶自动化的工程性能，劳动替代与周转必须分别比较。

支持劳动集中后，资本可以调整每名人员服务的车队范围、异常响应时限和轮班。较少人员配合更多自动车辆，是劳动购买与设备使用关系的变化；如果异常负担和待命随之提高，便同时发生劳动使用的变化。安全冗余则可能要求增加人员。应将服务规模、单位支持工时、工资和责任联读，区分省去了驾驶劳动、提高了支持效率和压缩了安全余量，不能用一个人员车辆比概括三者。

\section{公共监督须有信息、纠错与限制运行的权限}

软件怎样更新、车队怎样调度、事故资料怎样解释，会影响共同道路上的人。集中协调可以帮助及时修复，也可能使外部监督离不开经营者提供的信息。因此，经营者须能协调，公共机构也须能取得资料、查明责任并要求纠错。只有义务名称，没有可执行权限，监督就难以完成。

2024年NHTSA对Cruise的同意令涉及2023年10月事故后移动信息的报告问题。文件记载初期叙述存在遗漏，也记载相关视频已经交付，因而不能简化为全部资料被隐藏。\href{https://www.nhtsa.gov/sites/nhtsa.gov/files/2024-09/cruise-consent-order-2024-web.pdf}{历史同意令}。该事件表明，事故后续过程的完整报告与责任履行需要单独评价，车队平均事故率不足以替代这些信息。

许可、测试与经营也应分别理解。California DMV的2026年公告报告了规定年度超过900万英里的公共道路测试，并明确脱离报告的比较限制。\href{https://www.dmv.ca.gov/portal/news-and-media/autonomous-vehicle-permit-holders-in-california-logged-more-than-9-million-test-miles-between-december-1-2024-and-november-30-2025/}{官方年度说明}。测试活动能够积累信息，但不等于同等规模的实际运输服务，也不直接测得社会净收益。

中国2023年的准入与上路通行试点框架，同样区分测试评估、产品准入、限定区域通行与运输经营要求。\href{https://ythxxfb.miit.gov.cn/ythzxfwpt/hlwmh/tzgg/xzxk/clsczr/art/2023/art_ae314a1df30e4ba1a49bd4ffa10186f6.html}{四部门历史通知}。本章将其作为制度阶段的例子，不据此推定后续每家企业的当前许可状态。

发现问题后，谁能限制运行范围、取得资料，乘客与非乘客怎样提出异议，都影响风险由谁承担。用户可以决定坐不坐，却不能替其他道路使用者作选择；只公布总里程，公众又算不出对应事件率。共同道路影响的人多于交易双方，公共规则因而须覆盖交易以外的承担者。

\section{安全、费用与公共影响构成不同评价}

安全、费用和道路影响构成不同评价。事件率须与运行条件和版本对应，费用须覆盖完整劳动与设备，道路评价须包括空驶、停靠和高峰影响。三类资料相互补充，不能用一种结果替代其余结果。



还须分别查监督、劳动参与与所有权。披露和责任让公共机构能监督道路，劳动者参与部署和转岗让其影响过渡安排。降价和覆盖扩大，先改善取得服务；生产者或公众若又有持续改变投资、劳动和剩余用途的权利，相应决定关系才改变。每项权利都有自己的范围，取得监督权并不等于取得全部资本所有权。

\section{运营和公共规则决定自动化怎样使人受益}

本章的核心判断是驾驶自动化与运输组织不能等同。限定车队的安全改善不直接证明全部服务降本，车队规模缩小也不直接证明道路负担减少。运营和公共规则通过调度、替代交通与责任分担，决定三类结果怎样联系。

车队希望多接可获利订单，道路、应急和监管却靠更广的人群维持，这就给基本矛盾分析提供了具体关系。竞争可以让乘客分享部分降本，服务和责任规则可以改善覆盖、费用分担。要进一步判断生产条件由谁支配，还须查投资、运行范围、资料与剩余的实际权限，不能把服务改善直接写成占有关系转化。

交通案例将收益范围扩大到交易之外：乘客的服务改善与其他道路使用者的负担必须同时评价。第十一章则检验服务进入之前的分离：患者有健康需要，也有有效疗法，却未必具备完成治疗的支付和诊疗条件。两章分别补上交易外部影响与未实现需要。

\section*{本章参考文献}

1. Kusano, K. D., Scanlon, J. M., Chen, Y.-H., et al. 2024. Comparison of Waymo rider-only crash data to human benchmarks at 7.1 million miles. \emph{Traffic Injury Prevention}, 25(sup1): S66–S77. \href{https://doi.org/10.1080/15389588.2024.2380786}{DOI:10.1080/15389588.2024.2380786}；本章读取\href{https://arxiv.org/html/2312.12675v3}{接受稿v3}.
2. Scanlon, J. M., Kusano, K. D., Fraade-Blanar, L. A., et al. 2024. Benchmarks for retrospective automated driving system crash rate analysis using police-reported crash data. \emph{Traffic Injury Prevention}, 25(sup1): S51–S65. \href{https://doi.org/10.1080/15389588.2024.2380522}{DOI:10.1080/15389588.2024.2380522}；本章读取\href{https://arxiv.org/html/2312.13228v2}{接受稿v2}.
3. Di Lillo, L., Gode, T., Zhou, X., et al. 2024. \emph{Do Autonomous Vehicles Outperform Latest-Generation Human-Driven Vehicles? A Comparison to Waymo's Auto Liability Insurance Claims at 25 Million Miles}. 本章使用2024年12月发布的作者公开稿，不推定其为已核对正式期刊版本。\href{https://waymo.com/research/do-autonomous-vehicles-outperform-latest-generation-human-driven-vehicles-25-million-miles/}{发布方研究页面及PDF}.
4. Kalra, N., \& Paddock, S. M. 2016. \emph{Driving to Safety: How Many Miles of Driving Would It Take to Demonstrate Autonomous Vehicle Reliability?} RAND RR-1478-RC. \href{https://doi.org/10.7249/RR1478}{DOI:10.7249/RR1478}.
5. International Transport Forum. 2015. \emph{Urban Mobility System Upgrade: How Shared Self-Driving Cars Could Change City Traffic}. OECD Publishing. \href{https://doi.org/10.1787/5jlwvzdk29g5-en}{DOI:10.1787/5jlwvzdk29g5-en}.
6. NHTSA. 2026. \emph{Standing General Order on Crash Reporting}. 官方动态页面，首次核查2026-10-07、复查2026-10-08；页面仍显示数据截至2026-08-17。\href{https://www.nhtsa.gov/laws-regulations/standing-general-order-crash-reporting}{原文及限制说明}.
7. NHTSA. 2024. \emph{Consent Order: Cruise; Standing General Order 2021-01 Reporting}. Announced 30 September 2024. \href{https://www.nhtsa.gov/sites/nhtsa.gov/files/2024-09/cruise-consent-order-2024-web.pdf}{官方PDF}.
8. California Department of Motor Vehicles. 2026. \emph{Autonomous Vehicle Permit Holders in California Logged More Than 9 Million Test Miles Between December 1, 2024 and November 30, 2025}. 20 February 2026. \href{https://www.dmv.ca.gov/portal/news-and-media/autonomous-vehicle-permit-holders-in-california-logged-more-than-9-million-test-miles-between-december-1-2024-and-november-30-2025/}{官方说明}.
9. Waymo. 2026. \emph{Advice, Not Control: The Role of Remote Assistance in Waymo’s Operations}. From the Road, 17 February 2026. \href{https://waymo.com/blog/shorts/advice-not-control-the-role-of-remote-assistance/}{公司原文}.
10. Office of Senator Edward J. Markey. 2026. \emph{Remote Back Seat Operators: Revealing the Autonomous Vehicle Industry’s Reliance on Human Remote Assistance Operators}. \href{https://www.markey.senate.gov/imo/media/doc/remote_assistance_investigation_report.pdf}{调查报告}.
11. U.S. Department of Transportation. 2021. \emph{Driving Automation Systems in Long-Haul Trucking and Bus Transit: Preliminary Analysis of Potential Workforce Impacts}. \href{https://www.transportation.gov/av/workforce/report}{官方发布页及报告}.
12. 工业和信息化部、公安部、住房和城乡建设部、交通运输部. 2023. 关于开展智能网联汽车准入和上路通行试点工作的通知. 工信部联通装〔2023〕217号，2023-11-17。\href{https://ythxxfb.miit.gov.cn/ythzxfwpt/hlwmh/tzgg/xzxk/clsczr/art/2023/art_ae314a1df30e4ba1a49bd4ffa10186f6.html}{官方全文}.
13. Waymo. 2024. \emph{Fleet Response: Lending a Helpful Hand to Waymo’s Autonomously Driven Vehicles}. 21 May 2024. \href{https://waymo.com/blog/2024/05/fleet-response/}{公司原文}.

\chapter{基因编辑：疗效突破与患者实际受益的分离}\label{ch:11}

\section{疗效扩大治疗可能，支付与服务组织决定覆盖谁}

基因编辑的患者收益由疗效和服务进入共同形成。本章以严重镰状细胞病的体细胞治疗为例，区分限定临床结果、完整疗程容量及支付与风险安排。疾病负担给出需要的尺度，评估、等待、退出和完成记录则说明需要怎样成为实际治疗。

Jinek等2012年的研究证明RNA引导Cas9的可编程能力，为后续应用提供基础。\href{https://doi.org/10.1126/science.1225829}{原论文}与具体疾病的临床结果属于不同证据。2023年12月8日FDA公告批准的Casgevy使用CRISPR/Cas9编辑，同日公告的Lyfgenia采用另一种遗传改造途径。\href{https://www.fda.gov/news-events/press-announcements/fda-approves-first-gene-therapies-treat-patients-sickle-cell-disease}{历史批准公告}确认具体应用进展，不能把两种治疗及适用条件合并。

医疗案例补充服务进入的分析：技术有效，患者仍可能因资格、支付或服务容量未完成疗程。后文先确定临床证据范围，再比较容量与预算，最后讨论责任和资料权限。不同层面的结论分别以相应材料支持。

\section{疗效比例须保留患者分母、终点与随访期限}

Frangoul等2024年发表的第三期单组、开放标签研究报告44人接受exa-cel；进入主要疗效分析的30人中，29人在至少连续12个月内没有严重血管闭塞危象，约为97\%，原论文95\%置信区间为83\%—100\%。主要疗效人群要求输注后至少16个月随访，这与12个月终点窗口是两个时间条件。\href{https://www.astct.org/LinkClick.aspx?fileticket=1_e_CLAGAN8\%3D\&portalid=0}{正式论文统计分析与表2}，\href{https://doi.org/10.1056/NEJMoa2309676}{DOI}

该试验纳入12—35岁患者，要求筛查前两年每年至少出现两次严重血管闭塞危象；这限定了研究样本的年龄与病情范围，不能把其响应比例推广给全部患病者（原论文第1650页，Methods）。

表11-1把接受治疗、进入分析和达到终点的人数列开。它们分处研究流程的不同位置，不能换一个分母便宣布同一项结果变了。

\begingroup\footnotesize\setlength{\tabcolsep}{4pt}\renewcommand{\arraystretch}{1.25}
\begin{longtable}{>{\raggedright\arraybackslash}p{\dimexpr 0.21875\textwidth-8.0pt\relax}>{\raggedright\arraybackslash}p{\dimexpr 0.37500\textwidth-8.0pt\relax}>{\raggedright\arraybackslash}p{\dimexpr 0.40625\textwidth-8.0pt\relax}}
\toprule
\textbf{数字或条件} & \textbf{准确含义} & \textbf{在本章中的用途} \\
\midrule\endfirsthead
\toprule
\textbf{数字或条件} & \textbf{准确含义} & \textbf{在本章中的用途} \\
\midrule\endhead
\bottomrule\endfoot
44人 & 接受研究治疗的人数 & 描述已接受干预的规模 \\
30人 & 达到主要疗效分析要求的人群 & 作为29名响应者的分母 \\
29人 & 在该分析中满足所列主要终点 & 计算对应响应比例 \\
29/30\ensuremath{\approx}96.67\% & 舍入后报告为97\% & 限定人群和观察窗口的结果 \\
至少16个月随访 & 主要疗效分析所需条件 & 与终点本身的连续12个月分开 \\
95\%CI：83\%—100\% & 原研究报告的不确定范围 & 保留样本量与估计精度信息 \\
\end{longtable}\endgroup


\textbf{表11-1　临床证据的流程与分母。} 来源：Frangoul等（2024），研究资料截止2023年6月14日；本表没有重新分析个体数据，也未把单组结果变成随机对照效应。

如果直接将29除以44，得到的65.91\%并不是可以替换主要疗效结果的新估计。输注人群与主要疗效分析人群的随访要求不同，尚未进入该分析的人不能仅根据人数差额被判为治疗失败或成功。临床比例的解释因此须同时保留患者进入分析的标准、观察期限和终点定义。

该比例确认观察窗口内的疗效进展，持续性与相对其他治疗的效果仍需相应研究。连续12个月终点不等于终身结果，单组设计也不提供随机对照差异。监管与论文的不同截止版本分别保留，避免混用分母。

开发企业参与研究设计和分析，原文披露了这些角色。评价资料时须查谁生成、谁核验、是否完整报告，不能仅凭企业参与便否定观察结果，也不能把开发者报告当作已获独立核验。进一步判断要靠透明报告、完整随访与独立复核。本章只使用该研究公开的终点和患者资料。

\section{制造、诊疗与随访共同限制治疗容量}

EMA产品概览说明，此类产品按患者个体制造，并依赖具备相应能力的治疗中心、专业人员及后续监测。\href{https://www.ema.europa.eu/en/medicines/human/EPAR/casgevy}{官方概览}。其中“一次给药”描述输注步骤的频次，完整治疗则仍包含多个阶段的专业劳动与机构衔接。

流程还包括造血干细胞动员与采集、个体产品制备、输注前的白消安（busulfan）清髓预处理，以及输注后的恢复支持和监测。预处理与恢复阶段需要专业诊疗并具有相应风险，治疗中心、床位、人员和持续随访因而都是完成疗程的必要投入，治疗容量不能只按编辑产品数量推算（Frangoul等，2024，第1650页；上述EMA概览的使用与风险说明）。

完整疗程需要制造、临床与随访能力衔接。制造支持60例而临床只支持40例时，增加产品不会使完整治疗同比增加；临床扩容而其他环节不足时亦然。容量评价针对限制疗程的环节，不按新增投资额判断效果。

设已经符合临床条件且愿意在本期接受治疗的人数为\(D\)，制造、临床和随访的容量分别折算为每年可支持的完整治疗人次\(K_m,K_c,K_f\)，当年可用于所列治疗支出的预算为\(B\)，每人财务支出为\(c_t>0\)。在同质患者、相同期间及既定服务流程的简化下，完成治疗的人数\(q\)满足

\[
q\leq \min\left\{D,K_m,K_c,K_f,\left\lfloor B/c_t\right\rfloor\right\}.
\tag{11-1}
\]

向下取整表示人数为整数。式（11-1）是能力上界，达到它还需阶段之间及时衔接、患者能够完成流程及人员设备可用。它没有假定现实医疗可以像标准件一样无摩擦排产，也没有给任何真实患者作资格判定。

其中随访容量应先扣除既有队列责任，再按可为当年新增患者预留的完整路径折算；不能把全部随访资源视为都可分配给新患者。年度预算B与每人支出对应本期现金需要，后续年份的持续服务还需另有资金安排。

容量要按能支持多少完整疗程来算。中心能评估许多人，不等于能同样多地完成治疗与随访；患者需要的人员、时间和设施也会不同。实际核算须沿各阶段记录资源占用与衔接，中心数量不能直接相加作治疗能力。本节的简化关系只用来说明，价格或某项容量改变之后，哪项投入仍挡住完整疗程。

\subsection{基因编辑重组长期诊疗劳动，疗效不等于工资节省}

基因编辑可能减少反复治疗需要，却把一部分劳动集中到筛查、细胞采集、个体制造、预处理、输注和随访。对资本经营的药品制造与医疗服务，相关劳动力购买构成应分别考察的投入；公共医院预算、家庭照护与患者支出则属于不同账户，不能全部合成某家制药企业的可变资本。

同一患者的完整路径可以检验净劳动变化：新疗程的全部阶段工时，与原治疗在相同随访期限、健康结果和风险要求下的工时比较。一次输注不意味着一次劳动，减少未来住院也不意味着当前工资立即下降。专业人员保留并接诊其他患者时，减少的是每位患者所需劳动，机构工资总额可以不变；扩大治疗覆盖时还可能增加。

本章式（11-1）进一步说明，制造和临床瓶颈会限制这种扩张。实验人员、生产技术人员和临床人员不能仅按人数互换，稀缺技能及训练周期会影响报酬和可用工时。研发费用也不能全额等同于科研工资，治疗标价更不能代替各阶段工资。现有产品流程与容量资料支持确认互补劳动，尚未给出完整工资节省额。

长期健康改善还可能减少家庭照护和收入中断，这是社会再生产收益，不应只折成企业少付工资。若公共部门资助治疗、另一机构以后减少支出，付款与收益的分离也会改变采用激励。技术扩展了可以完成的治疗，资本对各类专业劳动力的购买怎样调整，仍需患者路径、机构工资和跨期付款相互核对。

治疗标准化也可能改变专业劳动的组织。相同人员照护更多患者，可能来自流程改进，也可能来自减少每位患者的沟通、恢复支持或随访。前者可能提高效率，后者可能牺牲服务质量；不能只以人均治疗量评价。对资本经营机构，工资购买同患者量和完整服务工时结合，才能解释可变资本变化；对公共服务机构，则按公共预算和质量责任考察。健康产出与劳动使用的分离，限制了将医疗技术一概写成节省人工的办法。

\section{降价减轻单次支出，扩容补服务短缺}

设一个假想体系每年有70人符合所列条件并愿意进入服务，制造能力60人、临床能力40人、随访能力80人；年度治疗预算为100个预算单位，每人完整的本期财务支出为2单位。所有费用都是假设尺度，与任何真实产品报价无关。

\begingroup\footnotesize\setlength{\tabcolsep}{4pt}\renewcommand{\arraystretch}{1.25}
\begin{longtable}{>{\raggedright\arraybackslash}p{\dimexpr 0.31707\textwidth-8.0pt\relax}>{\raggedright\arraybackslash}p{\dimexpr 0.17073\textwidth-8.0pt\relax}>{\raggedright\arraybackslash}p{\dimexpr 0.17073\textwidth-8.0pt\relax}>{\raggedright\arraybackslash}p{\dimexpr 0.17073\textwidth-8.0pt\relax}>{\raggedright\arraybackslash}p{\dimexpr 0.17073\textwidth-8.0pt\relax}}
\toprule
\textbf{假设变化} & \textbf{每人财务支出} & \textbf{临床容量} & \textbf{预算可支持人数} & \textbf{完整治疗能力上界} \\
\midrule\endfirsthead
\toprule
\textbf{假设变化} & \textbf{每人财务支出} & \textbf{临床容量} & \textbf{预算可支持人数} & \textbf{完整治疗能力上界} \\
\midrule\endhead
\bottomrule\endfoot
基线 & 2.0 & 40 & 50 & 40 \\
仅降低产品净价，使完整支出降至1.5 & 1.5 & 40 & 66 & 40 \\
临床容量扩大至65，其他费用不变 & 2.0 & 65 & 50 & 50 \\
降价与临床扩容同时完成 & 1.5 & 65 & 66 & 60 \\
\end{longtable}\endgroup


\textbf{表11-2　服务链约束的假设演算。} 制造60、随访80和需求70保持不变；取式（11-1）的上界。扩容建设费用需另行筹集，本表只考察扩容完成后的运行约束，不将新增设施当作免费资源。

仅降价，预算可支持人数由50升至66，临床却仍只支持40。少付钱有益，但本例上界没提高。仅扩容，临床限制放松，预算又把上界卡在50。两项一起改，上界才到60，接下来受制造能力限制。这里列的是能做到多少的上界，尚未证明实际完成了多少。

表11-2区分预算改善与容量改善。降价未增加治疗上界，不表明价格无作用，而表明当前临床约束未解除；扩容后的预算约束也不能靠床位解决。进入、等待、退出与完成记录，是识别措施作用位置的必要资料。

建中心还要能持续开起来。人员、病例或预算不足，名义容量就未必成为治疗。集中服务可以共用专业能力，却可能让患者多花路费和照护时间；分散服务可以近些，却可能提高固定费用，难以维持运营。须同时查质量、距离、规模与公平，比较哪种组织能持续服务需要的人，不能以中心越多便越好。

\section{健康需要须经过资格、支付与供给才能形成治疗}

GBD2021研究估计，2021年全球镰状细胞病患病人数约774万，95\%不确定区间为651万—920万。这是基于模型与资料综合的疾病负担估计，提供人口需要的尺度。\href{https://repository.up.ac.za/server/api/core/bitstreams/fd57c857-9ec2-4897-b38c-ff955abb4bb0/content}{正式论文e588—e589页}说明了方法；期刊随后更正图3、图5的部分国家着色及附录，本章不据相关地图估计覆盖，也不把更正理解为宣布全球人数发生变化。\href{https://doi.org/10.1016/S2352-3026(23)00215-6}{更正通知}

患病人数与式（11-1）的\(D\)并不相同。临床资格、替代治疗、患者意愿、信息、转诊与支付共同影响实际接受治疗的人群。即使需要很大，实际购买和治疗人数仍可能较小；将全球患病人数直接乘以产品标价，既不是现实预算，也不是企业能够实现的市场收入。

可以用一个条件概率链表示进入过程。设某地区已按具体标准确定的合格人群为\(N_e\)，完成评估、取得支付、选择治疗和完成流程的条件概率依次为\(p_1,p_2,p_3,p_4\)，则在所定义的同一人群与时期中，预期完成数为

\[
E[N_{\mathrm{complete}}]=N_e\,p_1p_2p_3p_4.
\tag{11-2}
\]

每个\(p_j\)都以已经通过前面步骤为条件，因而这个连乘不要求各阶段独立。现实容量不足会影响相应条件概率；如果还单独施加式（11-1），须明确两种表示怎样对应，避免重复计算同一障碍。

式（11-2）的各阶段概率须使用相应的患者分母。只知道最终治疗人数，无法区分评估不足、资格不符、支付失败、患者选择和服务中断。不同障碍要求不同措施：信息和转诊改善针对前端，公共采购与保障针对支付，医疗扩容与家庭支持针对后续。实际治疗需求的规模由这些准入、支付和供给安排共同形成，单独以市场规模不足概括上述障碍，无法解释其具体来源。

需要与有效支付的分离还会进入研发激励。预期回报、失败风险和资金期限影响疾病选择，公共保障或采购可以改变收益预期并支持原先无力支付者。扩大购买改变需求，公众能否影响研发投向则涉及另一类权限，二者分别调查。

\section{财务支付、资源投入与健康收益分别核算}

支付方付出的钱，进入销售方收入；全额价格不能再当作社会耗掉的资源。里面可能有研发回收、制造费用、风险补偿和经营收益，须用实际资料拆开。没有净价格、成本与合同，不能只看标价就推算利润，更不能把全部销售收入都算作超额所得。

疗程资源包括医护、设施、交通、照护与随访。支付减少可能来自投入节约，也可能来自转由家庭承担；厂家退款改变付款归属，并不返还已经使用的医疗劳动。因此，财务账户与资源账户按相同疗程和替代方案分别核算。

\begingroup\footnotesize\setlength{\tabcolsep}{4pt}\renewcommand{\arraystretch}{1.25}
\begin{longtable}{>{\raggedright\arraybackslash}p{\dimexpr 0.18919\textwidth-8.0pt\relax}>{\raggedright\arraybackslash}p{\dimexpr 0.37838\textwidth-8.0pt\relax}>{\raggedright\arraybackslash}p{\dimexpr 0.43243\textwidth-8.0pt\relax}}
\toprule
\textbf{观察主体} & \textbf{主要财务流量} & \textbf{另需记录的资源与结果} \\
\midrule\endfirsthead
\toprule
\textbf{观察主体} & \textbf{主要财务流量} & \textbf{另需记录的资源与结果} \\
\midrule\endhead
\bottomrule\endfoot
产品提供者 & 销售净收入、许可费用、退款 & 开发、制造和质量保障投入 \\
医疗机构 & 治疗与随访支付 & 专业工时、设施占用、持续服务能力 \\
保险或公共支付方 & 前期支付、后期节省与返还 & 预算时间、其他服务的机会成本 \\
患者及家庭 & 自付、交通住宿和收入变化 & 照护时间、健康、生活质量与选择能力 \\
\end{longtable}\endgroup


\textbf{表11-3　治疗经济评价的账户边界。} 表中为核算框架，未用未经核实的费用给真实产品填值。

健康价值还需与预算能力区分。用\(\Delta G\)表示相对于明确替代方案增加的健康结果，\(\Delta C\)表示同一评价视角下的增量费用。若健康结果采用质量调整生命年，即将生存时间与健康相关生活质量结合的度量，在给定评价权重\(\lambda\)下，可以写出净货币收益

\[
NMB=\lambda\Delta G-\Delta C.
\tag{11-3}
\]

\(\lambda\)是经济评价的资源权衡尺度，不是患者生命的市场售价。NICE方法手册讨论增量比较、成本效果与健康收益表达，本章借此说明方法关系，没有移用其具体阈值为本案例作新的推荐。\href{https://www.nice.org.uk/process/pmg36/chapter/economic-evaluation-2}{官方方法手册第4章}

即使式（11-3）在某组假设下为正，大规模前期支付仍可能超出当年预算。反过来，治疗价格低且容易支付，也不保证其相对其他医疗用途具有更大价值。健康增益、资源机会成本与当期现金需求，分别对应健康评价、资源配置和预算约束，须使用相应资料比较。患者的劳动市场收入更不能作为健康价值的唯一权重，否则原本在就业市场不利的人会在医疗评价中再次处于不利位置。

\section{当前付款、以后由别人省钱，可能影响采用}

有的治疗先付大笔费用，多年后才少花其他医疗费。患者后来换了支付体系，当前付款者便未必取得后续全部节省。同一疗程，放到较长时期可能有利，放到原支付方账户却可能难以收回前期支出。时间和账户不一致，可以影响采用动力；要知道是否发生，还须查转保、付款与后续节省的实际记录。

设治疗在未来第\(t\)年减少可比医疗支出\(S_t\)，折现率为\(r\)。若患者每年继续由当前支付方承担责任的条件概率为\(s\)，并为说明机制暂设这一概率不随年份改变，则原支付方预期取得的节省现值为

\[
PV_{\mathrm{payer}}=\sum_{t=1}^{T}\frac{s^t S_t}{(1+r)^t}.
\tag{11-4}
\]

在对医疗支出节省采用相同定义、并暂不计死亡及其他退出机制的简化下，覆盖所有后续支付方的金额现值为\(\sum_tS_t/(1+r)^t\)。前者较小，并不是治疗健康价值下降，而是财务收益跨主体转移。

举一个只为说明时序的假设例子：每年节省1个预算单位，持续10年，折现率3\%。若始终由原支付方负责，现值约为8.53；若每年延续概率为80\%，原支付方可归属的预期现值约为3.20。前期费用若不变，原支付方的采用动力就会不同。这里的概率与节省均未以真实患者或保险数据校准。

分期缓解当前现金约束，跨支付方分担重新配置责任，均不自动减少治疗资源。患者转保后的保障、后续付款与承诺执行须明确，否则初期支付降低仍可能伴随服务中断。

\section{承诺退款须有记录与执行，疗效风险才会转移}

NICE在2025年发布的TA1044中，对规定条件下的exa-cel采用有管理准入安排，并配套后续资料收集。\href{https://www.nice.org.uk/guidance/ta1044/chapter/1-recommendations}{推荐页面}，\href{https://www.nice.org.uk/guidance/ta1044/resources/managed-access-agreement-pdf-15250968973}{管理准入协议}。它将取得服务与继续评价联系起来，制度是否改善覆盖与长期结果，仍须观察实施资料。

美国CMS的CGT Access Model由相关州和厂家自愿参加，协助谈判结果挂钩安排，初始聚焦镰状细胞病基因治疗。\href{https://www.cms.gov/priorities/innovation/innovation-models/cgt}{官方项目页面}说明其设计与目标；启动和参加名单本身不是效果估计。该安排为分析风险如何重新分担提供具体制度背景。

该项目将协调写成具体义务：CMS与厂家谈判折扣和结果返还，州选择是否签约，参加后须执行约定的准入政策；厂家还须承担规定范围的生育力保存服务及其交通、住宿支持。治疗之外这些费用谁出，会影响患者能否走完整个流程。公开设计说明各方承诺做什么，能否增加完成治疗、减轻家庭负担，还要查实施资料。（CMS，上述项目页面，Design及Gene Therapy for Sickle Cell Disease）

考虑一份简化合同：先支付\(P\)，若未达到规定结果且满足退款条件，返还\(R\)。令合同所定义结果成功的概率为\(p_s\)，未成功后能够满足记录与执行条件并得到退款的概率为\(d\)。则预期净支付为

\[
E[P_{\mathrm{net}}]=P-(1-p_s)dR.
\tag{11-5}
\]

\(p_s\)是合同特定终点的概率，不能直接填入任何一项不同终点或不同人群的临床比例。\(d\)反映随访记录与退款执行对经济结果的影响：书面退款金额相同，记录完整性或执行程度不同，实际转移给提供者的风险也会不同。

假设\(P=100\)、\(R=50\)、\(p_s=0.8\)。退款执行充分时，预期净支付90；若\(d=0.5\)，则为95。全部是预算单位的假设算例，与真实合同无关。两种情况下，前期仍先支付100，若退款在以后发生，当年现金约束不会自动降到90或95。

结果挂钩合同只转移约定且可执行的部分财务风险。它不能替代治疗容量或家庭支持，也不消除患者已承受的健康风险。终点选择、随访信息与争议权限决定返还是否实现，书面退款金额不足以说明实际风险分担。

疗程结束后，患者、护理人员和机构还要持续提供信息，后续疗效判断与知识积累靠这些工作。资料可以改善公共判断，也可能集中在少数机构手里。贡献了信息的人能否知道用途、查错并提出纠正，公共机构能否核验，须落实到访问和使用权限。共同提供资料，没有自动形成共同监督。

\section{放开许可与建设设施，分别解除不同限制}

治疗供给由知识发现、临床研究、个体制造和专业诊疗共同形成，产品则通过具体许可与准入安排投入使用。排他安排可以支持开发费用回收，也可能限制替代提供者；其作用取决于期限、可替代技术及公共投入。WHO2021年建议将知识产权、登记、国际协作和公众参与列为治理议题。\href{https://www.who.int/publications/i/item/9789240030381}{官方概览}提供相关权利事项，实际影响仍由执行决定。

制造与诊疗已具备，采购或许可挡住使用，解除这些限制可能较快增加治疗。若缺人员、床位与随访，扩大使用资格也不能立刻完成更多疗程，仍须建设并维持运营。要把知识变成治疗，设施本来就是必要条件。由此既要查权利有没有放开，也要查放开后有没有真实能力使用。

公共出资建了什么、承担什么责任，也须查取得什么权利。建设费形成设施，运行保障支持人员和随访，稳定采购支持供给延续，作用不同。公共多担风险，却不能查资料、监督准入或影响用途，支持可能主要改善提供者回报。若接收方能维护流程、替换供应、培训人员并参与改进，才可能积累持续能力。能否替换还要看技术、费用和许可，不能只凭合同写了选项。

采购增加当前服务机会，能力建设扩大持续供给与替代的范围。两者应以完成治疗、运营持续性和替代提供者能否接手分别评价，项目数和降价不能代替这些结果。

\section{患者参与须能影响终点、资料用途与服务安排}

患者接受服务，也参与形成临床知识。评估花时间，治疗有经历，长期随访要持续提供资料，这些都有实际付出。但承认贡献，不等于把所有权利换算成股份。身体自主、研究同意、资料用途和经营收益各有对象，须分别规定，不能用一项经济补偿代替其他权利。

患者能决定自己是否接受治疗，未必能决定公共预算；公共机构能配置预算，也不能因此越过患者同意。个人选择和集体配置各有权限。服务组织要把两者接起来，使患者在了解风险与选择的条件下决定，同时公开公共资源怎样分配。

参与有没有作用，要查意见怎样进入决定。研究终点、交通支持、随访与争议处理是否因意见而改变，代表如何产生、资料怎样取得，都应有记录；办了多少场咨询，只说明举行过活动。患者对风险、时间与支持的看法还会不同，不能由少数代表发言便推定所有人同意。

体细胞治疗案例也不能代替可遗传编辑的治理判断。WHO的治理框架将这些范围区分讨论。\href{https://www.who.int/publications/i/item/9789240030060}{官方框架概览}。不同干预影响的群体和时间范围不同，现有患者的同意及支付不能代表未来所有受影响者。本章临床比例与治疗容量模型的适用范围，限于所列体细胞治疗及相应患者群体。

服务改善与治理变化分别评价。报销和扩容可以增加治疗机会，患者监督资料用途或公众影响研发投资则需要明确权限。实际进展应按已经改变的范围确认，避免将咨询活动等同于决策权转移。

\section{疗效成为患者受益，须接通支付、服务与责任}

本章说明疗效、支付和服务容量共同形成患者收益。临床比例限定于研究人群与随访，容量和预算模型说明两类约束不能替代，支付时序与合同执行解释部分财务风险怎样分担。治疗机会改善与研发、投资权限变化仍属不同结果。



患者提供资料，公共设施支持诊疗，长期知识积累支持产品，但准入、资料用途和研发投向仍由具体组织决定。这些权限与共同贡献怎样对应，就是本章要查的生产关系。需要的人得不到治疗，先查缺容量、缺预算，还是使用被排他条件挡住，再分别讨论扩容、保障与权利调整。进一步改变占有关系，须看谁能持续决定生产条件和目标，降价或举行咨询均不能单独证明。

医疗案例说明，有效技术与实际服务之间存在进入条件，支付和容量不能互相替代。第十二章进一步把“缺少能力”与“无权取得已有能力”分开：工序短缺限制制造，企业合同和国家许可规定交付资格。四个行业的比较由此形成对不同限制来源的辨析。

\section*{本章参考文献}

1. Jinek, M., Chylinski, K., Fonfara, I., Hauer, M., Doudna, J. A., \& Charpentier, E. (2012). A programmable dual-RNA-guided DNA endonuclease in adaptive bacterial immunity. \emph{Science, 337}(6096), 816–821. \href{https://doi.org/10.1126/science.1225829}{DOI}.
2. U.S. Food and Drug Administration. (2023, December 8). FDA approves first gene therapies to treat patients with sickle cell disease. \href{https://www.fda.gov/news-events/press-announcements/fda-approves-first-gene-therapies-treat-patients-sickle-cell-disease}{公告}.
3. Frangoul, H., Locatelli, F., Sharma, A., et al. (2024). Exagamglogene autotemcel for severe sickle cell disease. \emph{New England Journal of Medicine, 390}(18), 1649–1662. \href{https://doi.org/10.1056/NEJMoa2309676}{DOI}；\href{https://pubmed.ncbi.nlm.nih.gov/38661449/}{摘要}.
4. European Medicines Agency. (2024; webpage accessed 2026-10-07). Casgevy: EPAR medicine overview. \href{https://www.ema.europa.eu/en/medicines/human/EPAR/casgevy}{官方概览}.
5. GBD 2021 Sickle Cell Disease Collaborators. (2023). Global, regional, and national prevalence and mortality burden of sickle cell disease, 2000–2021: A systematic analysis from the Global Burden of Disease Study 2021. \emph{The Lancet Haematology, 10}(8), e585–e599. \href{https://doi.org/10.1016/S2352-3026(23)00118-7}{DOI}；\href{https://repository.up.ac.za/server/api/core/bitstreams/fd57c857-9ec2-4897-b38c-ff955abb4bb0/content}{作者机构所存正式论文}；\href{https://pubmed.ncbi.nlm.nih.gov/37331373/}{摘要}.

Correction to Lancet Haematol 2023; 10: e585–99. (2023). \emph{The Lancet Haematology, 10}(8), e574. \href{https://doi.org/10.1016/S2352-3026(23)00215-6}{DOI}.
6. National Institute for Health and Care Excellence. (2025, February 26). Exagamglogene autotemcel for treating severe sickle cell disease in people 12 years and over. Technology appraisal guidance TA1044. \href{https://www.nice.org.uk/guidance/ta1044/chapter/1-recommendations}{推荐}.
7. National Institute for Health and Care Excellence. (2025). TA1044: Managed access agreement. \href{https://www.nice.org.uk/guidance/ta1044/resources/managed-access-agreement-pdf-15250968973}{协议}.
8. Centers for Medicare \& Medicaid Services. (2026, September 30 update). Cell and Gene Therapy Access Model. \href{https://www.cms.gov/priorities/innovation/innovation-models/cgt}{项目}.
9. World Health Organization. (2021a). \emph{Human genome editing: Recommendations}. ISBN 978-92-4-003038-1. \href{https://www.who.int/publications/i/item/9789240030381}{出版概览}.
10. World Health Organization. (2021b). \emph{Human genome editing: A framework for governance}. ISBN 978-92-4-003006-0. \href{https://www.who.int/publications/i/item/9789240030060}{出版概览}.


National Institute for Health and Care Excellence. (2022, dynamic manual accessed 2026-10-08). NICE technology appraisal and highly specialised technologies guidance: the manual (PMG36), Chapter 4 Economic evaluation. \url{https://www.nice.org.uk/process/pmg36/chapter/economic-evaluation-2} 。本章引用增量经济评价方法，未引用具体当期支付阈值。

\chapter{半导体：专业协作与关键接口控制}\label{ch:12}

\section{专业协作形成芯片能力，关键接口影响制造选择}

半导体依靠设计、制造、封装测试与维护的专业协作。分工减少重复建设，互补工序限制交付，企业合同和国家许可规定取得条件。本章区分工序能力、费用与准入权限，分析专业化怎样同时形成效率和依赖。

本文以EUV光刻体系为重点。EUV使用13.5纳米波长，ASML的NXE与EXE系列数值孔径分别为0.33和0.55。波长、光学参数和商业节点各有含义，更先进光学支持更细图形，但具体产品还需要稳定工艺和经济可行性。\href{https://www.asml.com/en/products/euv-lithography-systems}{ASML技术说明}为设备能力提供依据，不能将13.5纳米当成芯片节点下限。

工序模型比较产能、良率和订单，成本与议价模型比较固定投入及外部选择，原始资料定位具体依赖。技术不可行、容量不足和权利受限需要不同解释，即使同时出现在一次断供中，也不能合并成单一原因。

\section{光刻更快，其他工序接不上，交付仍会受限}

设研究对象是一条在一个星期内运行的、生产同类产品的简化工艺链。第j道工序可提供K\_j个工序操作单位，每片投产晶圆需要r\_j个该工序单位。K\_j可以分别采用曝光次数、处理批次或测试次数，但这些数不能直接比较。先将它们换算为等效晶圆投产能力L\_j=K\_j/r\_j，单位才统一为片/周。假定工艺配方固定、排程已经协调且不考虑库存暂时缓冲，则全链最大投产量为：

\[
B=\min_j\{L_j\}=\min_j\{K_j/r_j\}.\tag{12-1}
\]

式（12-1）是互补投入的最小值约束，不是对现有晶圆厂的估计。在所列假设下，全链持续投产受到处理能力最低的必要工序约束。光刻能力增长，若刻蚀或检测已成为限制，新增能力可能先表现为闲置、排程弹性或较短等待，而不是交付增长。只有在光刻是唯一瓶颈且其他环节尚有余量时，少量提高L\_光刻才有dB/dL\_光刻=1；非瓶颈对应的局部导数为0。两项能力相等时，函数存在折点，不能在这个位置用一个光滑导数描述所有方向。

进一步区分投产数量与合格交付，设每片晶圆平均包含d个可测试晶粒，完成全链后的晶粒合格比例为y，0<y\ensuremath{\leq}1。潜在合格能力是ydB个/周；若该产品订单上限为D个/周，则本期实际投产与实际交付分别为：

\[
X=\min\{B,D/(yd)\},\qquad Q=ydX.\tag{12-2}
\]

这里将当期合格产品全部用于交付，并略去成品库存；实际企业可能在需求不足时积累库存，也可能采用不同运行策略。工程投产能力B、质量转换yd与实际交付Q须分别核算。不能把设备吞吐量叫作良率，也不能把合格能力叫作销量。合格率还要明确分母：测试图形、电学器件、晶粒与封装后产品不是同一种y，不能将局部测试比例逐项相乘，除非确实具有相应条件概率和完整流程资料。

下表只为演示式（12-1）—（12-2），各工序数字已经换成等效晶圆投产能力，取d=100个/片、y=0.8。它不对应ASML设备、客户工厂或具体节点。

\begingroup\footnotesize\setlength{\tabcolsep}{4pt}\renewcommand{\arraystretch}{1.25}
\begin{longtable}{>{\raggedright\arraybackslash}p{\dimexpr 0.16667\textwidth-8.0pt\relax}>{\raggedright\arraybackslash}p{\dimexpr 0.16667\textwidth-8.0pt\relax}>{\raggedright\arraybackslash}p{\dimexpr 0.16667\textwidth-8.0pt\relax}>{\raggedright\arraybackslash}p{\dimexpr 0.16667\textwidth-8.0pt\relax}>{\raggedright\arraybackslash}p{\dimexpr 0.16667\textwidth-8.0pt\relax}>{\raggedright\arraybackslash}p{\dimexpr 0.16667\textwidth-8.0pt\relax}}
\toprule
\textbf{解释情景} & \textbf{光刻L（片/周）} & \textbf{其他必要工序最低L（片/周）} & \textbf{B（片/周）} & \textbf{订单D（个/周）} & \textbf{实际Q（个/周）} \\
\midrule\endfirsthead
\toprule
\textbf{解释情景} & \textbf{光刻L（片/周）} & \textbf{其他必要工序最低L（片/周）} & \textbf{B（片/周）} & \textbf{订单D（个/周）} & \textbf{实际Q（个/周）} \\
\midrule\endhead
\bottomrule\endfoot
原生产链 & 100 & 80 & 80 & 10000 & 6400 \\
仅光刻提升 & 150 & 80 & 80 & 10000 & 6400 \\
瓶颈工序提升 & 100 & 100 & 100 & 10000 & 8000 \\
需求不足 & 100 & 100 & 100 & 6000 & 6000 \\
\end{longtable}\endgroup


局部提速并非没有收益。更快完成曝光可能减少返工、释放维护窗口，或让相同设备支持更复杂产品；这些作用在固定配方的最小值公式中尚未体现。若改善使r\_j下降，而非K\_j上升，也会提高L\_j。若减少缺陷而使y提高，则即使B不变也会增加合格能力。技术的产出效果须分别对应配方、吞吐量、停机时间和最终质量，不能以一个笼统的生产率比例替代这些机制。

imec于2025年9月22日发布的High NA EUV进展，涉及特定金属线间距、单次曝光和电学测试，依赖ASML—imec实验室及NanoIC设施。它说明局部图形、材料和工艺整合确有推进，不能直接替换式（12-2）中整颗商业芯片的y，更不能据该试验推算全球Q。\href{https://www.imec-int.com/en/press/imec-achieves-new-milestones-single-patterning-high-na-euv-lithography-both-damascene-and}{原始研发发布，正文及致谢}

\section{固定费用要靠合格产出分摊，还须有订单}

对$X>0$的实际投产，设同一时期固定设备服务、必要研发分摊与固定维护费用合计$F_M$，每片投产晶圆变动费用为$v$。这里的$v$是单位变动成本，包含哪些材料、能源与劳动须按成本口径列明；它不同于第二章的可变资本$v_M$。$F_M$是本期成本，不是全部建厂现金支出；假定无其他产品分担费用，则合格产品平均成本为：

\[
AC=\frac{F_M+vX}{ydX}=\frac{v}{yd}+\frac{F_M}{ydX}.\tag{12-3}
\]

保持产品和y不变，增加实际投产可以摊薄固定成本：\ensuremath{\partial}AC/\ensuremath{\partial}X=\ensuremath{-}F\_M/(ydX\ensuremath{^2})。保持X不变，提高最终合格率则有\ensuremath{\partial}AC/\ensuremath{\partial}y=\ensuremath{-}(v+F\_M/X)/(dy\ensuremath{^2})。后一式同时影响变动成本和固定成本的分母，因此质量改善具有实际经济价值。不过两种导数都是其他条件不变时的局部比较；改善良率所需的额外检测、材料或研发，仍要加入分子，不能只提高分母后宣称净成本必然下降。

在一组未以行业资料校准的成本假设中，若F\_M=100万元/周，v=1000元/片，d=100，y=0.8，实际投产80片时，AC=168.75元/个；投产100片时为137.50元/个。若能力仍为100片但订单只够75片，AC反而是179.17元/个。设备更先进与每件实际售出产品更便宜并不等价，需求不足会使固定投入无法充分分摊。这里的成本没有设定商品价格，也不能单独给出利润或消费者获益。

ASML年报披露了公司层面的收入与研发费用。2025财年公司总净销售约327亿欧元，研发费用约47亿欧元，采用US GAAP；年报于2026年2月25日发布。按报告显示的取整数字计算，研发费用/净销售约为14.4\%。这只是同一公司同一财年的费用强度，不能称为EUV独立研发比例，也不是研发投入社会收益率。费用包括当期会计确认的项目，并不等于本年全部现金支付或资本支出。\href{https://www.asml.com/investors/annual-report/2025/financials}{ASML财务页}，\href{https://investor.asml.com/annual-reports}{年报发布与会计口径}

\begingroup\footnotesize\setlength{\tabcolsep}{4pt}\renewcommand{\arraystretch}{1.25}
\begin{longtable}{>{\raggedright\arraybackslash}p{\dimexpr 0.22642\textwidth-8.0pt\relax}>{\raggedright\arraybackslash}p{\dimexpr 0.32075\textwidth-8.0pt\relax}>{\raggedright\arraybackslash}p{\dimexpr 0.15094\textwidth-8.0pt\relax}>{\raggedright\arraybackslash}p{\dimexpr 0.30189\textwidth-8.0pt\relax}}
\toprule
\textbf{真实数据或演算} & \textbf{时期与分母} & \textbf{结果} & \textbf{它解释的环节} \\
\midrule\endfirsthead
\toprule
\textbf{真实数据或演算} & \textbf{时期与分母} & \textbf{结果} & \textbf{它解释的环节} \\
\midrule\endhead
\bottomrule\endfoot
ASML总净销售 & 2025财年，公司全部业务，US GAAP & 327亿欧元 & 企业已经实现的总体销售收入 \\
ASML研发费用 & 同财年，公司R\&D costs & 47亿欧元 & 维持与开发技术能力的费用规模 \\
47/327\ensuremath{\times}100 & 上述公布取整金额之比 & 约14.4\% & 研发费用强度，不是投资回报率 \\
AC：80\ensuremath{\to}\allowbreak{}100片投产 & 原创情景，固定y及相同一周成本边界 & 168.75\ensuremath{\to}\allowbreak{}137.50元/个 & 交付规模使固定成本摊薄 \\
AC：100\ensuremath{\to}\allowbreak{}75片投产 & 原创情景，由订单不足造成 & 137.50\ensuremath{\to}\allowbreak{}179.17元/个 & 工程能力与成本实现之间的需求门槛 \\
\end{longtable}\endgroup


公司研发费用按会计范围确认，不能全部视为某一技术的固定准备成本。设计、工具、验证与售后属于不同阶段，资料复制也不能替代现场运行。网络核算中，上游投入若已通过设备售价转移，再与下游购机费用相加可能重复计量。

规模经济可以降低社会所需的重复投入，也可以提高新进入者需要跨越的销量门槛。若预期单件收入扣除变动成本后的贡献为g，采用准备费用为F\_A，且暂不考虑融资和风险，则需要预期合格销售量至少达到F\_A/g；g\ensuremath{\leq}0时，销量增加也无法回收准备费。真实采用决策须用现金流折现并处理试验失败，而不是把这个简化阈值当作产业预测。公共试验线、多项目流片或兼容工具若能降低F\_A，就可能降低低销量主体的采用门槛；这些主体能否继续扩大商业生产，还取决于需求、质量和运行费用。

\subsection{先进制程提高设备密度，专业劳动力购买未必减少}

半导体制造体现大量设备同专业劳动的互补。自动搬运和过程控制可以减少重复操作，同时更复杂的工艺、维护与良率管理可能增加对工程师和技术人员的需求。工资支付不仅取决于操作人数，还取决于技能组合、班次、训练和工厂实际扩产。设备先进化因而可能与可变资本绝对增加并存；在同口径不变资本增加更快时，其相对比重仍可下降。

工艺瓶颈还会限制劳动节约的实现。只提高光刻速度而其他工序不变，合格交付不增加，既有支持班组也未必缩减；最终良率提高，则可能减少每个合格芯片分摊的劳动，即使全年工资保持不变。式（12-3）中的单位变动费用$v$不提供这笔工资的独立数值，更不能用它替代$v_M$讨论价值增殖。

同样，ASML披露的研发费用含多种投入，不能全部计作研发人员工资；晶圆厂购机价格包含上游生产劳动与其他费用，却不是晶圆厂直接购买这些上游劳动力的支出。企业比较要区分本厂与供应商，生产网络比较则对各方工资分别求和，避免把设备采购和其中工资重复累计。

先进工厂的高设备投入还可能提高连续运行要求，停机损失、工艺变更和质量责任影响劳动组织。高技能工资上升与人员对关键系统依赖加深可以同时发生。因此，资本购买价格、可变资本比重与劳动者权限应分别观察；设备昂贵或工程师收入高，都不能独自说明劳动者已决定投资、工艺节奏和经营剩余的用途。

高技能劳动的增加还可能成为机器体系扩大运行的条件。工艺知识积累在企业数据库、设备和协作流程中，工程师虽可能取得较高工资，其工作仍可能受排产、良率目标和持续值班安排约束。需要分别考察知识和资格能否迁移、异常处置是否有停机权限、目标怎样修改。技术提高了专业劳动的生产作用，不自动把使用决定权交给专业人员；这使工资增长同劳动对资本组织的依赖加深能够同时出现。

\section{许可与交付权影响使用者的选择}

并非每个瓶颈都由权利人制造。罕见材料、长期训练或暂时不足的设备能力，也可以形成约束。但当取得关键环节必须经过少数可以决定许可、交付或维护的主体时，技术依赖便与组织权利结合。ASML2025年战略报告披露，Carl Zeiss SMT是关键光学部件的唯一供应者，并说明排他安排和产能限制风险。其意义是确认这家公司所报告的一项具体依赖；并不能据此确定所有半导体市场的法律边界。\href{https://ourbrand.asml.com/m/8ab959d4926657b/original/asml-2025-annual-report-strategic-report-section.pdf}{年报风险段，印刷第69页、摘录PDF第68页}

该披露指出了具体接口：专业研发与制造形成光学能力，排他供给安排又规定客户能从谁取得。供应者确有技术贡献，也可能取得排他优势。要分清各自带来多少收益，须查实际费用、合同和可用替代，不能把唯一供应者的全部收入都称为租金。

供应者的收费空间取决于使用者相对可行替代方案取得的额外收益，唯一供应者身份不足以单独确定收费水平。对一笔简化交易，设买方使用该节点、尚未支付接入费时，在同一时期所得净价值为A；采用最佳可行替代所得为B。两者已扣除相应其他成本。令G=A\ensuremath{-}B为使用该关系比退出多出的收益，供应方为完成该项供给需要真实成本k。若接入支付为t，买方相对退出的增益为G\ensuremath{-}t，供应方所得超过供给成本的部分为t\ensuremath{-}k。自愿可行交易要求：

\[
k\leq t\leq G,\qquad G\geq k.\tag{12-4}
\]

若G<k，单靠在双方间调整收费无法使交易同时合算，需要减少成本、改变需求或由其他主体提供支持。若双方按给定权重\ensuremath{\theta}分配可转移净收益，最大化(t\ensuremath{-}k)\textasciicircum{}\ensuremath{\theta}(G\ensuremath{-}t)\textasciicircum{}(1\ensuremath{-}\ensuremath{\theta})，则在0<\ensuremath{\theta}<1及G>k时得到：

\[
t=k+\theta(G-k).\tag{12-5}
\]

这里暂将供应者拒绝本次交易的净外部收益归一为零；真实供给成本k不包含另一客户可能提供的替代利润。若供应者有正外部收益，应从合作剩余另行扣除，并相应提高其接受交易的门槛。

对该对数目标求导可得：\ensuremath{\theta}/(t\ensuremath{-}k)=(1\ensuremath{-}\ensuremath{\theta})/(G\ensuremath{-}t)，整理即得式（12-5）。该式使用给定权重，既未估计ASML与Zeiss合同，也未将\ensuremath{\theta}视为供应者可自行决定的参数。其比较涉及三个不同机制：更可靠技术提高A，可能扩大共同剩余；替代能力改善提高B，可能降低节点的收费空间；合同、认证和退出条件影响权重，可能改变共同剩余如何分配。若G=100、k=20、\ensuremath{\theta}=0.5，则t=60；替代改善令G降至60时，t降至40。两次供应方真实供给成本相同，改变的是外部选择及可分配收益。

支付t是买方费用、供应方收入，双方合并核算后只剩k这类真实投入。把全部接入费当作社会耗费，会遗漏其收入去向；把收费全部看成无成本租金，又会抹去维护、培训和后续研发。实践中还可能有不同产品的交叉补贴、按使用量收费、知识产权约束或容量配给，此时要重写交易，不能将本节单笔议价公式直接放入第十四章固定成本的数量竞争均衡。

买了设备，能否改维护软件、取得工艺资料或再许可，还要查合同；供应商有知识产权，也不因此包办客户设计和现场决定。专用合作可以保护投入、提高质量，排他安排也可能增加换供应者的费用。客户数量少，还可能反过来约束供应方。究竟是哪种关系起作用，须查谁定条件、谁能修改、第二来源多久能认证、维护怎样响应、失败由谁承担，收入或研发总额均不能替代这些证据。

\section{知识取得与现场交付能力分别形成}

全球专业化节约了每家企业自行重建全部能力的费用。OECD2025年价值链研究把设计、晶圆制造、组装测试封装与材料设备联系起来；贸易分析主要覆盖2012—2022年，部分投入产出图使用2018年资料。分类无法完全分离所有半导体用途，出口地域也不是所有权地域。该报告可用于识别价值链关系，分析时仍须保留各图表的实际观察期，不能将2025出版日期作为全部数据年份。\href{https://www.oecd.org/content/dam/oecd/en/publications/reports/2025/06/mapping-the-semiconductor-value-chain_5ba52971/4154cdbf-en.pdf}{报告§2—3及附录C—D}

专业知识通过规则、接口与反复验证结合成产品。复制设计文件可以降低信息取得费用，现场能力仍需设备、训练与检测。生产社会化体现于这些投入的相互依赖，不仅体现于企业跨越的国家数量。

TSMC2025财年20-F的Design and Technology Platforms披露EDA流程、经硅验证的模块、PDK及Open Innovation Platform，并介绍多项目晶圆服务。其经济意义在于设计与制造之间可以共享准备、验证和工具合作，部分小规模原型不必独自承担整片晶圆全部成本。公司的“Open”不能替代客户合同，但也不应因商业平台有收费便否认它降低部分进入费用的可能。\href{https://investor.tsmc.com/sites/ir/sec-filings/2025_20F\%20Report.pdf}{公司原始申报，印刷第20页、PDF第26页}

SkyWater与Google的SKY130公开PDK构成另一种安排：部分工艺信息可以直接获得，从而降低学习和试验入口；项目状态文档同时将公开版本标为实验预览，不保证生产用途。教育设计、可靠量产和先进产品是不同目标，不能把130纳米开放项目与High NA流程作无条件成本排名。\href{https://www.skywatertechnology.com/open-source-asics-take-a-giant-leap-forward-with-first-ever-open-foundry-pdk/}{2020年原始发布}，\href{https://skywater-pdk.readthedocs.io/en/main/status.html}{项目状态}

开放了哪些资料，还要查使用者随后能做什么。设计准备费降了，若没有流片、封装或测试渠道，式（12-1）的其他限制仍挡住交付。兼容接口若让更多合格供应者接入，式（12-4）的外部选择B才可能提高。成果换个地方须重做全部验证，文件虽能分享，退出也未必更容易。因此，应沿资料取得、功能验证、制造申请、实际交付与成果保留，查具体少了哪项费用、解除了哪道限制。

SEMI的信息与控制委员会说明共同命令、数据采集和设备接口的范围。共同数据格式有助于信息交换，语义定义与校准影响数据解释，适配测试影响设备间协同，权限安排则规定使用者可执行的操作。标准由此可能减少协调成本，但各参与者实际取得的操作与修改权限仍取决于组织安排。本文只读取委员会与标准目录，未取得全部付费标准文本，不据此判断具体工厂符合性。\href{https://www.semi.org/en/products-services/standards/information_and_control}{SEMI官方目录}

接口可读之后，测试设备、人员与认证仍可能限制进入。共享测试和参考实现是否扩大选择，应由新增合格参与者、实际迁移与可用第二来源衡量。复制工厂、共享验证和改善恢复具有不同投入结构，需要分别比较。

\section{补贴转移费用，共享设施还可能减少重复投入}

共享试验、训练和测量，可以让多家企业少做重复准备；收益却未必全回到出资者。这是公共支持共同能力的一项理由。给商业工厂补贴，则先减少该企业自己付的钱，社会是否增加能力，还须查新增交付、可用替代与持续知识。两种支持都可能有价值，但只看受支持企业营业收入，分不清各自增加了什么。

Sandia2001年的发布记录国家实验室与EUV LLC的原型合作，列出Sandia、Lawrence Livermore和Lawrence Berkeley，并将1997至2002年初协议的项目研究经费列为产业出资。\href{https://newsreleases.sandia.gov/partners-unveil-first-extreme-ultraviolet-chip-making-machine/}{当年发布}表明公共实验室的人员设施与产业项目资金共同参与，不能将其中一方抹去。imec对NanoIC公共支持的致谢则确认另一阶段的共同试验条件；这些记录支持协作事实，没有给出各方全部历史贡献或利润分配比例。

美国商务部2024年11月向TSMC Arizona授予最高66亿美元直接支持，按里程碑安排。\href{https://www.nist.gov/news-events/news/2024/11/biden-harris-administration-announces-chips-incentives-award-tsmc-arizona}{当期公告}说明历史支持上限，不能当作已付款或当前交付。补贴降低私人固定支出，不自行减少社会建设资源；若由此形成新的互补能力与恢复路径，社会收益仍可能为正。共享试验设施则还可能降低多个项目的实际验证费用，两者应按不同作用核算。

共享能力还取决于时段申请、结果保留、维护和失败责任。公共支持若未改善访问与监督，可能主要改变受支持企业的财务条件；可迁移知识与持续准入则可能扩大其他生产者的选择。出资身份与项目名称不能替代运行和权限证据。

\section{国家许可改变可交付范围，须与企业产权分开}

商业上愿意交易，不意味着法律上可以交付。荷兰2025年1月的修订替换先进半导体制造设备管制附件，第III条规定2025年4月1日生效并设置既有许可过渡。附件涉及特定设备、软件和技术；当期官方解释说明相关出口至欧盟以外目的地采用逐案许可，许可要求本身并非全面禁运。本章将它作为已核原文的历史制度案例。\href{https://zoek.officielebekendmakingen.nl/stcrt-2025-1894.pdf}{正式法规第1页及附件}，\href{https://www.government.nl/latest/news/2025/01/15/klever-export-controls-on-advanced-semiconductor-manufacturing-equipment-to-be-tightened}{当期官方说明}

国家许可直接影响企业可取得的设备、软件和技术范围。若某种投入无法取得，工厂面对的并不只是一个更高价格，而可能是式（12-1）某项可用能力降低；若替代流程可行，则r\_j、良率、验证和研发费用都会变化。库存可以在有限时期缓冲，重新设计可能改变产品目标，本地维护可能改善恢复，但这些反应都有资源和时间成本。存在本地替代路线并不免除调整费用，短期供应瓶颈也不足以证明该投入将永久不可替代。

国家许可的效果需要连接规则、申请、决定与企业调整。出口变化还可能受需求、价格和周期影响，法规日期与产出同时变化不构成因果证明。本文只分析已核历史制度，不据此判断尚未完成后续法规核查的当前交易资格。

国家可能考虑安全、就业与供给中断，企业可能考虑客户、知识保护与长期收入；目标可以一致，也可以冲突。补贴和许可可以改变选址与风险承担，却没有自动把私人董事会的剩余使用权交给公众。工厂在哪、资产最终由谁控制、维护与技术仍依赖谁、公共担什么责任，要分别查，才能说明关系怎样变化。

\section{第二来源须真能接手，退出才成为实际选择}

遇到中断，替代体系须在需要的时间内交出合格产品，才有实际作用。两座工厂若仍用同一设备、材料或维护许可，也可能一起受影响；备用厂房缺人员和认证，也接不了订单。第二来源能运行，工艺与验证知识能保留并使用，才可能形成退出与恢复的选择。供应地点多了，本身没有证明这些条件已经具备。

备用能力具有成本。若其当期资本服务、训练和维护费用为C\_R，减少的预期中断损失为E[L\_0\ensuremath{-}L\_R]，其他可核查的进入收益为E\_J，则应比较E[L\_0\ensuremath{-}L\_R]+E\_J与C\_R，而不能把最大可能损失视为一定取得的收益。社会收益可能分散在客户及公共部门，单个提供者又未必愿意承担全部费用，这为共同出资和长期订单提供了协调理由。

共享测试可以少做重复准备，公开接口与可携带记录可以少花迁移、重验证费用，由此既可能增加供给，也可能改善替代选择。但有限设施若长期被少数老客户占满，小项目仍进不来。质量、费用、排队、申诉与规则修改权须一起查；公布文件、建成厂房各只完成一个环节。

专业协作怎样分配选择权，可以通过具体安排改变。进入和退出真的更容易，共享知识与设施便可能让更多生产者能做事；单一接口能持续单方改条件，多方形成的能力也可能受制于它。须看实际权限和可用替代，企业有多大、链条有多长，均没有单独说明谁能决定生产。

\section{工序、交易与国家许可形成不同限制}

本章以三类分析区分限制来源：工序互补规定可交付能力，合格产出与订单规定费用分摊，外部选择与合同规定交易空间。国家许可另行规定法律可交付范围。它们可以在同一接口交汇，仍具有不同证据和调整条件。

原始材料把关系落到具体事项：ASML披露单一光学供应与排他安排；公开PDK和多项目流片提供不同入口；公共实验室与产业资金共同支持技术；国家许可又规定法律上能交付什么。它们支持上述组织事实，尚未估计全产业的垄断收益或政策净效应。国家限制须与企业产权分开，产能移到本地，也没有自动让当地生产者取得决定权。

改善须对着实际限制做。缺验证与运行能力，可建设共享条件；缺替代，可查第二来源与迁移权；公共投入没有对应访问和责任，则要调整这些安排。进一步改变占有关系，还须看谁能决定设施、技术与剩余的用途。补贴多了、本地工厂多了，均不能单独说明支配权已改变。

四个行业共同说明，技术收益受交付条件、交易外部影响、服务进入和关键接口的不同限制。共同设施不足首先是生产问题；限制来自可排他的合同或许可时，才进一步涉及特定的权利关系。第十三章将比较从当期交付延长到持续生产：照护、训练和维护决定人和设施能否在下一期继续投入。

\section*{本章参考文献}

1. ASML. (2025财年；2026年发布). \emph{Annual Report 2025}. \href{https://www.asml.com/investors/annual-report/2025/financials}{官方Financials}，\href{https://ourbrand.asml.com/m/8ab959d4926657b/original/asml-2025-annual-report-strategic-report-section.pdf}{US GAAP战略报告摘录}。两入口属于同份年报。
2. ASML. (无年). \emph{EUV lithography systems}. \href{https://www.asml.com/en/products/euv-lithography-systems}{官方技术说明}。动态页面以访问日期定位。
3. imec. (2025, September 22). \emph{Imec achieves new milestones in single patterning High NA EUV lithography for both damascene and direct metal etch metallization processes}. \href{https://www.imec-int.com/en/press/imec-achieves-new-milestones-single-patterning-high-na-euv-lithography-both-damascene-and}{研发参与机构原始发布}。
4. OECD. (2025). \emph{Mapping the semiconductor value chain: Working towards identifying dependencies and vulnerabilities}. OECD Science, Technology and Industry Policy Papers, No. 182. Paris: OECD Publishing. \href{https://doi.org/10.1787/4154cdbf-en}{DOI：10.1787/4154cdbf-en}。
5. Taiwan Semiconductor Manufacturing Company Limited. (2025财年；2026年申报). \emph{Annual Report on Form 20-F for the year ended December 31, 2025}. \href{https://investor.tsmc.com/sites/ir/sec-filings/2025_20F\%20Report.pdf}{公司原始申报全文}。
6. SkyWater Technology. (2020, November 16). \emph{Open Source ASICs take a Giant Leap Forward with the First Ever Open Foundry PDK}. \href{https://www.skywatertechnology.com/open-source-asics-take-a-giant-leap-forward-with-first-ever-open-foundry-pdk/}{原始发布}。
7. SkyWater与Google项目文档. (无年). \emph{Current Status—Experimental Preview}. \href{https://skywater-pdk.readthedocs.io/en/main/status.html}{原始状态说明}。不将文档当前状态与基础工艺商业使用混合。
8. SEMI. (无年). \emph{Standards Information and Control}. \href{https://www.semi.org/en/products-services/standards/information_and_control}{委员会说明与标准目录}。
9. Sandia National Laboratories. (2001). \emph{Partners Unveil First Extreme Ultraviolet Chip-making Machine}. \href{https://newsreleases.sandia.gov/partners-unveil-first-extreme-ultraviolet-chip-making-machine/}{当时机构发布}。
10. U.S. Department of Commerce / NIST. (2024, November 15). \emph{Biden-Harris Administration Announces CHIPS Incentives Award with TSMC Arizona to Secure U.S. Leadership in Advanced Semiconductor Technology}. \href{https://www.nist.gov/news-events/news/2024/11/biden-harris-administration-announces-chips-incentives-award-tsmc-arizona}{当时官方公告}。
11. Minister voor Buitenlandse Handel en Ontwikkelingshulp. (2025). \emph{Regeling van 13 januari 2025, BZ2411160, tot wijziging van de Regeling geavanceerde productieapparatuur voor halfgeleiders}. \emph{Staatscourant}, 2025, 1894. \href{https://zoek.officielebekendmakingen.nl/stcrt-2025-1894.pdf}{正式法规}。
12. Government of the Netherlands. (2025, January 15). \emph{Klever: export controls on advanced semiconductor manufacturing equipment to be tightened}. \href{https://www.government.nl/latest/news/2025/01/15/klever-export-controls-on-advanced-semiconductor-manufacturing-equipment-to-be-tightened}{当期官方解释}。与上项共同构成一个政策案例。

原证据核查日期为2026年10月7日；本轮2026年10月8日重新读取ASML官方财务页及年报发布入口。不同证据类型、全文定位及未完成事项见原核查台账与《12与13改写说明》。



\chapter{社会再生产：参与能力的形成与共同费用的承担}\label{ch:13}

\section{时间、现金与共同服务决定劳动者能否持续参与}

生产能力要持续发挥，就需要学习、休息、照护、维护和废弃物处置；转换期间还要有生活现金。本章将这些必要投入纳入技术收益核算。\textbf{企业费用下降，可能来自完整投入减少，也可能来自家庭、服务人员或公共部门新增的承担。前者是资源节约，后者是承担转移，两者还可能同时发生。}

远程工具少了通勤，劳动者可能多出时间；设备、维修和无酬等待交给接单者，企业则可能只是少了付款。灵活排班让部分人更方便，也可能让照护者来不及安排接替。企业账上都可出现节省，生活结果却不同。因此，除了谁领收益，还要查谁花时间、出现金、维持后续工作，才能说明生产关系在生活中怎样起作用。

照护使人能够继续工作，但提供了什么服务、花了多少时间、取得多少收入，是不同的量。本章按时间和实际付款追踪这些活动，不把全部照护都算作某家企业的剩余价值，也不用一个假定工资冲掉健康和质量差别。需要解释的是：生产依靠这些投入，付出者能否取得收入，又能否影响照护和工作怎样安排。

跨境生产进一步分开收入取得与长期费用承担的地点。国家均值描述总体差异，资产、劳动身份、照护责任和合同则影响个人能否进入、持续参与或退出。下文从家庭时间和预算出发，依次考察技能、交易服务及环境责任，解释这些条件怎样扩大或限制实际选择。

\subsection{再生产技术改变劳动力供给，费用转移不等于价值降低}

照护、学习和休息维持劳动力，但由此不能把全部再生产投入都算作企业可变资本。家庭无酬照护没有发生资本购买劳动力的支付；公共服务支出与资本经营的照护服务也有不同关系。商业照护机构雇员工资要按其生产关系分析，财政服务工资则按公共预算分析，家庭时间仍保留为时间账。区分这些账户，才能查明企业购买劳动力的支出依靠什么外部条件。

技术和公共服务可以减少必要家务时间、接替照护、降低学习门槛，使更多人能够就业或接受训练。新增劳动力供给对工资的作用仍取决于岗位需求、技能和制度；不能从参与人数增加直接推出工资下降。雇主也可能因劳动力更容易持续工作而扩大生产，使工资支出增加。供给改善、企业购买和劳动者所得是相接的三个环节。

费用转移尤其需要同劳动力价值的降低分开。企业少付培训或把设备维护交给家庭，必要劳动并未因此消失；若由家庭新增时间或公共预算补上，只是承担变化。只有在既定生活和能力标准下，相关生产条件确实减少再生产所需的社会劳动，才支持讨论相应的价值降低。即使现金价格降低，也须查服务质量、税收和家庭新增时间，不能直接认定必要劳动时间同幅减少。

线上培训还可能改变技能形成费用与劳动供给结构。资料复制便宜，实践、反馈、认证和学习期间的生活仍有成本。雇主减少新人训练而要求个人预先取得资格，会减少自身当期支出，却可能增加劳动者进入前的承担。由谁支付、形成了什么能力、能否跨雇主使用，需要分别查证，才知道技术节约改善了再生产，还是改变了劳动力出售的条件。

再生产技术释放出的时间，也有不同用途。照护被接替以后，劳动者可以休息、学习或进入有酬工作，雇主也可能据此要求更长可用时间或更灵活排班。企业工资购买的扩大与家庭负担的变化因而不能直接相互抵销。应追踪同一家庭的有酬工作、照护、休息和学习时间，并查谁决定新增可用时间的安排，才知道释放的能力怎样成为劳动者收益或资本可使用的劳动供给。

\section{照护责任限制就业和学习的可用时间}

以互斥的活动分类建立时间约束，可以明确不同活动对技术学习的限制。以一周为观察期，总时间T分为有酬工作h\_paid、必要家务和照护h\_care、学习h\_training与休息及其他活动h\_rest：

\[
T=h_{paid}+h_{care}+h_{training}+h_{rest}.\tag{13-1}
\]

式（13-1）是本文的简化分区，不是国家时间利用调查的分类表。通勤、求职及接单等待应在其他活动内记录，或另设活动类别，以保持时间核算的完整性。带薪培训在这里记入学习时间、收入另行记录，不再同时放入工作小时。实际同时做两件事时，要选定主活动或另报并行活动，不能把重复时间相加超过T。

设维持健康与必要生活的休息下限为h\_rest,min，既有家务照护需求为C，已被其他家庭成员或合格服务实际接替的时间为x，则h\_care=C\ensuremath{-}x，0\ensuremath{\leq}x\ensuremath{\leq}C。在有酬工作小时暂不变、服务质量保持的条件下，可安排学习的总时间上限为：

\[
h_{training}\leq T-h_{paid}-(C-x)-h_{rest,min}.\tag{13-2}
\]

增加一小时可靠的照护接替，至多释放一小时学习时间；接送、沟通和新增行政事务可能占用部分释放时间。若培训需要连续两小时，式（13-2）满足也可能仍无法参加：每天多段碎片空闲并不等于一个可用连续时段。可以再要求某次学习所需时段l不超过家庭能够取得的最长连续可用时段u，即l\ensuremath{\leq}u。累计时间与连续时段分别构成数量约束和排程约束，二者均须满足。

式（13-2）区分三条改善路径：减少必要照护投入$C$，增加可靠接替$x$，或减少有酬工作$h_{paid}$。前者节约劳动，第二条重新分担劳动，第三条还要解决收入减少的问题。压缩必要休息也能在账面上腾出时间，却可能损害健康，不能与上述改善等量评价。家庭协商、雇主排班和服务供给决定哪条路径实际可行。

ILO2024年照护简报估计，2023年全球15岁及以上人口中，7.48亿人因照护责任处于劳动力之外，女性7.08亿、男性4000万。资料基础为125国住户调查，覆盖全球非劳动力人口约71\%，并对缺失国家和部分原因类别插补。按公布人数计算，女性占这组人口约94.7\%；分母是7.48亿照护原因非劳动力人口，绝不是全部女性。\href{https://www.ilo.org/sites/default/files/2024-10/GEDI-STAT\%20brief_formatted_28.10.24_final.pdf}{原简报第1、3及14—15页}

这一估计说明照护责任所关联的劳动市场可用性存在大规模性别不对称，但并未估计提供替代照护能够新增7.48亿岗位。岗位需求、技能、偏好和可行服务仍会限制进入。非劳动力状态也不等于没有劳动，ILO统计的是市场劳动力身份，家庭劳动并未因此消失。2018与2024版本估计方法和资料变化，不能将两者相减宣称照护退出增长了多少。

中国第三次全国时间利用调查在2024年5月11—31日开展，覆盖31个省区市及新疆生产建设兵团，抽中3.85万户、10.7万人；对象为6岁及以上常住成员。日记记录一个工作日和一个周末日，按5/7与2/7加权成一周七天的日均。第二号公报报告无酬劳动居民平均2小时5分钟，参与者平均2小时45分钟。第三号分性别数字为男性参与者1小时52分钟、女性参与者3小时29分钟，其对应参与率为67.5\%与83.9\%。\href{https://www.stats.gov.cn/xxgk/sjfb/zxfb2020/202410/t20241031_1957217.html}{样本方案}，\href{https://www.stats.gov.cn/xxgk/sjfb/zxfb2020/202410/t20241031_1957216.html}{总体及定义}，\href{https://www.stats.gov.cn/xxgk/sjfb/zxfb2020/202410/t20241031_1957215.html}{参与者数据及参与率}

\begingroup\footnotesize\setlength{\tabcolsep}{4pt}\renewcommand{\arraystretch}{1.25}
\begin{longtable}{>{\raggedright\arraybackslash}p{\dimexpr 0.22500\textwidth-8.0pt\relax}>{\raggedright\arraybackslash}p{\dimexpr 0.28750\textwidth-8.0pt\relax}>{\raggedright\arraybackslash}p{\dimexpr 0.28750\textwidth-8.0pt\relax}>{\raggedright\arraybackslash}p{\dimexpr 0.20000\textwidth-8.0pt\relax}}
\toprule
\textbf{真实统计及本文计算} & \textbf{年份、样本和分母} & \textbf{可复算的解释} & \textbf{进入本章哪项约束} \\
\midrule\endfirsthead
\toprule
\textbf{真实统计及本文计算} & \textbf{年份、样本和分母} & \textbf{可复算的解释} & \textbf{进入本章哪项约束} \\
\midrule\endhead
\bottomrule\endfoot
因照护处于劳动力之外748百万人，其中女性708百万人 & ILO全球2023年估计，15岁及以上，125国资料及插补 & 708/748约94.7\%，只说明这一组的构成 & 需要实际照护安排才能进入市场 \\
无酬劳动112与209分钟/日 & 中国2024年调查，分别为男性、女性活动参与者均值 & 参与者均值相差97分钟；不是每对夫妻分工差 & 照护与家务时间的分布不能省略 \\
相应参与率67.5\%与83.9\% & 同调查，分母分别为调查代表的男性与女性居民 & 两组参与者构成不同，不能只凭97分钟估计全部人口效应 & 申请与未申请群体必须一起观察 \\
居民125与参与者165分钟/日 & 同调查，全部居民与当日该领域参与者 & 分母变化解释两种均值并存 & “有人参与”与“参与者投入”不同 \\
\end{longtable}\endgroup


无酬劳动领域还含购物、无酬实习和志愿活动，其统计范围大于式（13-1）中的狭义照护。参与者差值也不是某项技术释放时间的估计。这些资料支持时间分布值得进入经济分析，而不是用统计均值替代每个家庭的C。实际检验须连接家庭构成、主活动、连续时段和排班，观察技术改变后约束是否松动。

\section{学费免了，学习时间、生活费和准备现金仍要解决}

时间可行还不够。设维持基本生活所需实际支出为C\_min，购买合格照护的单位价格为p\_c，购买量为x；网络、业务用电、交通与课程等当期支出合计f，贷款本期偿付为d\_t。设有酬工作小时工资为w，其他已取得现金收入为Y，培训期间实际得到的补助为g。一个不靠新增借款或提取储蓄维持的当期预算要求：

\[
C_{min}+p_cx+f+d_t\leq wh_{paid}+Y+g.\tag{13-3}
\]

若需要一次设备或课程准备支出F\_T，则还要满足进入当时的流动性约束：

\[
F_T\leq A_0+L_{approved}+G_0.\tag{13-4}
\]

A\_0是可使用现金，L\_approved是确实获准且可提款的信贷，G\_0是已可取得的一次支持。式（13-3）有意考察常态持续支付；若暂时动用储蓄，应另记提款和余额，不能同时把A\_0用于交费又用于整段生活。获贷也不会凭空创造家庭净财富，其后利息与还款进入d\_t。未来预期工资不能拿来付今天尚未获得的现金。

下面的两人算例只演示入口约束，不代表任何国家住户。二人每周T=168小时，有酬工作40小时，休息及其他必要活动84小时，工资30元/小时，基本现金支出1100元/周；培训要求8小时/周，课程免费，但开始需要400元设备准备。甲的家务照护为28小时、可用现金600元，乙为42小时、现金100元，不提供信贷或补助。

\begingroup\footnotesize\setlength{\tabcolsep}{4pt}\renewcommand{\arraystretch}{1.25}
\begin{longtable}{>{\raggedright\arraybackslash}p{\dimexpr 0.36364\textwidth-8.0pt\relax}>{\raggedright\arraybackslash}p{\dimexpr 0.15909\textwidth-8.0pt\relax}>{\raggedright\arraybackslash}p{\dimexpr 0.15909\textwidth-8.0pt\relax}>{\raggedright\arraybackslash}p{\dimexpr 0.31818\textwidth-8.0pt\relax}}
\toprule
\textbf{条件或计算} & \textbf{甲} & \textbf{乙} & \textbf{为什么结果不同} \\
\midrule\endfirsthead
\toprule
\textbf{条件或计算} & \textbf{甲} & \textbf{乙} & \textbf{为什么结果不同} \\
\midrule\endhead
\bottomrule\endfoot
培训前可用小时：168\ensuremath{-}40\ensuremath{-}84\ensuremath{-}C & 16小时 & 2小时 & 免费不增加乙的可用时间 \\
完成8小时课程需新增接替x & 0小时 & 至少6小时 & 乙需要改变照护或其他排程 \\
准备费现金余额：A\ensuremath{_0}\ensuremath{-}400 & 200元 & \ensuremath{-}300元 & 乙还受到一次支付门槛 \\
若服务20元/小时，最低新增周支出 & 0元 & 120元 & 现金结余原本仅1200\ensuremath{-}1100=100元 \\
扣除最低照护后周结余 & 100元 & \ensuremath{-}20元 & 仅解决准备费仍不能持续参加 \\
\end{longtable}\endgroup


乙的困难不是课程回报较低，而是参加条件没有满足：缺6小时照护接替、300元准备现金，购买接替后每周还缺20元。一次设备补助只解决准备费，持续照护和生活预算仍有缺口。可靠照护与带薪学习可以同时放松时间和预算限制。只观察已报名者，会遗漏被这些条件挡在课程之外的人。

不同解决办法还会改变负担的去向。少做有酬工作降低式（13-3）的收入，请家人接替增加另一成员的劳动；工具若在同等质量下减少家务或重复沟通，则可能同时节约时间和现金。比较各人的实际劳动、付款与服务质量，才能辨别改善来自投入减少还是重新分担。

能参加，还要查参加以后能得什么。未来新增工资，应扣新增工作费、放弃的训练期工资、持续照护和融资支出，再将折现收入与准备费比较。学费只是一项。货币回报为正，也没有单独说明健康、照料和自主性改善。本章不将无偿时间乘统一工资记入实得收益，避免把假定估价当作已经拿到的钱。

\section{照护接替可以扩大机会，新增时间由谁安排须另查}

照护技术可以减轻搬运、提醒、预约和重复登记，可靠服务可以提高式（13-2）的x。两者都能够创造真实收益。它们需要保持健康、安全和尊严的目标：减少必要陪伴或忽视异常，并不能因分钟数下降被叫作效率。评估应比较相同质量目标下的完整投入，包括家庭等待、沟通与服务人员工作。

Ajayi、Dao与Koussoubé的布基纳法索随机评价，在36个城市公共工程地点中随机选择18处提供社区托育，最初抽选2160名女性，2019年12月基线完成2150户，2021年2月随访1990人。结果支持若干经济和儿童发展改善，但家庭决策自主、性别态度及家庭内分工未见显著变化。\href{https://documents1.worldbank.org/curated/en/099215111282254210/pdf/IDU09ee54b690865904f71089a60d76ad621601c.pdf}{2022年原工作论文，§3及结果表}

这一实验支持照护服务改善特定人群的经济机会，却未检出家庭自主性显著变化。两种结果并不冲突：服务解除时间限制，未必改变家庭收入和活动安排的决定方式。服务地点的随机安排与个人实际使用也不同，疫情停开和差异失访限制了推断范围。

家庭腾出了时间，还要看接替者怎样工作。若照护工靠低薪、持续加班完成接替，家庭减负未必是社会少用了劳动；若设备省下服务工时并保持质量和合理排班，两方可以一起受益。低收费也有不同来源：效率提高、财政补贴、工作人员多担劳动，都可能出现。须查投入与承担者，才能解释价格降低意味着什么。

家庭总收入掩盖不了成员之间的差异。新增工资由谁使用，夜间照护由谁接替，不安全的工作安排由谁有权拒绝，都会改变家庭成员的实际收益。因此，收入分析要与成员各自的时间和处置权限相接，才能说明谁减轻了负担、谁增加了依赖。

\section{技能可携带与生活缓冲共同影响退出选择}

OECD2025年报告使用2023年成人技能调查，成人学习障碍分析限25—65岁、过去12个月非正规且与工作相关的学习。参加成员国非加权平均中，约24\%的人希望参加或进一步参加却遇到障碍；在报告障碍者中，48\%把工作或家庭时间不足列为首要原因。两个比例的分母不同，不能相乘后当世界人口训练需求。\href{https://www.oecd.org/en/publications/trends-in-adult-learning_ec0624a6-en/full-report/why-is-participation-not-more-common_73893b53.html}{第三章图3.6—3.7及注释}

自报时间不足，与本章约束相符，却没有估计解除障碍后能多大程度参加。要报告障碍，通常先要知道课程并希望参加；长期被挡住的人，也可能不再表达愿望。应沿不知课程、未申请、未获选、未完成和完成后未就业分别观察。只比结业者收入，可能把原有资产、技能与时间优势一并记为培训收益。

在线课程可以降低搜寻、地域和练习费用，模拟训练可以减少耗材，灵活学习能够适合部分照护者。这些真实能力收益之后，还要经过技能认可和岗位需求。企业内部工具成绩未必能向外部雇主展示；可验证作品和资格能够携带，才可能提高第四章分配模块中的外部机会。这一分析说明技能认可影响外部机会的可能机制；证书数量与谈判参数b之间仍需建立经验关系。

Alfonsi等人的乌干达双侧实验，涉及约1700名青年和1500家企业，实际符合资格申请1714人，比较六个月职业学校与补贴企业培训。主要培训发生在2013—2014年前后，追踪24、36和48个月，2020年发表。两条路径的启动率和收益轨迹不同，学校路径的较长持续改善与企业路径较快的早期改善，支持研究技能如何被其他雇主识别。\href{https://www.lse.ac.uk/economics/Assets/Documents/personal-pages/robin-burgess/tackling-youth-unemployment.pdf}{原论文§2—3及图1—2}

该实验比较的是两套培训安排，课程、证书和场所没有分别随机化。长期差异因此不能全部归因于证书，获邀培训的效果也不能直接当作实际参加者的效果。它为技能认可和参加门槛提供研究线索，不提供人工智能课程的预期回报参数。

能换工作，至少要有真实报价，还要撑得过转换期。新合同以后收入更多，若先要无薪等待、重新买设备或失去保障，今天也未必能选。现金缓冲、可携带资格和合理转换规则，可以补这些条件；昂贵单一认证、订阅许可或不认可已有经验，又可能增加新门槛。因此，技能提高之后，实际可选合同和可改任务仍须另查，不能保证选择权一路增加。

\section{数字连接和支付改善机会，使用权限仍需检验}

ITU《事实与数字2025》对2025年互联网使用作即时估计，综合会员国调查、模型、插补和人口加权。高收入组94\%、低收入组23\%，全球城市85\%、农村58\%；收入组依世界银行分类。这些不是同一天普查，也不是同一人面对两个政策的实验结果。\href{https://www.itu.int/itu-d/reports/statistics/2025/10/15/ff25-internet-use/}{互联网使用}，\href{https://www.itu.int/itu-d/reports/statistics/2025/10/15/ff25-internet-use-in-urban-and-rural-areas/}{城乡使用}，\href{https://www.itu.int/itu-d/reports/statistics/2025/10/15/ff25-methodology/}{估计方法}

\begingroup\footnotesize\setlength{\tabcolsep}{4pt}\renewcommand{\arraystretch}{1.25}
\begin{longtable}{>{\raggedright\arraybackslash}p{\dimexpr 0.35714\textwidth-8.0pt\relax}>{\raggedright\arraybackslash}p{\dimexpr 0.19048\textwidth-8.0pt\relax}>{\raggedright\arraybackslash}p{\dimexpr 0.45238\textwidth-8.0pt\relax}}
\toprule
\textbf{ITU2025估计的比较} & \textbf{相同指标的组间差} & \textbf{解释边界} \\
\midrule\endfirsthead
\toprule
\textbf{ITU2025估计的比较} & \textbf{相同指标的组间差} & \textbf{解释边界} \\
\midrule\endhead
\bottomrule\endfoot
高收入94\%，低收入23\% & 71个百分点 & 收入组人口使用率差，不是联网的工资效应 \\
全球城市85\%，农村58\% & 27个百分点 & 城乡人口使用率差，与收入组存在交叠 \\
\end{longtable}\endgroup


差值使用的是比例百分点，既不是增长率，也不是需要新增多少设备的估计。一次上网不保证可靠工作连接，独立设备、可负担数据、电力和支付资格还会进入式（13-3）—（13-4）。对照护者而言，网络能否在实际可用时段保持稳定，还关系到连续时间约束。全国使用率上升，可以与这些内部差异并存。

连接改善确有机会收益。Hjort与Poulsen2019年研究利用海底光缆渐次到达与陆地网络布局，以12国三套数据作双重差分，报告就业改善。Roodman2025年9月第三版复核则指出地理编码、处理范围与卫星更替敏感性；本章此次仍限官方摘要和版本说明，未复算微观结果，不把原系数作为通用投资回报。\href{https://www.aeaweb.org/articles?id=10.1257/aer.20161385}{原研究}，\href{https://arxiv.org/abs/2401.13694v3}{复核v3}

线路建了，效果仍要对应具体地点、处理范围和测量。线路附近原有商业优势，或位置编码错了，都可能改变比较；远程教育、客户接触和企业协调，又仍是可检验的渠道。研究应查稳定连接是否少了费用、多了可交付任务、改善了就业，不能以使用率上升代替这些结果。

支付接入也可能产生可测量的生活改善。Suri与Jack（2016）利用肯尼亚2008—2014年住户调查，比较2008—2010年住宅周围1公里内M-Pesa代理网点增加程度不同的家庭。2014年访得1,608户，贫困回归实际使用1,593户。表1中，女性户主家庭在接入增幅第75分位与第25分位之间的极端贫困概率差为\ensuremath{-}9.2个百分点，门槛是研究使用的每日人均1.25美元；这不是全国贫困率下降9.2个百分点。职业结果则以个人为单位，不能只归给女性户主家庭。（Suri与Jack，2016，第1288—1291页，表1、表3，\href{https://doi.org/10.1126/science.aah5309}{DOI：10.1126/science.aah5309}，\href{https://www.jefftk.com/suri2016.pdf}{正式论文公开镜像}）

作者进一步外推约194,000户退出极端贫困、185,000名女性转入商业或销售。补充材料说明，这些人数由回归系数、平均接入增幅与人口分母推算，不是逐户逐人观察的转换总数；原样本非内罗毕地区流失35\%，作者使用低流失社区及重新加权检查。识别依赖早期网点增长与未观测经济趋势不相关，接入增幅不是随机分配，也不是个人注册钱包的处理效应。职业结果还有强弱区别：表3女性商业或销售项的多重检验调整后p值为0.15，不能与女性退出农业的较稳健证据等量解释。（同上，表3；Suri与Jack，2016，\href{https://www.jefftk.com/suri2016-supp.pdf}{原始补充材料第6—8页、表S7—S8及S11—S12}）

移动支付可能通过及时收付款减轻流动性困难，使家庭少为生活支出停业或出售工具；职业转换也可能扩大收入来源。这是依据结果提出的机制解释，原研究没有将贫困改善全部分解为付款时点的效应。家庭资金量、到款时间与经营机会仍是三项需要分别观察的条件。

在研究识别假设成立的条件下，肯尼亚结果支持低收入家庭能够获得实际技术收益；与此同时，仍有另一项制度问题：家庭更容易使用支付网络，不等于家庭参与制定费用、账户准入和申诉规则。生活改善与决定权变化可以沿不同路径发生。比较其他地区时，应核对代理能否可靠兑付、交易费用与替代支付渠道，而不能仅按钱包注册率预测脱贫人数。对照护约束较强者，还须分别观察职业收入和无酬劳动时间；换入商业并不自身证明照护负担减少。

外部云服务可以将高额自建费用换成小额订阅，帮助小企业开始使用；资料能带走、工具兼容、替代可靠，又可以使退出更容易。若身份、支付、评价和业务格式带不走，初期方便便可能伴随以后转换更贵。应在同等质量下查能否换服务、能否协商规则，不能单凭服务器国籍判定收益与自主性。

\section{收费与无酬业务时间改变跨境任务的实际回报}

ILO2021年平台劳动报告将在线任务与地点型平台分开。其执行摘要说明，所跟踪的五家主要在线平台任务需求主要来自发达国家，劳动供给主要来自发展中国家，并讨论标注、转录、审核和业务外包。资料综合调查、访谈与平台记录，不能当作所有跨境劳动的概率样本或2026年全市场地图。\href{https://www.ilo.org/sites/default/files/wcmsp5/groups/public/\%40dgreports/\%40dcomm/documents/publication/wcms_771672.pdf}{官方执行摘要第2—6页}

跨境平台让一些不能迁移的人也能接客户、拿收入，小企业也可能少花购买服务的费用。但客户更容易找到劳动者，不意味着劳动者同样容易找到可信客户、追回拒付或带走评价。比较跨境与本地工作，须扣必要支出并算全业务时间，才能从报价走到实际回报。

设同一时期任务标价r元/计费小时，计费小时h\_b，必要而无酬的测试、搜寻、沟通和返工小时h\_u，平台费用比例s\_p，实际运行费用C\_op。忽略税费时，业务净现金除以完整业务时间得到：

\[
w_e=\frac{(1-s_p)rh_b-C_{op}}{h_b+h_u}.\tag{13-5}
\]

这一比率衡量完整业务时间的现金回报，家务照护与个人休息仍由独立时间账户记录。税费、拒付或停工造成的实际付款变化应计入分子；若账单按件计价，应直接用真实已付任务收入替换rh\_b。例：r=30元、h\_b=20小时、h\_u=10小时、s\_p=10\%、C\_op=60元，则净现金480元，完整业务小时回报16元，而非标价30元。它是原创算例，不代表某个平台观测均值。

在任务收入固定且质量相同的条件下，工具减少h\_u会提高w\_e，也能释放时间；若客户据工具提速调低单价，或要求更快响应，实际结果需要重新计算。工具减少业务小时还可能降低家中设备被占用时间，改善其他成员学习，这类收益可以直接测量。相反，短截止期和夜间响应可能挤占h\_rest或提高照护接替支出。业务净收入增长与社会再生产负担下降属于需要分别观察的结果。

跨境接单的长期价值取决于能否积累可携带的技能和客户关系。简单标注可能提供首份收入；长期只处理片段、缺少目标和反馈，则可能限制职业升级。有酬训练与完整项目参与提供另一条路径。沿同一劳动者的任务、客户集中度、付款和转出经历追踪，才能判断协作扩大了机会还是固化了位置。

外包、最终服务与软件采购是不同交易，不能全相加作新增产出。本地出口增加，还须拆出增加值、劳动报酬和所有者所得，各留在哪里。承包商、正式雇员与临时接单者可能分得不同，家庭要撑过付款延迟，年度宏观总额也未必显出这个困难。

\section{采购收益与健康负担可能分属不同主体}

技术要用材料，还可能需要污染控制、医疗和替代生计。本节沿采购查：谁得收入，谁维持受影响者以后的生活。投入便宜了，可能是真正少用资源，也可能有一部分修复费用尚无人付。后一判断须有现场暴露、损害与责任资料，不能仅看材料价格或能源装机。

Banza Lubaba Nkulu等人的刚果（金）钴开采研究，在2014年11月和2015年5月于科卢韦齐开展横截面调查，便利抽样纳入122人，其中矿区居民72人、挖掘者25人、对照区居民25人。研究发现矿区居民钴暴露较高，儿童差异尤其明显，并检验相关生物指标。它提供具体暴露路径的证据，但不是全国患病率调查，也不能把生物标记变化直接换成确定的终身经济损失。\href{https://pmc.ncbi.nlm.nih.gov/articles/PMC6166862/}{原论文，Methods、Results与Limitations}

矿区风险超出雇佣边界。儿童可能因居住地而暴露，家庭陪诊、照料和搬迁又未必计入矿产品报价。由此提出的机制是：采购和加工取得收益，而修复缺少持续资金时，家庭与公共服务可能承担后续费用。矿区暴露研究支持追查这一渠道，负担大小仍需健康、收入和责任记录相接。

在家庭时间和预算约束中，暴露后新增的陪诊与照料可能增加C，医疗、搬迁及替代生计费用可能提高实际C\_min。即使矿产品交易使另一端的企业成本下降，本地家庭的学习和工作集合也可能因此收缩。这里没有用生物标记估计新增小时或支出，所提出的是需要将健康与家庭记录连接的渠道。可靠修复和公共服务则可能减轻同一约束，收入与生态改善也可以并存。

钴案例不能代表全部矿区，也不能把所有锂电池都当作含钴。材料路线、开采与治理不同，工艺改进、减少需求、回收和监管可能减负；原论文也未证明大规模工业开采当然更好。买方排除手工采购，可以改变自己记录的风险，但当地人改做什么、污染是否减少，还须追踪，才知道生计与环境有没有改善。

资源需求也可能带来工资、设施与技术学习。须查工资和税费实际付了没有、收入是否支持当地服务、污染是否监测、居民能否拒绝或协商、投资结束后谁继续修复。税进入预算，还要经财政与服务才能改善健康；所有权名称没有替代这些环节。资源丰富地区的获益与负担，应按承担者分别报告。

采购方的价格收益与受影响者的健康损失可能分属不同地区和时期。比较应保留受损者、医疗服务者与财政承担者的结果；无法可靠货币化的健康和生态变化，直接报告指标及不确定性，避免任意折价掩盖严重损害。

\section{维护与处置责任决定设备的完整持续成本}

UNCTAD《数字经济报告2024》将开采、制造、使用和废弃置于同一生命周期，并指出矿物贸易与加工价值分布的地区差异。其国家分组不等于世界银行收入组，报告的宏观描述也不能据此识别某个家庭损失来自哪家企业。本章据此要求分析跨境投入和责任的对应，而不把“发展中国家”当作单一经济位置。\href{https://unctad.org/system/files/official-document/der2024_en.pdf}{原报告第二章第三节及官方报告入口}，\href{https://unctad.org/publication/digital-economy-report-2024}{发布机构入口}。

废弃物数据进一步揭示测量边界。《全球电子废弃物监测2024》估计，2022年全球电子废弃物为6200万吨，有记录的正规收集回收比例为22.3\%。这一总量涵盖多类电气电子设备，并非人工智能设备废弃量，也不能全部归因于跨境数字服务。未被记录为正规回收的部分，不等于全部被非法出口，更不等于全部没有任何处理。\href{https://ewastemonitor.info/the-global-e-waste-monitor-2024/}{ITU、UNITAR报告发布与数据说明}

维修与报废，还会改变家庭以后花多少钱、地方能做什么工作。延长可用寿命，可能少买新设备并保留经验；可维修产品，也可能支持当地技术岗位。配件或软件支持过早停止，则可能让有实体功能的设备难以使用。须分清物理故障、安全更新不足与主动换机，不能把所有更新一律称为蓄意报废。

设备生命周期还会改变采用时的流动性要求和后续运行预算。过早需要换机，会增加式（13-4）的准备现金需求；配件与软件支持若延长安全可用寿命，可以减少重复F\_T，并降低业务运行费用。一次便宜购买和长期可负担使用须分开，特别是现金紧张者可能只能购买当期负担得起、维护却更贵的设备。这个机制需要价格、维修和寿命记录，不能由全球废弃物总量直接推算家庭损失。

远程协调可以少跑重复路，维护更准可以延长寿命，信息服务也可以改善收集处理。有没有形成节约，须查完整使用周期和新增加的需求。单位省了，总量可能因使用更多而上升；替代高投入活动，也可能使总负担下降。应按要回答的问题统一边界，不能只选有利于既定结论的环节。

公共部门的实际负担取决于设施供给、处理标准、监管执行及收费规则，不能仅由部门属性判断。地方政府可以通过公共设施和监管改善回收，也可能因为治理不足而让污染扩大；私人维修与回收可以提供高质量服务，也可能在利润压力下回避低价值材料。评估应同时记录收费、处理标准、工作人员安全和最终去向，检验收入取得与处置责任的对应关系。收费机制若使低收入家庭无法交付废弃设备，形式完备的制度也可能产生实际排除。

\section{国家内部的资产和责任差异塑造调整能力}

国际位置与国内分配共同塑造采用条件。低收入国家中有资产、可靠服务和支付资格的人可以使用新工具，高收入国家中缺照护和生活缓冲的人仍可能无法参加。设施与许可规定共同条件，资产、劳动身份和责任分配再决定这些条件对谁可用。

同一出口项目，承包商可能拿管理收入，临时劳动者担订单波动，家庭担突发排班。矿区土地与许可持有者，和靠日工资、公共医疗生活的人，位置也不同。国家总收入涨了，不能宣布人人获益；一部分人受损，也不能抹去当地获益者与能力升级。须分清谁得什么、谁担什么。

性别与这些位置共同作用。ILO和中国时间统计指向明显分布差异，却不能用性别单独决定每个人的C或A\ensuremath{_0}。高收入、有服务支持的女性可能得到远程机会，男性主要照护者也可能受同样时段约束；移民、残障、年龄和语言会继续影响服务、认证与支付。研究要识别身份如何通过责任和规则起作用，而不将身份本身当作同质损失组。

比较应对准条件：工具质量与岗位需求相近时，照护、现金缓冲、可携带技能与合同，能否解释采用、退出和持续所得差异。缺资本者用低固定费服务取得好机会，应承认其入口改善；给了时间和现金仍没变化，应再查技能、需求、歧视或课程质量。具体约束解释须能接受这种反证，不能把一切结果都收进同一个原因。

\section{保障提供缓冲，共同服务扩大实际选择}

ILO《世界社会保障报告2024—2026》同版估计，2023年全球52.4\%人口获得至少一种现金社会保障待遇，高收入组85.9\%、低收入组9.7\%，同版差值76.2个百分点。分母为相应组人口，“至少一种”不是所有风险获得足额保障，也不能与医疗保护直接合并。方法与资料修订使它不宜机械连接旧版趋势。\href{https://www.ilo.org/sites/default/files/2024-09/SOCPRO_WSPR_2024_Executive_Summary_EN_Web.pdf}{执行摘要图ES.1及方法说明，第2—3页}

保障如何进入本章约束，可以具体说明。训练期实际收到现金支持g，提高式（13-3）右侧；设备支持改善式（13-4）；共同照护改变x；转换期间可携带的保障减少退出损失。四种作用不能由一个覆盖率推算金额或工资效应，待遇资格、等待期、持续性和实际申领都会影响结果。形式上有制度而无法获得支付，不能算家庭已具备缓冲。

在给定工作与可分配收入下，外部机会改善可能提高劳动谈判所得，但并不是保障支出自动创造产出。税收、缴费或支持资金须有实际来源；若同一家庭缴纳费用，还要记录其净支付。公共服务更需要合格工作人员、设施和排班，现金补助在服务极度短缺时可能主要推高价格，未充分提高可靠接替x。技术节约也只有实现为可征收入或真实资源后，才可能支持持续支出。

公共服务也可以先形成能力。可靠照护、健康与教育，让人能持续工作，可能少花以后训练和协调的费用；收益跨雇主、跨时期，单个企业未必能独自承担全部投入。共同组织因而可能更有持续性。质量差、资源只给容易完成者、申请过繁，又可能削弱作用。应查实际使用者和未获服务者，花钱增加没有单独证明机会增加。

服务的所有制名称不能代替结果评价。减负效果由质量、收费和可用容量判断，治理效果由资格制定、劳动安排、纠正程序和实际执行判断。现金补助改变预算，共同服务扩大可行活动，参与治理改变决定程序；三类作用各有独立的检验依据。

\section{承担变化与完整投入区分费用转移和资源节约}

时间与预算约束说明哪些方案可行，分工和组织解释为何有人具备条件、有人受到限制。家庭、合同与服务安排决定谁接照护、谁付现金、谁能培训；环境损害还会增加必要支出。进一步查明谁能修改这些安排，才能把再生产投入与生产支配联系起来。

进入比较应保留未申请者和退出者，考察免学费、照护接替与带薪学习分别改变了哪些障碍；退出比较应连接可携带技能、转换期现金和真实外部报价。业务收益比较则核对实际付款、收费与全部必要业务时间，不能以标价替代净回报，也不能把照护重复并入业务时薪。

要证明费用外移，须把同一活动前后的投入和承担者接上。企业少付，没有证明家庭多担；家庭多做，也可能因服务扩大或质量提高。完整投入、质量、现金与责任一起比较，才能分清节约、新增投入与付款转移。环境责任还须对应采购、暴露、医疗与修复的地点和时期，废弃物宏观总量没有单独指出某家企业应负多少责任。



成本转移以可行基准下的承担变化为依据。原有照护、教育或维护不能因企业受益便全部算作采用新增的转移；应比较同等质量下各方投入、付款和补偿的变化。企业减少的承担与家庭新增、未获补偿的承担相对应，才支持转移判断。双方投入都减少时，应承认真实节约；财政代付与公共服务减少重复投入，也分别记录。

\section{持续生产需要将家庭和公共投入纳入核算}

本章将技术收益的观察期从当前任务延长到持续生产。时间、生活现金、照护、训练和维护，决定劳动者能否进入、坚持和退出。培训、托育与支付案例显示，解除具体限制可以产生真实改善；免费工具或更高报价却未必同时解除其他限制。

这使费用评价发生变化。企业少付款，并不等于完整生产少用投入；其他主体原有的贡献，也不能全部记为技术采用新增的费用转移。只有同一活动中企业减少的承担与其他主体新增、未补偿的承担能够对应，费用转移判断才有依据。公共服务若减少重复劳动、保持质量并扩大选择，则产生相应社会收益。

再生产由此补上前文缺少的一环：共同能力依靠持续投入维持，投入者却未必取得相应补偿或决定权。第十四章据此分别核算企业利润、其他领取者和真实资源投入，并检验收益分化的条件。模型没有覆盖的家庭与长期健康费用，仍保留为经验研究的缺口。

\section*{本章参考文献}

1. ILO. 2024. \emph{The Impact of Care Responsibilities on Women’s Labour Force Participation}. Statistical Brief, October. \href{https://www.ilo.org/sites/default/files/2024-10/GEDI-STAT\%20brief_formatted_28.10.24_final.pdf}{官方PDF}。
2. 国家统计局. 2024.《第三次全国时间利用调查公报》第一、第二、第三号，10月31日。\href{https://www.stats.gov.cn/xxgk/sjfb/zxfb2020/202410/t20241031_1957217.html}{第一号}、\href{https://www.stats.gov.cn/xxgk/sjfb/zxfb2020/202410/t20241031_1957216.html}{第二号}、\href{https://www.stats.gov.cn/xxgk/sjfb/zxfb2020/202410/t20241031_1957215.html}{第三号}。
3. Ajayi, K. F., A. Dao, and E. Koussoubé. 2022. \emph{The Effects of Childcare on Women and Children: Evidence from a Randomized Evaluation in Burkina Faso}. World Bank Policy Research Working Paper 10239. \href{https://documents1.worldbank.org/curated/en/099215111282254210/pdf/IDU09ee54b690865904f71089a60d76ad621601c.pdf}{原论文}。
4. OECD. 2025. \emph{Trends in Adult Learning: New Data from the 2023 Survey of Adult Skills}. Getting Skills Right. \href{https://doi.org/10.1787/ec0624a6-en}{DOI:10.1787/ec0624a6-en}。
5. Alfonsi, L., O. Bandiera, V. Bassi, R. Burgess, I. Rasul, M. Sulaiman, and A. Vitali. 2020. “Tackling Youth Unemployment: Evidence from a Labor Market Experiment in Uganda.” \emph{Econometrica}, 88(6):2369–2414. \href{https://doi.org/10.3982/ECTA15959}{DOI:10.3982/ECTA15959}。
6. ITU. 2025. \emph{Measuring Digital Development: Facts and Figures 2025}. \href{https://www.itu.int/itu-d/reports/statistics/facts-figures-2025/}{官方在线报告}。
7. Hjort, J., and J. Poulsen. 2019. “The Arrival of Fast Internet and Employment in Africa.” \emph{American Economic Review}, 109(3):1032–1079. \href{https://doi.org/10.1257/aer.20161385}{DOI:10.1257/aer.20161385}。
8. Roodman, D. 2025. “The Arrival of Fast Internet and Employment in Africa: Comment.” arXiv:2401.13694v3，9月17日修订；初版2024。此处作为复核预印本引用。\href{https://arxiv.org/abs/2401.13694v3}{明确版本}。
9. ILO. 2021. \emph{World Employment and Social Outlook 2021: The Role of Digital Labour Platforms in Transforming the World of Work}. \href{https://www.ilo.org/sites/default/files/wcmsp5/groups/public/\%40dgreports/\%40dcomm/documents/publication/wcms_771672.pdf}{执行摘要}。
10. Banza Lubaba Nkulu, C., L. Casas, V. Haufroid, T. De Putter, N. D. Saenen, T. Kayembe-Kitenge, P. Musa Obadia, D. Kyanika Wa Mukoma, J.-M. Lunda Ilunga, T. S. Nawrot, O. Luboya Numbi, E. Smolders, and B. Nemery. 2018. “Sustainability of Artisanal Mining of Cobalt in DR Congo.” \emph{Nature Sustainability}, 1:495–504. \href{https://doi.org/10.1038/s41893-018-0139-4}{DOI:10.1038/s41893-018-0139-4}，\href{https://pmc.ncbi.nlm.nih.gov/articles/PMC6166862/}{开放原稿}。
11. UNCTAD. 2024. \emph{Digital Economy Report 2024: Shaping an Environmentally Sustainable and Inclusive Digital Future}. UNCTAD/DER/2024，7月10日。\href{https://unctad.org/publication/digital-economy-report-2024}{官方入口}。
12. Baldé, C. P., et al. 2024. \emph{The Global E-waste Monitor 2024}. ITU and UNITAR. \href{https://ewastemonitor.info/the-global-e-waste-monitor-2024/}{发布机构报告入口}。
13. ILO. 2024. \emph{World Social Protection Report 2024–26: Universal Social Protection for Climate Action and a Just Transition}. \href{https://doi.org/10.54394/ZMDK5543}{DOI:10.54394/ZMDK5543}，\href{https://www.ilo.org/sites/default/files/2024-09/SOCPRO_WSPR_2024_Executive_Summary_EN_Web.pdf}{执行摘要}。

14. Suri, T., \& Jack, W. (2016). The long-run poverty and gender impacts of mobile money. \emph{Science, 354}(6317), 1288—1292. \href{https://doi.org/10.1126/science.aah5309}{DOI：10.1126/science.aah5309}，\href{https://www.jefftk.com/suri2016.pdf}{正式论文公开镜像}，\href{https://www.jefftk.com/suri2016-supp.pdf}{原始补充材料公开镜像}。

\section{资料读取范围与未完成检验}

原证据核验日期为2026年10月7日；本轮2026年10月8日重读ILO照护简报、国家统计局第三号公报、OECD第三章及ITU使用与城乡页面。时间利用公报、ILO照护简报方法段、OECD第三章、两项培训/托育原论文相关设计与结果段，以及ITU在线方法与使用指标已读取。宽带原研究和复核此次读取官方摘要及版本说明，未重跑其微观复现；钴研究读取原稿索引返回的Methods、Discussion、Limitations等相关部分，未取得其微观原始数据；UNCTAD完整PDF直读受限，采用官方索引返回的第二章第三节及官方入口；电子废弃物采用发布机构的报告介绍，未逐表复算。具体页码、样本差异和限制见《13-社会再生产与国际不平等核查》。本文未开展新的住户调查、跨国因果估计或健康风险估值，原创条件机制不作为已经验证的经验结果。



新增移动货币案例于2026年10月8日读取正式论文及原始补充材料的公开镜像，实际范围和版本差异见《新增案例-制造与普惠》；未重新估计作者微观模型。

\chapter{技术收益与制度选择的条件模型：采用、分配与政治反馈}\label{ch:14}

\section{采用激励与各方收益需要分别解释}

企业采用，先比自己改与不改能赚多少；大家都改以后，市场价格和各方所得又会变化。\textbf{本章将采用激励与收益去向分开：前者看企业面对的选择，后者看竞争、费用与分配。} 技术能力提高，不能保证企业利润、劳动所得和社会收益同向或同比增加。

本章依次处理四个问题：数量竞争和固定费如何形成采用激励，给定工时与收入时劳动者怎样分享所得，给定支出规则时现金和商品如何闭合，以及补偿和政治成本如何改变自动化与制度选择。四个模块分别对应市场竞争、劳动分配、需求实现和制度反馈，为第四章命题提供条件推导；最后以效率和风险扩展检验政治结论的适用边界。

四个模块使用不同的给定条件，尚未构成联合均衡。竞争模块固定工资和需求，谈判模块固定工时和收入，现金模块固定价格、工资及支出规则，制度选择模块给定净收入、劳动参与约束和政治成本。把它们连接起来，需要同时求解工资、成本、产量、收入和支出，不能将相近的算例数值当作同一经济的连续变化。

参数是公开假设，未用前文案例联合校准。实证研究指出现实应查哪些变量，数值相近却不能据此选定企业数或谈判权重。代码、参数、输出与核验均保留，图画的是既定条件下的关系。谁能定这些条件、权限怎样变，仍须查制度和决策资料。

反例的用途是排除无条件推论：个别采用有利，并不保证普遍采用后企业利润提高；企业利润下降，并不保证局部社会收益下降；收入总额增加，并不保证劳动者所得增加。满足所列假设的反例足以否定必然关系，现实发生的频率和规模则由数据回答。

\section{劳动节约与接入收费共同决定企业成本}

设企业生产一单位产品需要$1/a$工时，$a>0$为每工时产出，时薪$w>0$。其他投入的单位货币成本为$m$，按产品数量支付的接入费为$\rho$。在当前工资和投入价格下，企业单位变动成本为：

\[
c=\frac{w}{a}+m+\rho,\qquad
\frac{\partial c}{\partial a}=-\frac{w}{a^2}.
\tag{14-1}
\]

\(a\) 的单位是产品/工时，\(w\) 的单位是货币/工时，后三项相加后均为货币/产品。这一成本分解区分了劳动投入节约与接入费用支付：即使技术本体免费，使用云服务、专用设备接口、基础设施或某些知识产权仍可能产生 \(\rho\)。接入费用提高，会抵消一部分企业看到的技术节约，但它究竟是社会资源耗费还是供应者收入转移，需要另做核算。

从 \(a_0\) 提高到 \(a_1>a_0\)，在工资及其他投入价格不变时，单位成本下降：

\[
\delta=c_0-c_1=w\left(\frac1{a_0}-\frac1{a_1}\right)>0.
\tag{14-2}
\]

若接入费用同时提高 \(\Delta\rho\)，则成本节约改为上述劳动节约减去 \(\Delta\rho\)；若额外能源、维护或材料投入提高 \(\Delta m\)，也需扣除。因此，把某项技术的实验速度提升直接当成商业成本同比下降，会遗漏收费和互补投入。前文能源、平台及芯片章节讨论的融资、电网、迁移和瓶颈，在这里分别进入固定费用、单位成本或产量约束，而不是被假定为不存在。

采用另需本期费用化的调试、培训和流程转换支出$F$，单位为货币/企业/期。它与现金模块的实体投资品数量$I$不同。本期全额费用化使采用博弈保持清楚；多年设备则应比较经营增量收益的折现现值与初始投入：

\[
F_{\text{初始}}\leq
\sum_{t=1}^{T}\frac{\Delta \pi_t^{\text{经营}}}{(1+r)^t},
\tag{14-3}
\]

式（14-3）还需结合残值、更新费用、风险和融资条件解释。采购时的现金支付、当期会计折旧与长期投资回报分别回答付款、费用确认和投资价值，单期采用费不能同时代表三者。

盈利条件与采用可行性是两道门槛。预计增量收益为正，企业仍可能缺预付款或无法取得必要服务；对手反应又会改变预计收益。因此，采用解释要在利润比较之外记录资金与服务准入。

\section{成本差异经数量竞争，改变价格与销量}

设一个局部产品市场存在 \(n\) 家企业，出售同质产品，各自在给定对手产量下选择自己的数量。行业总量为 \(Q=\sum_iq_i\)，逆需求为：

\[
p=\alpha-\beta Q,\qquad \alpha>0,\quad\beta>0.
\tag{14-4}
\]

“逆需求”表示市场交付越多，能够实现的单位价格越低。它把消费者支付意愿概括为一条线，不包含居民现金账户，也没有自动解决收入分配对市场需求的反馈。曲线有效区间为 \(0<Q<\alpha/\beta\)，不在负价格区间外推。

企业 \(i\) 的经营利润为 \((p-c_i)q_i\)。采用决定一旦作出，固定费用不会改变当期最优产量的一阶条件，只会改变是否采用的前期选择。把对手产量记为 \(Q_{-i}\)，企业求解：

\[
\max_{q_i\geq0}
[\alpha-\beta(q_i+Q_{-i})-c_i]q_i.
\]

内点一阶条件和二阶条件分别为：

\[
\alpha-c_i-\beta Q_{-i}-2\beta q_i=0,\qquad
\frac{\partial^2\pi_i}{\partial q_i^2}=-2\beta<0.
\tag{14-5}
\]

因此最优反应为
\(q_i=\max\{0,(\alpha-c_i-\beta Q_{-i})/(2\beta)\}\)。
非负产量约束允许高成本企业停止生产，因此内点公式须逐家检验适用条件。

把内点一阶条件对所有企业相加，得到
\(n\alpha-\sum_i c_i-\beta(n+1)Q=0\)。
进一步代回单企条件：

\[
Q^*=\frac{n\alpha-\sum_i c_i}{\beta(n+1)},\quad
p^*=\frac{\alpha+\sum_i c_i}{n+1},\quad
q_i^*=\frac{\alpha+\sum_j c_j-(n+1)c_i}{\beta(n+1)}.
\tag{14-6}
\]

需要逐家检验 \(q_i^*>0\) 和正价格。由于一阶条件还给出 \(p^*-c_i=\beta q_i^*\)，内点经营利润可写为
\(\pi_i^{\text{经营}}=\beta(q_i^*)^2\)。
上述利润表达式由凹利润函数的一阶条件导出，其成立以相应内点条件为前提。

对称成本 \(c_i=c\) 时，式（14-6）化为：

\[
q^*=\frac{\alpha-c}{\beta(n+1)},\quad
Q^*=\frac{n(\alpha-c)}{\beta(n+1)},\quad
p^*=\frac{\alpha+nc}{n+1}.
\tag{14-7}
\]

其内点条件是 \(\alpha>c\)。同期技术变化之外，企业数固定、没有进入退出、没有产品质量差异、没有价格歧视和产能约束。这些假设适合展示数量竞争的一条机制，不保证十项技术的所有市场都具有同一结构。古诺模型的历史来源可核验于Cournot（1838）的国家图书馆原书目录；本章线性需求与成本扩展由上面方程自行推导，不声称照搬原书全部结论。\href{https://catalogue.bnf.fr/ark:/12148/cb30280488q}{法国国家图书馆原版书目}

\section{单件少用劳动，销量扩张仍可能增加总工时}

在对称均衡中，单位成本变化的价格传导率为：

\[
\frac{\partial p^*}{\partial c}=\frac{n}{n+1},\qquad
\frac{\partial Q^*}{\partial c}=-\frac{n}{\beta(n+1)}.
\tag{14-8}
\]

若两家企业同步降低单位成本20，价格下降 \(20\times2/3=13.333\)，而不是必然下降20。成本节约的一部分通过更低价格实现，另一部分改变企业利润和销量。固定企业数的这一传导率，也不是在任意需求曲率或市场结构下都成立的定律。

把技术对成本的影响代入：

\[
\frac{dQ^*}{da}=\frac{nw}{\beta(n+1)a^2}>0,\qquad
\frac{dp^*}{da}=-\frac{nw}{(n+1)a^2}<0.
\tag{14-9}
\]

技术提高同时增加交付和降低价格。给定同样产品，直接劳动需要下降；但是行业总工时需要用新的交付量计算。令 \(k=m+\rho\)，则：

\[
H^*=\frac{Q^*}{a}
=\frac{n}{\beta(n+1)}
\left[\frac{\alpha-k}{a}-\frac{w}{a^2}\right],
\]

\[
\frac{dH^*}{da}
=\frac{n}{\beta(n+1)a^2}
\left[\frac{2w}{a}-(\alpha-k)\right].
\tag{14-10}
\]

内点要求 \(a>w/(\alpha-k)\)。在这一域内，工时可能先升后降：当 \(a<2w/(\alpha-k)\)，需求扩张增加的工时超过单位产品节约；越过该点，节约效应较强。式（14-10）也可以写成
\(d\ln H/d\ln a=d\ln Q/d\ln a-1\)。
销量对效率的弹性大于1，总工时就增加；小于1则减少。这使“生产效率提高多少就减少多少工时”的说法成为可以核对的条件命题。

本章主算例使用以下参数，选择它们是为了让相反机制可见，不是因为已估计出现实需求截距120或真实工资60。

\begingroup\footnotesize\setlength{\tabcolsep}{4pt}\renewcommand{\arraystretch}{1.25}
\begin{longtable}{>{\raggedright\arraybackslash}p{\dimexpr 0.36364\textwidth-8.0pt\relax}>{\raggedright\arraybackslash}p{\dimexpr 0.21212\textwidth-8.0pt\relax}>{\raggedright\arraybackslash}p{\dimexpr 0.42424\textwidth-8.0pt\relax}}
\toprule
\textbf{参数} & \textbf{取值} & \textbf{单位与含义} \\
\midrule\endfirsthead
\toprule
\textbf{参数} & \textbf{取值} & \textbf{单位与含义} \\
\midrule\endhead
\bottomrule\endfoot
\(\alpha,\beta\) & 120，1 & 需求截距为货币/产品；斜率为货币/产品数量平方 \\
\(n\) & 2 & 两家数量竞争企业 \\
\(w\) & 60 & 货币/工时，竞争阶段给定 \\
\(m,\rho\) & 10，10 & 货币/产品，真实其他投入与接入收费分列 \\
\(a_0,a_1\) & 1，1.5 & 产品/工时 \\
\(F\) & 300 & 货币/企业，本期采用费用 \\
局部社会账户 \(b_s,\kappa,s\) & 30，0，1 & 劳动机会成本/工时、接入真实成本/产品、固定费真实资源份额 \\
\end{longtable}\endgroup


双方未采用时 \(c_0=80\)，价格93.333，总交付26.667，总工时26.667。双方采用后 \(c_1=60\)，价格80，总交付40，总工时仍为26.667，工资总额均为1600。效率与销量都提高50\%，因而总工时恰好不变。沿连续效率变化计算，工时在 \(a=1.2\) 达到约27.778，再下降；这说明端点“没有变化”并不表示中间导数始终为零。

\begin{center}\begin{tikzpicture}\begin{axis}[width=.31\textwidth,height=5.6cm,
title={价格},xlabel={生产率 a},ylabel={货币/件},
tick label style={font=\scriptsize},label style={font=\small},title style={font=\small},
legend style={font=\scriptsize,draw=none,at={(0.5,-0.3)},anchor=north},legend columns=1,
axis lines=left,grid=major,grid style={gray!15},scaled ticks=false,xtick={1,1.2,1.5,2},]
\addplot+[blue!70!black,mark=none,thick,] coordinates {(1,93.3333333) (1.01,92.9372937) (1.02,92.5490196) (1.03,92.1682848) (1.04,91.7948718) (1.05,91.4285714) (1.06,91.0691824) (1.07,90.7165109) (1.08,90.3703704) (1.09,90.030581) (1.1,89.6969697) (1.11,89.3693694) (1.12,89.047619) (1.13,88.7315634) (1.14,88.4210526) (1.15,88.115942) (1.16,87.816092) (1.17,87.5213675) (1.18,87.2316384) (1.19,86.9467787) (1.2,86.6666667) (1.21,86.3911846) (1.22,86.1202186) (1.23,85.8536585) (1.24,85.5913978) (1.25,85.3333333) (1.26,85.0793651) (1.27,84.8293963) (1.28,84.5833333) (1.29,84.3410853) (1.3,84.1025641) (1.31,83.8676845) (1.32,83.6363636) (1.33,83.4085213) (1.34,83.1840796) (1.35,82.962963) (1.36,82.745098) (1.37,82.5304136) (1.38,82.3188406) (1.39,82.1103118) (1.4,81.9047619) (1.41,81.7021277) (1.42,81.5023474) (1.43,81.3053613) (1.44,81.1111111) (1.45,80.9195402) (1.46,80.7305936) (1.47,80.5442177) (1.48,80.3603604) (1.49,80.1789709) (1.5,80) (1.51,79.8233996) (1.52,79.6491228) (1.53,79.4771242) (1.54,79.3073593) (1.55,79.1397849) (1.56,78.974359) (1.57,78.8110403) (1.58,78.649789) (1.59,78.490566) (1.6,78.3333333) (1.61,78.1780538) (1.62,78.0246914) (1.63,77.8732106) (1.64,77.7235772) (1.65,77.5757576) (1.66,77.4297189) (1.67,77.2854291) (1.68,77.1428571) (1.69,77.0019724) (1.7,76.8627451) (1.71,76.7251462) (1.72,76.5891473) (1.73,76.4547206) (1.74,76.3218391) (1.75,76.1904762) (1.76,76.0606061) (1.77,75.9322034) (1.78,75.8052434) (1.79,75.679702) (1.8,75.5555556) (1.81,75.4327808) (1.82,75.3113553) (1.83,75.1912568) (1.84,75.0724638) (1.85,74.954955) (1.86,74.8387097) (1.87,74.7237077) (1.88,74.6099291) (1.89,74.4973545) (1.9,74.3859649) (1.91,74.2757417) (1.92,74.1666667) (1.93,74.0587219) (1.94,73.95189) (1.95,73.8461538) (1.96,73.7414966) (1.97,73.6379019) (1.98,73.5353535) (1.99,73.4338358) (2,73.3333333)};
\addplot+[red!70!black,mark=none,thick,only marks,mark=*] coordinates {(1,93.3333333) (1.2,86.6666667) (1.5,80)};
\end{axis}\end{tikzpicture}\hfill\begin{tikzpicture}\begin{axis}[width=.31\textwidth,height=5.6cm,
title={交付},xlabel={生产率 a},ylabel={件},
tick label style={font=\scriptsize},label style={font=\small},title style={font=\small},
legend style={font=\scriptsize,draw=none,at={(0.5,-0.3)},anchor=north},legend columns=1,
axis lines=left,grid=major,grid style={gray!15},scaled ticks=false,xtick={1,1.2,1.5,2},]
\addplot+[blue!70!black,mark=none,thick,] coordinates {(1,26.6666667) (1.01,27.0627063) (1.02,27.4509804) (1.03,27.8317152) (1.04,28.2051282) (1.05,28.5714286) (1.06,28.9308176) (1.07,29.2834891) (1.08,29.6296296) (1.09,29.969419) (1.1,30.3030303) (1.11,30.6306306) (1.12,30.952381) (1.13,31.2684366) (1.14,31.5789474) (1.15,31.884058) (1.16,32.183908) (1.17,32.4786325) (1.18,32.7683616) (1.19,33.0532213) (1.2,33.3333333) (1.21,33.6088154) (1.22,33.8797814) (1.23,34.1463415) (1.24,34.4086022) (1.25,34.6666667) (1.26,34.9206349) (1.27,35.1706037) (1.28,35.4166667) (1.29,35.6589147) (1.3,35.8974359) (1.31,36.1323155) (1.32,36.3636364) (1.33,36.5914787) (1.34,36.8159204) (1.35,37.037037) (1.36,37.254902) (1.37,37.4695864) (1.38,37.6811594) (1.39,37.8896882) (1.4,38.0952381) (1.41,38.2978723) (1.42,38.4976526) (1.43,38.6946387) (1.44,38.8888889) (1.45,39.0804598) (1.46,39.2694064) (1.47,39.4557823) (1.48,39.6396396) (1.49,39.8210291) (1.5,40) (1.51,40.1766004) (1.52,40.3508772) (1.53,40.5228758) (1.54,40.6926407) (1.55,40.8602151) (1.56,41.025641) (1.57,41.1889597) (1.58,41.350211) (1.59,41.509434) (1.6,41.6666667) (1.61,41.8219462) (1.62,41.9753086) (1.63,42.1267894) (1.64,42.2764228) (1.65,42.4242424) (1.66,42.5702811) (1.67,42.7145709) (1.68,42.8571429) (1.69,42.9980276) (1.7,43.1372549) (1.71,43.2748538) (1.72,43.4108527) (1.73,43.5452794) (1.74,43.6781609) (1.75,43.8095238) (1.76,43.9393939) (1.77,44.0677966) (1.78,44.1947566) (1.79,44.320298) (1.8,44.4444444) (1.81,44.5672192) (1.82,44.6886447) (1.83,44.8087432) (1.84,44.9275362) (1.85,45.045045) (1.86,45.1612903) (1.87,45.2762923) (1.88,45.3900709) (1.89,45.5026455) (1.9,45.6140351) (1.91,45.7242583) (1.92,45.8333333) (1.93,45.9412781) (1.94,46.04811) (1.95,46.1538462) (1.96,46.2585034) (1.97,46.3620981) (1.98,46.4646465) (1.99,46.5661642) (2,46.6666667)};
\addplot+[red!70!black,mark=none,thick,only marks,mark=*] coordinates {(1,26.6666667) (1.2,33.3333333) (1.5,40)};
\end{axis}\end{tikzpicture}\hfill\begin{tikzpicture}\begin{axis}[width=.31\textwidth,height=5.6cm,
title={直接工时},xlabel={生产率 a},ylabel={工时},
tick label style={font=\scriptsize},label style={font=\small},title style={font=\small},
legend style={font=\scriptsize,draw=none,at={(0.5,-0.3)},anchor=north},legend columns=1,
axis lines=left,grid=major,grid style={gray!15},scaled ticks=false,xtick={1,1.2,1.5,2},]
\addplot+[blue!70!black,mark=none,thick,] coordinates {(1,26.6666667) (1.01,26.7947587) (1.02,26.9127259) (1.03,27.0210827) (1.04,27.1203156) (1.05,27.2108844) (1.06,27.2932242) (1.07,27.3677468) (1.08,27.4348422) (1.09,27.4948798) (1.1,27.5482094) (1.11,27.5951627) (1.12,27.6360544) (1.13,27.6711828) (1.14,27.700831) (1.15,27.7252678) (1.16,27.7447483) (1.17,27.7595149) (1.18,27.769798) (1.19,27.7758162) (1.2,27.7777778) (1.21,27.7758805) (1.22,27.7703126) (1.23,27.7612532) (1.24,27.7488727) (1.25,27.7333333) (1.26,27.7147896) (1.27,27.6933887) (1.28,27.6692708) (1.29,27.6425696) (1.3,27.6134122) (1.31,27.5819202) (1.32,27.5482094) (1.33,27.51239) (1.34,27.4745675) (1.35,27.4348422) (1.36,27.3933103) (1.37,27.350063) (1.38,27.305188) (1.39,27.2587685) (1.4,27.2108844) (1.41,27.1616116) (1.42,27.1110229) (1.43,27.0591879) (1.44,27.0061728) (1.45,26.9520412) (1.46,26.8968537) (1.47,26.8406682) (1.48,26.7835403) (1.49,26.7255229) (1.5,26.6666667) (1.51,26.6070202) (1.52,26.5466297) (1.53,26.4855397) (1.54,26.4237927) (1.55,26.3614291) (1.56,26.2984878) (1.57,26.2350062) (1.58,26.1710196) (1.59,26.1065622) (1.6,26.0416667) (1.61,25.9763641) (1.62,25.9106843) (1.63,25.8446561) (1.64,25.7783066) (1.65,25.7116621) (1.66,25.6447477) (1.67,25.5775873) (1.68,25.5102041) (1.69,25.4426199) (1.7,25.3748558) (1.71,25.306932) (1.72,25.2388679) (1.73,25.1706817) (1.74,25.1023913) (1.75,25.0340136) (1.76,24.9655647) (1.77,24.8970602) (1.78,24.8285149) (1.79,24.759943) (1.8,24.691358) (1.81,24.622773) (1.82,24.5542004) (1.83,24.485652) (1.84,24.4171393) (1.85,24.348673) (1.86,24.2802636) (1.87,24.211921) (1.88,24.1436547) (1.89,24.0754738) (1.9,24.0073869) (1.91,23.9394022) (1.92,23.8715278) (1.93,23.803771) (1.94,23.7361392) (1.95,23.6686391) (1.96,23.6012772) (1.97,23.53406) (1.98,23.4669932) (1.99,23.4000825) (2,23.3333333)};
\addplot+[red!70!black,mark=none,thick,only marks,mark=*] coordinates {(1,26.6666667) (1.2,27.7777778) (1.5,26.6666667)};
\end{axis}\end{tikzpicture}\end{center}


\textbf{图14-1　竞争实现与工时反转。} 固定工资、企业数和线性需求；曲线来自101个公开参数点，标注采用前后与导数变号点。工时只含本竞争模块直接生产劳动，不含旧模型工程维护工时。纵轴分别为价格、交付和工时，不能当作同一福利指标。

\subsection{单位工资下降与预付资本节约属于不同结果}

本章竞争模型中的$W=wH$是本期工资支出。在所有所列工时均为同一生产范围的受雇劳动这一设定下，它反映该期购买劳动力的货币额；模型没有工资支付和销售回款日期，因而没有求出预付可变资本，更没有直接估计劳动力价值或剩余价值率。纳什模块中的$V$是可分配净收入，也不表示可变资本。

数值结果显示，单位工资成本由$w/a=60$降至40，而双方采用前后的行业工资总额同为1600。产量扩张恰好吸收了每单位省下的工时。因此，该算例支持“单位劳动力购买费用降低”，不支持“行业可变资本实际使用额减少”。单方采用和普遍采用还影响企业间的订单份额，企业自身减支与行业减支须分开。

周转则可用另一个独立假设展示。设全年实际使用的可变资本货币额为1200，等额资本每年完整周转6次，所需预付为200；若实际使用额不变而周转增至12次，预付为100。年度工资可以保持1200，减少的是为维持它所需的预付额。以上只是式（2-V2）的条件演算，未由竞争模块求得，也不以延迟付薪为缩短周期的手段。实际企业是否能取得这一结果，要看工资支付与货款回流的匹配。

\section{抢先采用包含不能由全行业同时取得的份额优势}

企业的单边采用比较以对手行为给定为条件，其收益不同于全行业同时改变技术后的平均利润。设原本所有企业成本为 \(c_0\)，一家采用后降为 \(c_1=c_0-\delta\)。由不对称均衡式：

\[
q_A=\frac{\alpha+(n-1)c_0-nc_1}{\beta(n+1)}
=q_B+\frac{n\delta}{\beta(n+1)},
\]

\[
q_R=\frac{\alpha+c_1-2c_0}{\beta(n+1)}
=q_B-\frac{\delta}{\beta(n+1)},
\quad q_B=\frac{\alpha-c_0}{\beta(n+1)}.
\tag{14-11}
\]

采用者扩大数量，竞争对手缩减数量。对手仍有正产量要求
\(\delta<\alpha-c_0\)，等价于 \(\alpha+c_1-2c_0>0\)。
若节约非常大，其他企业可能退出，式（14-11）不再是全部企业内点均衡。本算例 \(\delta=20<40\)，四种采用状态均通过内点检查。

抢先采用增加的经营利润为：

\[
G_{\text{先}}=\beta(q_A^2-q_B^2)
=\frac{2nq_B\delta}{n+1}
+\frac{n^2\delta^2}{\beta(n+1)^2}.
\tag{14-12}
\]

而全行业同步降成本后，单企数量只增加
\(\delta/[\beta(n+1)]\)，单企经营利润变化为：

\[
G_{\text{同}}=
\frac{2q_B\delta}{n+1}
+\frac{\delta^2}{\beta(n+1)^2}.
\tag{14-13}
\]

当 \(n>1,\delta>0,q_B>0\)，式（14-12）大于式（14-13）。抢先采用的优势包含从对手那里转移来的销量，由市场份额转移形成的相对优势，无法由所有企业同时取得。把同步收益代入企业的单家采用条件，会漏掉这一战略收益。

主算例中，独家采用者数量为26.667，对手为6.667，市场价格86.667。采用者经营利润711.111，而原本为177.778，故抢先采用阈值 \(G_{\text{先}}=533.333\)。支付300以后，它本期净利润411.111，仍高于原本177.778。全行业同步采用时，每家经营利润只有400，增加222.222；支付同样300以后，每家净利润100，低于原本水平。

这一结果来自两项变化：抢先者取得的相对销量优势不能由全部企业同时保有，普遍采用又使每家支付固定费用。能力提高和购买价格下降因而可以与企业净利润减少并存。原销量乘单位节约，只能给出忽略竞争反应的局部比较。

\section{抢先与跟随都须比较，固定费决定采用均衡}

为了确定采用是否会扩散，设两家企业先同时选择采用或不采用，再按既定成本进行数量竞争。四种状态对应的第二阶段均衡先求出，第一阶段再比较扣除 \(F\) 的收益。这样处理不同于把新技术成本直接填进一个全行业对称式：它明确保留了自身不采用、对手采用时的收益，完整刻画单边选择的比较基准。

\begingroup\footnotesize\setlength{\tabcolsep}{4pt}\renewcommand{\arraystretch}{1.25}
\begin{longtable}{>{\raggedright\arraybackslash}p{\dimexpr 0.11429\textwidth-8.0pt\relax}>{\raggedright\arraybackslash}p{\dimexpr 0.20000\textwidth-8.0pt\relax}>{\raggedright\arraybackslash}p{\dimexpr 0.11429\textwidth-8.0pt\relax}>{\raggedright\arraybackslash}p{\dimexpr 0.11429\textwidth-8.0pt\relax}>{\raggedright\arraybackslash}p{\dimexpr 0.22857\textwidth-8.0pt\relax}>{\raggedright\arraybackslash}p{\dimexpr 0.22857\textwidth-8.0pt\relax}}
\toprule
\textbf{企业1状态、企业2状态} & \textbf{\(q_1,q_2\)} & \textbf{价格} & \textbf{\(H\)} & \textbf{经营利润1、2} & \textbf{扣本期采用费后的净利润1、2} \\
\midrule\endfirsthead
\toprule
\textbf{企业1状态、企业2状态} & \textbf{\(q_1,q_2\)} & \textbf{价格} & \textbf{\(H\)} & \textbf{经营利润1、2} & \textbf{扣本期采用费后的净利润1、2} \\
\midrule\endhead
\bottomrule\endfoot
都不采用 & 13.333，13.333 & 93.333 & 26.667 & 177.778，177.778 & 177.778，177.778 \\
仅企业1采用 & 26.667，6.667 & 86.667 & 24.444 & 711.111，44.444 & 411.111，44.444 \\
仅企业2采用 & 6.667，26.667 & 86.667 & 24.444 & 44.444，711.111 & 44.444，411.111 \\
双方采用 & 20，20 & 80 & 26.667 & 400，400 & 100，100 \\
\end{longtable}\endgroup


对手未采用时，采用收益增加 \(533.333-F\)。对手已经采用时，自己若不采用只能赚44.444，若也采用经营利润400，所以跟随采用收益增加：

\[
G_{\text{跟}}-F=(400-44.444)-F=355.556-F.
\tag{14-14}
\]

当 \(F=300\)，抢先采用净增233.333，跟随采用净增55.556。无论对手选择什么，采用都严格提高自己的收益，因此双方采用是唯一严格均衡；但均衡中每家赚100，低于都不采用时的177.778。这是一个有限条件下的采用竞争困境，不是所有技术史都具有这种结局的断言。

三条阈值还让结论可以反转。\(F<222.222\) 时，双方同步采用也能使各自净利润高于旧技术；\(222.222<F<355.556\) 时，双方采用仍为严格均衡，却都比双方不采用少赚；\(355.556<F<533.333\) 时，对手不采用则愿意采用、对手已采用则不愿跟随，两个单家采用状态成为严格纯策略均衡；\(F>533.333\) 时，双方不采用是严格均衡。中间区间还存在对称混合策略均衡：双方各以 \(x=(G_{\text{先}}-F)/(G_{\text{先}}-G_{\text{跟}})\) 的概率采用，使对手两种行动的期望收益相等；例如F=400时x=0.75。这是策略随机化概率，不是现实采用率估计。在抢先或跟随阈值处，至少一项行动比较无差异，不能宣称唯一严格均衡。同步利润阈值222.222只表示双方采用与双方不采用的跨状态利润相等；此时抢先和跟随净收益仍严格为正，双方采用仍是唯一严格均衡。固定费不改变这四种给定成本状态下的第二阶段内点产量，但会改变哪些状态在第一阶段实现。

\begin{center}\begin{tikzpicture}\begin{axis}[width=.53\textwidth,height=7cm,
title={采用收益与固定费},xlabel={固定费 F},ylabel={净收益变化},
tick label style={font=\scriptsize},label style={font=\small},title style={font=\small},
legend style={font=\scriptsize,draw=none,at={(0.5,-0.3)},anchor=north},legend columns=1,
axis lines=left,grid=major,grid style={gray!15},scaled ticks=false,]
\addplot+[blue!70!black,mark=none,thick,] coordinates {(0,533.333333) (5,528.333333) (10,523.333333) (15,518.333333) (20,513.333333) (25,508.333333) (30,503.333333) (35,498.333333) (40,493.333333) (45,488.333333) (50,483.333333) (55,478.333333) (60,473.333333) (65,468.333333) (70,463.333333) (75,458.333333) (80,453.333333) (85,448.333333) (90,443.333333) (95,438.333333) (100,433.333333) (105,428.333333) (110,423.333333) (115,418.333333) (120,413.333333) (125,408.333333) (130,403.333333) (135,398.333333) (140,393.333333) (145,388.333333) (150,383.333333) (155,378.333333) (160,373.333333) (165,368.333333) (170,363.333333) (175,358.333333) (180,353.333333) (185,348.333333) (190,343.333333) (195,338.333333) (200,333.333333) (205,328.333333) (210,323.333333) (215,318.333333) (220,313.333333) (225,308.333333) (230,303.333333) (235,298.333333) (240,293.333333) (245,288.333333) (250,283.333333) (255,278.333333) (260,273.333333) (265,268.333333) (270,263.333333) (275,258.333333) (280,253.333333) (285,248.333333) (290,243.333333) (295,238.333333) (300,233.333333) (305,228.333333) (310,223.333333) (315,218.333333) (320,213.333333) (325,208.333333) (330,203.333333) (335,198.333333) (340,193.333333) (345,188.333333) (350,183.333333) (355,178.333333) (360,173.333333) (365,168.333333) (370,163.333333) (375,158.333333) (380,153.333333) (385,148.333333) (390,143.333333) (395,138.333333) (400,133.333333) (405,128.333333) (410,123.333333) (415,118.333333) (420,113.333333) (425,108.333333) (430,103.333333) (435,98.3333333) (440,93.3333333) (445,88.3333333) (450,83.3333333) (455,78.3333333) (460,73.3333333) (465,68.3333333) (470,63.3333333) (475,58.3333333) (480,53.3333333) (485,48.3333333) (490,43.3333333) (495,38.3333333) (500,33.3333333) (505,28.3333333) (510,23.3333333) (515,18.3333333) (520,13.3333333) (525,8.33333333) (530,3.33333333) (535,-1.66666667) (540,-6.66666667) (545,-11.6666667) (550,-16.6666667) (555,-21.6666667) (560,-26.6666667) (565,-31.6666667) (570,-36.6666667) (575,-41.6666667) (580,-46.6666667) (585,-51.6666667) (590,-56.6666667) (595,-61.6666667) (600,-66.6666667)};
\addlegendentry{抢先}
\addplot+[red!70!black,mark=none,thick,] coordinates {(0,355.555556) (5,350.555556) (10,345.555556) (15,340.555556) (20,335.555556) (25,330.555556) (30,325.555556) (35,320.555556) (40,315.555556) (45,310.555556) (50,305.555556) (55,300.555556) (60,295.555556) (65,290.555556) (70,285.555556) (75,280.555556) (80,275.555556) (85,270.555556) (90,265.555556) (95,260.555556) (100,255.555556) (105,250.555556) (110,245.555556) (115,240.555556) (120,235.555556) (125,230.555556) (130,225.555556) (135,220.555556) (140,215.555556) (145,210.555556) (150,205.555556) (155,200.555556) (160,195.555556) (165,190.555556) (170,185.555556) (175,180.555556) (180,175.555556) (185,170.555556) (190,165.555556) (195,160.555556) (200,155.555556) (205,150.555556) (210,145.555556) (215,140.555556) (220,135.555556) (225,130.555556) (230,125.555556) (235,120.555556) (240,115.555556) (245,110.555556) (250,105.555556) (255,100.555556) (260,95.5555556) (265,90.5555556) (270,85.5555556) (275,80.5555556) (280,75.5555556) (285,70.5555556) (290,65.5555556) (295,60.5555556) (300,55.5555556) (305,50.5555556) (310,45.5555556) (315,40.5555556) (320,35.5555556) (325,30.5555556) (330,25.5555556) (335,20.5555556) (340,15.5555556) (345,10.5555556) (350,5.55555556) (355,0.555555556) (360,-4.44444444) (365,-9.44444444) (370,-14.4444444) (375,-19.4444444) (380,-24.4444444) (385,-29.4444444) (390,-34.4444444) (395,-39.4444444) (400,-44.4444444) (405,-49.4444444) (410,-54.4444444) (415,-59.4444444) (420,-64.4444444) (425,-69.4444444) (430,-74.4444444) (435,-79.4444444) (440,-84.4444444) (445,-89.4444444) (450,-94.4444444) (455,-99.4444444) (460,-104.444444) (465,-109.444444) (470,-114.444444) (475,-119.444444) (480,-124.444444) (485,-129.444444) (490,-134.444444) (495,-139.444444) (500,-144.444444) (505,-149.444444) (510,-154.444444) (515,-159.444444) (520,-164.444444) (525,-169.444444) (530,-174.444444) (535,-179.444444) (540,-184.444444) (545,-189.444444) (550,-194.444444) (555,-199.444444) (560,-204.444444) (565,-209.444444) (570,-214.444444) (575,-219.444444) (580,-224.444444) (585,-229.444444) (590,-234.444444) (595,-239.444444) (600,-244.444444)};
\addlegendentry{跟随}
\addplot+[green!50!black,mark=none,thick,] coordinates {(0,222.222222) (5,217.222222) (10,212.222222) (15,207.222222) (20,202.222222) (25,197.222222) (30,192.222222) (35,187.222222) (40,182.222222) (45,177.222222) (50,172.222222) (55,167.222222) (60,162.222222) (65,157.222222) (70,152.222222) (75,147.222222) (80,142.222222) (85,137.222222) (90,132.222222) (95,127.222222) (100,122.222222) (105,117.222222) (110,112.222222) (115,107.222222) (120,102.222222) (125,97.2222222) (130,92.2222222) (135,87.2222222) (140,82.2222222) (145,77.2222222) (150,72.2222222) (155,67.2222222) (160,62.2222222) (165,57.2222222) (170,52.2222222) (175,47.2222222) (180,42.2222222) (185,37.2222222) (190,32.2222222) (195,27.2222222) (200,22.2222222) (205,17.2222222) (210,12.2222222) (215,7.22222222) (220,2.22222222) (225,-2.77777778) (230,-7.77777778) (235,-12.7777778) (240,-17.7777778) (245,-22.7777778) (250,-27.7777778) (255,-32.7777778) (260,-37.7777778) (265,-42.7777778) (270,-47.7777778) (275,-52.7777778) (280,-57.7777778) (285,-62.7777778) (290,-67.7777778) (295,-72.7777778) (300,-77.7777778) (305,-82.7777778) (310,-87.7777778) (315,-92.7777778) (320,-97.7777778) (325,-102.777778) (330,-107.777778) (335,-112.777778) (340,-117.777778) (345,-122.777778) (350,-127.777778) (355,-132.777778) (360,-137.777778) (365,-142.777778) (370,-147.777778) (375,-152.777778) (380,-157.777778) (385,-162.777778) (390,-167.777778) (395,-172.777778) (400,-177.777778) (405,-182.777778) (410,-187.777778) (415,-192.777778) (420,-197.777778) (425,-202.777778) (430,-207.777778) (435,-212.777778) (440,-217.777778) (445,-222.777778) (450,-227.777778) (455,-232.777778) (460,-237.777778) (465,-242.777778) (470,-247.777778) (475,-252.777778) (480,-257.777778) (485,-262.777778) (490,-267.777778) (495,-272.777778) (500,-277.777778) (505,-282.777778) (510,-287.777778) (515,-292.777778) (520,-297.777778) (525,-302.777778) (530,-307.777778) (535,-312.777778) (540,-317.777778) (545,-322.777778) (550,-327.777778) (555,-332.777778) (560,-337.777778) (565,-342.777778) (570,-347.777778) (575,-352.777778) (580,-357.777778) (585,-362.777778) (590,-367.777778) (595,-372.777778) (600,-377.777778)};
\addlegendentry{同步比较}
\addplot+[orange!80!black,mark=none,thick,dashed] coordinates {(0,0) (600,0)};
\addlegendentry{零线}
\end{axis}\end{tikzpicture}\hfill\begin{minipage}[b]{0.40\textwidth}\small\centering 给定 $F=300$ 的收益表\par\medskip\resizebox{\linewidth}{!}{\begin{tabular}{c|cc} & 对手不采用 & 对手采用\\\hline 不采用 & $177.778,177.778$ & $44.444,411.111$\\ 采用 & $411.111,44.444$ & $100,100$\end{tabular}}\par\medskip 行：企业1；列：企业2。\par 数字顺序：企业1、企业2。\par\medskip 同步利润比较不是第三项行动。\end{minipage}\end{center}


\textbf{图14-2　抢先、跟随与同步利润的不同阈值。} 左侧121个固定费点，零线以上为相应比较的净收益增加；右侧完整两企业收益表，数字顺序为企业1、企业2。同步利润线是行业共同变化的比较基准，并非企业面对给定对手时的第三项可选行动。

采用阈值规定盈利选择，预付款和服务资格规定选择能否实行。固定费减免、融资和准入改革因此作用于不同障碍。此外，重复配置可能提供冗余、可靠性或质量收益，本例没有计入这些结果，利润困境本身不足以证明部署浪费。

\section{企业少赚，不排除消费者和其他领取者多得}

仅以产品企业利润合计评价采用，会遗漏消费者、劳动者及接入供应者的收益；将接入费 \(\rho Q\) 当成已经消失的社会资源，又会重复扣除供应者收入。需要明确谁在核算边界内，以及哪些付款对应真实机会成本。

在线性需求与准线性货币评价下，消费总支付意愿为
\(U(Q)=\alpha Q-\beta Q^2/2\)，消费者剩余为
\(CS=U(Q)-pQ=\beta Q^2/2\)。
这里的支付意愿并不等于基本需要，也不自动给不同收入群体相同福利权重。用于局部总收益核算时，它只提供一项透明的评价约定。

令 \(b_s\) 为劳动每工时真实机会成本，\(\kappa\) 为接入方每产品真实资源成本，且 \(0\leq\kappa\leq\rho\)。固定费用中真实资源份额为 \(s\)，采用企业数为 \(A_N\)。把消费者、劳动者、产品企业及接入/许可供应者都放进边界，局部净收益写为：

\[
\Omega=U(Q)-b_sH-(m+\kappa)Q-sFA_N.
\tag{14-15}
\]

\(wH\) 在产品企业账户中构成工资成本，在社会收益账户中同时构成劳动收入。减去 \(b_sH\) 是机会成本约定，剩余 \((w-b_s)H\) 是劳动所得超过外部替代价值的部分。接入收费的收入转移为 \((\rho-\kappa)Q\)；固定许可等转移为 \((1-s)FA_N\)。因此同一收益也能由领取者分项表达：

\[
\Omega=CS+\sum_i\pi_i^{\text{净}}
+(w-b_s)H+(\rho-\kappa)Q+(1-s)FA_N.
\tag{14-16}
\]

式（14-15）和（14-16）在代码中逐状态核对。核算分别记录收入转移、真实资源投入与消费者剩余：工资和接入收入不另构成新增商品，固定费中的资源消耗与许可方收入也须避免重复计入。若接入设施确有维护和能源消耗，就以 \(\kappa\) 登记；若费用仅对应产权许可，其领取者收入仍须保留在分配账户中。

在本例 \(b_s=30,\kappa=0,s=1\) 的约定下：

\begingroup\footnotesize\setlength{\tabcolsep}{4pt}\renewcommand{\arraystretch}{1.25}
\begin{longtable}{>{\raggedright\arraybackslash}p{\dimexpr 0.13208\textwidth-8.0pt\relax}>{\raggedright\arraybackslash}p{\dimexpr 0.16981\textwidth-8.0pt\relax}>{\raggedright\arraybackslash}p{\dimexpr 0.16981\textwidth-8.0pt\relax}>{\raggedright\arraybackslash}p{\dimexpr 0.16981\textwidth-8.0pt\relax}>{\raggedright\arraybackslash}p{\dimexpr 0.16981\textwidth-8.0pt\relax}>{\raggedright\arraybackslash}p{\dimexpr 0.18868\textwidth-8.0pt\relax}}
\toprule
\textbf{采用状态} & \textbf{消费者剩余} & \textbf{两企业净利润合计} & \textbf{劳动所得减机会成本} & \textbf{接入方收入减真实成本} & \textbf{局部净收益 \(\Omega\)} \\
\midrule\endfirsthead
\toprule
\textbf{采用状态} & \textbf{消费者剩余} & \textbf{两企业净利润合计} & \textbf{劳动所得减机会成本} & \textbf{接入方收入减真实成本} & \textbf{局部净收益 \(\Omega\)} \\
\midrule\endhead
\bottomrule\endfoot
都不采用 & 355.556 & 355.556 & 800 & 266.667 & 1777.778 \\
仅一家采用 & 555.556 & 455.556 & 733.333 & 333.333 & 2077.778 \\
双方采用 & 800 & 200 & 800 & 400 & 2200 \\
\end{longtable}\endgroup


双方采用相对都不采用，局部净收益增加422.222，尽管产品企业净利润合计减少155.556。消费者剩余增加444.444，接入收入增加133.333，劳动收益分项不变，四项变化相加正好等于422.222。这个例子说明私人采用竞争、产业利润变化与社会收益不能互相替代，却也不证明现实采用的社会净收益必为正。若固定费主要是转移，若真实接入资源成本更高，若增加产量带来污染或照护负担，分项和边界都应改写。

本表也不支持“降低接入费一定等额增加社会收益”。在给定产量下把一笔纯转移换领取者，总量不变；但是降低企业面对的 \(\rho\) 会改变成本、价格和交付，因而可能经由行为产生真正的总量效果。该行为效应需要重新求解均衡，仅调整既有表格的收费分项不足以确定结果。若接入供应者在国外，本国收入边界与包括国外供应者的全球边界还会不同，不能用同一 \(\Omega\) 回答二者。

\section{给定工时与收入时，外部机会和谈判权决定劳动所得}

竞争模块固定 \(w\)，能够追踪成本与市场实现，却没有解释工资为什么是60，也没有解释采用收益怎样由合同分配。为此另建一个给定工时 \(H\)、可分配净收入 \(V\) 的谈判子问题。\(V\) 是双方待分配的收入总额，已扣除谈判之外的其他投入与约定支出，应与未扣除这些支出的销售总额区别，避免扩大工资和利润的可分配范围。

谈判失败时，劳动者获得外部所得 \(bH\)，企业获得外部收益 \(\bar\pi\)。在这一子问题中，\(b\) 描述劳动者的破裂所得，与上节社会影子机会成本 \(b_s=30\) 不强制相等。企业外部收益也不是一个保证无风险实现的利润，现实中可能来自替代劳动、停产选择或其他业务，须有证据估计。可合作剩余为：

\[
S=V-bH-\bar\pi.
\tag{14-17}
\]

只有 \(S>0\)，双方才能同时比破裂点更好。令劳动谈判权重为 \(0<\eta<1\)，使用货币线性效用和可转移所得，最大化加权净收益乘积：

\[
\max_{bH<W<V-\bar\pi}
(W-bH)^\eta(V-W-\bar\pi)^{1-\eta}.
\tag{14-18}
\]

Nash（1950）原研究给出净效用增益乘积的谈判框架；本章采用的非对称权重与线性货币分配，是进一步公开设定的子问题，不把 \(\eta\) 宣称为1950论文替当下工会估计出的系数。\href{https://doi.org/10.2307/1907266}{原论文DOI}

取对数，目标为
\(\eta\ln(W-bH)+(1-\eta)\ln(V-W-\bar\pi)\)。
一阶条件是
\(\eta/(W-bH)=(1-\eta)/(V-W-\bar\pi)\)；
其二阶导数为两个严格负项之和，所以内点解唯一。交叉相乘并整理：

\[
\eta(V-W-\bar\pi)=(1-\eta)(W-bH),
\]

\[
W=bH+\eta(V-bH-\bar\pi)=bH+\eta S.
\tag{14-19}
\]

企业所得为 \(V-W=\bar\pi+(1-\eta)S\)。所得分配分别取决于劳动外部所得、企业外部所得与合作剩余，不能仅以销售收入乘以预定份额代替工资决定。权重取0或1时按边界极限理解；没有正剩余时当前内点公式不适用。

给定 \(H,V\)，比较静态为：

\[
\frac{\partial W}{\partial V}=\eta,\quad
\frac{\partial W}{\partial \eta}=S,\quad
\frac{\partial W}{\partial b}=(1-\eta)H,\quad
\frac{\partial W}{\partial\bar\pi}=-\eta.
\tag{14-20}
\]

可分配收入增加、谈判权重提高或劳动外部机会改善，会增加劳动所得；企业破裂所得更高则压低这一子问题中的劳动份额。工资增长是否跟随收入增长，不是生产率一个参数独自决定的。若生产率变化还改变 \(H\)，则需要全微分
\(dW=\eta\,dV+(1-\eta)b\,dH+(1-\eta)H\,db-\eta\,d\bar\pi+S\,d\eta\)。
即使 \(\eta\) 不变，工时、外部机会和企业替代方案也会影响结果。

示例设 \(H=20,V=2000,b=20,\bar\pi=200,\eta=0.4\)，则 \(S=1400,W=960\)，企业所得1040，平均每工时劳动所得48。将 \(V\) 单独增至2400，劳动所得增加160、企业所得增加240；将 \(b\) 单独提高到40，劳动所得变为1200；将 \(\eta\) 单独提高到0.6，劳动所得变为1240。相同收入规模可以对应不同分配，劳动所得改善也不必全部来自产出增加。

\noindent\begin{minipage}{\textwidth}
\begin{center}\begin{tikzpicture}\begin{axis}[width=.44\textwidth,height=7cm,
title={局部社会账户},xlabel={采用状态},ylabel={货币单位},
tick label style={font=\scriptsize},label style={font=\small},title style={font=\small},
legend style={font=\scriptsize,draw=none,at={(0.5,-0.3)},anchor=north},legend columns=1,
axis lines=left,grid=major,grid style={gray!15},scaled ticks=false,ybar stacked,ymin=0,bar width=12pt,xtick={0,1,2},xticklabels={均未,单家,双方},]
\addplot+[blue!70!black,mark=none,thick,] coordinates {(0,355.555556) (1,555.555556) (2,800)};
\addlegendentry{消费者剩余}
\addplot+[red!70!black,mark=none,thick,] coordinates {(0,800) (1,733.333333) (2,800)};
\addlegendentry{劳动净收益}
\addplot+[green!50!black,mark=none,thick,] coordinates {(0,177.777778) (1,411.111111) (2,100)};
\addlegendentry{企业1}
\addplot+[orange!80!black,mark=none,thick,] coordinates {(0,177.777778) (1,44.4444444) (2,100)};
\addlegendentry{企业2}
\addplot+[violet,mark=none,thick,] coordinates {(0,266.666667) (1,333.333333) (2,400)};
\addlegendentry{接入收益}
\end{axis}\end{tikzpicture}\hfill\begin{tikzpicture}\begin{axis}[width=.44\textwidth,height=7cm,
title={给定条件的劳动所得},xlabel={谈判权重 \ensuremath{\eta}},ylabel={货币单位},
tick label style={font=\scriptsize},label style={font=\small},title style={font=\small},
legend style={font=\scriptsize,draw=none,at={(0.5,-0.3)},anchor=north},legend columns=1,
axis lines=left,grid=major,grid style={gray!15},scaled ticks=false,]
\addplot+[blue!70!black,mark=none,thick,] coordinates {(0,400) (0.01,414) (0.02,428) (0.03,442) (0.04,456) (0.05,470) (0.06,484) (0.07,498) (0.08,512) (0.09,526) (0.1,540) (0.11,554) (0.12,568) (0.13,582) (0.14,596) (0.15,610) (0.16,624) (0.17,638) (0.18,652) (0.19,666) (0.2,680) (0.21,694) (0.22,708) (0.23,722) (0.24,736) (0.25,750) (0.26,764) (0.27,778) (0.28,792) (0.29,806) (0.3,820) (0.31,834) (0.32,848) (0.33,862) (0.34,876) (0.35,890) (0.36,904) (0.37,918) (0.38,932) (0.39,946) (0.4,960) (0.41,974) (0.42,988) (0.43,1002) (0.44,1016) (0.45,1030) (0.46,1044) (0.47,1058) (0.48,1072) (0.49,1086) (0.5,1100) (0.51,1114) (0.52,1128) (0.53,1142) (0.54,1156) (0.55,1170) (0.56,1184) (0.57,1198) (0.58,1212) (0.59,1226) (0.6,1240) (0.61,1254) (0.62,1268) (0.63,1282) (0.64,1296) (0.65,1310) (0.66,1324) (0.67,1338) (0.68,1352) (0.69,1366) (0.7,1380) (0.71,1394) (0.72,1408) (0.73,1422) (0.74,1436) (0.75,1450) (0.76,1464) (0.77,1478) (0.78,1492) (0.79,1506) (0.8,1520) (0.81,1534) (0.82,1548) (0.83,1562) (0.84,1576) (0.85,1590) (0.86,1604) (0.87,1618) (0.88,1632) (0.89,1646) (0.9,1660) (0.91,1674) (0.92,1688) (0.93,1702) (0.94,1716) (0.95,1730) (0.96,1744) (0.97,1758) (0.98,1772) (0.99,1786) (1,1800)};
\addlegendentry{基准}
\addplot+[red!70!black,mark=none,thick,] coordinates {(0,800) (0.01,810) (0.02,820) (0.03,830) (0.04,840) (0.05,850) (0.06,860) (0.07,870) (0.08,880) (0.09,890) (0.1,900) (0.11,910) (0.12,920) (0.13,930) (0.14,940) (0.15,950) (0.16,960) (0.17,970) (0.18,980) (0.19,990) (0.2,1000) (0.21,1010) (0.22,1020) (0.23,1030) (0.24,1040) (0.25,1050) (0.26,1060) (0.27,1070) (0.28,1080) (0.29,1090) (0.3,1100) (0.31,1110) (0.32,1120) (0.33,1130) (0.34,1140) (0.35,1150) (0.36,1160) (0.37,1170) (0.38,1180) (0.39,1190) (0.4,1200) (0.41,1210) (0.42,1220) (0.43,1230) (0.44,1240) (0.45,1250) (0.46,1260) (0.47,1270) (0.48,1280) (0.49,1290) (0.5,1300) (0.51,1310) (0.52,1320) (0.53,1330) (0.54,1340) (0.55,1350) (0.56,1360) (0.57,1370) (0.58,1380) (0.59,1390) (0.6,1400) (0.61,1410) (0.62,1420) (0.63,1430) (0.64,1440) (0.65,1450) (0.66,1460) (0.67,1470) (0.68,1480) (0.69,1490) (0.7,1500) (0.71,1510) (0.72,1520) (0.73,1530) (0.74,1540) (0.75,1550) (0.76,1560) (0.77,1570) (0.78,1580) (0.79,1590) (0.8,1600) (0.81,1610) (0.82,1620) (0.83,1630) (0.84,1640) (0.85,1650) (0.86,1660) (0.87,1670) (0.88,1680) (0.89,1690) (0.9,1700) (0.91,1710) (0.92,1720) (0.93,1730) (0.94,1740) (0.95,1750) (0.96,1760) (0.97,1770) (0.98,1780) (0.99,1790) (1,1800)};
\addlegendentry{外部机会提高}
\addplot+[green!50!black,mark=none,thick,] coordinates {(0,400) (0.01,418) (0.02,436) (0.03,454) (0.04,472) (0.05,490) (0.06,508) (0.07,526) (0.08,544) (0.09,562) (0.1,580) (0.11,598) (0.12,616) (0.13,634) (0.14,652) (0.15,670) (0.16,688) (0.17,706) (0.18,724) (0.19,742) (0.2,760) (0.21,778) (0.22,796) (0.23,814) (0.24,832) (0.25,850) (0.26,868) (0.27,886) (0.28,904) (0.29,922) (0.3,940) (0.31,958) (0.32,976) (0.33,994) (0.34,1012) (0.35,1030) (0.36,1048) (0.37,1066) (0.38,1084) (0.39,1102) (0.4,1120) (0.41,1138) (0.42,1156) (0.43,1174) (0.44,1192) (0.45,1210) (0.46,1228) (0.47,1246) (0.48,1264) (0.49,1282) (0.5,1300) (0.51,1318) (0.52,1336) (0.53,1354) (0.54,1372) (0.55,1390) (0.56,1408) (0.57,1426) (0.58,1444) (0.59,1462) (0.6,1480) (0.61,1498) (0.62,1516) (0.63,1534) (0.64,1552) (0.65,1570) (0.66,1588) (0.67,1606) (0.68,1624) (0.69,1642) (0.7,1660) (0.71,1678) (0.72,1696) (0.73,1714) (0.74,1732) (0.75,1750) (0.76,1768) (0.77,1786) (0.78,1804) (0.79,1822) (0.8,1840) (0.81,1858) (0.82,1876) (0.83,1894) (0.84,1912) (0.85,1930) (0.86,1948) (0.87,1966) (0.88,1984) (0.89,2002) (0.9,2020) (0.91,2038) (0.92,2056) (0.93,2074) (0.94,2092) (0.95,2110) (0.96,2128) (0.97,2146) (0.98,2164) (0.99,2182) (1,2200)};
\addlegendentry{可分配收入提高}
\end{axis}\end{tikzpicture}\end{center}

\textbf{图14-3　社会收益分项与给定条件下的劳动所得。} 左图使用竞争算例及局部机会成本；右图使用独立的20工时、2000或2400收入谈判算例，两边没有自动联立。接入费用作为领取者收入保留，固定费真实资源份额默认1。右图三条曲线改变的是谈判条件，不是已经估计的制度改革效果。
\end{minipage}

\section{传导与分享须对应真实交易，数值相近不能校准模型}

推导产生的是条件命题，原始研究则帮助判断现实是否值得沿这些条件寻找证据。成本传导的相关证据来自Ganapati、Shapiro和Walker（2020）：研究美国1972—1997年六类较同质制造产品的工厂资料，涵盖纸箱、面包、水泥、混凝土、汽油炼制和胶合板，使用能源价格变化识别成本传导。原论文概括，短中期内能源变化引致单位边际成本增加1美元，产品价格平均约增加0.70美元，而行业存在明显异质性。观察期不是2020，平均也不是所有制造业的统一常数。\href{https://www.sganapati.com/files/Ganapati_Shapiro_Walker_Energy_Cost_Pass-Through_AEJ_AE_2020.pdf}{原始论文及数据方法，pp.305—306、315—318}

这个经验平均与本例两企业传导率 \(2/3\) 数量上接近，却不能据此认定这些行业恰有两家企业，也不能将能源成本上涨的局部估计当成AI提高效率后的降价幅度。研究本身讨论不完全竞争、投入替代和需求性质。因此，实测传导率若偏离 \(n/(n+1)\)，应检查需求曲率、质量、产能和市场结构，企业数的确定应依据实际市场结构，不能仅以传导率数值接近作为选择依据。前文光伏成本与终端价之间的差别，在这个环节转化为明确的识别问题：成本变化的具体项目，以及各交易层级价格对该变化的响应。

收益分配的相关证据来自Kline、Petkova、Williams和Zidar（2019）。其原研究把美国约2001—2011年的公开专利申请与企业税务、劳动者收入记录连接，比较首次申请中初次获准和初次被拒的企业。论文报告，对其识别的有价值专利，劳动者平均通过收入增加获得约每1美元新增剩余中的0.30美元；申请年就在企业的劳动者份额约为此的两倍，收益在不同群体之间分布不均。这个结果说明创新租金可以传给劳动者，同时也存在企业内部差异。\href{https://zidar.princeton.edu/sites/g/files/toruqf3371/files/zidar/files/kpwz.pdf}{原论文pp.1343—1347}

0.30不能直接填作本章 \(\eta\)。经验中的剩余构造、员工进入退出、收入变动与专利冲击不同于固定 \(H,V,b,\bar\pi\) 的静态分配；专利获准还改变排他权，不是单纯提高生产效率。它支持“需要测量谁获得收益”，不能单独证明加权Nash是唯一真实工资机制，更不能从该样本推出当下所有AI劳动者获得相同比例。

两项经验研究分别识别成本向价格传导、创新所得向劳动收入传递的一段关系。检验完整机制，还要将采用费、对手反应、人员流动和可分配收入接到同一市场与时期。跨行业结果提供比较依据，不能拼成同一笔技术收益的连续去向。

\section{有支出意愿，还须有商品交付与可用现金}

竞争模块的需求曲线把收入支出背景概括掉，谈判模块又固定可分配收入。要追问某些收入变动能否支持全经济持续交付，需要明确居民和企业的期初现金、当期收入支出及期末余额。这里保留原单期核算，其任务是检查支出与资源能否闭合，尚未将前述竞争与谈判结果联立为宏观均衡。Godley和Lavoie（2007）提供存量流量一致方法背景；当前保留的是最小现金和商品账户。\href{https://doi.org/10.1057/9780230626546}{正式出版DOI}

现金模块包含劳动收入居民、资本收入居民、企业聚合部门和预算平衡的转移政府。价格、时薪、投资及消费比例外生给定；两类居民期初现金各120，企业200，总量440，发行者对应负债固定为440。本期不含新增货币、银行信贷或外部资金。技术能力与实现量分别记为：

\[
Q=C+I,\quad H=h_0+Q/a,\quad Q_{\max}=a(\bar H-h_0),
\quad h_0=J+M.
\tag{14-21}
\]

\(J\) 为安装培训工时，\(M\) 为维护工时。这里的 \(H\) 比竞争模块多了工程维护劳动，\(Q_{\max}\) 是可产量而非库存。\(I\) 是与消费品同单位的真实投资品数量；在企业聚合部门内，投资品买方付款与卖方收入配对，交付只计一次。工资与工程支出不能各自再增加一项最终产出。

工资 \(W=wH\)，毛营业盈余 \(\Pi=pQ-W\)。它没有扣除折旧、利息和未设立的其他投入，不能与前面包含 \(m,\rho,F\) 的净利润直接比较。政府把正毛盈余比例 \(\tau\) 征收并全部转给劳动群体，企业把税后正盈余比例 \(d\) 分红：

\[
T=\tau\max(\Pi,0),\qquad
D=d\max(\Pi-T,0),\qquad pQ=W+\Pi.
\tag{14-22}
\]

消费支出意愿为
\(E_w=c_w(W+T)+A_w\)、\(E_k=c_kD+A_k\)。
\(A_w,A_k\) 来自各自期初现金，不能超提；\(c_w,c_k\) 是支出行为假设，不是所有权定义。意愿交付满足
\(Q_d=I+[E_w(Q_d)+E_k(Q_d)]/p\)，实现量为
\(Q=\min(Q_d,Q_{\max})\)。
能力不足时用同一比例
\(r=p(Q-I)/(E_w+E_k)\)
配给两群体消费，得到 \(C_w=rE_w/p,C_k=rE_k/p\)。

部门现金变化为：

\[
\Delta B_w=W+T-pC_w,\quad
\Delta B_k=D-pC_k,\quad
\Delta B_f=\Pi-T-D-pI.
\tag{14-23}
\]

由商品与收入恒等式，三者合计为零。企业中的 \(pI\) 抵掉已经包含在 \(pQ\) 中的部门内投资收入，留下来自居民消费的净收款。政府税与转移严格相等；转移变领取者，只有经支出和生产反馈才可能间接增加交付。现金下降不能由未说明的贷款自动补足，企业期末现金及预付工资能力均受约束。

正毛盈余分支令
\(\theta=c_w\tau+c_kd(1-\tau)\)，则：

\[
Q_d=
\frac{pI+(c_w-\theta)wh_0+A_w+A_k}
{p(1-\theta)-(c_w-\theta)w/a}.
\tag{14-24}
\]

若亏损，税和分红归零，须改用亏损分支；既有程序使用原始含最大值函数括定和二分。令 \(u=w/(pa)\)，两段支出反馈斜率分别为 \(c_wu\) 与
\(\theta+(c_w-\theta)u\)。
它们均小于1，且支出非负时，交付减意愿购买连续严格递增，有限固定点唯一。所有收入消费且全部分红可能导致给定投资和提用下没有有限无约束解；当前模型拒绝这一闭合，不能把数值截断当成均衡。

现金核算补上付款来源、商品约束和账户对应，尚未同时确定需求曲线与工资。联立还需明确生产部门、工资决定时点、最终与中间投入、购买行为，以及信贷和外部部门。账户闭合是联合解释的必要条件，不能代替这些行为关系。

\section{支出规则改变交付，产能与现金限制实现}

旧基准设置 \(a=1,J=0,M=2,I=20,w=1,p=2,\bar H=100\)，采用共同设置 \(a=1.5,J=6,M=4,I=24\)。居民支出比例0.95、0.5，分红比例0.8，期初提用7、3。与竞争算例相比，旧价格没有由产量竞争决定，工资分享和降价是外生合同规则。因此，它呈现的是给定规则后的闭合结果，不能倒过来当作竞争式已被验证。

\begingroup\footnotesize\setlength{\tabcolsep}{4pt}\renewcommand{\arraystretch}{1.25}
\begin{longtable}{>{\raggedright\arraybackslash}p{\dimexpr 0.16327\textwidth-8.0pt\relax}>{\raggedright\arraybackslash}p{\dimexpr 0.18367\textwidth-8.0pt\relax}>{\raggedright\arraybackslash}p{\dimexpr 0.16327\textwidth-8.0pt\relax}>{\raggedright\arraybackslash}p{\dimexpr 0.16327\textwidth-8.0pt\relax}>{\raggedright\arraybackslash}p{\dimexpr 0.16327\textwidth-8.0pt\relax}>{\raggedright\arraybackslash}p{\dimexpr 0.16327\textwidth-8.0pt\relax}}
\toprule
\textbf{旧核算情景} & \textbf{可产量} & \textbf{实际交付} & \textbf{实际工时} & \textbf{劳动群体购买力 \((W+T)/p\)} & \textbf{毛营业盈余} \\
\midrule\endfirsthead
\toprule
\textbf{旧核算情景} & \textbf{可产量} & \textbf{实际交付} & \textbf{实际工时} & \textbf{劳动群体购买力 \((W+T)/p\)} & \textbf{毛营业盈余} \\
\midrule\endhead
\bottomrule\endfoot
基准 & 98.000 & 78.615 & 80.615 & 40.308 & 76.615 \\
工资价格不变 & 135.000 & 76.200 & 60.800 & 30.400 & 91.600 \\
时薪分享，w=1.5 & 135.000 & 101.923 & 77.949 & 58.462 & 86.923 \\
价格传导，p=2/1.5 & 135.000 & 109.615 & 83.077 & 62.308 & 63.077 \\
缩时且提时薪 & 85.000 & 85.000 & 66.667 & 50.000 & 70.000 \\
盈余税转移，\ensuremath{\tau}=0.35 & 135.000 & 106.777 & 81.185 & 63.757 & 132.370 \\
只缩时 & 85.000 & 76.200 & 60.800 & 30.400 & 91.600 \\
\end{longtable}\endgroup


工资价格不变时，可产量扩张但实际交付减少约3.07\%；劳动群体支出减少没有被其他支出充分补足。它与新竞争模块交付增加并不矛盾：新模块固定一条有效需求曲线，旧核算让收入领取和消费规则改变需求。两者恰好标明了需要实证检验的条件差别，其结果应按照各自的需求设定和闭合规则解释，不构成可互换的经验反事实。

旧降价情景的工时83.077高于基准，也与新模块式（14-10）的可能扩张机制相呼应；但降价规则、采用工程劳动和支出闭合不同，不能要求两个算例数值相同。时薪分享使交付增加、工时略减和劳动购买力提高，构成默认需求不足结果的反例。缩时且提时薪受产能限制，实际消费配给约0.917，说明收入权利与商品供给需要同时核对。

只缩时并未改变实现量和工时，因为原本实际需求的60.8小时低于新上限66.667；改变上限没有自动改变工资和分配。这一结果不能当作工时制度无效的实证结论，但足以要求讨论区分总工时、每人工作时间与人数。就业人数 \(N\) 还需每人工时 \(h\)，通过 \(H=Nh\) 定义。当前两个数量模块都没有估计 \(N\)，不能直接报告技术造成失业多少。

投资闭合提供另一项反转。保持旧采用工资价格不变，在正毛盈余未触及产能时，要使交付不低于旧基准 \(Q_0\)，需：

\[
I^*=
\frac{Q_0[p(1-\theta)-(c_w-\theta)w/a]
-(c_w-\theta)wh_0-A_w-A_k}{p}.
\tag{14-25}
\]

默认阈值25.006，采用投资24；只增加约1.006真实投资品，在可融资条件下即可使交付不再下降。投资不是凭空的新需求，它对应真实产品和企业付款；也不因为能增加当期交付就保证设备有用。取消一次性工程劳动并恢复投资20的独立反事实交付62.640、工时45.760，不能当作上期现金递推产生的未来路径。

两类支出比例相同且全部正盈余分红时，\(\theta=c\) 不随转移税率变化，原程序的再分配中性例给出相同交付。这补充了谈判模块：改变领取者能否改变需求，还取决于领取者是否以不同方式支出。控制描述只改标签而不改方程，结果也不变；这表示模型没有描述控制效应，不表示现实改变采购、准入或盈余决策权无效。

现金模块另保留两组$81\times81$敏感性网格，共13122点，分别比较效率与价格传导、工资分享与税转移。交付变化范围约为[-4.156\%,87.476\%]和[-3.072\%,71.722\%]；网格面积不代表现实发生概率，$g=1$全列的零值属于同一基准的退化边界。原图和输出保存在模型目录，原三图集中呈现竞争、采用与谈判关系，新增图14-P1单列制度选择阈值。

\section{制度选择将补偿责任反馈到自动化决策}\label{sec:political-automation}

\subsection{从主体选择推进到制度结果}

Acemoglu、Gitmez与Shadmehr（2026）的\emph{Automation and Repression}把任务生产、劳动者集体行动和资本方政策选择相连接：自动化的分配后果改变政治约束，国家再比较管制、再分配与镇压。其互补性结论依赖资本方目标、集体行动及政策成本等设定。\href{https://doi.org/10.3386/w35336}{NBER工作论文35336}、\href{https://economics.mit.edu/sites/default/files/2026-06/Automation\%20and\%20Repression.pdf}{2026年6月5日作者稿}。本文借鉴“假设—选择—约束—收益—比较静态—反例”的论述顺序，构造下列约化模型；后文单独重算原文政策收益的效率敏感性，不声称复现其全球博弈、资本积累或政变模型。

原论文用政治风险解释再分配与镇压的比较，本文则以可执行的劳动参与约束刻画协商补偿，并允许限制政治行动的边际成本随部署增加。两种模型的成本函数与政治过程不同，推论不能相互替代。这个扩展的任务是明确现稿中谁选择技术、谁承担补偿，以及制度成本怎样反过来改变选择；它不预测任何具体国家的镇压概率。

原论文的分散均衡还保留了一个关键区别：劳动供给$L$固定时，工资总额为$wL=L$，资本收入为$rK=K$，资本扩大使劳动份额$L/(K+L)$下降，而不要求工资总额同时下降（原稿命题2，印刷第11页）。政治分配压力因而可能来自相对所得的变化，可变资本绝对减少不是其必要前提。本文下述净工资节约系数及参与要求仍是另外的约化设定，不能用原论文这项结果为其数值背书。

\subsection{净收入、劳动力购买与参与约束}

给定资本可用规模指数$K\geq0$，资本方能够选择被研究子流程的自动化部署强度$x\in[0,\bar x]$及新增劳动密集型任务规模$n\in[0,\bar n]$。$x=1$仅表示这个子流程达到部署上界，不表示全社会不再需要劳动。各参数均对应同一期、同一计价物；本节的$K$不是第二章价值构成中的$c_M$。令净收入在扣除所列非工资投入后、支付工资与政治费用前为

\[
Z(x,n;K)=Z_0(K)+\mu(K)x-\frac{d_x}{2}x^2
                 +gn-\frac{d_n}{2}n^2,
\quad d_x,d_n>0.
\tag{14-P1}
\]

$\mu(K)$衡量部署的初始净收入收益，$d_x$概括部署成本及边际收益递减；$g,d_n$分别刻画新增任务收益及其递减。该式是局部约化关系，不是已估计的技术生产函数，也不是社会福利。闲暇、失业风险和政治权利损失尚未全部计入，不能把最大化$Z$称为社会最优。

设原岗位平均时薪为$w_0$，原付酬工时为$H_0$，单位部署净节省工时为$h_x$；新增任务的平均时薪和单位工时为$w_n,h_n$。记$W_0=w_0H_0$、$\lambda=w_0h_x>0$、$\nu=w_nh_n>0$，则

\[
H=H_0-h_xx+h_nn,\qquad W=W_0-\lambda x+\nu n,
\qquad P=Z-W.
\tag{14-P2}
\]

这里$P$是扣除工资后的资本经营余额，不是剩余价值$s_M$的直接估计。$\lambda>0$预设所考察期间仍有净付酬劳动节约；需求扩张若已抵消节约，就应重估该参数，不能把本式套到所有市场。原岗位工时要求$H_0\geq h_x\bar x$，工资与收入按实际归属期计量。确认所列雇佣劳动的生产职能以后，$W$才为可变资本使用的货币记录提供相应指标，预付额仍须另查周转。

设劳动群体具有可执行的参与或实施约束，其在协商方案中要求的总所得为$\bar W(K)=W_0+q(K)$，其中$q(K)\geq0$。所需补偿为

\[
T(x,n;K)=[\bar W(K)-W(x,n)]_+
       =[q(K)+\lambda x-\nu n]_+.
\tag{14-P3}
\]

这是假设的集体参与约束，不是劳动力价值的定义，也不是客观确定的政治行动收益。$q$需由具体协议、公共保障和受影响者的替代条件识别。$T$为资本方向劳动群体的转移，不额外创造产出。可执行的补偿制度要求$P\geq T$；若资金不足，协商方案不可按现有参数实施，不能靠式（14-P3）凭空支付。

本节P1—P6采用线性预期经济收益，只比较协商补偿与镇压两个分支。协商方案满足参与约束后，假设本期实施不再被该集体行动阻止；经济风险没有另行随机化。这种确定协商与风险镇压的不对称须单独检验，后文以凹效用限定其适用范围。选择顺序是：先给定技术、资本规模与制度约束，再由资本方选择制度及其部署方案；劳动所得、补偿和政治费用随后按方案实现。国家按资本方收益作选择，是本节的比较假设，不是对所有国家目标的经验概括。劳动者在协商分支中通过可执行约束影响选择，政治失败在另一分支中压缩为外生风险；本文没有推导逐个劳动者的协调博弈。

\subsection{私人采用与协商补偿产生不同部署强度}

先令$n=0$，并记$A(K)=Z_0(K)-W_0$。忽略政策反馈的私人采用基准最大化$P=A+(\mu+\lambda)x-d_xx^2/2$；这只是参照选择，不宣称已构成多个企业的政治均衡。协商补偿分支$N$在$T>0$时的资本方收益为

\[
U_N=P-T=A(K)-q(K)+\mu(K)x-\frac{d_x}{2}x^2.
\tag{14-P4}
\]

\textbf{命题P1：补偿约束严格生效、融资可行且最优选择为内点时，私人采用基准高于协商部署强度。} 对目标求导，$P_x=\mu+\lambda-d_xx$，$(U_N)_x=\mu-d_xx$，二阶导数均为$-d_x<0$，故唯一解为

\[
x_{\mathrm{私}}=\frac{\mu+\lambda}{d_x},\qquad
x_N=\frac{\mu}{d_x},\qquad
x_{\mathrm{私}}-x_N=\frac{\lambda}{d_x}>0.
\tag{14-P5}
\]

工资节省在私人参照中增加经营余额；协商补偿却将这部分节省返还给劳动群体，资本方不能将其全部保留。因此，$x_N$恰好也是本模型最大化$Z$的部署强度。这是转移与目标函数的条件结果，不表明现实中的集体协商总能实现最优生产。若式（14-P5）越过选择区间，应比较端点或将选择投影到$[0,\bar x]$；后续内点收益公式不能继续直接套用。

\subsection{镇压偏好取决于补偿、风险和政治成本}

设限制劳动者政治行动的分支为$R$。其在既有产权延续条件下由资本方承担的成本为$C_R(x)=R_0+\kappa x$，其中$R_0,\kappa\geq0$；$\kappa$可概括随着部署扩大而增加的执行、问责或其他制度成本。镇压后既有资本方仍失去收益请求权的概率为$\chi\in[0,1)$。沿用仅在产权延续状态支付该费用的比较设定，其期望收益为

\[
U_R=(1-\chi)[P-R_0-\kappa x].
\tag{14-P6}
\]

若政治失败，原资本方收益记为零，不表示产品或收入凭空消失；新的收入归属没有在本节展开。$P\geq C_R$是该方案的融资条件。$U_R$只表示资本方的预期经济收益，没有把政治伤害与其他群体损失全部货币化，不能以它较高判断镇压具有社会正当性。若费用无论成败均须先支付，则目标函数和阈值均要重推。

\textbf{命题P2：在上述内点条件下，镇压与较高部署强度的互补性取决于边际政治成本。} 因为$(U_R)_x=(1-\chi)(\mu+\lambda-\kappa-d_xx)$，且$(U_R)_{xx}=-(1-\chi)d_x<0$，所以

\[
x_R=\frac{\mu+\lambda-\kappa}{d_x},\qquad
x_R-x_N=\frac{\lambda-\kappa}{d_x}.
\tag{14-P7}
\]

当$\kappa<\lambda$，限制政治行动省去的边际补偿超过相应边际制度成本，$x_R>x_N$；当$\kappa\geq\lambda$，这一排序不再成立。$\kappa=0$时$x_R=x_{\mathrm{私}}$，也不等于原论文的完整技术排序。本模型把政治成本设为固定项加线性项，不能照搬其按产出比例收费等条件下的全部结论。

\textbf{命题P3：制度选择须比较各分支分别优化后的收益，而不是在同一部署水平上比较费用。} 将两项最优选择代入，在内点及融资条件成立时有

\[
\begin{aligned}
U_N^*&=A-q+\frac{\mu^2}{2d_x},\\
U_R^*&=(1-\chi)\left[A-R_0+
                  \frac{(\mu+\lambda-\kappa)^2}{2d_x}\right],\\
\Delta^*&=U_R^*-U_N^*
 =q-\chi A-(1-\chi)R_0+
 \frac{(1-\chi)(\mu+\lambda-\kappa)^2-\mu^2}{2d_x}.
\end{aligned}
\tag{14-P8}
\]

资本方在$\Delta^*>0$时选择$R$，在$\Delta^*<0$时选择$N$，等于零时无差异。等价地，定义

\[
R_0^{\mathrm{临}}=A+\frac{(\mu+\lambda-\kappa)^2}{2d_x}
-\frac{A-q+\mu^2/(2d_x)}{1-\chi},
\quad R\text{被选择}\ \Longleftrightarrow\ R_0<R_0^{\mathrm{临}}.
\tag{14-P9}
\]

若临界值非正，非负的$R_0$不能满足该严格条件。基于包络关系，$\partial\Delta^*/\partial R_0=-(1-\chi)<0$、$\partial\Delta^*/\partial\kappa=-(1-\chi)x_R<0$；在镇压分支净收益为正时，提高$\chi$同样降低$\Delta^*$。这些导数分别对应制度成本和剩余风险，不能由一个笼统的“控制能力”指标代替。提高$q$则增加所需补偿并提高绕过约束的资本方收益差，这解释了为什么补偿要求与权利执行需要共同考察。

\subsection{资本扩大只在特定区间提高镇压收益差}

\textbf{命题P4：资本规模增加是否推动制度转换，取决于补偿要求、基准收益与部署收益的联合变化。} 令$A,\mu,q$随$K$变化，保持$\lambda,d_x,R_0,\kappa,\chi$不变，在内点分支上对式（14-P8）求导得到

\[
\frac{d\Delta^*}{dK}=q'(K)-\chi A'(K)
 +\frac{\mu'(K)}{d_x}
 \left[(1-\chi)(\lambda-\kappa)-\chi\mu(K)\right].
\tag{14-P10}
\]

即使$\mu'>0$使资本扩大与自动化收益互补，右侧也未必为正。新增可索取收入使政治失败的损失扩大，或边际制度成本足够高时，资本方仍可能选择协商。若在一个可行区间内该导数严格为正，$\Delta^*$至多穿越零线一次；再有区间两端异号，才存在唯一转换阈值。没有端点异号就不能宣称已经出现转换，超出内点或融资范围则需分段求解。

\textbf{推论P4.1：静态阈值不是长程制度收敛的证明。} 只有另外说明$K$怎样积累、路径是否越过阈值、政治与融资条件能否持续，才能把比较静态接成时间路径。本节不引入资本运动方程，不把一次$K$比较解释为资本主义必然转向镇压。原论文的长程分析属于另有积累与存续条件的模型，本节只借鉴其先求选择、再讨论转换的逻辑。

\subsection{新增劳动任务可以降低制度转换的收益差}

恢复$n$。协商分支在$T>0$时，$\nu n$增加的工资相应减少补偿；镇压分支则由资本方直接承担该工资。\textbf{命题P5：若$g>\nu>0$，新增任务的选择为内点、补偿与融资约束满足，则协商分支选择更多新增劳动任务。} 对两个目标分别求导，二阶导数为$-d_n$和$-(1-\chi)d_n$，从而

\[
n_N=\frac{g}{d_n},\qquad n_R=\frac{g-\nu}{d_n},\qquad
n_N-n_R=\frac{\nu}{d_n}>0.
\tag{14-P11}
\]

劳动收入参与补偿责任，使新增任务在协商分支中兼有扩大净收入与减轻转移支出的作用。该解释依赖参与约束严格生效；当新任务工资已超过所得要求时，补偿归零，不能继续把新增工资全部冲抵转移。此时应回到式（14-P3）的分段原始目标。$g\leq\nu$或最优选择超界时也需检查角点。

\textbf{命题P6：在同一内点分支中，允许新增劳动任务降低镇压相对于协商的资本方收益差，但未必改变制度选择。} 两分支分别优化后，新增项为

\[
\Delta_{\mathrm{新任务}}^*=
\Delta_{n=0}^*+
\frac{(1-\chi)(g-\nu)^2-g^2}{2d_n}
<\Delta_{n=0}^*.
\tag{14-P12}
\]

不等式来自$0<g-\nu<g$与$0<1-\chi\leq1$。它不能证明新增任务必然消除镇压偏好；原来的正收益差若足够大，新增项不足以反转。也不能将本结果冒充原论文新任务扩展的证明，因为本文的参与约束和任务费用另有设定。新的劳动机会是否具有足够需求、工资和可执行权利，才是将该推论接到部门案例的条件。

\subsection{反事实算例展示阈值及反转条件}

公开设定$A=60+2K$、$\mu=12+4K$、$q=8+6K$、$d_x=40$、$W_0=40$、$\lambda=8$、$R_0=10$、$\kappa=0$、$\chi=0.05$、$\bar x=1$；新增任务使用$g=20$、$\nu=10$、$d_n=20$、$\bar n=2$。收入为假设货币单位，$K$为指数，未按任何国家、行业或原论文图形校准。$q$随$K$增长是明确的参与约束设定，不能从资本扩大本身推定现实劳动群体提出了同幅补偿要求。

\begingroup\footnotesize\setlength{\tabcolsep}{4pt}\renewcommand{\arraystretch}{1.25}
\begin{longtable}{>{\raggedright\arraybackslash}p{\dimexpr 0.25\textwidth-8pt\relax}>{\raggedright\arraybackslash}p{\dimexpr 0.06\textwidth-8pt\relax}>{\raggedright\arraybackslash}p{\dimexpr 0.08\textwidth-8pt\relax}>{\raggedright\arraybackslash}p{\dimexpr 0.08\textwidth-8pt\relax}>{\raggedright\arraybackslash}p{\dimexpr 0.13\textwidth-8pt\relax}>{\raggedright\arraybackslash}p{\dimexpr 0.10\textwidth-8pt\relax}>{\raggedright\arraybackslash}p{\dimexpr 0.10\textwidth-8pt\relax}>{\raggedright\arraybackslash}p{\dimexpr 0.20\textwidth-8pt\relax}}
\toprule
\textbf{假设情景} & \textbf{$K$} & \textbf{$x_N$} & \textbf{$x_R$} & \textbf{$n_N/n_R$} & \textbf{$U_N^*$} & \textbf{$U_R^*$} & \textbf{资本方选择} \\
\midrule\endfirsthead
\toprule
\textbf{假设情景} & \textbf{$K$} & \textbf{$x_N$} & \textbf{$x_R$} & \textbf{$n_N/n_R$} & \textbf{$U_N^*$} & \textbf{$U_R^*$} & \textbf{资本方选择} \\
\midrule\endhead
\bottomrule\endfoot
基准，未提供新增任务 & 0 & 0.30 & 0.50 & 0.00/0.00 & 53.800 & 52.250 & 协商补偿$N$ \\
资本规模提高 & 1 & 0.40 & 0.60 & 0.00/0.00 & 51.200 & 56.240 & 镇压$R$ \\
边际政治成本提高至16 & 1 & 0.40 & 0.20 & 0.00/0.00 & 51.200 & 50.160 & 协商补偿$N$ \\
剩余政治风险提高至0.4 & 1 & 0.40 & 0.60 & 0.00/0.00 & 51.200 & 35.520 & 协商补偿$N$ \\
新增劳动任务可供选择 & 1 & 0.40 & 0.60 & 1.00/0.50 & 61.200 & 58.615 & 协商补偿$N$ \\
\end{longtable}\endgroup

\textbf{表14-P1　分别优化后的自动化与制度选择。} 各行是独立反事实；收益只对应资本方，未表示社会福利。补偿、融资及内点条件逐行核查，数值来自本文模型程序。

基准中资本方选择协商；$K$增至1后，其最优制度收益转向镇压。但提高$\kappa$或$\chi$、允许新增劳动任务，均使同一资本规模下的选择回到协商。这些反例否定从“资本增加”直接推出单一制度结果。两个分支的部署强度、收入和工资要分别计算，不能只保持$x$相同而挑选较低政策费用。

在$K=1$、没有新增任务时，协商分支工资为36.8、补偿17.2，劳动总所得为54；镇压分支工资为35.2。允许新增任务后，协商分支工资46.8、补偿7.2，总所得仍为54，$n_N=1$；镇压分支$n_R=0.5$。工资、转移、任务与资本方制度收益分别记录，不能把补偿计成新增工资或真实产出。

在$K\in[0,2]$的已核查可行内点区间，基准$\Delta^*$严格增加，并在$K\approx0.2349$越过零线；新增任务使转换阈值推迟至$K\approx1.3931$。这是本模型的两组独立反事实，不是观测到的政治时间线，也不是将新任务选项当作同一实际经济已经发生的后续阶段。

\begin{center}\begin{tikzpicture}
\begin{axis}[width=0.95\textwidth,height=7.2cm,
xlabel={资本可用规模指数 $K$},ylabel={资本方收益差 $U_R^*-U_N^*$},
xmin=0,xmax=2,ymin=-10,ymax=13,grid=major,
legend style={font=\footnotesize,at={(0.02,0.98)},anchor=north west},
tick label style={font=\small},label style={font=\small}]
\addplot[thick,blue!70!black] coordinates {(0.000000,-1.550000000) (0.100000,-0.890100000) (0.200000,-0.230400000) (0.300000,0.429100000) (0.400000,1.088400000) (0.500000,1.747500000) (0.600000,2.406400000) (0.700000,3.065100000) (0.800000,3.723600000) (0.900000,4.381900000) (1.000000,5.040000000) (1.100000,5.697900000) (1.200000,6.355600000) (1.300000,7.013100000) (1.400000,7.670400000) (1.500000,8.327500000) (1.600000,8.984400000) (1.700000,9.641100000) (1.800000,10.297600000) (1.900000,10.953900000) (2.000000,11.610000000)};
\addlegendentry{基准，无新增任务}
\addplot[thick,green!45!black] coordinates {(0.000000,-9.175000000) (0.100000,-8.515100000) (0.200000,-7.855400000) (0.300000,-7.195900000) (0.400000,-6.536600000) (0.500000,-5.877500000) (0.600000,-5.218600000) (0.700000,-4.559900000) (0.800000,-3.901400000) (0.900000,-3.243100000) (1.000000,-2.585000000) (1.100000,-1.927100000) (1.200000,-1.269400000) (1.300000,-0.611900000) (1.400000,0.045400000) (1.500000,0.702500000) (1.600000,1.359400000) (1.700000,2.016100000) (1.800000,2.672600000) (1.900000,3.328900000) (2.000000,3.985000000)};
\addlegendentry{新增劳动任务可供选择}
\addplot[thick,orange!85!black] coordinates {(0.000000,-6.110000000) (0.100000,-5.602100000) (0.200000,-5.094400000) (0.300000,-4.586900000) (0.400000,-4.079600000) (0.500000,-3.572500000) (0.600000,-3.065600000) (0.700000,-2.558900000) (0.800000,-2.052400000) (0.900000,-1.546100000) (1.000000,-1.040000000) (1.100000,-0.534100000) (1.200000,-0.028400000) (1.300000,0.477100000) (1.400000,0.982400000) (1.500000,1.487500000) (1.600000,1.992400000) (1.700000,2.497100000) (1.800000,3.001600000) (1.900000,3.505900000) (2.000000,4.010000000)};
\addlegendentry{较高边际政治成本}
\addplot[dashed,gray,forget plot] coordinates {(0,0) (2,0)};
\end{axis}\end{tikzpicture}\end{center}

\textbf{图14-P1　补偿、政治成本与新增任务改变制度选择阈值。} 正值表示在该假设模型中资本方偏好镇压分支，负值偏好协商补偿。曲线为分别优化所得的收益差，并非镇压发生概率；零线是无差异边界。图内数据嵌入源稿，独立PNG、SVG与完整CSV另存。


\subsection{数学验证与经验对应分别完成}

符号核验检查一阶条件、严格凹性、最优收益差和资本规模导数；原始分段目标另由一维有界优化及二维约束优化求解。100组随机可行内点参数中，解析部署与数值解最大差异为$4.72\times10^{-8}$，收益最大差异$5.69\times10^{-14}$，有限差分导数误差不超过$2.64\times10^{-9}$。工资、转移与资本余额的最大账户残差为$2.85\times10^{-14}$；另查了部署角点和补偿归零的分支。两项转换阈值由解析收益差与数值优化结果分别求根，结果一致。

这些核验确认给定模型的计算，不识别政治偏好或制度效果。经验对应需要同时记录部署净收益$\mu,d_x$，同一群体工资与任务变化$\lambda,\nu$，补偿资格和执行$q,T$，实际政治约束及成本$R_0,\kappa,\chi$。不能以算法监控、低劳动份额或一次冲突代替全部参数，更不能用$U_R>U_N$证明实际国家代表同一种资本方目标。

模型代码、参数、情景表、曲线和核验结果保存于独立文件，旧竞争、谈判与现金模型未被覆盖。政治扩展与原三个模块尚未联立：原模型给定价格和工资，本节给定净收入及政治约束。加入共同工资、产品需求、税收融资和劳动者逐人协调，才可能进一步构成可估计的联合体系。

\section{公共品效率与政治风险限定制度选择}\label{sec:political-robustness}

\subsection{成本不对称是条件设定，政策方向仍须检验}

前节说明补偿责任与政治成本能够改变部署选择，但其预期收益目标仍为线性，协商成功又被设为确定。要把这一解释推进到实际政治回应，还须检验公共服务效率、两种政策的失败风险及决策者的风险偏好。刘镇锐对\emph{Automation and Repression}的\href{https://www.zhihu.com/question/2083147887706366153/answer/2088453310651032052}{知乎评论}提出了公共品效率固定和政治选择风险中性两项问题。本文依据用户提供的评论文本分析，将其作为假设检验的线索；以下推导和算例另行核对，不把该评论当作实证结果。

原论文的镇压总成本为$RY$，并非固定货币额；固定的是成本占产出的比例$R$。资本收入份额$s_K$增加时，其相对负担$RY/(s_KY)=R/s_K$下降。再分配则把税收收入用于公共品，每单位支出带来的劳动者效用由固定系数$\alpha$表示。这两项设定支持原文的政策比较，却没有证明实际镇压成本和公共品效率总按这样的比例变化。（Acemoglu等，2026，印刷第13—18页）

廉价消费品增加，也不能直接等同于$\alpha$提高。该系数比较公共品支出与工资收入的效用；生产端的真实产出变化已经进入$Y$。若要另设效率增长，必须说明公共服务的生产率、相对价格、质量和覆盖如何变化，避免重复计入同一项技术收益。医疗、教育和照护的效率能否提高，及改善能否被失去工资的人实际取得，均需对应资料。

\subsection{公共品效率提高可以增加最优税率}

先在原文的再分配目标中固定劳动收入份额$s_L$，并保持$\alpha>2$及政权存续概率小于1的内部解条件。税率最大化$(s_L(1-\tau)+\alpha\tau)(1-\tau)$，因而

\[
\tau^*(s_L,\alpha)=\frac{\alpha-2s_L}{2(\alpha-s_L)},\qquad
V(s_L,\alpha)=\frac{\alpha^2}{4(\alpha-s_L)}.
\tag{14-S1}
\]

这里$V$是代入最优税率后的收益系数，并非劳动者效用总量。式（14-S1）对应原稿Lemma 2和式（23）；以下导数由本文核算：

\[
\frac{\partial\tau^*}{\partial\alpha}
=\frac{s_L}{2(\alpha-s_L)^2}\geq0,\qquad
\frac{\partial\tau^*}{\partial s_L}
=-\frac{\alpha}{2(\alpha-s_L)^2}<0.
\tag{14-S2}
\]

\textbf{命题P7：在给定劳动份额和内部解条件下，公共品效率提高不降低资本方最优再分配税率；劳动份额为正时严格提高。达到同一稳定程度的最低支出则可以减少。} 两者的区别在目标。对固定存续概率$p$，原模型给出$p\beta s_K=s_L+(\alpha-s_L)\tau$，所以可行目标的所需税率为$(p\beta s_K-s_L)/(\alpha-s_L)$；当分子为正且目标可达时，效率提高降低这一税率。资本方却不固定$p$，而是在少缴税与减少政治失败之间选择。稳定更容易取得时，资本方可以选择增加购买，而不是只保留原来的稳定程度。

因此，公共品效率提高不能通过“最优税率必然下降”来反驳原文。不过它可以提高再分配的间接收益，削弱镇压的相对优势。只要$s_L/\alpha\to0$且内部解条件保持，式（14-S1）的税率仍趋于$1/2$；仅令$\alpha$上升并不会自动取消这一极限。若效率提高使存续概率达到1，则应加入概率上界并重新求解，不能继续使用内部解的税率。

\subsection{效率增长能够抵消镇压的相对收益增长}

为检验原文的政策排序，本小节单独保留其CES生产结构，令原模型实物资本为$K_A$，与前节的规模指数$K$分开。劳动供给$L$固定，替代弹性$\sigma>0$且$\sigma\ne1$；以资本收入份额$s_K$重写产出：

\[
\mathcal Y(s_K;K_A)=
\left[\frac{s_K^\sigma}{K_A^{\sigma-1}}
      +\frac{(1-s_K)^\sigma}{L^{\sigma-1}}\right]^{-1/(\sigma-1)}.
\tag{14-S3}
\]

允许$\alpha=\alpha(K_A)$、$R=R(K_A)$和$\chi=\chi(K_A)$平滑变化，保持行动收益参数$\beta$固定。每个给定$K_A$下，这些参数不随政策所选的自动化水平另外变化。两种政策分别优化后的资本方预期收益为

\[
\begin{aligned}
\mathcal U_r^*(K_A)
 &=\max_{s_K\in[R,1]}(1-\chi)(s_K-R)\mathcal Y(s_K;K_A),\\
\mathcal U_s^*(K_A)
 &=\max_{s_K\in[0,1]}
   \frac{\alpha^2\mathcal Y(s_K;K_A)}{4\beta(\alpha-1+s_K)}.
\end{aligned}
\tag{14-S4}
\]

式（14-S4）的第二式仍要求优化解处存续概率小于1。本节只重算这两项政策目标，不重新证明劳动者全球博弈或资本积累均衡。设最优资本份额为$s_K^r,s_K^s$，并记$s_L^s=1-s_K^s$。定义固定份额下的产出资本弹性

\[
\eta(s_K;K_A)=
\frac{s_K^\sigma K_A^{1-\sigma}}
     {s_K^\sigma K_A^{1-\sigma}+(1-s_K)^\sigma L^{1-\sigma}}.
\tag{14-S5}
\]

\textbf{命题P8：固定公共品效率下的收益单交叉，不能直接推广到效率和政治成本随资本变化的情形。} 在最优解唯一、可微、融资约束不绑定且两种收益为正的区间，利用包络定理和式（14-S1），得到

\[
\begin{aligned}
\frac{d}{dK_A}\log\frac{\mathcal U_r^*}{\mathcal U_s^*}
={}&\frac{\eta(s_K^r;K_A)-\eta(s_K^s;K_A)}{K_A}
 -\frac{\chi'}{1-\chi}-\frac{R'}{s_K^r-R}\\
 &-\frac{\alpha-2s_L^s}{\alpha(\alpha-s_L^s)}\alpha'.
\end{aligned}
\tag{14-S6}
\]

撇号均表示对$K_A$求导。原文固定$\alpha,R,\chi$时，只剩第一项；其证明利用镇压下较高的资本份额建立正号。效率提高、镇压成本比例上升或镇压失败概率增加，都可能抵消这项收益优势。其中$\alpha>2$保证效率导数前的系数为正。只有右边小于零，才能说镇压的相对收益在该区间下降；只有收益比进一步越过1，才发生政策转换。角点或概率上界生效时需要重新比较最优值，不能沿用这一导数。

给定$L=10$、$\sigma=0.8$、$R=0.8$、$\chi=0.05$、$\beta=10$，分别优化每一种政策。固定效率的参照与效率增长路径$\alpha(K_A)=0.3K_A$在$K_A=10$处重合，在$K_A=20$处分开：

\begingroup\small\setlength{\tabcolsep}{4pt}\renewcommand{\arraystretch}{1.2}
\begin{longtable}{>{\raggedright\arraybackslash}p{\dimexpr .43\textwidth-8pt\relax}>{\centering\arraybackslash}p{\dimexpr .18\textwidth-8pt\relax}>{\centering\arraybackslash}p{\dimexpr .18\textwidth-8pt\relax}>{\centering\arraybackslash}p{\dimexpr .21\textwidth-8pt\relax}}
\toprule
情景 & $\mathcal U_r^*$ & $\mathcal U_s^*$ & 资本方选择 \\
\midrule\endfirsthead
\toprule
情景 & $\mathcal U_r^*$ & $\mathcal U_s^*$ & 资本方选择 \\
\midrule\endhead
\bottomrule\endfoot
$K_A=10,\ \alpha=3$ & 1.967 & 1.847 & 镇压 \\
$K_A=20,\ \alpha=3$，效率固定 & 3.870 & 2.586 & 镇压 \\
$K_A=20,\ \alpha=6$，效率提高 & 3.870 & 4.786 & 再分配 \\
\end{longtable}\endgroup

\textbf{表14-S1　公共品效率增长可以反转政策排序。} 表中数据由式（14-S3）和（14-S4）重算，不取自原论文，也没有经验校准。沿效率增长路径，资本方从镇压转向再分配。三行再分配存续概率约为0.402、0.273、0.488，税率约为0.368、0.412、0.466；政策反转与税率提高可以同时发生。

这个反例说明，固定效率是原结论适用范围的重要条件；它没有推翻原假设下的Lemma 4，也没有证明现实效率增长已经足够快。$\alpha(K_A)$刻画资本积累过程中伴随的供给效率改善，不等同于直接设定$\alpha(I)$。若效率随国家所选自动化水平变化，还会改变技术选择的一阶条件，应另行建模。现实检验须估计效率变化，不能把算例中的翻倍设定当作已观察到的趋势。

\subsection{风险厌恶的方向由两种政策的收益分布决定}

原文政治比较使用预期货币回报，动态部分则以$\log C+\phi'\log B$选择消费与遗赠，并沿政权存续的路径分析资本运动（印刷第22—24页）。这种处理没有把同一个凹效用函数用于政治失败风险。要检验风险厌恶，应先交代失败后可维持多少消费，再比较两种政策的完整分布。

令政策$j$的可行部署与转移方案为$z\in\mathcal Z_j(K)$，存续概率为$p_j$，两个状态下的消费为$c_j^S,c_j^F>0$。风险偏好一致的政策目标应写为

\[
\mathcal E_j^*(K)=\max_{z\in\mathcal Z_j(K)}
\left[p_j(z,K)u(c_j^S(z,K))
 +(1-p_j(z,K))u(c_j^F(z,K))\right].
\tag{14-S7}
\]

存续和失败后的收入均按实际消费来源界定；不能把资本请求权归零直接当作消费为零，然后代入对数效用。若失败后确实没有任何正消费，期望对数效用为负无穷，原来的有限收益比较不再适用。

\textbf{命题P9：风险厌恶本身不能确定资本方更偏好镇压还是再分配。} 为说明其方向，设再分配以0.7概率获得消费140、以0.3概率获得20；镇压以0.95概率获得105、以0.05概率获得20。两者期望消费为104与100.75，风险中性者选择再分配；期望对数效用却约为4.358与4.571，排序转向镇压。该例只是两状态分布的逻辑反例，没有从生产或集体行动模型推导这些概率。它说明风险厌恶可能偏好失败概率较低的方案，不能只扣减其中一种政策的风险收益。

“资本增加使赌注变大”也不足以单独改变选择。对数效用满足$E[\log(ac)]=\log a+E[\log c]$；两种政策所有状态的消费同比放大，效用差不变。常数相对风险厌恶效用同样保留这种排序。固定绝对风险厌恶、固定失败后消费底线或风险概率随资本变化，则是不同设定，需要各自求解。

\subsection{本文的确定协商假设使风险作用具有单向边界}

前节参与约束模型另有一项不对称：协商满足约束后，本期收益确定；镇压存在失败风险。令两种方案共有的、不会因政治失败丧失的消费资源为$B_0>0$，将前节单期余额视为本期可消费所得，则协商消费为$B_0+U_N$，镇压存续消费为$B_0+P-C_R$、失败消费为$B_0$。此处是显式新增的消费假设，不把经营余额在现实中一律当作消费。

\textbf{命题P10：在协商收益确定、消费来源一致且效用递增凹的条件下，风险厌恶不能把期望货币收益已经劣于协商的镇压方案变为更优。} 若前节分别优化的$U_R^*\leq U_N^*$，则任一可行镇压方案由Jensen不等式满足

\[
\begin{aligned}
E[u(c_R)]&\leq u(E[c_R])=u(B_0+U_R)\\
 &\leq u(B_0+U_N^*)=\mathcal E_N^*.
\end{aligned}
\tag{14-S8}
\]

协商下递增效用保留原有货币收益最大化方案；镇压下则必须重新优化，不能只把原来的最优收益代入效用后当作新的最优值。若$U_R^*>U_N^*$，风险厌恶仍可能使协商更有利，但是否反转取决于曲率、失败概率和消费底线。因而，前节P1—P6属于线性预期收益下的结果，不能宣称其阈值已经对任意风险偏好稳健。协商同样具有失败风险时，应回到式（14-S7），P10的单向结论不能继续套用。

本节将制度判断落实为三项比较：同样支出能够提供多少有效保障，两种政策分别承担多少成本，以及政治失败如何改变消费与收益请求权。它们共同限制“自动化推动镇压”的适用范围。效率提高与风险厌恶都不是自动通向某种制度的力量，关键在于其变化是否足以改变分别优化后的收益排序。

数值核验程序保存为\texttt{模型/公共品效率与政治风险.py}，情景CSV与核验JSON另存数据目录。税率、效率导数和产出弹性经过符号核对；三行算例各检查100001个份额点及边界，并以原始自动化和税率联合优化独立复核，最大收益差小于$5\times10^{-13}$。独立优化显式限制存续概率不超过1。计算核验支持上述假设反例，不识别现实的政策成本、风险偏好或制度效果。

\section{数学与账户核验通过，行为假设仍须查实际数据}

新模型的验证首先比较数学解析解与另一路数值求解。四个主采用状态及120组固定随机种子生成的内点参数，共124组，独立逐家最优反应迭代与解析数量最大差异约 \(1.12\times10^{-11}\)，最大一阶条件残差约 \(4.97\times10^{-14}\)。每家在1001个候选偏离数量上的利润，没有超过均衡利润。数值最优反应不是把解析解直接重写为断言，因此能发现成本代入、一阶条件和索引错误；有限网格偏离检查只提供补充验证，严格最优性仍来自负二阶导数。

代码拒绝负需求斜率、非有限成本与被错误套入内点式的退出情形。对于成本[20,110]、需求截距120和斜率1的角点，独立数值解为[50,0]，而不是保留解析内点式中的负数量。主采用博弈四格最小数量6.667、最大交付40，均满足内点及正价格要求。新模型固定费改变的是第一阶段收益，完整博弈逐格比较了抢先与跟随选择，没有只检查同步利润就宣布采用均衡。

对价格、交付和总工时的导数，用步长 \(10^{-4}\) 的中心有限差分在四个效率点核对，最大绝对误差约 \(7.05\times10^{-7}\)。该检查特别覆盖工时导数变号附近，不只画出预设向上或向下趋势。谈判公式另用对数目标的一维黄金分割最大化检验五个内点权重，所得货币分配与解析解最大差异约 \(3.64\times10^{-5}\)，并检查双方外部所得及收入相加等于 \(V\)。无正合作剩余的例子被拒绝，权重端点只按边界处理。

局部社会收益由“总支付意愿减真实机会成本”与“各领取者剩余相加”两种方式计算，四格最大差异约 \(4.55\times10^{-13}\)。另将固定费真实资源份额设为0、0.5、1，接入真实成本设为0、5、10，在四种采用状态下核对36个账户，避免只有全资源默认例通过而纯转移例遗漏领取者。121个固定费点也由完整收益表逐格检查纯策略采用均衡，而不是只按阈值给曲线贴标签。这些检验不验证机会成本数值真实，F、\ensuremath{\rho}的现实构成仍需项目账和合同支持。

旧模型及其文件保持原内容。已有7情景、6个有效边界和9项无效输入拒绝继续有效；现金最大残差约 \(2.13\times10^{-14}\)，敏感性网格会计最大残差约 \(5.68\times10^{-14}\)，最小部门期末现金115.941，最大反馈斜率0.84。新脚本运行前后对旧模型目录文件散列检查通过，避免把新版输出覆盖成旧核验的假延续。

解析推导、最优反应与账户完整性已有相应核验，现实是否符合这些行为则仍要数据。成本与成交价查传导，采用时点、固定投入与对手反应查抢先收益，合同、外部岗位与替代成本查分配，住户支出、企业现金、投资与外需查需求。各项须对应同一市场、主体与时期，才能从计算正确走到现实解释。

\section{收益经竞争与合同分配，谁能改条件仍须查权限}

本章得到四项相互区别的结论。\textbf{企业采用有利，不保证普遍采用后企业利润提高。} 抢先采用包含订单转移的相对收益，同行跟进后这种优势收窄；固定支出决定双方采用是否仍优于最初状态。主算例中采用是严格占优选择，两家最终各获100，低于原有177.778。较低固定费能够反转利润结果，因此问题落在竞争位置和投入条件，而非技术名称本身。

\textbf{企业利润下降，也不保证社会收益下降。} 同一算例中价格下降、交付增加，劳动总工时保持不变，消费者剩余与接入方收入增加。在公开的机会成本和资源边界下，局部净收益增加422.222。其经济意义是不同领取者的收益可以分化，不能以产品企业利润代表全部社会结果；支付意愿评价和遗漏的环境、照护等成本，又限制了该总量的福利解释。

\textbf{增加的可分配收入与劳动者实际取得的收益仍须分别解释。} 条件谈判表明，劳动外部所得、企业替代收益和权重共同影响劳动所得。现金核算进一步说明，收入由谁领取以及怎样支出，会影响给定条件下的交付；工资分享、降价和投资又能改变默认情景。竞争、谈判与支出分别提供一部分解释，联立模型还须同时确定工资、产量、收入及资金流，当前演算没有完成这一经验或理论闭合。



\textbf{制度成本能改变资本方技术偏好，资本扩大也不必然带来同一种政治回应。} 协商参与约束把工资节省接到补偿责任，镇压分支则由政治成本和失败风险限制。两分支分别优化以后才有制度比较；新增劳动任务还可改变所得与补偿。在给定算例中，资本规模提高使镇压更有利，提高政治成本、剩余风险或提供新增任务又使协商更有利。命题P1—P6证明线性预期收益下的条件关系；P7—P10进一步区分最优税率与固定稳定目标、效率增长与收益排序、以及两种方案的风险结构。效率增长反例使选择从镇压转向再分配，风险厌恶则没有脱离收益分布的统一方向。这些结果未估计现实政体选择。

本章分别检验经济收益与制度选择：采用动机、企业利润、社会收益和劳动所得没有无条件的同向关系，补偿责任和政治成本又可以反馈到部署与任务创造。模型明确了改变结果的条件，但未把现实所有权与政治权限估计为参数。第十五章因此结合费用、需求、谈判和可执行的制度约束比较措施，不用构造收益排序代替制度效果。

\section*{本章参考文献}

1. Cournot, A.-A. (1838). \emph{Recherches sur les principes mathématiques de la théorie des richesses}. Paris: L. Hachette。\href{https://catalogue.bnf.fr/ark:/12148/cb30280488q}{BnF原版目录}。本轮核验原书元数据，未声称通读法文原著；本章方程自行推导。
2. Nash, J. F., Jr. (1950). The bargaining problem. \emph{Econometrica, 18}(2), 155–162. \href{https://doi.org/10.2307/1907266}{DOI:10.2307/1907266}。读取\href{https://www.haverford.edu/sites/default/files/Nash1950.pdf}{大学公开原论文}谈判定义、净效用增益与乘积解相关段；非对称权重为本章假设。
3. Ganapati, S., Shapiro, J. S., \& Walker, R. (2020). Energy cost pass-through in US manufacturing: Estimates and implications for carbon taxes. \emph{American Economic Journal: Applied Economics, 12}(2), 303–342. \href{https://doi.org/10.1257/app.20180474}{DOI:10.1257/app.20180474}。读取作者公开正式论文的摘要、引言、数据方法及传导讨论相关页；未重跑受限普查微观数据。
4. Kline, P., Petkova, N., Williams, H., \& Zidar, O. (2019). Who profits from patents? Rent-sharing at innovative firms. \emph{Quarterly Journal of Economics, 134}(3), 1343–1404. \href{https://doi.org/10.1093/qje/qjz011}{DOI:10.1093/qje/qjz011}。核验正式出版元数据，读取作者机构公开论文摘要及引言样本/结果相关页；未宣称复现全部方法与微观记录。
5. Godley, W., \& Lavoie, M. (2007). \emph{Monetary economics: An integrated approach to credit, money, income, production and wealth}. Palgrave Macmillan. \href{https://doi.org/10.1057/9780230626546}{DOI:10.1057/9780230626546}。沿用前轮方法背景与书目核验，不声称当前模型复现其完整体系。
6. Acemoglu, D., Gitmez, A. A., \& Shadmehr, M. (2026). \emph{Automation and Repression}. NBER Working Paper 35336. \href{https://doi.org/10.3386/w35336}{DOI}；\href{https://economics.mit.edu/sites/default/files/2026-06/Automation\%20and\%20Repression.pdf}{6月5日作者稿}。政策比较见第13—20页，动态效用见第22—24页，Lemma 4证明见附录A-13—A-15页。
7. 刘镇锐（发表日期未核实）。对\emph{Automation and Repression}的评论［知乎回答］。\href{https://www.zhihu.com/question/2083147887706366153/answer/2088453310651032052}{原回答链接}。依据用户提供的评论文本分析；网页直接访问未成功，不将评论视为实证研究或同行评审论文。

经典方法来源超过十年，其用途是理论定位；两项实证研究的观察期与当下市场不同，本章没有用它们给2026行业参数背书。自有参与约束模型、代码和情景图由本项目构建；效率敏感性使用已标明来源的原文生产与政策收益公式，导数、反例及风险扩展另行核算。AI协助推导检查、写作与核验，真实研究数据未由本报告重新采集。原文只作必要概括，不复制原论文长段、表格或图像。全报告伦理、资助、利益冲突及作者贡献声明由总稿统一列示。


\chapter{制度条件与比较路径：经济约束与决策权限的调整}\label{ch:15}

\section{收益分享、采用机会与生产决定权需要不同制度}

多得收入、能够换服务、能够共同定规则，是不同结果。\textbf{改善分配，要改收益与负担怎样承担；扩大采用，要补进入和迁移条件；改变占有关系，则须调整生产条件归属与决定权。} 三类措施可以配合，本章按各自作用比较，不能用其中一项改善代替另外两项。

成本式$c=w/a+m+\rho$、固定采用支出$F$和分配式$W=bH+\eta(V-bH-\bar\pi)$，提供了比较费用、采用门槛与谈判条件的语言。分配式适用于合作剩余为正且工时、收入给定的情况。真实制度会同时影响多个条件：共同设施可以降低准备费用并扩大准入，劳动者参与可以改变报酬，也可以重新选择任务和技术。投资审批、经营目标及剩余用途超出这些参数，须由产权与实际决策资料判断。

制度效果以可信的比较基准判断。同期繁荣可能抬高收入，改革也可能只减少了原本会发生的损失。公共、私人、合作和开放安排均要说明实际解除的限制、费用承担与冲突时的决策程序，组织名称不能代报结果。

本文比较制度，明确追求三项结果：劳动者与使用者实际受益，必要费用不被无偿留给别的承担者，受影响者能够改变相关生产安排。利润和效率方程不能替我们选定这些目标。目标之间有冲突，就须分别列质量、服务、时间、财政与权限，说明先改什么、代价由谁担。先定评价目标，再查有效约束，改革选择才有具体依据。

\subsection{补偿的可执行性与政治约束影响技术路径}

第十四章的新模型将制度回应接到技术选择：补偿资格实际能够约束方案实施时，资本方需将工资节省同转移责任一起比较；限制政治行动的成本与失败风险，又影响能否绕过这一约束。因此，提出补偿标准、保证付款和维护参与权限是不同的制度任务。只有标准而不能执行，不能视为模型中的有效参与约束。公共服务的效率还影响同样财政支出能够兑现多少保障：税收金额、服务质量和实际覆盖要分别记录。第十四章的效率反例说明，这种改善可以改变资本方政策排序，但算例中的翻倍效率不是政策成效估计。

模型还说明，新岗位需要同时具备净收益、报酬和真实任务内容。在新增任务可供选择的反事实中，劳动工资增加并减少所需转移，协商方案重新成为资本方较有利的选择；它不支持只靠岗位名称或临时补贴稳定制度。较高政治成本也须对应可执行的法律、监督和问责，不是将任意数字写进模型就算实现了治理改革。

这些收益比较解释资本方在特定条件下如何选择，公共制度评价还要直接查政治权利、劳动者所得、安全与持续服务。原模型按资本方目标行事的假设，不能转为改革的公共目标；所有权、实际表决和申诉执行仍是第十五章判断决定权的依据。

\section{竞争与价格传导：进入和迁移需要可用资源}

在给定单位成本的数量竞争模型中，更多竞争者会扩大总供应、降低价格；单位成本下降向价格传导的比例为$n/(n+1)$。这些结果说明竞争可以使能力收益进入购买条件，但企业数在该模块中是给定的。以制度解释市场进入，还需纳入采用支出、认证要求、销售渠道和退出条件。一个有潜在技术能力的生产者，如果无法融资、无法证明质量，就没有成为实际竞争者。

降低进入支出和约束排他收费作用于不同位置。前者可能使原先无法采用者进入，后者使既有使用者的$\rho$下降。两者可以互补，也可能相互抵消。例如许可费用下降后，上游提高不可替代的维护费用，下游总成本未明显改善；终端平台数量增加，但都使用相同关键设施，也未必增加相应上游的替代选择。应按完整链条确定市场边界。

迁移权的实际价值由完整转换过程决定。资料导出之后能否被新系统使用，重新验证能否获客户认可，停工和培训费用能否承担，决定退出是否可行。制造与半导体案例中的认证适配，因此既是资源问题，也是使用者能否取得和保留必要资料的问题。

兼容可以少做不必要重复、减少单一依赖，安全、质量与维护的必要劳动仍须完成。共同标准可能互认测试，也可能规定只有哪些路径可进入。应在同一质量目标下比方案：费用补偿了真实验证，还是因为资料与认可只能从原提供者取得而增加限制。没有这个比较，认证贵本身无法分清必要投入与排他收费。

补贴效果取决于是否改变采用行为。第十四章$F=300$时，两家企业已都会采用；若财政各补200，使私人有效支出降至100，既定企业数下的产量和价格不变，企业净利润各增加200。公共部门支付400，真实准备投入没有因此减少。这一比较支持付款承担变化，不提供新增产出或资源节约的证据；其他质量、创新和进入效果需另行检验。



即使用户能换平台，劳动者也未必同样能换工作。多个平台也可能采用相近监控规则；企业争客户，也可能争熟练人员而改善劳动外部机会。哪种路径发生，须各自查。价格、劳动收入与工作自主性是不同结果，不能因一项改善就宣布整套制度有效。

\section{共同设施减少重复准备，财政代付改变付款承担}

共同实验室、计量设施、培训和测试接口可以让多个生产者复用准备投入，降低真实$F$。第八章的参考测量和资格核验、第十二章的原型平台表明，采用所需条件超出单台设备。共享的优势来自投入能够多次使用，成立范围由服务质量、容量与适配费用决定。

设每家自建相同能力的设施需600，共同设施也需600且可供两家使用。忽略新增协调费用，每家分摊300，准备资源由1200降为600。若两家仍分别自建，财政各补贴300，企业同样支付300，真实准备却仍为1200。\textbf{共享与代付可以产生相同私人门槛，社会资源结果不同。} 现实比较还需计入人员、材料、维护、排队和用户适配；免费服务可能通过等待与配额分配稀缺容量。

融资解决付款时点和风险承担，共享设施减少重复准备。贷款附带还本付息，担保可能使公共部门承担违约损失。第九章算例说明融资条件如何改变采用回报；实际评价仍需列明利息、违约、追偿、残值以及需求不足时的损失承担。

设施共同建设，还要规定谁能用。按先后来、按支付能力、给小组织留名额，各有目标和机会成本，须公开优先次序、等待、利用与复议。公共所有也不自动等于使用者共同决定；扩容、收费、技术方向仍由管理者单定，就尚未取得这些权限。谁能监督预算、修改规则，才说明治理改到哪一步。

\section{协商、缩时和保障作用于不同分配条件}

给定$H,V$的分配式显示三条不同路径。提高劳动权重$\eta$，增加劳动方取得的合作剩余；提高外部机会$b$，改善协议破裂时的所得；降低企业外部机会$\bar\pi$，改变企业拒绝协议的收益。这些参数并非现实制度的直接度量；集体协商、可迁移技能、其他岗位和生活保障，则提供了研究其变化的具体机制。

沿用第十四章的20工时、$V=2000$、$b=20$、$\bar\pi=200$例子，将$\eta$从0.2提高到0.4，劳动所得由680升至960，增加280；收入总额未变，这是分配变化。保持$\eta=0.4$而将$b$提高到40，劳动所得由960升至1200，变化通过外部机会发生。两种路径对应不同条件，数字不代表任何工会或保障政策的效果估计。

保障能否支持拒绝不利工作，要查离开后还能拿到什么。在职才有、退出即失去的待遇，未必改善破裂所得；资格、等待期与地区服务也会影响实际可用性。培训须形成带得走的能力，不能只发证书。企业能否迁移、外包或自动化，又影响另一方选择，双方条件须一起看。

有薪缩时涉及工时和收入两个变量。若产出保持、生产率提高，所需$H$下降；若每小时工资不变，工资总额随之下降。要保持原收入，就需要提高每小时报酬或建立其他支付，这会改变企业成本，必须重新比较价格、销量和采用。降低原模型的可用工时上限只改变生产约束，同等收入的可支付性仍需工资、产出和现金账户支持。第十四章将各种独立比较分开核算。

缩时须算全劳动。少了直接操作，多了待命、核验、家庭工作或即时响应，自由时间便未必增加。核心员工平均收入保持，外围和外包者却可能多担损失。第六、十三章的时间费用账户须进入协议，把这些劳动一起列清，才能知道谁真正分享。

协商时点限定可改变的事项。采购后主要讨论补偿，沉没的适配支出可能使替换方案更贵；采购前讨论质量、安全、培训和任务目标，则可能改变技术选择。早期参与需要可用资料、讨论时间和独立技术支持，会议次数不能代表这些条件已经具备。

OECD《就业展望》2023年第七章基于七国制造和金融调查讨论咨询与报告结果。咨询企业的劳动者更常报告健康安全改善，但原文明确不能确定因果，咨询本身也可能改变感知；该资料与Lane等的调查共享基础，不能算两次独立验证。（Lane等，2023，\href{https://doi.org/10.1787/ea0a0fe1-en}{调查来源}；OECD，2023，Figure 7.3及其后识别限制段，\href{https://www.oecd.org/en/publications/2023/07/oecd-employment-outlook-2023_904bcef3/full-report/social-dialogue-and-collective-bargaining-in-the-age-of-artificial-intelligence_5828691a.html}{章节原文}）它支持进一步研究咨询的过程，而不是将满意度直接换成健康的因果收益。

协商的效力由具体权限判断：知情提供依据，提议增加方案，谈判要求达成协议，否决或共同审批改变最终决定程序。协议还需核验、争议处理及运行资源，使权限在冲突时能够执行。席位比例、会议次数和满意度不能直接替代谈判权重或实际决定权。

\subsection{保障劳动力再生产，协商技术节约的实际用途}

可变资本分析把劳动政策的对象进一步分清。最低报酬和集体协商作用于劳动力购买条件；保薪缩时作用于省下时间的用途；培训和社会保障维持转岗期间的能力与生活；付款保障限制拖欠工资形成的劳动者垫付。这些改变可以改善所得与再生产，不能仅因企业工资支出上升就判断技术失效。

技术部署还需协商完整岗位和责任，而非只约定裁员人数。维护、核验和待命若被加到原人员身上，人数稳定也可能伴随强度和责任增加；外包后本企业工资减少，也可能伴随保障和训练责任转移。公开岗位工时、工资口径和付款周期，并让劳动者参与缩时、培训与安全安排，才把收益分享落到可执行事项。若进一步要求占有关系转化，还须改变投资、生产条件与剩余用途的决定权。

\section{代表席位不自动转化为工资收益，治理效应取决于具体权限}

Jäger、Schoefer和Heining利用德国1994年改革形成的设立日期差异：在改革覆盖范围内，不超过500名员工的新设股份公司取消三分之一监事会代表要求，既有相关公司保留。处理是代表要求的差异，而非实际席位随机分配。正式研究未检出工资、工资结构、租金分享和劳动份额效果，也未支持资本形成减少。（Jäger等，2021，\href{https://doi.org/10.1093/qje/qjaa038}{正式DOI}；\href{https://jaeger.scholar.princeton.edu/document/91}{2020年10月20日作者稿，§2及结果概述}）它约束了代表席位必然提高工资的判断；对管理监督、投资和不同群体的作用，还需相应权限及执行资料。

Gorton与Schmid研究的是另一处理：1989—1993年较大德国非金融上市公司中，一半与三分之一劳动代表的比较。其模型估计前者股权市场价值对账面价值比平均约低31\%；表4年度半对数系数按exp(系数)\ensuremath{-}1转换，五年平均约\ensuremath{-}30.66\%。这不是原始均值之比，更不是社会收益下降31\%的估计。（Gorton与Schmid，2004，pp.873、878—880；\href{https://doi.org/10.1162/1542476042782260}{DOI}，\href{https://spinup-000d1a-wp-offload-media.s3.amazonaws.com/faculty/wp-content/uploads/sites/20/2020/12/Capital-Labor-and-the-Firm-A-Study-of-German-Codetermination.pdf}{正式全文}）股东估值反映预期收入与风险，变化既可能涉及资源损耗，也可能涉及收益或权限转移，须与工资、投资、质量和风险分别比较。

两项研究分别比较代表要求与代表比例，观察期和结果也不同，不能合成参与制度的统一效果。向平台、制造或医疗外推，应比较派单争议、投资审批和服务优先顺序等具体权限及其执行。本章未重新估计原数据，也未将代表比例或估计系数填作$\eta$。

\section{财政分享需要接通付款、服务与持续运行}

收益分享可以通过企业协议，也可以通过财政覆盖企业之外的承担者。税收与转移改变可支配收入，公共购买提供设施和服务；二者不同。假设一项收入税取得$T$，其中转移为$R$、公共购买为$G$，在无借款和储备提用的当前比较中，必须满足$T=R+G$。财政付款有收款者，实际服务还有设备、劳动和材料，不能将两者重复当成新增资源。

给定总收入和其他条件，把企业所得的1000转给劳动家庭，若两者消费倾向分别为0.5和0.9，消费增加400。这个局部关系来自第四章，说明分配可能改变已有收入用途；它不证明技术能力增加，也不保证实际交付增加。若能力受限、进口增加或价格响应，支出效果会通过其他路径实现。企业投资和政府购买还需分别记录。

税收对采用的影响取决于征收和扣除规则。在一个假设税制中，经营收益与采用费用完全对称处理，税后采用差额为$(1-\tau)(\Delta\pi^{op}-F)$，$0\leq\tau<1$，符号不变。若采用支出不能抵扣、亏损不能同样处理，或税改变融资现金，就不能沿用这个结果。这里没有描述任何辖区现行税法，公式只是说明为什么“征税影响投资”还需要具体规则。

付款之外，采购还可能增加实际能力。建网络、聘临床人员、提供测量，可补其他主体的必要条件，同时也占用设备、劳动和材料。须查设施有没有用、质量怎样、少了哪些重复费用，再将资金来源与资源供给对应。第九章并网储能和第十一章疗程，都需要沿组织之间记录付款、交付和承担者。

公共项目的持续性取决于建设与运行两段安排。初始预算需要接上维护、更新和故障准备，并按使用期记录税源、融资与风险。年度拨款是付款流量，长期资产的当期资源消耗另需核算；当前赤字也不足以独立判断项目长期可持续性。

谁先得到支持，也要查。统一设备补贴，可能先由有配套资金者领走；压低公共采购价，又可能将未补偿费用压给供应商与劳动者。支付链之外的人也可能承担，应保留未进入合同者。公平可以采用不同权重，但须先列收入、服务、风险和实际费用，才能讨论公共利益。

\section{合作治理改变内部决定权，成员边界与外部依赖限定其范围}

合作治理的区别，在于成员能共同选择投入、经营目标与剩余用途。多分工资尚未单独说明有这种权限。现有竞争模型按利润最大化求解，没有同时求成员收入、工时、服务与积累的共同选择。因此，要看章程和实际程序，不能给模型换个所有制标签就称为合作治理。

ILO第193号建议书以自主、共同所有和民主控制界定合作社，并反对借合作形式规避劳动规范或建立伪装雇佣关系。这是关于成员关系的规范文本，现实执行须另行核查。（ILO，2002，第2与8(1)(b)段；\href{https://webapps.ilo.org/public/english/standards/relm/ilc/ilc90/pdf/rec-193.pdf}{原始文本}）成员资格决定谁分享所得、承担风险并参与表决。非成员雇员和外包人员可能处于不同规则下；应记录加入条件、资本支付、退出补偿及其劳动权利，内部成员改善不能代表全部参与者。

成员作主，也要解决钱和跨期取舍。自筹可能把家庭财产与经营风险放在一起，借款要还，外部股权又可能改投票和剩余权。现在多分、少留，可能影响设备更新、培训与后来成员。治理要让这些选择可讨论、可查、可问责，并列失败由谁承担，资源限制并不会因此消失。

Pencavel的2012年选择性综述分别讨论收入就业波动、效率、存续和民主治理，并区分收益激励与参与作用。本章使用其比较框架，没有独立核查综述中的全部案例。（Pencavel，2012，Introduction及Conclusion；\href{https://docs.iza.org/dp6932.pdf}{原工作稿}）效率结果不能单独验证治理权限，民主程序也需要检验其经营条件。

成员内部能作主，外部仍可能有单一客户认证、平台或融资合同限制。共同标准与设施可以扩展选择，实际能否使用却须验证。应同时肯定已形成的内部权限，并列仍受谁限制，才能界定共同治理改了哪些关系、还差哪些条件。

\section{开放扩大使用选择，维护和责任保证持续能力}

在许可允许相应用途的条件下，公开设计和模型权重可以扩大学习、修改与运行的选择，并降低部分$F$、$\rho$和重复研发投入。设备、材料、验证与维护仍决定交付能力。文档、测试和兼容需要持续劳动，免费文件只降低了特定取得费用。

开放项目可以由公共资助、共同出资或收费维持，各有预算与决定规则。能复制文件，未必能延续服务；出资后能选择维护方向，才取得相应治理权。长期停机、运行失败与成员外劳动的遗漏，都须在公共、合作和开放安排中同样查，不能因名称有利就省略责任。

医疗、交通与高责任制造还须查保护和核验。限制数据、道路权限或工艺资料，应说明对应什么风险、怎样验证、谁监督，再按同一质量目标比少些限制的方案。必要质量要求与认证资料带不走，是不同问题：前者可以保护使用者，后者可能增加对原提供者的依赖。

预防、监测、停用与赔偿，都要有能执行的主体。风险不可逆，事后付钱便不能替代事前预防；谁能暂停、谁能取资料、谁能获赔，会影响运行。核验若靠劳动者无偿完成，企业效率可能包含费用外移。补偿这些劳动并规定相应停机权限，才能把质量、承担和决定接起来。

\section{制度组合与改革顺序应依据部门瓶颈确定}

措施要对着实际限制。平台查收费和等待，制造查验证能力，储能查时序收入与融资，医疗查完整疗程，半导体查工艺、设备和知识配合。一环改善，另一环仍可能接不上。下表列的是机制比较与可能失败，不是新增政策效果估计。

\begingroup\footnotesize\setlength{\tabcolsep}{4pt}\renewcommand{\arraystretch}{1.25}
\begin{longtable}{>{\raggedright\arraybackslash}p{\dimexpr 0.13433\textwidth-8.0pt\relax}>{\raggedright\arraybackslash}p{\dimexpr 0.41791\textwidth-8.0pt\relax}>{\raggedright\arraybackslash}p{\dimexpr 0.44776\textwidth-8.0pt\relax}}
\toprule
\textbf{优先约束} & \textbf{直接改变的条件} & \textbf{配套与可能失败} \\
\midrule\endfirsthead
\toprule
\textbf{优先约束} & \textbf{直接改变的条件} & \textbf{配套与可能失败} \\
\midrule\endhead
\bottomrule\endfoot
采用准备与融资 & 共享验证和培训降低真实$F$；信用改变现金与折现条件 & 维护容量不足、补贴替代私人付款却未减少真实准备；强使用者优先获益 \\
关键入口费用 & 替代服务、可迁移合同和接口约束$\rho$与退出支出 & 终端竞争共用上游瓶颈；文件迁移不能完成质量和技能迁移 \\
劳动分配与选择 & 协商、保障与可携带技能影响$\eta$、$b$及企业外部机会 & 只咨询无决定效力；工时减少转为无偿核验；外围劳动者不被代表 \\
公共需求与服务 & 税、转移、采购改变支出用途；设施放松容量 & 资源与财政重复计算；建设有钱、运行无预算；价格降低但覆盖不扩展 \\
投资与剩余决定 & 可执行共同审批、信息和问责改变方案选择与用途 & 席位未形成实际权限；内部治理仍受单一客户、融资与认证限制 \\
\end{longtable}\endgroup


改革顺序应比较当前服务收益与未来转换费用。先采用专用系统可能积累难以收回的适配投入，先建立标准和培训则可能推迟交付。延迟成本、迁移成本、质量和风险在各部门不同，因此没有由本框架直接推出的统一顺序。

将来有收益，今天仍要撑过去。参加训练先花时间、生活和照护，换服务先要停用、重核验。只给席位和文件，可能留下承担不起的过渡费用。第十三章的家庭时间与保障应从方案开始就列入，才能说明机会实际给了谁。

推广评价应纳入未进入者、退出者和失败项目。资源充足的试点可能低估参加门槛，临时补贴可能掩盖运行缺口。规模扩大后，维护和监督容量能否增长也影响制度效果，技术复制费用下降不保证组织费用同步下降。

共同安排要逐项列权限。公共付建设费、私人定技术、劳动者谈工作、用户能带资料，各可能改善，也可能留下关键限制。冲突时谁最后决定、谁能暂停和复议、谁付损失，应有可核验规则。资金与运营混合，没有自动说明技术、责任和剩余怎样决定。

\section{改善分配有独立价值，占有关系转化还须改变生产条件支配}

竞争、保障、协商和公共服务，可以分别改善收入、工时、服务与费用承担，不必等所有权一并改变才有价值。共享让实际准备少重复，财政代付让私人少付款，也须按各自结果评价。比较的对象是哪些人实际得了什么、负担怎样变，不能由组织形式代报答案。

稳定收入可以撑住拒绝不利安排，缩时可以腾出学习与协商时间，但知情、提议、审批和否决仍要有协议、章程与执行。劳动者能共同改安全或培训规则，说明相关权限已变；投资目标与剩余仍由资本定，就说明改变的范围。分配与决定可以相互支持，应逐项说明已经做到哪里。

\textbf{占有关系的转化要求同时考察生产条件的归属与支配：谁掌握生产资料，谁选择投资和劳动组织，谁决定社会剩余的用途。} 提高$\eta$、降低$\rho$或改变$F$分别描述收入、费用和门槛，不能代替上述治理判断。公共与合作所有提供不同权利基础，其实质效力还受问责、成员范围、运行资源和外部依赖影响。

制度比较最终回答两组问题：收入、服务和负担改善了多少，相关生产事项由谁决定。前一组沿付款、容量和分配查结果，后一组沿资料、替代选择和可执行权限查程序。两者可以相互支持，却不是收入不断提高便自动走向共同占有的阶梯。第十六章据此归纳全文的经济结论和权利判断。

\section*{本章参考文献}

1. Lane, M., Williams, M., \& Broecke, S. (2023). The impact of AI on the workplace: Main findings from the OECD AI surveys of employers and workers. OECD Social, Employment and Migration Working Papers, No. 288. \href{https://doi.org/10.1787/ea0a0fe1-en}{DOI}.
2. OECD. (2023). \emph{OECD Employment Outlook 2023: Artificial Intelligence and the Labour Market}. 第7章Social dialogue and collective bargaining in the age of artificial intelligence. \href{https://doi.org/10.1787/08785bba-en}{DOI}；\href{https://www.oecd.org/en/publications/2023/07/oecd-employment-outlook-2023_904bcef3/full-report/social-dialogue-and-collective-bargaining-in-the-age-of-artificial-intelligence_5828691a.html}{章节原文}.
3. Jäger, S., Schoefer, B., \& Heining, J. (2021). Labor in the boardroom. \emph{The Quarterly Journal of Economics, 136}(2), 669–725. \href{https://doi.org/10.1093/qje/qjaa038}{DOI}.
4. Gorton, G., \& Schmid, F. A. (2004). Capital, labor, and the firm: A study of German codetermination. \emph{Journal of the European Economic Association, 2}(5), 863–905. \href{https://doi.org/10.1162/1542476042782260}{DOI}.
5. International Labour Organization. (2002). Promotion of Cooperatives Recommendation, 2002 (No. 193). Adopted 20 June 2002. \href{https://webapps.ilo.org/public/english/standards/relm/ilc/ilc90/pdf/rec-193.pdf}{原始文本}.
6. Pencavel, J. H. (2012). Worker cooperatives and democratic governance. IZA Discussion Paper No. 6932. 本章使用2012工作稿，不与2013书章视为独立研究。\href{https://docs.iza.org/dp6932.pdf}{原稿}.

本章对第06—14章的引用指向其具体证据与模型，不将重复提及作为新增来源数量。部门比较表、制度路径、权利地图与判准为报告原创综合分析；没有采集合同、访谈或个体行政数据。

\chapter{结论、局限与研究议程：技术收益的实现与分配条件}\label{ch:16}

\section{技术扩大可能，生产关系规定其支配与用途}

本文对新技术与资本主义基本矛盾的回应是：\textbf{技术进步能够改善劳动和生活，但不会使共同参加生产的人自动取得共同决定生产的权利。在本文讨论的资本主导组织中，新增生产能力仍通过资本的投资、用工和经营安排使用。} 基本矛盾的具体表现，因而不能只用利润或工资的增减判断，还包括共同形成的能力怎样受到支配。

全文的四项判断涉及劳动过程、生产组织、社会再生产和制度反馈：技术节约经需求、竞争与报酬形成收入和时间收益；知识与协作扩展仍可能受关键条件限制；家庭与公共投入需要完整核算；补偿责任、公共品效率和政治风险又能改变资本方技术选择。它们共同解释能力扩大与生产支配延续为何可以并存，并使制度回应成为待解释的结果。资本收益增加不是任何一种政治回应的充分条件。

证据对任务表现和局部收益的支持强于对决定权变化的支持。因此，本文解释技术改善与既有资本支配并存的条件，未证明所有新技术都强化控制。制度回应也分为三个层次：分享改变所得，准入改变可用条件，所有权与可执行治理改变决定权限。

\section{单项技术有效，还要看整个工作能否完成}

客服研究在一家企业5,172名人员的分批部署中发现，每小时成功解决问题数平均提高约15\%；咨询实验则同时观察到适用任务中的改善与越界任务中的正确率下降。它们确认局部能力进步，也要求把核验和交付纳入任务范围。固定任务量下13.04\%的工时节约是本文按客服结果换算，原研究并未据此估计就业下降。（Brynjolfsson等，2025；Dell’Acqua等，2026；\hyperref[ch:5]{第五章}）

2026年10月的数学发布进一步表明，部分研究论证可以通过组织模型计算批量生成。OpenAI于10月6日公开首批722篇数学稿件；撤回3篇后，10月8日目录为719篇、372组相关成果。稿件数不等于独立解决的难题数，全部论证也尚未经本文独立确认。结合推理训练和算法研发证据，可提出的判断是：\textbf{当生成和验证达到要求，追加研究产出的约束可能从专家逐项处理转向模型接入、算力和验证设施。} 问题选择、理解与协作仍需分别考察，现有材料尚不能证明所有智力工作都将完成这种转化。（\hyperref[ch:7]{第七章7.1—7.2}）

采用还取决于完整投入能否由预期服务覆盖。增材制造节省模具准备，优势主要在适当材料和小批量范围；物联网和数字孪生需要测量、整合与维护。光伏融资算例中，同一设备与电量在不同资金回收条件下具有不同经济可行性。参考测量、质量记录和网络设施因此是生产能力的组成部分，文件传送和单台设备表现未必提供完整交付。（\hyperref[ch:8]{第八章}；\hyperref[ch:9]{第九章表9-1及演算}）

医疗和先进制造亦需区分局部结果与持续服务。exa-cel研究44人接受输注，可评价30人中29人达到所列连续12个月终点，其疗效范围受入组与随访条件限定；EUV成像和电学试验转入量产还需材料、检测及稳定工艺。（Frangoul等，2024；\hyperref[ch:11]{第十一章}；\hyperref[ch:12]{第十二章}）技术收益由这些互补条件共同形成，准入、支付与容量决定哪些人能够使用。真实资源不足、私人排他和国家限制具有不同来源，应分别判断。

\section{竞争传递价格收益，也可能压低采用者净利润}

第十四章给出个别采用激励与行业利润分离的条件反例。统一参数下，抢先与跟随采用的经营利润增量分别为533.333和355.556，$F=300$使两家企业均愿升级；双方升级后各自净利润却由177.778降为100。能力提高与采用者获利减少可以并存，原因是竞争降价与采用支出改变了私人回报。较低$F$、需求扩张或新产品可以改变结果，这个未校准模型未估计现实企业平均损失。

价格收益和收入实现还经过时序、容量与交易安排。光伏发电量加权价格不同于全年市场均价，储能调整交付时段也产生设备和运行费用。中国新型储能同版资料显示，2021—2025年功率由4.00升至135.93GW、能量容量由8.00升至351.12GWh，全国利用小时2023—2025年由611升至1195。部署与使用扩大是实际变化，项目利润和终端电费仍需各自资料。独立与配建储能的分组利用差异也受地区、年龄与接入选择影响。（\hyperref[ch:9]{第九章表9-2、9-3}）

平台收费、模型接口与制造认证又影响节约的去向。廉价接口可以降低小组织的固定负担，开放权重可以提供替代方案；当替代不具备适用质量和容量时，使用者仍依赖原提供者。半导体专业协作可以支持创新，也可能使扩产和迁移受关键投入限制。下游价格竞争与上游依赖能够同时存在，不能只凭企业数量或设备规模判断控制。（\hyperref[ch:6]{第六章}；\hyperref[ch:7]{第七章}；\hyperref[ch:12]{第十二章}）

由此，竞争的两种后果可以同时出现：购买者以低价取得收益，采用者因同行反应和固定支出失去部分利润。社会收益评价还需纳入劳动者、供应者和公共部门的所得与负担，不能以任一领取者的利润代替完整结果。

价值构成的利润率压力与经营中的销售约束还须分开。第二章的$r=e/(1+k)$说明资本构成和剩余价值率如何共同作用，第十四章的竞争算例则说明售价、数量和采用支出如何改变货币利润。两者并不证明相同的因果过程。利润率下降可以与利润总额增长并存，生产能力扩张也可以与订单不足并存。基本矛盾在此表现为共同生产能力扩大，而其使用仍须服从各资本分别实现增殖的条件；现实危机的程度与形态，尚需进一步连接收入、部门比例、投资和信用。

\section{劳动者能否分享节约取决于需求、报酬和退出条件}

单位任务省时向工资与就业的传递并非同比。美国机器人地区研究观察到就业人口比和工资的负向效应，德国研究则发现制造业置换由服务业新增抵消，调整更多发生在年轻进入者的职业入口；跨国行业研究未检出总工时显著下降。不同结果支持同时考察置换、需求扩张和部门再配置，而非使用统一替代率。（Acemoglu与Restrepo，2020；Graetz与Michaels，2018；Dauth等，2021；\hyperref[ch:5]{第五章}）丹麦截至2024年的调查与行政记录还显示，AI相关任务调整、职业迁移关联可以与短期收入和登记工时近零的平均估计并存；平均数尚未明显变化，不等于劳动配置保持原状。（Humlum与Vestergaard，2026年3月13日\href{https://www.rfberlin.com/wp-content/uploads/2026/03/26078.pdf}{固定工作稿}）

关系$H=Q/a$明确了需求的作用。统一竞争算例中，生产率与产量均提高50\%，总工时保持$80/3$，固定时薪下工资总额保持1600。销量增幅小于生产率增幅时总工时下降，反之增加；人数还取决于每人时长。劳动所得又受分配条件影响：在$H=20$、$b=20$、$\bar\pi=200$时，$V=2000$、$\eta=0.4$对应所得960；若$V$升至2400而$\eta$降至0.2，所得为760。这是保存情景中的组合比较，说明收入增长仍可能伴随劳动所得下降，不是经验预测。竞争与谈判模块使用不同给定条件，尚未联立求解。（\hyperref[ch:14]{第十四章}；谈判输出）

平台案例把分配关系落实到收费、派单和等待：匹配省时可以与收费增加并存，必要无偿劳动、退出和外包又会改变实际回报。职业收益还包括跨期能力，新人当前表现改善不能独立证明工具撤除后的水平。训练、反馈与完整责任的取得，决定部分当前收益能否延续为职业机会。

劳动者同时可能取得消费和服务收益。所列Waymo运行范围中的有利安全结果、低价产品和有效治疗均有实际价值，特定岗位损失和新增家庭承担也需保留。（Kusano等，2024；\hyperref[ch:10]{第十章}）评价应连接同一群体的收入、可支配时间、风险和服务，使技术改善最终体现为哪些人的收益，而非让单项平均数替代全部劳动后果。

\subsection{技术重组可变资本，部门结果不能归结为统一替代率}

跨行业分析得到的共同结论是：\textbf{技术重组劳动力购买，也改变对已购劳动力的使用；可变资本总额由实际产量、岗位报酬和购买范围决定，预付额另受周转约束。}单位用工减少、工资支出减少、相对比重下降与预付节约是四个不同结果。生活资料生产率对劳动力价值的影响，又不能由上述任何单项变化直接代替。

\begingroup\footnotesize\setlength{\tabcolsep}{4pt}\renewcommand{\arraystretch}{1.25}
\begin{longtable}{>{\raggedright\arraybackslash}p{\dimexpr 0.14\textwidth-8pt\relax}>{\raggedright\arraybackslash}p{\dimexpr 0.32\textwidth-8pt\relax}>{\raggedright\arraybackslash}p{\dimexpr 0.30\textwidth-8pt\relax}>{\raggedright\arraybackslash}p{\dimexpr 0.24\textwidth-8pt\relax}}
\toprule
\textbf{领域} & \textbf{劳动力购买与使用的变化} & \textbf{抵消或扩大支出的条件} & \textbf{现有证据的范围} \\
\midrule\endfirsthead
\toprule
\textbf{领域} & \textbf{劳动力购买与使用的变化} & \textbf{抵消或扩大支出的条件} & \textbf{现有证据的范围} \\
\midrule\endhead
\bottomrule\endfoot
机器人 & 部分操作劳动减少，招聘与岗位结构调整 & 扩产、新任务、技能工资及转岗安排 & 地区、行业与人员结果不同；非统一资本账户 \\
平台 & 等待、派单、报酬与付款来源重组 & 报酬标准、任务规模、经营成本和合同关系 & 纽约报送付款与时间可核；雇佣范围仍须辨明 \\
AI知识工作 & 起草与编码可能省时，核验和部署劳动变化 & 更多项目、质量责任、上游工资与专业技能 & 局部任务证据较多；完整工资效应有限 \\
增材与孪生 & 制模、返工与巡检可能减少，新技术劳动增加 & 后处理、认证、维护、订单及回款周期 & 工艺与采用资料不能代替工资节省 \\
光伏与储能 & 制造、建设和运维阶段劳动重组 & 新建规模、运行任务、供应链及交付电量 & 装机和成本可核；分阶段工资资料不足 \\
自动驾驶 & 驾驶劳动减少，支持与维护劳动保留 & 服务扩张、工资差异、轮班及现场任务 & 已确认支持任务；全年工资节省尚未测出 \\
基因编辑 & 劳动集中到制造、临床与长期随访 & 治疗扩容、专业技能、既有人员再利用 & 疗效和流程可核；完整工资账户缺失 \\
半导体 & 重复操作减少，工程和良率管理可能增加 & 工厂扩产、工艺复杂度、班次和技能报酬 & 设备与研发费用不能直接估计可变资本 \\
社会再生产 & 照护与学习改变劳动力供给和进入条件 & 岗位需求、服务质量、家庭及公共承担 & 参与和负担证据须同资本购买分开 \\
\end{longtable}\endgroup

\textbf{表16-V1　技术影响可变资本的部门路径与证据边界。} 本文机制综合，依据第五至十三章案例和第二章范畴辨析；各行不是可比的因果效应估计。表中“可能”表示需要相应条件和工资记录检验，未以设备或技术表现代替资本观测。

这些差异深化了基本矛盾的劳动表现。共同知识与设施扩大生产能力，资本却决定省下的劳动用于扩产、减员还是缩时，也决定购买什么技能、将哪些责任交给外部。工资改善和绝对可变资本增加可以同这种决定权延续并存。因而，本文既不能用减员证明一切技术强化资本控制，也不能用工资增长证明生产者已共同决定技术收益的用途。

参考视频所提示的机械化与劳动强化关系，跨行业分析还区分了三个结果：少购买劳动力，更多使用留任人员，降低生活资料生产所需劳动。它们分别涉及用工数量、劳动密度和必要劳动时间，不能互相替代。工资、任务产出和资本构成的变化也不足以独立测出剩余价值率。本文能够明确这些传导条件和部分组织环节，尚未建立覆盖各行业的价值量估计与动态危机模型。

\section{继续工作所需的时间和费用，须有人承担}

社会再生产是技术收益实现的参与条件。照护、学习、生活现金与维护，决定劳动者能否进入、持续工作和退出。第十三章算例表明，免学费并不解除照护和准备费限制；跨境任务的标价也须扣收费并计入无酬业务时间，才形成完整业务回报。

放松这些限制能够带来实际好处。布基纳法索的托育实验改善了若干就业、财务和儿童发展结果，家庭自主性等指标未检出显著变化；肯尼亚移动支付研究在其识别条件下支持贫困改善。所得和机会确有改善，就应承认改善；参与者能否改变生产和服务规则，还要另找权限的证据。（\hyperref[ch:13]{第十三章}）

生产需要照护、训练和公共服务，这些费用不会因为企业没有支付就消失。企业少付的钱，如果由家庭或公共部门新增而未获补偿的投入顶上，就包含费用转移；如果工具和服务在同等质量下减少了整个过程的劳动，才是相应范围的真实节约。判断须查同一项活动原来怎样做、现在怎样做、谁新增投入、谁取得补偿。费用发生在企业外面，本身还不能证明转嫁。

材料采购和废弃物处置，还把生产的影响带到别的地点和以后时期。一方取得采购收益，健康损害、维修费用和公共服务负担却可能由另一方承担。工资或出口增加了，不能因此假定受损者已获补偿。要把地点、时期、监测、付款和修复责任联系起来，也不能由一处暴露或一个全球总量推断所有技术的平均损失。

收入分配还可能改变需求实现。收入转向消费倾向较高的领取者，在其他支出未抵消时增加消费，实际交付仍受产能、价格、投资和现金限制。第十四章现金模块检验这些条件，没有建立全经济动态危机模型。再生产分析因而连接两件事：持续投入如何维持能力，所得和责任如何支持继续使用这些能力。

\section{改善所得和改变生产支配，各有实际条件}

第十五章讨论三种不同的制度改变。竞争、劳动保障和公共服务可以改善收益与负担；共同决定改变某些事项由谁作主；生产条件共同占有和共同管理，进一步改变生产者与生产资料的关系。设施共用可以少做重复准备，记录能够迁移可以约束收费，保障可以改善退出条件，公共服务可以减轻照护负担。补贴本来就会采用的企业，则可能只增加利润，没有改变当期产量和价格。各办法改变什么，须分别核算；谁替谁付款，也要同真正省下投入分清。

共同决定是不是有效，要看实际决定了哪些事，程序能不能执行。德国三分之一监事会代表要求的研究未检出所考察的工资与租金分享效应，另一研究比较不同代表比例下的股东估值；两者不能替所有参与制度下同一个结论。（Jäger等，2021；Gorton与Schmid，2004）成员能否取得信息、独立核验、要求纠正争议，决定纸面权利能否用于投资和劳动安排。工资涨了，还不能据此说这些权利已经落实。

劳动者能够共同修改安全、工时或培训安排，相应关系就已改变。要进一步说明占有关系整体转化，还须查生产条件归谁，投资服务于什么目标，经营剩余由谁决定用途。公共所有和合作所有提供了不同的权利基础；成员范围、问责、运行资金和外部依赖，又影响权利能够实行到什么程度。本文没有证明新技术已经促成这种总体转化，而是说明要提出这一判断，必须查清哪些权利已经改变。（\hyperref[ch:15]{第十五章}）

\subsection{制度反馈补充了技术收益与决定权之间的连接}

新增政治经济模型将论证从技术省工推进到主体选择：净收入、工资购买、补偿责任和政治成本各自进入目标，分别优化后才得到制度排序。协商分支选择$x_N=\mu/d_x$，镇压分支选择$x_R=(\mu+\lambda-\kappa)/d_x$；两者的互补性取决于$\lambda>\kappa$，制度是否转换则取决于完整收益差。资本增加只在式（14-P10）的联合条件成立时提高镇压偏好，不是由设备增加直接推出的政治规律。

本节算例使“技术扩张必然压缩参与”有了明确反例。$K=1$时，资本方原本在镇压分支得到56.240，高于协商的51.200；可选择新增劳动任务后，协商收益为61.200，高于镇压的58.615。资本方收益不等于社会福利，但反转说明，任务创造与可执行参与约束能够共同改变私人选择。两组转换阈值约0.2349与1.3931均来自未校准情景，不是实际社会的临界资本量。

公共品效率扩展又区分了另一条路径：资本积累伴随的保障效率提高，可以使再分配收益超过镇压收益，即使最优税率同时提高。风险扩展则表明，风险厌恶只有结合两种方案的失败概率、消费底线和收益分布才有确定含义。我们模型中的确定协商假设支持P10的单向边界，不能替代原文两种政策均有失败风险的比较。

该扩展借鉴Acemoglu等工作论文的论证组织，区分自有参与约束模型与原文政策收益敏感性。修改假设后的反例限定结果的推广范围，没有推翻原假设内的定理。现实制度判断仍须核查国家目标、劳动者协调、资产归属及实际决定程序，不能由模型分支名称代替。

\section{案例比较支持条件解释，尚未形成统一因果估计}

九个案例把前文的问题落到具体事实。任务实验、市场估计、运行记录和数学成果发布，各有自己的对象，不能加成一个统一的技术收益分数。表16.1列出原有口径。比较时要查的是：哪一步生产能力变了，好处怎样传给各方，参加者实际能决定什么。

\begingroup\footnotesize\setlength{\tabcolsep}{4pt}\renewcommand{\arraystretch}{1.25}
\begin{longtable}{>{\raggedright\arraybackslash}p{\dimexpr 0.23810\textwidth-8.0pt\relax}>{\raggedright\arraybackslash}p{\dimexpr 0.34524\textwidth-8.0pt\relax}>{\raggedright\arraybackslash}p{\dimexpr 0.21429\textwidth-8.0pt\relax}>{\raggedright\arraybackslash}p{\dimexpr 0.20238\textwidth-8.0pt\relax}}
\toprule
\textbf{案例与观察范围} & \textbf{主要结果及其口径} & \textbf{可支持的讨论} & \textbf{仍需补充的证据} \\
\midrule\endfirsthead
\toprule
\textbf{案例与观察范围} & \textbf{主要结果及其口径} & \textbf{可支持的讨论} & \textbf{仍需补充的证据} \\
\midrule\endhead
\bottomrule\endfoot
Copilot标准任务，2022年 & 成功完成任务者平均耗时减少55.8\% & 辅助工具能改善特定任务流程 & 未完成者、完整项目与长期收入 \\
METR真实代码库任务，2025年初至中期 & 允许AI条件下，经模型调整的完成时间增加19\% & 完整任务中的核验与整合投入需要计量 & 其他人员、任务和后续工具版本 \\
纽约送餐平台，2023—2024年第一季度 & 报送小时总收入11.72升至19.26美元，总周工时169.9万降至146.8万 & 报酬规则、等待与接单机会需要一起观察 & 去重人员的净收入、退出与政策反事实 \\
法国制造业，2010—2015年 & 采用企业自身工时与竞争者工时的估计方向不同 & 企业扩张可伴随同行业务转移 & 采用选择及跨行业一般均衡 \\
助听器贸易，1995—2016年，2007年扩散节点 & 固定2019年工作稿估计出口相对基准增加约58\% & 增材技术可以扩大既有跨境生产网络 & 后续版本、生产量与劳动收益，以及识别条件 \\
CAISO运行记录，2024年 & 弃风弃光342.34万MWh，约92\%属于本地拥塞经济调度类 & 空间与时序约束影响能力的实际使用 & 节点反事实、设施费用与福利变化 \\
肯尼亚接电实验，2014年干预 & 免费报价处理组363／380户接受；就业效应区间仍较宽 & 连接选择、生活改善与市场收入须分开 & 更长期随访及互补投入的独立效应 \\
肯尼亚移动支付，2008—2014年 & 女性户主家庭接入增幅第75相对第25分位对应极端贫困概率下降9.2个百分点 & 交易服务可扩大低收入家庭可用收益 & 非随机布局、失访及收益分配机制 \\
OpenAI数学资料发布，2026年10月 & 首发722篇，撤回3篇后现存719篇、372组相关成果 & 模型参与研究生成，成果公开与生产能力取得分开 & 逐项正确性、原创性、完整投入及实际模型使用条件 \\
\end{longtable}\endgroup


\textbf{表16.1　案例结果、证据性质与讨论边界。} Copilot的55.8\%来自成功完成者条件样本；METR的19\%来自任务时间模型，两者不能直接作均值差比较。纽约数据为平台报送汇总，CAISO为可修订的运行数据，移动支付和助听器贸易为识别假设下的估计；数学发布按2026年10月8日目录快照记录，为开发者材料及修订记录，稿件数不等于已独立确认的解决问题数。来源、样本、置信区间及计算分别见\hyperref[ch:5]{第五章}、\hyperref[ch:6]{第六章}、\hyperref[ch:7]{第七章}、\hyperref[ch:8]{第八章}、\hyperref[ch:9]{第九章}和\hyperref[ch:13]{第十三章}。原有数值输入见新增案例数据说明，数学事件的读取范围见发布核查。



助听器工作稿中的出口增长构成“工艺进步必然使生产本地化”的反例：新工艺也能扩大既有跨境网络。结合半导体案例，组织边界取决于材料、认证、客户和设施的完整联系，不能仅由某道工序更易复制推断。生产地点和贸易变化又不同于劳动者决定权限的变化。

新增产业序列把数量与份额分开：UNIDO五国制造业常价增加值在2010—2024年均增长，GDP份额却下降；IFR汽车和电子电气安装量也在份额下降时扩大。它们共同约束了以份额单独判断生产萎缩的做法，未估计自动化对增加值或就业的净贡献。（\hyperref[ch:8]{第八章8.8}）

能源和交易服务的案例，要从能够接入查到实际所得。CAISO的本地拥塞分类说明，地点和时段可以限制电力使用；少弃多少电值得投入多少，还须比较设施、损耗和可靠性费用。肯尼亚接电实验确认了连接选择和若干生活改善，就业估计仍不确定；移动支付研究在其识别条件下支持低收入家庭贫困改善。改善发生在哪里，没有实现的好处卡在哪里，都应查清。研究设计不同的两个案例，不能给各种技术排出普遍的高低次序。

数学发布要分清能看到成果和能自主使用研究工具。稿件、部分证明代码和推理摘要公开了，其他研究者可以学习和核查；产生这些结果的模型在该次发布时却尚未开放，研究者不能仅凭稿件调用同一系统研究自己选的问题。AGMAI要求公平取得强大工具和足够计算资源，也明确咨询不包含企业决策权。\textbf{分享研究成果扩大了可用知识，能够自主调用和运行工具，才进一步关系到生产条件的支配。} 一些研究任务可以交给模型计算了，谁能调用、谁能调配算力、谁定规则，仍是实际问题。（\hyperref[ch:7]{第七章7.1—7.2}）

上述案例用于比较条件，并未合并估计共同因果效果。它们支持的联系是：共同生产能力能够扩大，实际收益与决定权却经不同安排形成。技术效果由对应任务和交付结果判断，占有关系变化由生产条件归属及实际权限判断，收益改善不能单独代替后者。

\section{后续研究需要连接收益轨迹与实际决定过程}

本文的主要限制在证据连接：定向文献比较覆盖了任务、市场、家庭和制度的不同环节，尚缺同一批主体的完整长期记录；行业合同、原始数据和决定过程也未统一取得。条件模型未经验校准，来源台账所列的访问与读取范围限制了各项判断。文献数量不能替代直接支持，也不能据此估计跨技术的平均总体效应。

后续研究首先要把采用、竞争、分配和制度回应接到同一组事实。现有竞争模块给定工资和需求，谈判模块给定工时和收入，制度模型又给定补偿和政治成本，而现实中的分配会改变成本、需求、投资与参与条件。须取得同一产品市场的采用支出、销量、价格、完整工时、工资和合同，比较其中确有依据的变化。还要跟踪离开者、新进入者和工具撤除后的任务表现，才能从当前省时查到职业机会与长期收入。

政治反馈还需要独立资料：补偿标准何时提出、哪些主体能要求执行、部署方案如何因此修改、政治约束是否实际改变、新增任务是否进入持续需求。还须对应记录同等支出的服务质量与覆盖、工资收入的实际购买能力、两种政策的失败概率及失败后的消费来源。先检验模型目标、约束和风险结构是否适用，再考察收益差；不能用劳动份额下降推定镇压，也不能用商品降价推定公共品效率必然提高。加入劳动者协调、资本运动与融资约束后，才可能讨论静态阈值以外的制度路径。

另一项研究任务是连接实际迁移与决策过程。记录转换平台、模型和认证所花的时间、费用及失败原因，检验替代是否可行；记录投资提案、审批、规则修订和争议裁定，检验权限是否执行。能源、交通和医疗还需将容量、付款、质量与承担者连接，才能比较服务改善和治理变化。

全文归结为一个判断：\textbf{技术进步扩大生产的可能，共同生产者是否支配这些可能，则取决于生产条件的归属和实际决定权。} 竞争降价、报酬改善和公共服务可以带来独立的生活收益，这些改善仍可能发生在资本决定投资、用工和剩余用途的组织中。占有关系转化要求共同生产者对相应生产条件取得实际支配，其已经实现的范围须由所有权与决策过程查证。

\section{本章参考文献与复核入口}

本章综合前述章节，不增加独立田野或因果估计。以下列出正文直接涉及的原研究；完整版本、样本、定位及其他技术来源分别见相应章末书目和证据台账。

Acemoglu, D., \& Restrepo, P. (2020). Robots and jobs: Evidence from US labor markets. \emph{Journal of Political Economy, 128}(6), 2188–2244. \href{https://doi.org/10.1086/705716}{DOI}

Graetz, G., \& Michaels, G. (2018). Robots at work. \emph{Review of Economics and Statistics, 100}(5), 753–768. \href{https://doi.org/10.1162/rest_a_00754}{DOI}

Dauth, W., Findeisen, S., Suedekum, J., \& Woessner, N. (2021). The adjustment of labor markets to robots. \emph{Journal of the European Economic Association, 19}(6), 3104–3153. \href{https://doi.org/10.1093/jeea/jvab012}{DOI}

Brynjolfsson, E., Li, D., \& Raymond, L. (2025). Generative AI at work. \emph{Quarterly Journal of Economics, 140}(2), 889–942. \href{https://doi.org/10.1093/qje/qjae044}{DOI}

Dell’Acqua, F., McFowland, E., III, Mollick, E., Lifshitz, H., Kellogg, K. C., Rajendran, S., Krayer, L., Candelon, F., \& Lakhani, K. R. (2026). Navigating the jagged technological frontier: Field experimental evidence of the effects of artificial intelligence on knowledge worker productivity and quality. \emph{Organization Science}, advance online publication. \href{https://doi.org/10.1287/orsc.2025.21838}{DOI}

Kusano, K. D., Scanlon, J. M., Chen, Y.-H., et al. 2024. Comparison of Waymo rider-only crash data to human benchmarks at 7.1 million miles. \emph{Traffic Injury Prevention}, 25(sup1): S66–S77. \href{https://doi.org/10.1080/15389588.2024.2380786}{DOI:10.1080/15389588.2024.2380786}；第十章读取\href{https://arxiv.org/html/2312.12675v3}{接受稿v3}.

Frangoul, H., Locatelli, F., Sharma, A., et al. (2024). Exagamglogene autotemcel for severe sickle cell disease. \emph{New England Journal of Medicine, 390}(18), 1649–1662. \href{https://doi.org/10.1056/NEJMoa2309676}{DOI}；\href{https://pubmed.ncbi.nlm.nih.gov/38661449/}{摘要}.

Ajayi, K. F., A. Dao, and E. Koussoubé. 2022. \emph{The Effects of Childcare on Women and Children: Evidence from a Randomized Evaluation in Burkina Faso}. World Bank Policy Research Working Paper 10239. \href{https://documents1.worldbank.org/curated/en/099215111282254210/pdf/IDU09ee54b690865904f71089a60d76ad621601c.pdf}{原论文}。

Jäger, S., Schoefer, B., \& Heining, J. (2021). Labor in the boardroom. \emph{The Quarterly Journal of Economics, 136}(2), 669–725. \href{https://doi.org/10.1093/qje/qjaa038}{DOI}.

Gorton, G., \& Schmid, F. A. (2004). Capital, labor, and the firm: A study of German codetermination. \emph{Journal of the European Economic Association, 2}(5), 863–905. \href{https://doi.org/10.1162/1542476042782260}{DOI}.

International Renewable Energy Agency. (2026). \emph{Renewable power generation costs in 2025}. \href{https://www.irena.org/Publications/2026/Jul/Renewable-Power-Generation-Costs-in-2025}{原来源}

国家能源局能源节约和科技装备司、国家发展和改革委员会创新和高技术发展司、电力规划设计总院. (2026). \emph{中国新型储能发展报告2026}. \href{https://www.nea.gov.cn/20260730/0a1e9823b7a14e8ba16e82c34b3bea2b/c.html}{原来源}

能源观测与融资演算：\hyperref[ch:9]{第九章}、融资参数、融资计算、2021—2025装机、2023—2025利用。观测来自同版官方报告，派生比率与融资结果为本文计算，二者不作因果互换。

原创机制及情景：\hyperref[ch:2]{第二章概念}、\hyperref[ch:4]{第四章H1—H9}、\hyperref[ch:14]{第十四章模块与演算}、\hyperref[ch:15]{第十五章制度比较}。

Humlum, A., \& Vestergaard, E. (2026). Still waters, rapid currents: Early labor market transformation under generative AI. RFBerlin Discussion Paper 078/26, March 13 working-paper version. \url{https://www.rfberlin.com/wp-content/uploads/2026/03/26078.pdf}

学院派Academia.（2021）。\emph{【硬核社会学】资本主义的本质是什么？}。Bilibili视频，2021年2月2日，BV1Jt4y1B72f，P1“资本主义三幻神”。\href{https://www.bilibili.com/video/BV1Jt4y1B72f/}{视频入口}。作为问题讨论材料，机器文字稿的定位、核对范围与使用限制见第二章，不作为马克思原文或行业效应的证明。

Acemoglu, D., Gitmez, A. A., \& Shadmehr, M. (2026). \emph{Automation and Repression}. NBER Working Paper 35336, June. \href{https://doi.org/10.3386/w35336}{DOI:10.3386/w35336}；\href{https://economics.mit.edu/sites/default/files/2026-06/Automation\%20and\%20Repression.pdf}{2026年6月5日作者稿}。工作论文；核读模型、政策比较、动态条件及Lemma 4证明。本文参与约束模型另有设定，效率扩展只重算政策收益的部分敏感性，不声明复现全部原论文。

\backmatter\chapter*{统一参考文献}\addcontentsline{toc}{chapter}{统一参考文献}


按实际章末书目归并，保留章节与版本核查信息；入口重复与发布方字段差异的处置见修订记录。实际读取范围见逐章台账，不宣称每条文献全文已读。

1. Acemoglu, D., \& Restrepo, P. (2019). Automation and new tasks: How technology displaces and reinstates labor. \emph{Journal of Economic Perspectives, 33}(2), 3–30. \url{https://doi.org/10.1257/jep.33.2.3}
引用章节：01, 05。

2. Bresnahan, T. F., \& Trajtenberg, M. (1992). \emph{General Purpose Technologies “Engines of Growth?”} NBER Working Paper 4148. \url{https://doi.org/10.3386/w4148} 。本章核查原工作稿摘要与发表信息，引用此版本，未把正式1995稿逐页核查与工作稿混同。
引用章节：01。

3. Brynjolfsson, E., Li, D., \& Raymond, L. (2025). Generative AI at work. \emph{Quarterly Journal of Economics, 140}(2), 889–942. \url{https://doi.org/10.1093/qje/qjae044}
引用章节：01, 05, 16。

4. Grossman, S. J., \& Hart, O. D. (1986). The costs and benefits of ownership: A theory of vertical and lateral integration. \emph{Journal of Political Economy, 94}(4), 691–719. \url{https://doi.org/10.1086/261404}
引用章节：01。

5. Ostrom, E. (1990). \emph{Governing the Commons: The Evolution of Institutions for Collective Action}. Cambridge University Press. \url{https://doi.org/10.1017/CBO9780511807763} 。本章只核查出版社书目与研究说明。
引用章节：01。

6. 恩格斯（1878）。《反杜林论》，第三编第二章。访问版本：Emile Burns英译，Progress Publishers，1947；\href{https://www.marxists.org/archive/marx/works/1877/anti-duhring/index.htm}{MIA版本说明与目录}。
引用章节：01, 02。

7. 马克思（1867）。《资本论》第一卷。访问版本：Samuel Moore与Edward Aveling英译，恩格斯编辑，1887；\href{https://www.marxists.org/archive/marx/works/1867-c1/index.htm}{英文版本说明}，\href{https://www.marxists.org/chinese/marx-engels/23/index.htm}{《马克思恩格斯全集》第二十三卷中文转录目录}。章号对应见正文及理论核查台账。
引用章节：01, 02。

8. 马克思（1885）。《资本论》第二卷，恩格斯整理出版。访问版本：I. Lasker英译，Progress Publishers，1956，依据1893年修订第二版；\href{https://www.marxists.org/archive/marx/works/1885-c2/index.htm}{版本说明}。
引用章节：02。

9. 马克思（1894）。《资本论》第三卷，恩格斯整理出版。访问文本为MIA英文转录，所列来源包括Institute of Marxism-Leninism，1959及International Publishers；\href{https://www.marxists.org/archive/marx/works/1894-c3/index.htm}{版本说明}。
引用章节：02。

10. 马克思（1857—1858年草稿；德文首次出版1939—1941）。《政治经济学批判大纲》。访问版本：Martin Nicolaus英译，1973；\href{https://www.marxists.org/archive/marx/works/1857/grundrisse/index.htm}{版本说明与草稿目录}。一般智力段落按网站分段定位，不将网页编号冒充作者正式章号。
引用章节：02。

11. 马克思（1875年写作；1891年刊布）。《哥达纲领批判》，第一部分。访问版本来自Progress Publishers，1970，《马克思恩格斯选集》第三卷英文转录；\href{https://www.marxists.org/archive/marx/works/1875/gotha/index.htm}{版本说明}。
引用章节：02。

12. Callaway, B., \& Sant’Anna, P. H. C. (2021). Difference-in-Differences with multiple time periods. \emph{Journal of Econometrics, 225}(2), 200–230. \url{https://doi.org/10.1016/j.jeconom.2020.12.001} 。核查作者2020-12-01公开稿第1—7页及定义；未在本报告实施估计。
引用章节：03。

13. OECD. (2026). \emph{Competition in the age of AI: Initial evidence from microdata}. OECD Artificial Intelligence Papers, No. 64. \url{https://doi.org/10.1787/6f88a1ea-en} 。核查公开PDF第9、52页的范围与讨论。
引用章节：03, 04, 07。

14. Textor, J. (2023). \emph{Drawing and Analyzing Causal DAGs with DAGitty}, version 3.1 manual, July 18. \url{https://www.dagitty.net/manual-3.x.pdf} 。核查第8—9页；本报告未运行DAGitty效应识别算法。
引用章节：03。

15. 国家能源局能源节约和科技装备司、国家发展和改革委员会创新和高技术发展司、电力规划设计总院（2026）。《中国新型储能发展报告2026》。\url{https://www.nea.gov.cn/20260730/0a1e9823b7a14e8ba16e82c34b3bea2b/202607300a1e9823b7a14e8ba16e82c34b3bea2b_270e42e93bb1a647609c851527f42364ee.pdf} 。图19、图24；本章计算为描述统计。
引用章节：03, 09, 16。

16. Nash, J. F., Jr. (1950). The bargaining problem. \emph{Econometrica}, 18(2), 155–162. \href{https://doi.org/10.2307/1907266}{DOI:10.2307/1907266}。\href{https://www.haverford.edu/sites/default/files/Nash1950.pdf}{Haverford公开原论文}，原刊pp.157—159核查破裂基点、公理与增益乘积解；权重、货币简式和劳动分配算例为本报告设定。访问日期：2026年10月8日。
引用章节：04, 14。

17. Acemoglu, D., \& Restrepo, P. (2020). Robots and jobs: Evidence from US labor markets. \emph{Journal of Political Economy, 128}(6), 2188–2244. \url{https://doi.org/10.1086/705716}
引用章节：05, 16。

18. Graetz, G., \& Michaels, G. (2018). Robots at work. \emph{Review of Economics and Statistics, 100}(5), 753–768. \url{https://doi.org/10.1162/rest_a_00754}
引用章节：05, 16。

19. Dauth, W., Findeisen, S., Suedekum, J., \& Woessner, N. (2021). The adjustment of labor markets to robots. \emph{Journal of the European Economic Association, 19}(6), 3104–3153. \url{https://doi.org/10.1093/jeea/jvab012}
引用章节：05, 16。

20. Noy, S., \& Zhang, W. (2023). Experimental evidence on the productivity effects of generative artificial intelligence. \emph{Science, 381}(6654), 187–192. \url{https://doi.org/10.1126/science.adh2586}
引用章节：05。

21. Dell’Acqua, F., McFowland, E., III, Mollick, E., Lifshitz, H., Kellogg, K. C., Rajendran, S., Krayer, L., Candelon, F., \& Lakhani, K. R. (2026). Navigating the jagged technological frontier: Field experimental evidence of the effects of artificial intelligence on knowledge worker productivity and quality. \emph{Organization Science}, advance online publication. \url{https://doi.org/10.1287/orsc.2025.21838}
引用章节：05, 16。

22. Autor, D., Chin, C., Salomons, A., \& Seegmiller, B. (2024). New frontiers: The origins and content of new work, 1940–2018. \emph{Quarterly Journal of Economics, 139}(3), 1399–1465. \url{https://doi.org/10.1093/qje/qjae008}
引用章节：05。

23. Acemoglu, D. (2025). The simple macroeconomics of AI. \emph{Economic Policy, 40}(121), 13–58. \url{https://doi.org/10.1093/epolic/eiae042}
引用章节：05。

24. Humlum, A., \& Vestergaard, E. (2026). Still waters, rapid currents: Early labor market transformation under generative AI. \emph{RFBerlin Discussion Paper 078/26}, March 13 version, working paper. \url{https://www.rfberlin.com/wp-content/uploads/2026/03/26078.pdf}
引用章节：05, 16。

25. Peng, S., Kalliamvakou, E., Cihon, P., \& Demirer, M. (2023). \emph{The impact of AI on developer productivity: Evidence from GitHub Copilot}. arXiv:2302.06590v1, February 13, preprint. \href{https://doi.org/10.48550/arXiv.2302.06590}{DOI：10.48550/arXiv.2302.06590}; \href{https://arxiv.org/pdf/2302.06590v1}{本章所读固定版本}。访问日期：2026-10-08。
引用章节：05。

26. Becker, J., Rush, N., Barnes, E., \& Rein, D. (2025). \emph{Measuring the impact of early-2025 AI on experienced open-source developer productivity}. arXiv:2507.09089v2, July 25, preprint. \href{https://doi.org/10.48550/arXiv.2507.09089}{DOI：10.48550/arXiv.2507.09089}; \href{https://arxiv.org/pdf/2507.09089v2}{本章所读固定版本}。论文题页将Barnes署名列为Beth，arXiv书目信息列为Elizabeth。访问日期：2026-10-08。
引用章节：05。

27. Becker, J., Rush, N., Cunningham, T., Rein, D., \& Mahamud, K. (2026). \emph{We are changing our developer productivity experiment design}. METR, February 24, research update. \href{https://metr.org/blog/2026-02-24-uplift-update/}{作者原始说明}。仅用于补充原研究区间与时效、选择限制，不替代2025年原论文结果。访问日期：2026-10-08。
引用章节：05。

28. International Labour Office. (2021). \emph{World Employment and Social Outlook 2021: The role of digital labour platforms in transforming the world of work}. Geneva: ILO. PDF ISBN 9789220319413. \href{https://www.ilo.org/publications/flagship-reports/role-digital-labour-platforms-transforming-world-work}{官方出版入口}。本章实际核查执行摘要与方法附录，总报告全文访问限制见台账。
引用章节：06, 13。

29. Milanez, A., Lemmens, A., \& Ruggiu, C. (2025). \emph{Algorithmic management in the workplace: New evidence from an OECD employer survey}. OECD Artificial Intelligence Papers, No. 31. \href{https://doi.org/10.1787/287c13c4-en}{DOI：10.1787/287c13c4-en}。
引用章节：06。

30. Kellogg, K. C., Valentine, M. A., \& Christin, A. (2020). Algorithms at work: The new contested terrain of control. \emph{Academy of Management Annals, 14}(1), 366—410. \href{https://doi.org/10.5465/annals.2018.0174}{DOI：10.5465/annals.2018.0174}。
引用章节：06。

31. Rosenblat, A., \& Stark, L. (2016). Algorithmic labor and information asymmetries: A case study of Uber’s drivers. \emph{International Journal of Communication, 10}, 3758—3784. \href{https://ijoc.org/index.php/ijoc/article/download/4892/1739}{期刊全文}。\href{https://doi.org/10.2139/ssrn.2686227}{DOI：10.2139/ssrn.2686227}对应SSRN版本，非期刊DOI。
引用章节：06。

32. Wood, A. J., Graham, M., Lehdonvirta, V., \& Hjorth, I. (2019). Good gig, bad gig: Autonomy and algorithmic control in the global gig economy. \emph{Work, Employment and Society, 33}(1), 56—75. \href{https://doi.org/10.1177/0950017018785616}{DOI：10.1177/0950017018785616}。
引用章节：06。

33. Baiocco, S., Fernández-Macías, E., Rani, U., \& Pesole, A. (2022). \emph{The algorithmic management of work and its implications in different contexts}. Background Paper No. 9, Joint EU–ILO Project “Building Partnerships on the Future of Work”. \href{https://www.ilo.org/publications/algorithmic-management-work-and-its-implications-different-contexts}{官方出版页}。
引用章节：06。

34. Rani, U., Pesole, A., \& González Vázquez, I. (2024). \emph{Algorithmic Management practices in regular workplaces: Case studies in logistics and healthcare}. Luxembourg: Publications Office of the European Union. EUR 31846 EN; JRC136063. \href{https://doi.org/10.2760/712475}{DOI：10.2760/712475}。
引用章节：06。

35. Hadwiger, F. (2022). \emph{Realizing the opportunities of the platform economy through freedom of association and collective bargaining}. ILO Working Paper 80. \href{https://doi.org/10.54394/VARD7939}{DOI：10.54394/VARD7939}。
引用章节：06。

36. New York City Department of Consumer and Worker Protection. (2024). \emph{Restaurant delivery app data: January–March 2024}. Q1 2024 report, July 2024 corrected version; Tables 1–3 and Technical Note. \href{https://www.nyc.gov/assets/dca/downloads/pdf/workers/Restaurant-Delivery-App-Data-Q1-2024.pdf}{官方报告}。观察期与出版、纠错时间分别保留。访问日期：2026-10-08。
引用章节：06。

37. OECD. (2026a). \emph{Artificial Intelligence markets: Recent developments and competition issues}. Published 10 July 2026. \href{https://doi.org/10.1787/d531d73f-en}{DOI:10.1787/d531d73f-en}.
引用章节：07。

38. Stanford HAI. (2026). \emph{The 2026 AI Index Report}, Chapter 4: Economy. \href{https://hai.stanford.edu/assets/files/ai_index_report_2026_chapter_4_economy.pdf}{官方章节PDF}.
引用章节：07。

39. International Energy Agency. (2026). \emph{Key Questions on Energy and AI}. Paris: IEA. Published 16 April 2026. \href{https://www.iea.org/reports/key-questions-on-energy-and-ai}{报告主页}.
引用章节：07。

40. Competition and Markets Authority. (2024). \emph{AI Foundation Models: Technical Update Report}. 16 April 2024. \href{https://assets.publishing.service.gov.uk/media/661e5a4c7469198185bd3d62/AI_Foundation_Models_technical_update_report.pdf}{官方PDF}.
引用章节：07。

41. Kapoor, S., Bommasani, R., Klyman, K., et al. (2024). \emph{On the Societal Impact of Open Foundation Models}. 本章使用arXiv固定v1作者稿。 \href{https://doi.org/10.48550/arXiv.2403.07918}{DOI:10.48550/arXiv.2403.07918}；\href{https://arxiv.org/html/2403.07918v1}{固定版本}.
引用章节：07。

42. Groeneveld, D., Beltagy, I., Walsh, P., et al. (2024). \emph{OLMo: Accelerating the Science of Language Models}. 本章使用arXiv固定v1作者稿，提交于1 February 2024。 \href{https://doi.org/10.48550/arXiv.2402.00838}{DOI:10.48550/arXiv.2402.00838}；\href{https://arxiv.org/html/2402.00838v1}{固定版本}.
引用章节：07。

43. DeepSeek-AI. (2024). \emph{DeepSeek-V3 Technical Report}. 本章使用arXiv固定v1技术报告。 \href{https://doi.org/10.48550/arXiv.2412.19437}{DOI:10.48550/arXiv.2412.19437}；\href{https://arxiv.org/pdf/2412.19437v1}{固定版本PDF}.
引用章节：07。

44. Open Source Initiative. (2024). \emph{The Open Source AI Definition—1.0}. \href{https://opensource.org/ai/open-source-ai-definition}{定义原文}.
引用章节：07。

45. Meta. (2024). \emph{Llama 3.1 Community License Agreement}. Version release date: 23 July 2024. \href{https://raw.githubusercontent.com/meta-llama/llama-models/main/models/llama3_1/LICENSE}{发布方许可文本}.
引用章节：07。

46. 国家发展改革委、国家数据局、中央网信办、工业和信息化部、国家能源局. (2023). 关于深入实施“东数西算”工程加快构建全国一体化算力网的实施意见. 发改数据〔2023〕1779号；签署2023年12月25日，公布2023年12月29日。\href{https://www.ndrc.gov.cn/xxgk/zcfb/tz/202312/t20231229_1363000.html}{官方文件}.
引用章节：07。

47. Guo, D., Yang, D., Zhang, H., et al. (2025). DeepSeek-R1 incentivizes reasoning in LLMs through reinforcement learning. \emph{Nature, 645}, 633–638. \href{https://doi.org/10.1038/s41586-025-09422-z}{DOI:10.1038/s41586-025-09422-z}；\href{https://www.nature.com/articles/s41586-025-09422-z}{正式论文}。本章引用正式稿中R1-Zero结果，不与2025年1月预印本的不同数字混用。
引用章节：07。

48. Snell, C., Lee, J., Xu, K., \& Kumar, A. (2024). \emph{Scaling LLM Test-Time Compute Optimally can be More Effective than Scaling Model Parameters}. arXiv:2408.03314v1，6 August 2024. \href{https://doi.org/10.48550/arXiv.2408.03314}{DOI}；\href{https://arxiv.org/html/2408.03314v1}{本章使用的固定v1}。未将其与改题后的ICLR 2025版本混作逐页核查。
引用章节：07。

49. Novikov, A., Vũ, N., Eisenberger, M., et al. (2025). \emph{AlphaEvolve: A coding agent for scientific and algorithmic discovery}. arXiv:2506.13131v1，16 June 2025. \href{https://doi.org/10.48550/arXiv.2506.13131}{DOI}；\href{https://arxiv.org/html/2506.13131v1}{固定v1白皮书}。系统开发团队的结果报告，不视为独立随机实验。
引用章节：07。

50. OpenAI. (2026a). \emph{Ten advances in mathematics and theoretical computer science}. 1 August 2026. \href{https://openai.com/index/ten-advances-in-mathematics/}{官方发布说明}。本章核对公开的过程披露，未独立验证十项数学证明。
引用章节：07。

51. OpenAI. (2026b). \emph{math: mathematical manuscripts and supporting proof artifacts}. \href{https://github.com/openai/math}{GitHub仓库及README}。访问日期2026年10月8日；目录为719篇稿件、372组相关成果。
引用章节：07。

52. OpenAI. (2026c). \emph{math/history.md}. 7 October 2026更新，\href{https://github.com/openai/math/blob/main/history.md}{撤回、修订与形式化记录}。本章首发722篇由当次目录719篇加已撤回3篇核算，保留版本差异。
引用章节：07。

53. OpenAI. (2026d). \emph{Sharing AI progress in mathematics}. 6 October 2026. \href{https://openai.com/index/sharing-ai-progress-in-mathematics/}{官方发布说明}。计算口径为ChatGPT Pro等效思考量，不是完整研发成本。
引用章节：07。

54. OpenAI. (2026e). \emph{Advisory Group on Mathematics and Artificial Intelligence}. 21 September 2026. \href{https://openai.com/index/advisory-group-on-mathematics-and-ai/}{成立说明}。
引用章节：07。

55. Advisory Group on Mathematics and Artificial Intelligence (AGMAI). (2026). \emph{Independence, Transparency, and Accountability; On OpenAI's Release of Mathematical Results}. 6 October 2026发布声明，\href{https://agmai.org/}{咨询组官网}。访问日期2026年10月8日。
引用章节：07。

56. Thomas, D. S., \& Gilbert, S. W. (2014). \emph{Costs and Cost Effectiveness of Additive Manufacturing: A Literature Review and Discussion}. NIST SP 1176. \href{https://doi.org/10.6028/NIST.SP.1176}{DOI：10.6028/NIST.SP.1176}。
引用章节：08。

57. Hernandez Korner, M. E., Lamban, M. P., Albajez, J. A., Santolaria, J., Ng Corrales, L. C., \& Royo, J. (2024). Cost Model Framework for Pieces Additively Manufactured in Fused Deposition Modeling for Low to Medium Batches. \emph{3D Printing and Additive Manufacturing, 11}(1), 287—298. \href{https://doi.org/10.1089/3dp.2022.0044}{DOI：10.1089/3dp.2022.0044}。
引用章节：08。

58. NIST. (2018，2025年更新). \emph{Additive Manufacturing Part Qualification}. \href{https://www.nist.gov/programs-projects/additive-manufacturing-part-qualification}{官方项目说明}。
引用章节：08。

59. ISO/ASTM. (2023). \emph{ISO/ASTM 52920:2023: Additive manufacturing—Qualification principles—Requirements for industrial additive manufacturing processes and production sites}. \href{https://www.iso.org/standard/76911.html?browse=tc}{官方摘要}。
引用章节：08。

60. Anderson, G. (2016). \emph{The Economic Impact of Technology Infrastructure for Advanced Manufacturing: An Overview}. NIST Economic Analysis Briefs 1. \href{https://www.nist.gov/system/files/documents/2024/12/05/NIST.EAB_.1_R_1.pdf}{全文}。
引用章节：08。

61. NIST. (2024). \emph{Additive Manufacturing Powder Metrology Laboratory}. \href{https://www.nist.gov/laboratories/tools-instruments/additive-manufacturing-powder-metrology-laboratory}{官方设施说明}。
引用章节：08。

62. Shao, G. (2021). \emph{Use Case Scenarios for Digital Twin Implementation Based on ISO 23247}. NIST AMS 400-2. \href{https://doi.org/10.6028/NIST.AMS.400-2}{DOI：10.6028/NIST.AMS.400-2}。
引用章节：08。

63. Kim, D. B., Shao, G., \& Jo, G. (2022). A Digital Twin Implementation Architecture for Wire + Arc Additive Manufacturing Based on ISO 23247. \emph{Manufacturing Letters, 34}, 1—5. \href{https://doi.org/10.1016/j.mfglet.2022.08.008}{DOI：10.1016/j.mfglet.2022.08.008}。
引用章节：08。

64. Thomas, D. (2024b). \emph{Economics of Digital Twins: Costs, Benefits, and Economic Decision Making}. NIST AMS 100-61. \href{https://doi.org/10.6028/NIST.AMS.100-61}{DOI：10.6028/NIST.AMS.100-61}。
引用章节：08。

65. World Economic Forum. (2025). \emph{Global Lighthouse Network: The Mindset Shifts Driving Impact and Scale in Digital Transformation}. \href{https://www.weforum.org/publications/global-lighthouse-network-the-mindset-shifts-driving-impact-and-scale-in-digital-transformation/}{出版入口}。同日\href{https://www.weforum.org/press/2025/01/global-lighthouse-network-2025-world-economic-forum-recognizes-companies-transforming-manufacturing-through-innovation/}{新批次发布说明}为同一项目资料，不计独立重复验证。
引用章节：08。

66. Unilever. (2025). \emph{Five ways Unilever’s new Lighthouse site applies AI for impact}. \href{https://www.unilever.com/news/news-search/2025/five-ways-unilevers-new-lighthouse-site-applies-ai-for-impact/}{公司案例}。
引用章节：08。

67. 工业和信息化部办公厅（2026）。《关于推广5G工厂典型应用实践的通知》，工信厅信管函〔2026〕58号。\href{https://www.miit.gov.cn/ztzl/rdzt/gyhlw/wjfb/art/2026/art_6b35987d001445d49487650795eac294.html}{官方原文}。
引用章节：08。

68. Thomas, D. (2024a). \emph{Annual Report on the U.S. Manufacturing Economy: 2024}. NIST AMS 600-16. \href{https://doi.org/10.6028/NIST.AMS.600-16}{DOI：10.6028/NIST.AMS.600-16}。
引用章节：08。

69. Acemoglu, D., Lelarge, C., \& Restrepo, P. (2020). Competing with Robots: Firm-Level Evidence from France. \emph{AEA Papers and Proceedings, 110}, 383—388. \href{https://doi.org/10.1257/pandp.20201003}{DOI：10.1257/pandp.20201003}，\href{https://economics.mit.edu/sites/default/files/publications/Competing\%20with\%20Robots\%20-\%20Firm-Level\%20Evidence\%20from\%20F.pdf}{作者公开正式稿}。
引用章节：08。

70. Freund, C., Mulabdic, A., \& Ruta, M. (2019). \emph{Is 3D Printing a Threat to Global Trade? The Trade Effects You Didn’t Hear About}. World Bank Policy Research Working Paper 9024, September. \href{https://documents1.worldbank.org/curated/en/152701569432061451/pdf/Is-3D-Printing-a-Threat-to-Global-Trade-The-Trade-Effects-You-Didnt-Hear-About.pdf}{本文采用的固定工作稿}。本轮核查表1及方法、结论相关段落，访问日期2026年10月8日。
引用章节：08。

71. Egli, F., Steffen, B., \& Schmidt, T. S. (2018). A dynamic analysis of financing conditions for renewable energy technologies. \emph{Nature Energy, 3}, 1084–1092. \url{https://doi.org/10.1038/s41560-018-0277-y}
引用章节：09。

72. Hittinger, E. S., \& Azevedo, I. M. L. (2015). Bulk energy storage increases United States electricity system emissions. \emph{Environmental Science \& Technology, 49}(5), 3203–3210. \url{https://doi.org/10.1021/es505027p}
引用章节：09。

73. International Energy Agency. (2024). \emph{Batteries and secure energy transitions}. \url{https://www.iea.org/reports/batteries-and-secure-energy-transitions}
引用章节：09。

74. International Energy Agency. (2026a). \emph{Global EV outlook 2026: Electric vehicle batteries}. \url{https://www.iea.org/reports/global-ev-outlook-2026/electric-vehicle-batteries}
引用章节：09。

75. International Energy Agency. (2026b). \emph{Electricity 2026}. \url{https://www.iea.org/reports/electricity-2026}
引用章节：09。

76. International Energy Agency. (2026c). \emph{Global critical minerals outlook 2026: Market overview}. \url{https://www.iea.org/reports/global-critical-minerals-outlook-2026/market-overview}
引用章节：09。

77. International Renewable Energy Agency. (2023). \emph{The cost of financing for renewable power}. \url{https://www.irena.org/Publications/2023/May/The-cost-of-financing-for-renewable-power}
引用章节：09。

78. International Renewable Energy Agency. (2025). \emph{Renewable power generation costs in 2024}. \url{https://www.irena.org/Digital-Report/Renewable-Power-Generation-Costs-in-2024}
引用章节：09。

79. International Renewable Energy Agency. (2026). \emph{Renewable power generation costs in 2025}. \url{https://www.irena.org/Publications/2026/Jul/Renewable-Power-Generation-Costs-in-2025}
引用章节：09, 16。

80. Jenkins, K., McCauley, D., Heffron, R., Stephan, H., \& Rehner, R. (2016). Energy justice: A conceptual review. \emph{Energy Research \& Social Science, 11}, 174–182. \url{https://doi.org/10.1016/j.erss.2015.10.004}
引用章节：09。

81. Walker, G., \& Devine-Wright, P. (2008). Community renewable energy: What should it mean? \emph{Energy Policy, 36}(2), 497–500. \url{https://doi.org/10.1016/j.enpol.2007.10.019}
引用章节：09。

82. 国家能源局. (2026a). \emph{2025年可再生能源并网运行情况}. \url{https://www.nea.gov.cn/20260212/742b8c6a078347b0b39de676c05c5d58/c.html}
引用章节：09。

83. California Independent System Operator. (2024). \emph{Wind and Solar Curtailment: December 31, 2024}. Posted January 1, 2025. \url{https://www.caiso.com/documents/wind-solar-real-time-dispatch-curtailment-report-dec-31-2024.pdf}
引用章节：09。

84. U.S. Energy Information Administration. (2025, May 28). \emph{Solar and wind power curtailments are increasing in California}. \url{https://www.eia.gov/todayinenergy/detail.php?id=65364}
引用章节：09。

85. Lee, K., Miguel, E., \& Wolfram, C. (2020). Experimental evidence on the economics of rural electrification. \emph{Journal of Political Economy, 128}(4), 1523–1565. \url{https://doi.org/10.1086/705417} \href{https://emiguel.econ.berkeley.edu/assets/miguel_research/74/REPP_2019-04-25-FINAL-All.pdf}{作者2019年4月最终稿及附录}。
引用章节：09。

86. Kusano, K. D., Scanlon, J. M., Chen, Y.-H., et al. (2024). Comparison of Waymo rider-only crash data to human benchmarks at 7.1 million miles. \emph{Traffic Injury Prevention}, 25(sup1): S66–S77. \href{https://doi.org/10.1080/15389588.2024.2380786}{DOI:10.1080/15389588.2024.2380786}；本章读取\href{https://arxiv.org/html/2312.12675v3}{接受稿v3}.
引用章节：10, 16。

87. Scanlon, J. M., Kusano, K. D., Fraade-Blanar, L. A., et al. (2024). Benchmarks for retrospective automated driving system crash rate analysis using police-reported crash data. \emph{Traffic Injury Prevention}, 25(sup1): S51–S65. \href{https://doi.org/10.1080/15389588.2024.2380522}{DOI:10.1080/15389588.2024.2380522}；本章读取\href{https://arxiv.org/html/2312.13228v2}{接受稿v2}.
引用章节：10。

88. Di Lillo, L., Gode, T., Zhou, X., et al. (2024). \emph{Do Autonomous Vehicles Outperform Latest-Generation Human-Driven Vehicles? A Comparison to Waymo's Auto Liability Insurance Claims at 25 Million Miles}. 本章使用2024年12月发布的作者公开稿，不推定其为已核对正式期刊版本。\href{https://waymo.com/research/do-autonomous-vehicles-outperform-latest-generation-human-driven-vehicles-25-million-miles/}{发布方研究页面及PDF}.
引用章节：10。

89. Kalra, N., \& Paddock, S. M. (2016). \emph{Driving to Safety: How Many Miles of Driving Would It Take to Demonstrate Autonomous Vehicle Reliability?} RAND RR-1478-RC. \href{https://doi.org/10.7249/RR1478}{DOI:10.7249/RR1478}.
引用章节：10。

90. International Transport Forum. (2015). \emph{Urban Mobility System Upgrade: How Shared Self-Driving Cars Could Change City Traffic}. OECD Publishing. \href{https://doi.org/10.1787/5jlwvzdk29g5-en}{DOI:10.1787/5jlwvzdk29g5-en}.
引用章节：10。

91. NHTSA. (2026). \emph{Standing General Order on Crash Reporting}. 官方动态页面，首次核查2026-10-07、复查2026-10-08；页面仍显示数据截至2026-08-17。\href{https://www.nhtsa.gov/laws-regulations/standing-general-order-crash-reporting}{原文及限制说明}.
引用章节：10。

92. NHTSA. (2024). \emph{Consent Order: Cruise; Standing General Order 2021-01 Reporting}. Announced 30 September 2024. \href{https://www.nhtsa.gov/sites/nhtsa.gov/files/2024-09/cruise-consent-order-2024-web.pdf}{官方PDF}.
引用章节：10。

93. California Department of Motor Vehicles. (2026). \emph{Autonomous Vehicle Permit Holders in California Logged More Than 9 Million Test Miles Between December 1, 2024 and November 30, 2025}. 20 February 2026. \href{https://www.dmv.ca.gov/portal/news-and-media/autonomous-vehicle-permit-holders-in-california-logged-more-than-9-million-test-miles-between-december-1-2024-and-november-30-2025/}{官方说明}.
引用章节：10。

94. Waymo. (2026). \emph{Advice, Not Control: The Role of Remote Assistance in Waymo’s Operations}. From the Road, 17 February 2026. \href{https://waymo.com/blog/shorts/advice-not-control-the-role-of-remote-assistance/}{公司原文}.
引用章节：10。

95. Office of Senator Edward J. Markey. (2026). \emph{Remote Back Seat Operators: Revealing the Autonomous Vehicle Industry’s Reliance on Human Remote Assistance Operators}. \href{https://www.markey.senate.gov/imo/media/doc/remote_assistance_investigation_report.pdf}{调查报告}.
引用章节：10。

96. U.S. Department of Transportation. (2021). \emph{Driving Automation Systems in Long-Haul Trucking and Bus Transit: Preliminary Analysis of Potential Workforce Impacts}. \href{https://www.transportation.gov/av/workforce/report}{官方发布页及报告}.
引用章节：10。

97. 工业和信息化部、公安部、住房和城乡建设部、交通运输部. (2023). 关于开展智能网联汽车准入和上路通行试点工作的通知. 工信部联通装〔2023〕217号，2023-11-17。\href{https://ythxxfb.miit.gov.cn/ythzxfwpt/hlwmh/tzgg/xzxk/clsczr/art/2023/art_ae314a1df30e4ba1a49bd4ffa10186f6.html}{官方全文}.
引用章节：10。

98. Waymo. (2024). \emph{Fleet Response: Lending a Helpful Hand to Waymo’s Autonomously Driven Vehicles}. 21 May 2024. \href{https://waymo.com/blog/2024/05/fleet-response/}{公司原文}.
引用章节：10。

99. Jinek, M., Chylinski, K., Fonfara, I., Hauer, M., Doudna, J. A., \& Charpentier, E. (2012). A programmable dual-RNA-guided DNA endonuclease in adaptive bacterial immunity. \emph{Science, 337}(6096), 816–821. \href{https://doi.org/10.1126/science.1225829}{DOI}.
引用章节：11。

100. U.S. Food and Drug Administration. (2023, December 8). FDA approves first gene therapies to treat patients with sickle cell disease. \href{https://www.fda.gov/news-events/press-announcements/fda-approves-first-gene-therapies-treat-patients-sickle-cell-disease}{公告}.
引用章节：11。

101. Frangoul, H., Locatelli, F., Sharma, A., et al. (2024). Exagamglogene autotemcel for severe sickle cell disease. \emph{New England Journal of Medicine, 390}(18), 1649–1662. \href{https://doi.org/10.1056/NEJMoa2309676}{DOI}；\href{https://pubmed.ncbi.nlm.nih.gov/38661449/}{摘要}.
引用章节：11, 16。

102. European Medicines Agency. (2024; webpage accessed 2026-10-07). Casgevy: EPAR medicine overview. \href{https://www.ema.europa.eu/en/medicines/human/EPAR/casgevy}{官方概览}.
引用章节：11。

103. GBD 2021 Sickle Cell Disease Collaborators. (2023). Global, regional, and national prevalence and mortality burden of sickle cell disease, 2000–2021: A systematic analysis from the Global Burden of Disease Study 2021. \emph{The Lancet Haematology, 10}(8), e585–e599. \href{https://doi.org/10.1016/S2352-3026(23)00118-7}{DOI}；\href{https://repository.up.ac.za/server/api/core/bitstreams/fd57c857-9ec2-4897-b38c-ff955abb4bb0/content}{作者机构所存正式论文}；\href{https://pubmed.ncbi.nlm.nih.gov/37331373/}{摘要}.
引用章节：11。

104. Correction to Lancet Haematol 2023; 10: e585–99. (2023). \emph{The Lancet Haematology, 10}(8), e574. \href{https://doi.org/10.1016/S2352-3026(23)00215-6}{DOI}.
引用章节：11。

105. National Institute for Health and Care Excellence. (2025, February 26). Exagamglogene autotemcel for treating severe sickle cell disease in people 12 years and over. Technology appraisal guidance TA1044. \href{https://www.nice.org.uk/guidance/ta1044/chapter/1-recommendations}{推荐}.
引用章节：11。

106. National Institute for Health and Care Excellence. (2025). TA1044: Managed access agreement. \href{https://www.nice.org.uk/guidance/ta1044/resources/managed-access-agreement-pdf-15250968973}{协议}.
引用章节：11。

107. Centers for Medicare \& Medicaid Services. (2026, September 30 update). Cell and Gene Therapy Access Model. \href{https://www.cms.gov/priorities/innovation/innovation-models/cgt}{项目}.
引用章节：11。

108. World Health Organization. (2021a). \emph{Human genome editing: Recommendations}. ISBN 978-92-4-003038-1. \href{https://www.who.int/publications/i/item/9789240030381}{出版概览}.
引用章节：11。

109. World Health Organization. (2021b). \emph{Human genome editing: A framework for governance}. ISBN 978-92-4-003006-0. \href{https://www.who.int/publications/i/item/9789240030060}{出版概览}.
引用章节：11。

110. National Institute for Health and Care Excellence. (2022, dynamic manual accessed 2026-10-08). NICE technology appraisal and highly specialised technologies guidance: the manual (PMG36), Chapter 4 Economic evaluation. \url{https://www.nice.org.uk/process/pmg36/chapter/economic-evaluation-2} 。本章引用增量经济评价方法，未引用具体当期支付阈值。
引用章节：11。

111. ASML. (2025财年；2026年发布). \emph{Annual Report 2025}. \href{https://www.asml.com/investors/annual-report/2025/financials}{官方Financials}，\href{https://ourbrand.asml.com/m/8ab959d4926657b/original/asml-2025-annual-report-strategic-report-section.pdf}{US GAAP战略报告摘录}。两入口属于同份年报。
引用章节：12。

112. ASML. (无年). \emph{EUV lithography systems}. \href{https://www.asml.com/en/products/euv-lithography-systems}{官方技术说明}。动态页面以访问日期定位。
引用章节：12。

113. imec. (2025, September 22). \emph{Imec achieves new milestones in single patterning High NA EUV lithography for both damascene and direct metal etch metallization processes}. \href{https://www.imec-int.com/en/press/imec-achieves-new-milestones-single-patterning-high-na-euv-lithography-both-damascene-and}{研发参与机构原始发布}。
引用章节：12。

114. OECD. (2025). \emph{Mapping the semiconductor value chain: Working towards identifying dependencies and vulnerabilities}. OECD Science, Technology and Industry Policy Papers, No. 182. Paris: OECD Publishing. \href{https://doi.org/10.1787/4154cdbf-en}{DOI：10.1787/4154cdbf-en}。
引用章节：12。

115. Taiwan Semiconductor Manufacturing Company Limited. (2025财年；2026年申报). \emph{Annual Report on Form 20-F for the year ended December 31, 2025}. \href{https://investor.tsmc.com/sites/ir/sec-filings/2025_20F\%20Report.pdf}{公司原始申报全文}。
引用章节：12。

116. SkyWater Technology. (2020, November 16). \emph{Open Source ASICs take a Giant Leap Forward with the First Ever Open Foundry PDK}. \href{https://www.skywatertechnology.com/open-source-asics-take-a-giant-leap-forward-with-first-ever-open-foundry-pdk/}{原始发布}。
引用章节：12。

117. SkyWater与Google项目文档. (无年). \emph{Current Status—Experimental Preview}. \href{https://skywater-pdk.readthedocs.io/en/main/status.html}{原始状态说明}。不将文档当前状态与基础工艺商业使用混合。
引用章节：12。

118. SEMI. (无年). \emph{Standards Information and Control}. \href{https://www.semi.org/en/products-services/standards/information_and_control}{委员会说明与标准目录}。
引用章节：12。

119. Sandia National Laboratories. (2001). \emph{Partners Unveil First Extreme Ultraviolet Chip-making Machine}. \href{https://newsreleases.sandia.gov/partners-unveil-first-extreme-ultraviolet-chip-making-machine/}{当时机构发布}。
引用章节：12。

120. U.S. Department of Commerce / NIST. (2024, November 15). \emph{Biden-Harris Administration Announces CHIPS Incentives Award with TSMC Arizona to Secure U.S. Leadership in Advanced Semiconductor Technology}. \href{https://www.nist.gov/news-events/news/2024/11/biden-harris-administration-announces-chips-incentives-award-tsmc-arizona}{当时官方公告}。
引用章节：12。

121. Minister voor Buitenlandse Handel en Ontwikkelingshulp. (2025). \emph{Regeling van 13 januari 2025, BZ2411160, tot wijziging van de Regeling geavanceerde productieapparatuur voor halfgeleiders}. \emph{Staatscourant}, 2025, 1894. \href{https://zoek.officielebekendmakingen.nl/stcrt-2025-1894.pdf}{正式法规}。
引用章节：12。

122. Government of the Netherlands. (2025, January 15). \emph{Klever: export controls on advanced semiconductor manufacturing equipment to be tightened}. \href{https://www.government.nl/latest/news/2025/01/15/klever-export-controls-on-advanced-semiconductor-manufacturing-equipment-to-be-tightened}{当期官方解释}。与上项共同构成一个政策案例。
引用章节：12。

123. ILO. (2024). \emph{The Impact of Care Responsibilities on Women’s Labour Force Participation}. Statistical Brief, October. \href{https://www.ilo.org/sites/default/files/2024-10/GEDI-STAT\%20brief_formatted_28.10.24_final.pdf}{官方PDF}。
引用章节：13。

124. 国家统计局. 2024.《第三次全国时间利用调查公报》第一、第二、第三号，10月31日。\href{https://www.stats.gov.cn/xxgk/sjfb/zxfb2020/202410/t20241031_1957217.html}{第一号}、\href{https://www.stats.gov.cn/xxgk/sjfb/zxfb2020/202410/t20241031_1957216.html}{第二号}、\href{https://www.stats.gov.cn/xxgk/sjfb/zxfb2020/202410/t20241031_1957215.html}{第三号}。
引用章节：13。

125. Ajayi, K. F., A. Dao, and E. Koussoubé. (2022). \emph{The Effects of Childcare on Women and Children: Evidence from a Randomized Evaluation in Burkina Faso}. World Bank Policy Research Working Paper 10239. \href{https://documents1.worldbank.org/curated/en/099215111282254210/pdf/IDU09ee54b690865904f71089a60d76ad621601c.pdf}{原论文}。
引用章节：13, 16。

126. OECD. (2025). \emph{Trends in Adult Learning: New Data from the 2023 Survey of Adult Skills}. Getting Skills Right. \href{https://doi.org/10.1787/ec0624a6-en}{DOI:10.1787/ec0624a6-en}。
引用章节：13。

127. Alfonsi, L., O. Bandiera, V. Bassi, R. Burgess, I. Rasul, M. Sulaiman, and A. Vitali. (2020). “Tackling Youth Unemployment: Evidence from a Labor Market Experiment in Uganda.” \emph{Econometrica}, 88(6):2369–2414. \href{https://doi.org/10.3982/ECTA15959}{DOI:10.3982/ECTA15959}。
引用章节：13。

128. ITU. (2025). \emph{Measuring Digital Development: Facts and Figures 2025}. \href{https://www.itu.int/itu-d/reports/statistics/facts-figures-2025/}{官方在线报告}。
引用章节：13。

129. Hjort, J., and J. Poulsen. (2019). “The Arrival of Fast Internet and Employment in Africa.” \emph{American Economic Review}, 109(3):1032–1079. \href{https://doi.org/10.1257/aer.20161385}{DOI:10.1257/aer.20161385}。
引用章节：13。

130. Roodman, D. (2025). “The Arrival of Fast Internet and Employment in Africa: Comment.” arXiv:2401.13694v3，9月17日修订；初版2024。此处作为复核预印本引用。\href{https://arxiv.org/abs/2401.13694v3}{明确版本}。
引用章节：13。

131. Banza Lubaba Nkulu, C., L. Casas, V. Haufroid, T. De Putter, N. D. Saenen, T. Kayembe-Kitenge, P. Musa Obadia, D. Kyanika Wa Mukoma, J.-M. Lunda Ilunga, T. S. Nawrot, O. Luboya Numbi, E. Smolders, and B. Nemery. (2018). “Sustainability of Artisanal Mining of Cobalt in DR Congo.” \emph{Nature Sustainability}, 1:495–504. \href{https://doi.org/10.1038/s41893-018-0139-4}{DOI:10.1038/s41893-018-0139-4}，\href{https://pmc.ncbi.nlm.nih.gov/articles/PMC6166862/}{开放原稿}。
引用章节：13。

132. UNCTAD. (2024). \emph{Digital Economy Report 2024: Shaping an Environmentally Sustainable and Inclusive Digital Future}. UNCTAD/DER/2024，7月10日。\href{https://unctad.org/publication/digital-economy-report-2024}{官方入口}。
引用章节：13。

133. Baldé, C. P., et al. (2024). \emph{The Global E-waste Monitor 2024}. ITU and UNITAR. \href{https://ewastemonitor.info/the-global-e-waste-monitor-2024/}{发布机构报告入口}。
引用章节：13。

134. ILO. (2024). \emph{World Social Protection Report 2024–26: Universal Social Protection for Climate Action and a Just Transition}. \href{https://doi.org/10.54394/ZMDK5543}{DOI:10.54394/ZMDK5543}，\href{https://www.ilo.org/sites/default/files/2024-09/SOCPRO_WSPR_2024_Executive_Summary_EN_Web.pdf}{执行摘要}。
引用章节：13。

135. Suri, T., \& Jack, W. (2016). The long-run poverty and gender impacts of mobile money. \emph{Science, 354}(6317), 1288—1292. \href{https://doi.org/10.1126/science.aah5309}{DOI：10.1126/science.aah5309}，\href{https://www.jefftk.com/suri2016.pdf}{正式论文公开镜像}，\href{https://www.jefftk.com/suri2016-supp.pdf}{原始补充材料公开镜像}。
引用章节：13。

136. Cournot, A.-A. (1838). \emph{Recherches sur les principes mathématiques de la théorie des richesses}. Paris: L. Hachette。\href{https://catalogue.bnf.fr/ark:/12148/cb30280488q}{BnF原版目录}。本轮核验原书元数据，未声称通读法文原著；本章方程自行推导。
引用章节：14。

137. Ganapati, S., Shapiro, J. S., \& Walker, R. (2020). Energy cost pass-through in US manufacturing: Estimates and implications for carbon taxes. \emph{American Economic Journal: Applied Economics, 12}(2), 303–342. \href{https://doi.org/10.1257/app.20180474}{DOI:10.1257/app.20180474}。读取作者公开正式论文的摘要、引言、数据方法及传导讨论相关页；未重跑受限普查微观数据。
引用章节：14。

138. Kline, P., Petkova, N., Williams, H., \& Zidar, O. (2019). Who profits from patents? Rent-sharing at innovative firms. \emph{Quarterly Journal of Economics, 134}(3), 1343–1404. \href{https://doi.org/10.1093/qje/qjz011}{DOI:10.1093/qje/qjz011}。核验正式出版元数据，读取作者机构公开论文摘要及引言样本/结果相关页；未宣称复现全部方法与微观记录。
引用章节：14。

139. Godley, W., \& Lavoie, M. (2007). \emph{Monetary economics: An integrated approach to credit, money, income, production and wealth}. Palgrave Macmillan. \href{https://doi.org/10.1057/9780230626546}{DOI:10.1057/9780230626546}。沿用前轮方法背景与书目核验，不声称当前模型复现其完整体系。
引用章节：14。

140. Lane, M., Williams, M., \& Broecke, S. (2023). The impact of AI on the workplace: Main findings from the OECD AI surveys of employers and workers. OECD Social, Employment and Migration Working Papers, No. 288. \href{https://doi.org/10.1787/ea0a0fe1-en}{DOI}.
引用章节：15。

141. OECD. (2023). \emph{OECD Employment Outlook 2023: Artificial Intelligence and the Labour Market}. 第7章Social dialogue and collective bargaining in the age of artificial intelligence. \href{https://doi.org/10.1787/08785bba-en}{DOI}；\href{https://www.oecd.org/en/publications/2023/07/oecd-employment-outlook-2023_904bcef3/full-report/social-dialogue-and-collective-bargaining-in-the-age-of-artificial-intelligence_5828691a.html}{章节原文}.
引用章节：15。

142. Jäger, S., Schoefer, B., \& Heining, J. (2021). Labor in the boardroom. \emph{The Quarterly Journal of Economics, 136}(2), 669–725. \href{https://doi.org/10.1093/qje/qjaa038}{DOI}.
引用章节：15, 16。

143. Gorton, G., \& Schmid, F. A. (2004). Capital, labor, and the firm: A study of German codetermination. \emph{Journal of the European Economic Association, 2}(5), 863–905. \href{https://doi.org/10.1162/1542476042782260}{DOI}.
引用章节：15, 16。

144. International Labour Organization. (2002). Promotion of Cooperatives Recommendation, 2002 (No. 193). Adopted 20 June 2002. \href{https://webapps.ilo.org/public/english/standards/relm/ilc/ilc90/pdf/rec-193.pdf}{原始文本}.
引用章节：15。

145. Pencavel, J. H. (2012). Worker cooperatives and democratic governance. IZA Discussion Paper No. 6932. 本章使用2012工作稿，不与2013书章视为独立研究。\href{https://docs.iza.org/dp6932.pdf}{原稿}.
引用章节：15。

146. Acemoglu, D., Gitmez, A. A., \& Shadmehr, M. (2026). \emph{Automation and Repression}. NBER Working Paper 35336, June. \href{https://doi.org/10.3386/w35336}{DOI:10.3386/w35336}；\href{https://economics.mit.edu/sites/default/files/2026-06/Automation\%20and\%20Repression.pdf}{2026年6月5日作者稿}。工作论文；核读模型、政策比较、动态条件及Lemma 4证明。本文参与约束模型另有设定，效率扩展只重算政策收益的部分敏感性，不声明复现全部原论文。
引用章节：1, 4, 14, 15, 16。

147. 刘镇锐（发表日期未核实）。对\emph{Automation and Repression}的评论［知乎回答］。\href{https://www.zhihu.com/question/2083147887706366153/answer/2088453310651032052}{原回答链接}。依据用户提供的评论文本分析；网页直接访问未成功，不将评论视为实证研究或同行评审论文。
引用章节：14。

\end{document}
